% 第17章：机器人导航与具身智能
\chapter{机器人导航与具身智能}
\label{ch:17}

\paragraph*{本章阅读锚点}
本章的锚点是机器人可用的地图。原始 3DGS 主要服务渲染；机器人方法还需要可定位、可规划、可更新，并能表达不确定性和可通行空间。


机器人导航与具身智能需要一种既能提供精确几何信息、又能支持语义理解和实时渲染的环境表示。传统方法使用占据栅格或稀疏点云，在几何精度和视觉保真度之间存在根本矛盾。3DGS 提供了一种统一解决方案：高斯基元既能编码几何（通过位置和协方差），又能承载语义（通过特征嵌入），同时支持照片级渲染——这使得 3DGS 成为机器人仿真训练、场景理解和导航规划的通用场景表示。

3DGS 在机器人领域的应用可以归纳为五条技术路线。第一条路线是\textbf{导航与路径规划}，将 3DGS 作为环境地图表示，支持安全导航、语义目标搜索和轨迹优化。第二条路线是\textbf{机械臂操控}，利用 3DGS 的可编辑渲染能力生成训练数据、预测操作结果和规划抓取动作。第三条路线是\textbf{具身感知与交互}，将 3DGS 与基础模型结合，构建多模态空间记忆，支持遥操作、人机协作和主动感知。第四条路线是\textbf{仿真与真实迁移}（sim-to-real），利用 3DGS 构建高保真数字孪生仿真环境，大幅降低机器人策略学习的真实世界数据需求。第五条路线是\textbf{强化学习与场景表示}，将 3DGS 作为强化学习智能体的环境表示，提升样本效率和泛化能力。

本章围绕上述五条技术路线，系统梳理了 170 篇相关论文。第~\ref{sec:17_1}~节讨论导航与路径规划（56 篇），覆盖 3DGS 地图构建、视觉导航、语义搜索、安全规划和无人机导航。第~\ref{sec:17_2}~节介绍机械臂操控（42 篇），涵盖抓取规划、操作策略学习、功能区域推理、可变形物体操控和多指灵巧手。第~\ref{sec:17_3}~节深入具身感知与交互（34 篇），探讨遥操作系统、数字孪生、多模态记忆和主动感知。第~\ref{sec:17_4}~节分析仿真与真实迁移（32 篇），覆盖 Real2Sim2Real 框架、数据增强和评估基准。第~\ref{sec:17_5}~节讨论强化学习与场景表示（6 篇）。第~\ref{sec:17_summary}~节总结全章。

\begin{figure}[htbp]
\centering
\resizebox{\textwidth}{!}{\begin{tikzpicture}[
  route/.style={rectangle, rounded corners=5pt, draw=#1, fill=#1!10,
    text width=3.2cm, minimum height=0.9cm, align=center, font=\small},
  method/.style={rectangle, rounded corners=3pt, draw=#1, fill=#1!6,
    font=\scriptsize, text width=3.0cm, align=center, minimum height=0.6cm},
  arr/.style={-{Stealth[length=4pt]}, techgray!60},
  >=Stealth
]

  \node[route=techgray, text width=4cm, font=\small\bfseries] (root) at (25.5, 5.6) {机器人导航与具身智能方法分类};
  \node[route=techblue] (b0) at (8.8, 3.1) {导航与路径规划};
  \node[method=techblue] (b0_l0) at (2.0, 0.6) {地图构建导航\\{\tiny ATLAS Navigator, NavGSim, NVSim}};
  \node[method=techblue] (b0_l1) at (6.6, 0.6) {语义探索安全\\{\tiny VISTA, SplatSearch, GSplatVNM}};
  \node[method=techblue] (b0_l2) at (10.9, 0.6) {无人机导航\\{\tiny UAV-GS, DroneNav-GS}};
  \node[route=techgreen] (b1) at (19.3, 3.1) {机械臂操控};
  \node[method=techgreen] (b1_l0) at (15.2, 0.6) {功能推理抓取\\{\tiny GraspGS, SplatGrasp, ManipGS}};
  \node[method=techgreen] (b1_l1) at (19.3, 0.6) {灵巧操控\\{\tiny DexGS, DextrousSplat}};
  \node[route=techorange] (b2) at (25.5, 3.1) {具身感知与交互};
  \node[method=techorange] (b2_l0) at (23.2, 0.6) {遥操作孪生\\{\tiny TeleGS, Gello-GS, MimicGS}};
  \node[method=techorange] (b2_l1) at (27.3, 0.6) {多模态感知\\{\tiny ConceptGS, SceneGraph-GS}};
  \node[route=techpurple] (b3) at (32.0, 3.1) {仿真与真实迁移};
  \node[method=techpurple] (b3_l0) at (32.0, 0.6) {Real2Sim2Real\\{\tiny SplatSim, R2D2-GS}};
  \node[method=techpurple] (b3_l1) at (36.7, 0.6) {Sim2Real\\{\tiny RoboGSim, GenSim-GS}};
  \node[route=techblue] (b4) at (40.6, 3.1) {强化学习};
  \node[method=techblue] (b4_l0) at (40.6, 0.6) {场景表示\\{\tiny RL-GS, NavRL-GS}};
  \node[method=techblue] (b4_l1) at (44.0, 0.6) {RL训练\\{\tiny GSRL, SceneRL-GS}};
  \draw[arr] (root) -- (b0);
  \draw[arr] (root) -- (b1);
  \draw[arr] (root) -- (b2);
  \draw[arr] (root) -- (b3);
  \draw[arr] (root) -- (b4);
  \draw[arr] (b0) -- (b0_l0);
  \draw[arr] (b0) -- (b0_l1);
  \draw[arr] (b0) -- (b0_l2);
  \draw[arr] (b1) -- (b1_l0);
  \draw[arr] (b1) -- (b1_l1);
  \draw[arr] (b2) -- (b2_l0);
  \draw[arr] (b2) -- (b2_l1);
  \draw[arr] (b3) -- (b3_l0);
  \draw[arr] (b3) -- (b3_l1);
  \draw[arr] (b4) -- (b4_l0);
  \draw[arr] (b4) -- (b4_l1);
\end{tikzpicture}}
\caption{机器人导航与具身智能方法分类}
\label{fig:ch17_robot_taxonomy}
\end{figure}

% ======================================================================
\section{导航与路径规划}
\label{sec:17_1}

导航是机器人最基础也最核心的能力之一。传统导航系统依赖激光雷达点云或占据栅格地图进行定位和路径规划，但这些表示缺乏视觉语义信息，难以支持语言引导的目标搜索和开放词汇导航。3DGS 为导航系统带来了新的可能性：它既能提供几何精确的环境模型用于碰撞检测和路径规划，又能通过照片级真实感渲染生成任意视角的图像，支持视觉导航策略的训练和部署。本节将 3DGS 导航研究分为五个子方向：3DGS 地图构建与定位、视觉目标导航、语义搜索与主动探索、安全规划与碰撞避免、以及无人机导航。

\begin{figure}[htbp]
\centering
\resizebox{\textwidth}{!}{\begin{tikzpicture}[
  box/.style={rectangle, rounded corners=4pt, draw=#1, fill=#1!10,
    text width=3.2cm, align=center, font=\small, minimum height=1cm},
  sbox/.style={rectangle, rounded corners=3pt, draw=#1!60, fill=#1!5,
    font=\scriptsize, text width=2.4cm, align=center, minimum height=0.5cm},
  arr/.style={-{Stealth[length=6pt]}, thick, techgray},
  darr/.style={-{Stealth[length=6pt]}, thick, techgray, dashed},
  >=Stealth
]

  \node[box=techblue] (n1) at (2.8, 0.6) {3DGS场景(已建图)};
  \node[box=techgreen] (n2) at (9.7, 0.6) {场景理解(语义/目标)};
  \node[box=techorange] (n3) at (16.5, 0.6) {路径规划(A*/RRT)};
  \node[box=techpurple] (n4) at (23.3, 0.6) {运动控制(避障/导航)};
  \node[box=techblue] (n5) at (29.8, 0.6) {执行(移动/抓取)};
  \node[font=\scriptsize, text=techgray] at (9.7, 1.8) {LLM指令};
  \draw[arr] (n1) -- (n2);
  \draw[arr] (n2) -- (n3);
  \draw[arr] (n3) -- (n4);
  \draw[arr] (n4) -- (n5);
  \draw[darr] (n5) to[bend right=30] node[midway, font=\scriptsize, text=techgray] {在线更新} (n1);
\end{tikzpicture}}
\caption{3DGS 导航 pipeline：从场景建图到路径规划与运动控制}
\label{fig:ch17_nav_pipeline}
\end{figure}

\subsection{地图构建与目标导航}
\label{ss:17_1a}

本节讨论两条技术路线：（1）3DGS 地图构建与定位；（2）视觉目标导航。

\subsubsection{3DGS 地图构建与定位}

本路线介绍将 3DGS 作为导航底层地图表示的构建与定位方法。将 3DGS 作为导航的底层地图表示，需要解决地图构建效率、定位精度和多机器人协同等问题。

\paragraph{在线语义地图构建}

ATLAS Navigator（2025.02）~\cite{2502_20386} 修改了 3DGS 的地图表示与规划环节——用语言嵌入的层次化子图替代标准全局 3DGS 地图、用层次化离散-连续双层规划替代单层路径搜索——3DGS 的核心可微渲染管线保持不变。

\textbf{基线与方法论}：传统导航系统依赖占据栅格进行碰撞检测，但缺乏语义信息支持语言引导的任务规划；而纯语义地图（如 ConceptGraphs）虽能支持高层推理，却无法提供稠密几何用于安全路径规划。ATLAS Navigator 的核心思想是构建\textbf{层次化语言嵌入高斯泼溅}——在单一 3DGS 地图中同时编码稠密几何和压缩语言特征，使地图既能用于碰撞检测，又能响应自然语言任务查询。其设计动机在于：机器人导航同时需要"安全"（避免碰撞）和"语义"（找到目标）两种能力，而此前的表示方法只能满足其一。ATLAS Navigator 通过在每个高斯上附加紧凑的语言特征向量，使同一张地图既可用于渲染图像、检测碰撞，也可用于查询"椅子在哪里""找到红色靠垫"等开放词汇任务。

\textbf{技术细节}：系统在地图构建环节采用增量式子图（submap）策略——将全局地图划分为以 6-DoF 锚点位姿为参考帧的局部子图，每个子图独立存储高斯参数。子图的设计动机是：随着机器人探索，高斯数量迅速增长至百万量级，若在全局地图上做碰撞检查，计算代价极高。通过仅加载机器人周围半径 $R_\text{loc}$ 内的子图（例如仅处理 369,513 个高斯而非全局 247 万个），碰撞检查的计算复杂度从 $O(N)$ 降至 $O(N_\text{local})$。此外，子图的锚点位姿可由外部 SLAM（Simultaneous Localization and Mapping，同时定位与建图）系统的位姿图优化（Pose Graph Optimization, PGO）结果直接更新，从而天然支持回环检测后的全局地图校正，而无需重新优化所有高斯参数。

每个高斯附加 $N_c=24$ 维压缩语言特征向量，用于语义嵌入。直觉上，这相当于为每个 3D 点赋予一个简短的"语义指纹"，使得地图不仅能回答"这里有多远"（几何），还能回答"这里是什么"（语义）。具体而言，系统通过增量主成分分析（Incremental PCA, IPCA）——一种可逐批处理数据的主成分分析方法，无需一次性加载全部训练样本——在 COCO 2017 数据集上离线拟合 512 维 CLIP 特征的主方向，运行时仅需将特征投影到前 24 个主成分上即可。这种离线-在线分离的设计保证了跨场景泛化能力：IPCA 基底在 COCO 上预训练一次后固定不变，新场景无需重新训练。

层次化规划将地图信息分三层组织。

第一层是\textbf{度量-语义聚类}。直觉上，如果两个高斯既在空间上接近、又在语义上相似（例如都属于同一把椅子），它们就应该被归为同一个对象。为此，系统定义了一个同时编码空间距离和语义相似度的距离度量：

\begin{equation}
q = q_e + \lambda(1 - q_s)
\end{equation}

其中 $q_e$ 为欧氏距离，$q_s$ 为压缩语言特征的余弦相似度，$\lambda$ 为权重参数。直观理解：$q_s$ 越大（语义越相似），$(1-q_s)$ 越小，总距离 $q$ 就越小，意味着这两个高斯更可能属于同一语义对象。基于此距离矩阵，系统在每个子图内进行凝聚聚类（agglomerative clustering，一种自底向上的层次聚类方法，先将每个点视为一类，再逐步合并最相似的类对），构建子图-区域-对象三层层次树。给定任务后，系统计算每个叶节点（对象）与任务嵌入的语义相关度，再沿层次树逐层向上求和传播，得到区域级和子图级的任务效用分数。

第二层是\textbf{离散规划}。离散规划器在子图级别的稀疏图上工作，每个节点对应一个子图的质心，边权为子图间距。它求解一个带旅行预算约束的优化问题：在总行程不超过 $\mathcal{B}$ 的条件下，寻找效用最大的视点序列，用于高层路径决策。

第三层是\textbf{连续轨迹规划}。连续规划器基于非完整轮式机器人（unicycle）动力学模型，在稠密高斯场中搜索动态可行轨迹。碰撞检测的核心思路是：将机器人和每个高斯点都建模为三维正态随机变量 $R \sim \mathcal{N}(\mu_\text{rob}, \Sigma_\text{rob})$、$G_i \sim \mathcal{N}(\mu_i, \Sigma_i)$，碰撞判定等价于两者的距离是否小于碰撞半径 $r_\text{coll}$。为了在百万量级高斯中高效计算碰撞概率，系统通过二分搜索找到最小局部半径 $R_\text{loc}$，使得半径外所有高斯的总碰撞概率不超过容差 $p_\text{tol}$——这保证了只需检查局部区域内的少量高斯，即可近似全局碰撞概率。

该系统在 Clearpath Jackal 平台上以 1 Hz 地图更新率和 1 Hz 规划频率运行，在室内杂乱环境和公里级户外环境中均实现了约 60\% 的导航成功率（竞争比率 $\sim$0.59，即实际路径长度约为最短路径的 1.7 倍）。

ATLAS Navigator 展示了语言嵌入 3DGS 地图在导航中的威力，但有两个关键局限：其一，它假设环境已经预先建图，无法在机器人首次进入的未知环境中运行；其二，它不量化地图的几何不确定性——即无法判断哪些区域尚未被充分观测、哪些区域还需要更多视角来完善重建。当目标物体不在初始视野范围内时，仅靠语义搜索而无几何引导可能导致探索效率急剧下降。VISTA（2025.07）~\cite{2507_01125} 正是在这一背景下提出的增量方法：它继承了 ATLAS Navigator 的语言嵌入建图思路（将 CLIP 语义特征蒸馏到 3DGS 中），但将整个范式从"在已有地图上搜索"推进到"边建图边搜索"——在未知环境中同步完成语义建图和目标导向的主动探索。

\textbf{在线语义建图}。VISTA 基于 SplatBridge（NerfBridge 的 3DGS 扩展）实现实时 3DGS 训练：机器人采集的 RGB-D 图像和位姿持续汇入流式数据集，3DGS 参数同步优化。语义建图方面，VISTA 通过与 3DGS 并行训练的多分辨率哈希网格 + MLP（多层感知机，Multi-Layer Perceptron）语义场 $f: \mathbb{R}^3 \to \mathbb{R}^l$，将 CLIP（Contrastive Language-Image Pre-training，对比语言-图像预训练）图像嵌入蒸馏到 3D 空间中的每个查询点——优化目标为预测嵌入与真实 CLIP 嵌入之间的均方误差（MSE）和余弦相似度联合损失。与 ATLAS Navigator 的离线 IPCA 压缩策略不同，VISTA 的语义场在线增量化训练，无需预训练基底，但语义场参数与 3DGS 参数同步更新。训练过程中，系统实时发布包含 RGB 颜色和语义嵌入的点云，再将其体素化为以机器人为中心的固定分辨率体素网格。通过从训练相机位姿出发的射线体素遍历（voxel traversal），空间被划分为\textbf{占据}（含点云信息）、\textbf{自由}（被射线穿过但无点云）和\textbf{未观测}（无射线到达）三类——这一设计将地图复杂度与高斯数量解耦，使其不随场景规模增长。

\textbf{VISTA-Score：信息增益量化}。核心挑战在于如何同时量化"几何上哪里还没看清"和"语义上哪里可能藏着目标"。为此，VISTA 提出视点-语义覆盖率（VISTA-Score）指标，融合几何视角多样性和任务语义相关性两个分量。

几何信息增益 $G_I(x)$ 的直觉是：如果一个候选视角的射线方向与该体素此前被观测到的所有方向都差异很大，则该视角能提供新的几何信息。具体而言，对候选相机位姿的每条射线 $d_x^n$，计算它与该体素存储的所有历史观测方向 $d_v$ 的点积，取最大值（即最近似方向），再归一化到 $[0,1]$：

\begin{equation}
g_I(d_x^n) = \frac{-\max_{d_v}(d_v^\top d_x^n) + 1}{2}
\end{equation}

值越接近 1 表示该方向越新颖（与所有历史方向越不同）。候选位姿的几何增益为所有像素增益的均值 $G_I(x) = \frac{1}{N_r}\sum_{n=1}^{N_r} g_I(d_x^n)$。语义信息增益 $G_S(x)$ 则通过体素遍历渲染语义值图并取像素均值得到。

沿路径 $\Psi$ 的 VISTA-Score 综合两个分量，并以折扣因子对远期航点降权：

\begin{equation}
G(\Psi) = \sum_{k=1}^{K} \gamma^{K-k} \cdot \bigl(c \cdot G_I(x_k) + G_S(x_k)\bigr)
\end{equation}

其中 $K$ 为路径中航点总数，$k$ 为航点索引，$\gamma$ 为折扣因子（近期航点权重更高），$c$ 为几何信息增益权重。

\textbf{MPC（Model Predictive Control，模型预测控制）探索规划}。探索规划采用 GMM（高斯混合模型）采样策略生成偏向高信息增益区域的候选轨迹。首先将 3D 体素网格沿高度方向压缩为 2D 表示，从中提取前沿单元（frontier，自由-未观测边界）和语义 top-$m$ 单元，分别拟合 GMM 分布作为目标位置采样器。再以 Dijkstra 最短路径连接机器人当前位置与各采样目标，生成候选轨迹集合。候选轨迹经 VISTA-Score 评分后，选择得分最高的路径执行。整个过程在 MPC 循环中以高频重规划。

值得注意的是，与 ATLAS Navigator 等方法的"探索-利用切换"模式不同，VISTA 始终同时考虑几何和语义信息增益，不做模式切换。随着 MPC 重规划轮次 $i$ 增加，几何权重按 $c \leftarrow \gamma^i \cdot c$ 衰减，自然地将探索重心从几何覆盖逐步过渡到语义搜索——这是一种平滑的注意力转移，而非硬切换。在四旋翼硬件实验中，VISTA 在目标物体被遮挡的困难场景下取得了 100\% 的成功率，而纯语义基线和纯几何基线均仅为 16.67\%，实现了 6 倍的成功率提升。

\paragraph{大规模导航地图}

NavGSim（2026.03）~\cite{2603_15186} 修改了 3DGS 的大规模场景渲染与碰撞表示环节——用层次化 3DGS（H3DGS）替代标准 3DGS、用高斯切片碰撞检测替代传统网格碰撞——高斯泼溅的核心渲染管线保持不变。

\textbf{动机与问题}。现有导航仿真器（如 Habitat、AI2-THOR）依赖重建或合成网格，在大规模场景（多房间乃至整层楼）中视觉渲染质量不足；而标准 3DGS 虽然渲染质量高，却无法高效处理数百平方米规模的场景，且缺乏碰撞检测所需的体积占据表示。NavGSim 的核心思路是：利用层次化 3DGS 解决大规模渲染的可扩展性问题，再通过解析的高斯切片技术从连续概率场中提取可导航区域，从而在单一框架内同时实现照片级真实感渲染和碰撞检测。

\textbf{层次化 3DGS 场景建模}。为处理数百至数千平方米的大规模场景，NavGSim 采用分治策略：将输入图像集按空间划分为多个块（chunk），每个块独立优化局部 3DGS 模型；局部模型训练完成后，通过合并、节点清理与对齐构建全局层次结构。该层次结构以树形组织：叶节点用密集各向异性高斯编码细粒度视觉细节，内部节点将子节点信息聚合为粗粒度表示，支持细节层次（Level-of-Detail, LOD）渲染——渲染时根据视点距离动态选择不同层次的高斯，并在过渡区域平滑插值。这种设计使系统能够利用多 GPU 并行处理覆盖约 $700\,\text{m}^2$ 的真实办公/工厂环境。

\textbf{高斯切片碰撞检测}。在获得可渲染的 3DGS 场景后，NavGSim 需要从中提取碰撞信息以支持路径规划。其核心思路是：不将 3D 高斯直接投影为 2D 椭圆（这会丢失高度维度的碰撞信息），而是沿机器人高度范围以固定间距（例如每 $10\,\text{cm}$）定义一组水平切片平面，每个平面与 3D 高斯场求交，解析地计算出切片处的 2D 高斯分布。

具体而言，给定 3D 高斯的定义：

\begin{equation}
G_{\text{3D}}(\mathbf{x}) = e^{-\frac{1}{2}(\mathbf{x}-\boldsymbol{\mu})^T \boldsymbol{\Sigma}_{\text{3D}}^{-1} (\mathbf{x}-\boldsymbol{\mu})}
\end{equation}

其中 $\mathbf{x}=(x,y,z)$，$\boldsymbol{\Sigma}_{\text{3D}}^{-1}$ 为 $3\times3$ 精度矩阵。直觉上，这个公式描述了以 $\boldsymbol{\mu}$ 为中心、沿三个轴方向按协方差矩阵展开的一个三维"模糊椭球"。在高度 $z$ 处的切片 2D 高斯为：

\begin{equation}
G_{\text{2D}}(\mathbf{x}_p, z) = e^{-\frac{1}{2}\frac{(z-\mu_z)^2}{\sigma^2}} \cdot e^{-\frac{1}{2}[\mathbf{x}_p - \tilde{\boldsymbol{\mu}}(z)]^T \boldsymbol{\Sigma}_{\text{2D}}^{-1} [\mathbf{x}_p - \tilde{\boldsymbol{\mu}}(z)]}
\end{equation}

其中 $\mathbf{x}_p=(x,y)$，$\sigma$ 由 3D 协方差矩阵元素计算得到，$\boldsymbol{\Sigma}_{\text{2D}}^{-1}$ 为切片后的 $2\times2$ 精度矩阵，$\tilde{\boldsymbol{\mu}}(z)$ 为高度相关的偏移中心。简而言之，公式左半部分的衰减项控制切片"有多厚"——当 $z=\mu_z$ 时取得最大切片，偏离越远切片越小；右半部分给出切片在水平面上的 2D 高斯分布。该方法受 4D-Rotor-GS 按时间平面切片 4D 高斯的启发，将"沿时间切片"推广为"沿高度切片"。由于 NavGSim 中高斯体积较小，碰撞体积误差被限制在 $5\,\text{cm}$ 以内。

\textbf{VLA 导航策略训练与 Sim-to-Real 迁移}。在 NavGSim 仿真器中，研究者收集了 4,500 条导航轨迹用于训练视觉-语言-动作（VLA）模型（基于 Uni-NaVid 扩展），并在 9 个语义目标上评估。实验表明，基于 H3DGS 构建的仿真器训练的模型在真实世界部署中保持高成功率；而用标准 3DGS（渲染质量较低）训练的模型，真实世界性能下降约 $81\%$。该系统在 Unitree Go2 四足机器人上以 $8\,\text{Hz}$ 闭环频率运行，验证了从仿真到真实的有效迁移。

NavGSim 证明了层次化 3DGS 可以在大规模场景中同时实现照片级渲染和碰撞检测，但其训练数据仍来自预先扫描的高质量多视角图像。实际机器人部署场景中，导航仿真器往往只能获取机器人沿行走轨迹的前视相机图像序列——地面区域的观测极为稀疏，导致 3DGS 重建在地板上产生大量浮游伪影（floater），直接破坏可通行区域的视觉可靠性。NVSim（2025.10）~\cite{2510_24335} 正是从这一实际约束出发，沿袭 NavGSim「分治大规模场景 + 从高斯场中提取导航信息」的基本思路，但在场景分解策略、地面伪影处理和碰撞检测机制三个环节做出根本性差异化设计。

\textbf{场景分解}。NVSim 同样采用分治策略应对大规模场景，但与 NavGSim 构建层次化 LOD 树（叶节点编码细节、内部节点聚合粗粒度表示、多 GPU 并行渲染）不同，NVSim 直接对相机轨迹做 Agglomerative Clustering 划分为若干独立子图（submap），子图之间通过空间邻域重叠（距离阈值 $\delta$ 内的邻域图像共享）确保边界处的多视角覆盖和几何连续性，各子图独立训练后直接合并渲染。这种设计舍弃了层次化 LOD 结构，但换来更简单的训练流程——无需构建全局树形索引和层次间对齐，适用于仅需拓扑导航图而不需要实时 LOD 渲染的仿真场景。

\textbf{楼层感知高斯泼溅}。这是 NVSim 相对于 NavGSim 最核心的新增模块。机器人前视相机的固定视角导致地面区域观测极为稀疏，标准 3DGS 在地板上产生浮游伪影和模糊斑块，严重影响导航环境的可靠性。NVSim 提出混合楼层分割策略融合语义分割（EVF-SAM，文本提示 ``floor''）与几何法线估计（GeoWizard）两种互补信息源——前者语义精确但覆盖不全，后者能捕获所有水平面但无法区分地板与桌面；先取两者的交集获得高置信度候选像素，再以其为点提示驱动 SAM2 在法线图（而非 RGB 图，从而获得光照不变性）上推理最终楼层掩码。在此基础上，NVSim 的核心设计是不再用高斯基元建模地面，转而用基于球谐函数的背景 MLP（受 Splatfacto-W 天空建模启发）表示低频纹理地面，并通过两个损失函数协同约束：Floor Suppression Loss 直接惩罚楼层掩码区域内的高斯 alpha 值推向零，阻止不稳定高斯生成；Background Floor Loss 引导背景模型在楼层区域忠实还原地面颜色和纹理。此外，NVSim 将 Splatfacto-W 的离散外观嵌入索引替换为以相机 3D 位置为输入的连续 MLP，使新视角下的外观插值更加合理。

\textbf{无网格可通行性检查}。NavGSim 通过解析的高斯切片在固定高度平面上计算 2D 碰撞概率场，属于基于几何的体积碰撞检测。NVSim 则走上完全不同的路线——直接从渲染结果推断通行性，无需任何几何切片操作。具体而言，算法以广度优先搜索从任意起始视点出发，在 8 个基本方向上逐步扩展可通行空间：在每个候选方向上渲染球面视角，提取累积 alpha 图和零样本法线估计（GeoWizard），在预定义关注区域（RoI）内判断两个条件——alpha 均值低于阈值 $\tau_{alpha}$（表示该方向无障碍物遮挡）且法线朝上分量高于阈值 $\tau_{normal}$（表示该方向存在可通行地面）；双条件同时满足时在该方向上放置新视点节点，最终自动构建拓扑导航图。这种方法将碰撞检测从 NavGSim 的解析体积计算转变为基于渲染外观的视觉推断，在无网格场景表示下绕过了显式几何重建的需求，以球面渲染的 alpha 通道和法线语义取代了高度切片和体积碰撞场。

\paragraph{多机器人地图融合}

上述方法各自构建单一机器人的 3DGS 地图，但在实际大规模场景中，单台机器人的电池续航和工作空间均有限，多机器人协同建图与地图融合成为必然需求。然而，多机器人场景下各子图之间通常没有共享的相机位姿、图像或相对变换信息，传统的配准方法（如 GaussReg 需要单独训练几何 Transformer 和 CNN，PhotoReg 需要访问相机关键帧）难以直接适用。

与 GaussReg 依赖额外训练的几何 Transformer 和 CNN、PhotoReg 需要访问相机关键帧不同，SIREN（2025.02）~\cite{2502_06519} 从根本上取消了所有初始化前提——无需相机位姿、源图像或地图间变换的初始估计，仅凭高斯基元自身的语义特征即可完成跨地图配准。其核心洞察在于：真实场景中蕴含丰富的语义信息（车辆、行人、植被等），这些语义可以作为跨地图配准的可靠锚点。具体而言，SIREN 设计了三阶段语义级联框架（Semantic Cascade Framework, SCF），逐步从粗到精完成配准，每一阶段都利用语义信息替代 GaussReg/PhotoReg 所依赖的外部数据源。

\textbf{阶段一：语义特征蒸馏与对应关系建立}。SIREN 同时训练一个多分辨率神经哈希网格语义场 $\psi: \mathbb{R}^3 \mapsto \mathbb{R}^d$ 与 GSplat 模型，利用 GSplat 自身提供的高精度深度估计将 CLIP（对比语言-图像预训练，一种将图像和文本映射到共享嵌入空间的基础模型）的 2D 语义特征反投影至 3D 空间并蒸馏到每个高斯基元，避免了传统 NeRF 方法中训练 proposal network 的额外计算和显存开销。随后，基于开放词汇查询的语义相关性评分识别地图中的语义丰富区域——这些区域的配准问题本身更好定义（better posed）——并在源和目标地图的高斯基元之间通过余弦相似度建立候选对应关系集合 $\mathcal{E}$。

\textbf{阶段二：高斯到高斯的几何粗配准}。以对应关系集合 $\mathcal{E}$ 为基础，SIREN 构建加权最小二乘优化问题，联合约束对应高斯的均值（位置）和协方差（形状），求解包含旋转、平移和缩放的非刚性相似变换。其中权重 $w_{ij}$ 定义为对应高斯基元间语义特征的余弦相似度，使语义一致性高的对应对优化结果产生更大影响，有效抑制离群匹配的干扰。该优化问题可通过 SVD 分解高效获得闭式最优解。

\textbf{阶段三：基于新视角合成的光度精配准}。在粗配准的基础上，SIREN 利用 GSplat 的新视角合成能力，在配准位姿处从两个子地图分别渲染候选图像。为避免低质量渲染干扰精配准，SIREN 引入基于语义的图像滤波器筛选可靠候选图像对，再通过图像级语义特征匹配建立初始对应，最终以局部特征匹配和运动恢复结构（Structure from Motion, SfM）方法计算高精度相对变换，生成高保真的融合地图。

在包括机械臂、无人机和四足机器人三种不同平台采集的真实数据上，SIREN 均展现出显著优于竞争基线的配准精度。在最具挑战性的四足机器人建图场景中，SIREN 实现了约 90 倍更小的旋转误差、300 倍更小的平移误差和 44 倍更小的缩放误差，而竞争方法在该场景下几乎无法收敛。

\paragraph{LiDAR 增强定位}

上述方法聚焦于从 3DGS 地图中实现导航功能，但地图本身的几何精度同样影响定位可靠性——当高斯点偏离真实表面时，基于地图的位姿估计也会产生偏差。GeomGS（2025.01）~\cite{2501_13417} 正是针对这一问题，聚焦于提升 3DGS 地图的几何精确性以服务于自动驾驶场景中的 LiDAR 定位。其核心动机在于：现有 LiDAR-3DGS 融合方法（如 DrivingGaussian）仅对所有高斯点施加统一的欧氏距离约束，要求高斯基元贴近 LiDAR 点云。然而，场景中大量高斯点对应的是 LiDAR 无法扫描到的区域（如天空、远处建筑），对它们施加同样的距离约束会严重扭曲 3DGS 的几何结构，反而降低重建质量。

为此，GeomGS 提出\textbf{几何置信度分数}（Geometric Confidence Score, GCS）$\gamma_i \in (0,1)$，作为每个高斯基元的新增可优化参数，用于衡量该点的结构可靠性。具体而言，GCS 通过非对称 Sigmoid 函数 $\sigma_{\text{asym}}(x) = 1/(1 + e^{k(x-d)})$ 将高斯基元到最近 LiDAR 点的距离映射为置信度——距离在阈值 $d$ 以内的点获得接近 1 的高置信度，而天空等远离 LiDAR 表面的点则被赋予接近 0 的低置信度。非对称设计确保了距离衰减的陡峭性，使置信度在阈值附近快速过渡。在此基础上，GeomGS 设计\textbf{概率距离约束}损失 $\mathcal{L}_{\text{prob}} = \frac{1}{N}\sum_{i}(\ln(1-\gamma_i) + d_i/(1-\gamma_i))$，其中 $d_i$ 为高斯基元到最近 LiDAR 点的距离。该损失的关键性质在于：当 $\gamma_i \to 1$ 时，$1/(1-\gamma_i)$ 急剧增大，距离项 $d_i$ 被强烈惩罚，迫使高置信度点紧贴 LiDAR 表面；而当 $\gamma_i \to 0$ 时，约束自然松弛，低置信度点无需被迫靠近 LiDAR，保留其对图像渲染质量的贡献。这种差异化约束策略使 GeomGS 在提升几何精度的同时不牺牲光度性能——在 KITTI-360 和 KITTI 数据集上，PSNR 和 SSIM 均优于使用相同 LiDAR 初始化的基线方法。

GeomGS 进一步将 GCS 的几何感知能力延伸到位姿估计，提出 GCS 加权 ICP 与光度精化的迭代定位框架。核心思路是利用 GCS 作为加权 ICP 的点对权重矩阵 $\mathbf{W}=\text{diag}(\gamma)$，使几何可靠的高斯点在位姿对齐中占据主导地位，抑制低置信度离群点的干扰，通过 SVD 分解求解最优旋转和平移。随后，在 ICP 给出的初始位姿处渲染图像，与真实图像进行光度对比以精化位姿。两个步骤交替迭代，逐步收敛至精确位姿。这种将 LiDAR 几何约束与 3DGS 渲染能力统一在同一置信度框架下的设计，在多个自动驾驶基准上实现了当时最优的几何与定位性能。

\subsubsection{视觉目标导航}

本路线介绍基于 3DGS 可渲染特性的视觉目标导航方法。视觉目标导航（visual goal navigation）要求机器人仅使用视觉信息（图像或语言描述）导航至目标位置。3DGS 的可渲染特性为这类任务提供了独特的优势——可以从已构建的 3DGS 地图中渲染任意视角的图像，用于目标匹配和路径规划。

\paragraph{图像目标导航}

IGL-Nav（2025.08）~\cite{2508_00823} 修改了 3DGS 的\textbf{地图构建与定位环节}——用\textbf{前馈单目增量重建}替代离线多视角优化，以\textbf{3D 卷积等价的体素匹配}替代逐点几何搜索——3DGS 的高斯表示与可微渲染保持不变。

\textbf{核心问题与动机。}图像目标导航（image-goal navigation）要求机器人仅凭一张目标照片导航至拍摄位置。传统方法将图像匹配局限在 2D 特征空间，无法处理目标照片与机器人相机视角、内参完全不同的"自由视角"设定。3DGS 的显式几何与可渲染性为这一问题提供了新思路：如果能实时构建环境的 3DGS 表示，就可以将目标图像嵌入 3D 空间，在几何层面直接搜索匹配位姿。然而，标准 3DGS 需要离线多视角优化，无法满足导航中流式输入的实时性要求；已有的前馈方法依赖多视角输入，同样不适用于机器人逐步探索的场景。IGL-Nav 正是为了弥合这一差距——设计首个支持单目 RGB-D 流式输入的前馈 3DGS 增量重建，并以此为基础实现粗到精的目标定位。

\textbf{增量式 3DGS 场景构建。}IGL-Nav 的第一个核心贡献是将 3DGS 从离线优化转为前馈单目增量预测。直觉上，传统 3DGS 如同一幅需要反复打磨的画作——用数百轮迭代精调每个高斯参数；而 IGL-Nav 则像一台快速相机——看一眼就给出合理的场景表示。具体而言，在时刻 $t$，网络接收当前帧的 RGB 图像 $I_t$ 和深度图 $D_t$，通过 UNet 编码器提取稠密单目场景嵌入 $E'_t$，再经高斯预测头 $H$ 输出像素对齐的 3DGS 参数：

\begin{equation}
(\Delta C_{2D}, \Delta D, \Sigma, c) = H(E'_t), \quad E'_t = E(I_t, D_t)
\end{equation}

其中 $\Delta C_{2D}$ 和 $\Delta D$ 是相对于像素坐标和深度的残差修正，$\Sigma$ 为协方差，$c$ 为球谐系数。利用相机内参 $M$ 和位姿 $T_t$，通过逆投影将 2D 像素提升至 3D 高斯位置：$\mu = \pi^{-1}(C_{2D} + \Delta C_{2D},\; D + \Delta D \mid M, T_t)$。场景表示 $G_t$ 随探索增量更新：$G_t = G_{t-1} \cup \{\mu_t, \Sigma_t, \alpha_t, c_t\}$，$E_t = E_{t-1} \cup E'_t$。当高斯数量过大时，基于不透明度和密度进行剪枝以控制内存。该前馈模型用离线 RGB-D 视频流训练，损失函数为 $\mathcal{L}_2$ 与 LPIPS 感知损失的线性组合。

\textbf{粗定位：3D 卷积等价的体素匹配。}增量构建场景表示后，下一步是在 3D 空间中定位目标照片的拍摄位姿。目标位姿包含 5 个自由度——位置 $(x, y, z)$ 加旋转 $(\phi, \theta)$——搜索空间极大。IGL-Nav 的关键洞察是：将这一搜索问题转化为 3D 卷积操作。直觉上，如果将场景的 3D 嵌入体素化，将目标的 3D 嵌入在不同旋转下也体素化，那么"在所有位置和旋转上寻找最佳匹配"就等价于"用目标体素作为卷积核在场景体素上做 3D 卷积"。

具体实现上，3D 空间被体素化为间距为 $v$ 的网格；球面旋转通过正二十面体的 $\zeta$-层细分（实验中 $\zeta = 3$）离散化为 $N$ 个方向。目标图像的 3D 嵌入 $E_g$ 在 $N$ 个旋转下分别体素化，形成形状为 $L \times L \times L \times C_{\text{in}}$ 的 3D 卷积核 $K^{L \times L \times L \times C_{\text{in}} \to C_{\text{out}}}$，其中 $C_{\text{out}} = N$。粗定位公式为：

\begin{equation}
\arg\max_{x,y,z,k}\; \mathcal{C}\bigl(f_1(\mathcal{V}(E_t)),\; f_2(K)\bigr)[x][y][z][k]
\end{equation}

其中 $\mathcal{C}$ 为 3D 卷积操作，$\mathcal{V}$ 将场景嵌入体素化，$f_1 / f_2$ 为可学习 MLP 将场景嵌入与卷积核投射到对齐的特征空间。输出激活图的形状为 $X \times Y \times Z \times N$，取最大值索引即得粗定位结果。为进一步提升效率，IGL-Nav 采用基于柱体（pillar）的体素化替代标准 3D 体素化。粗定位模块用 focal loss 监督激活图训练。

\textbf{精定位：匹配约束的可微渲染优化。}当粗定位引导机器人接近目标区域后，IGL-Nav 通过可微渲染精化 6-DoF 位姿。首先，利用当前帧的 3DGS $G_t$ 以目标相机内参渲染图像，通过 LoFTR 局部特征匹配判断目标是否出现在视野中——匹配点数超过阈值 $\eta$ 时触发精定位。与直接的光度损失优化不同，IGL-Nav 发现增量构建的 3DGS 质量有限，全局光度优化不稳定。其解决方案是\textbf{只关注高置信度的匹配点对}：先用 LoFTR 检测匹配像素对 $(x_g, x)$，再查询对应深度 $(d_g, d)$，反投影至 3D 空间得 $(X_g, X)$，优化目标为：

\begin{equation}
T^* = \arg\min_{T \in SE(3)}\; \frac{1}{Q}\sum_{i=0}^{Q-1} \| X_g^i - X^i \|_2
\end{equation}

其中 $Q$ 为匹配点对数。由于 LoFTR 的匹配是可微的，梯度可通过渲染的 RGB 和深度图回传，实现端到端的位姿优化。该策略有效利用了增量 3DGS 中最可靠的渲染区域，避免了低质量区域对优化的干扰。

\textbf{导航流程与实验表现。}整体导航分为探索与到达两阶段。探索阶段，粗定位的激活图投射至鸟瞰图（Bird's Eye View, BEV）得到得分图 $S_a$，与前沿探索得分图 $S_f$ 融合引导探索方向；到达阶段，渲染判别器检测目标出现后切换至精定位模块，用 Fast Marching Method（FMM）规划路径。在 Gibson 数据集上，IGL-Nav 在传统图像目标导航和更困难的自由视角设定下均大幅超越当时最优方法（零样本迁移性能甚至超过其他方法的监督训练结果），并成功部署于真实机器人平台，展示了从前馈 3DGS 增量重建到粗到精定位的完整闭环。

IGL-Nav 证明了前馈增量 3DGS 重建在导航中的可行性，但其场景表示始终是一张不断增长的全局体素地图——随探索范围扩大，体素化的存储与计算开销急剧上升，且前端单目前馈网络的重建精度受限。与 IGL-Nav 的单一大规模 3DGS 不同，YOPO-Nav（2025.12）~\cite{2512_09903} 采取了\textbf{全局-局部层次化}的架构：不再维护一张全局体素地图，而是将环境分解为多个小而精的局部 3DGS 模型，用拓扑图组织它们之间的空间关系。这一设计的关键差异在于场景表示的粒度——IGL-Nav 的体素化 3D 嵌入需要在整个探索空间上做 3D 卷积匹配，而 YOPO-Nav 将定位拆解为先用视觉地点识别（Visual Place Recognition, VPR）选择正确的局部 3DGS 节点，再在节点内部做精确位姿估计，从而兼具拓扑图的轻量可扩展性与 3DGS 的几何精确性。

\textbf{离线建图：轨迹级局部 3DGS 与跨轨迹图构建。}与 IGL-Nav 在线逐帧增量构建全局表示不同，YOPO-Nav 的建图完全在离线阶段完成。具体而言，采用 AnySplat 前馈网络从视频序列直接预测 3DGS 参数及相机位姿（无需已知位姿信息），每 55 帧（448$\times$336 分辨率）构建一个局部 3DGS 模型，对应轨迹的一段局部场景；随后经 SOG（Spatially Ordered Gaussians）压缩，将每个模型从 150--250 MB 降至 25--40 MB（压缩比约 6:1）。多个 3DGS 模型之间通过跨轨迹特征匹配建立图连接——帧间用 FAISS KNN 检索 top-5 候选，再用 XFeat+LighterGlue 精确匹配，保留对应点超过 800 的边（最优阈值从 500--900 的消融实验中确定），形成支持路径规划（如 Dijkstra 算法）的拓扑图。该设计的核心优势在于：每个 3DGS 模型仅覆盖局部区域，建图精度不受全局规模影响；图结构天然支持增量扩展——新增轨迹只需建图并链接至已有节点。

\textbf{运行时导航：VPR 粗定位与局部 3DGS 精定位。}运行时导航分为两阶段，其粗到精的思路与 IGL-Nav 类似，但实现机制截然不同。IGL-Nav 用 3D 卷积在体素空间中做粗定位，而 YOPO-Nav 的粗定位由 VPR 模块（YOPO-Loc）完成：以 DINOv3（ViT-L）为骨干网络、Bag of Queries（128 个可学习查询，输出 4096 维描述符）为聚合器，在覆盖 2.7 km$^2$ 的 GND 大规模校园数据集（55,317 张训练图像）上训练后迁移至目标校园，检索当前观测最可能所属的 3DGS 模型节点。精定位则在选定的局部 3DGS 模型内进行：从当前观测与目标帧提取 XFeat 特征对应，结合 3DGS 深度反投影建立 2D-3D 对应关系，通过 SQPNP-RANSAC（500 次迭代、5 像素重投影阈值、0.99 置信度）+Gauss-Newton 精化估计 6-DoF 位姿；估计位姿与目标帧位姿之差直接生成平移和旋转控制动作。为解决 3DGS 模型与真实世界的尺度差异，YOPO-Nav 通过 SAM-2 分割渲染图像下方 30\% 区域的地面像素，子采样至 20,000 点后 RANSAC 拟合地面平面，取所有已知位姿拟合结果的中位数，将 3DGS 相机高度与 Jackal 机器人固定相机高度（250 mm）之比作为尺度因子，将 3DGS 空间中的平移量映射为真实世界米制单位（旋转无需尺度校正）。

\textbf{实验表现与可解释的人机协作。}在 YOPO-Campus 数据集（35 条校园步道、总覆盖 6 km、4 小时视频、26,500 帧图像）的真实部署实验中，547 个 3DGS 模型经 SOG 压缩后总存储仅 17.8 GB。系统在无 GPU 的机器人端（i3-9100TE CPU）通过 QUIC 协议（基于 TLS 1.3 和后量子密码学的 s2n-quic 实现）连接远程 GPU 工作站实现实时推理。在图像目标导航中，YOPO-Nav 在短距离（1.5 m）到长距离（12 m）各阈值上均优于零样本基础模型 ViNT 和 NoMaD，且优势随轨迹长度增加而愈加显著——这表明基于环境特化的 3DGS 图表示在长程导航中优于通用策略模型。值得关注的是季节变化鲁棒性：在冬季测试时（数据采集于夏秋两季），YOPO-Nav 仍能成功导航，只要场景变化不占据主导地位。此外，YOPO-Nav 支持可解释的人机协作导航——GUI 界面实时显示机器人视角、VPR 匹配帧、鸟瞰图路径和 3DGS 渲染画面，允许操作员在必要时以离散动作（0.25 m 前后或 15$^\circ$ 旋转）介入纠偏，系统随后通过 YOPO-Loc 自动重定位并恢复自主导航。

上述方法解决的是通用图像目标导航——给定一张场景级参考图像，导航至大致匹配的位置。但在实际部署中，机器人常需要精确定位到某个特定物体实例，而非仅仅到达相似的视觉区域。以下方法进一步利用 3DGS 的可渲染性，支持实例级精度的图像目标导航。

\paragraph{实例图像目标导航}

与 IGL-Nav 和 YOPO-Nav 面向场景级定位不同，SplatSearch（2025.11）~\cite{2511_12972} 聚焦实例图像目标导航（Instance Image Goal Navigation, IIN）——目标不再是一张场景照片对应的大致位置，而是某个特定物体实例。这一升级带来两个关键挑战：目标参考图像可能来自任意视角（如鸟瞰），机器人需在视角不变的前提下识别目标；同时，在未知环境中在线构建的 3DGS 地图仅有稀疏视角覆盖，渲染的候选视角图像存在大量缺失区域。

SplatSearch 围绕上述挑战设计了六模块闭环架构。

\textbf{地图构建模块}（MGM）以 RGB-D 观测增量构建稀疏 3DGS 地图和占据栅格图。

\textbf{语义对象识别模块}（SOIM）用 YOLOv7 检测当前帧中与目标同类的候选对象，MobileSAM 分割实例掩码并映射至 3DGS 空间估计质心。

\textbf{新视角合成模块}（NVSM）围绕每个候选对象的 3D 质心，在球面上采样 $N_R$ 个相机位姿，通过 3DGS 地图渲染多视角图像。由于稀疏重建导致渲染图像存在大面积空洞，SplatSearch 训练了\textbf{视角一致图像补全网络}（VCICN）——一种多视角扩散模型，以渲染图像及其缺失掩码为输入，经 VAE 编码、前向扩散加噪、去噪网络预测（以 CLIP 文本编码的目标类别为条件）三个阶段，输出结构一致的完整视角图像。随后，\textbf{特征提取模块}（FEM）用 EfficientLoFTR 将补全后的多视角图像与目标图像进行特征匹配，最高匹配分数超过阈值 $\tau_m$ 时确认为目标实例。若未确认，匹配分数作为视觉上下文得分传入\textbf{探索规划器}（EP）。

EP 的前沿选择策略综合三类信号：距离偏好（归一化逆距离）、语义上下文得分（BLIP-2 编码的当前观测与目标图像的余弦相似度）和视觉上下文得分（FEM 的特征匹配分数），使机器人优先探索与目标实例语义相近且视觉特征匹配度高的前沿区域。其中语义上下文在远离目标时引导大范围搜索，视觉上下文在接近候选对象时实现实例级区分——两者协同构成了\textbf{语义-视觉联合前沿探索策略}。

实验在 Habitat 仿真器的 HM3D-val 和更困难的 HM3D-val-hard（目标图像改为鸟瞰等非标准视角）两个基准上进行。SplatSearch 在成功率和成功路径长度（SPL）上均优于 GaussNav（预建 3DGS）、IEVE（DRL 探索-验证切换）和 UniGoal（LLM 场景图）等基线方法，尤其在视角变化剧烈的 hard 集上优势更明显。消融实验验证了 NVSM、VCICN、语义得分和视觉得分各组件的必要性。在 35m$\times$14m 真实室内环境中，SplatSearch 直接从仿真迁移部署（无额外训练），在包含人员、电视、椅子、背包和笔记本电脑五种目标的 25 次试验中达到 SR$=$0.77、SPL$=$0.48，验证了方法对未见物体和任意视角目标的泛化能力。

上述方法（IGL-Nav、YOPO-Nav、SplatSearch）均将 3DGS 作为在线定位的地图表示——机器人在导航过程中需要持续地相对于 3DGS 进行自定位。SplatSearch 更是在线增量构建稀疏 3DGS 地图并实时渲染候选视角以匹配目标实例。GSplatVNM（2025.03）~\cite{2503_05152} 提出了一个根本性的范式转换：将 3DGS 从在线定位地图重新定义为\textbf{离线视角合成器}，完全不需要机器人在导航时进行度量定位——3DGS 仅在导航前用于离线渲染目标视角序列，导航过程中机器人不与 3DGS 交互。这一设计的动机源自视觉导航模型（VNM）的观察：NoMaD 等零样本策略已能在超过 100 小时的异构数据集上训练后，通过目标视角图像直接生成路径点，无需精确自定位或环境特定训练；但传统 VNM 依赖图像拓扑图（ITG）提供目标图像序列，而 ITG 在图像数据库稀疏时覆盖不足，且存储开销随图像数量线性增长。GSplatVNM 利用 3DGS 同时解决这两个瓶颈——3DGS 可合成任意视角图像填补稀疏数据库的间隙，且将图像信息压缩为高斯参数使存储与图像数量无关。

整个流程具体如下：首先通过 LPIPS 感知相似度与碰撞惩罚的联合损失 $\mathcal{L}(\mathbf{q}_*) = \mathcal{L}_\text{img}(I_*, I_\text{rendered}(\mathbf{q}_*)) + \mathbf{1}_\text{collision}[d_\text{min}(\mathbf{q}_*) \leq r]$（其中 $\mathbf{q}_* \in \mathrm{SE}(2)$ 为 2D 位姿，$r$ 为机器人半径），用 Tree-structured Parzen Estimator（Optuna 实现）黑盒优化在 3DGS 中估计起点和目标点的全局位姿；然后以 A* 算法在 3DGS 的椭球碰撞表示中规划无碰撞全局路径点序列，偏航角由相邻路径点的 $\mathrm{atan2}$ 计算，最终形成轨迹 $\mathcal{T} = \{\mathbf{q}_\text{start}, (x_0, y_0, \theta_0), \ldots, \mathbf{q}_\text{goal}\}$，沿轨迹渲染密集的目标视角图像序列 $\mathcal{V}_r = \{I_0^r(\mathbf{q}_0), \ldots, I_{M-1}^r(\mathbf{q}_{M-1}), I_\text{goal}\}$；最后以 NoMaD 进行局部规划——利用距离估计网络 $\phi$ 在 $\mathcal{V}_r$ 中定位最接近当前观测的图像 $I_t^r$，以次一帧 $I_{t+1}^r$ 作为子目标，条件化扩散策略（从 4 帧观测历史 $\mathcal{O}_\text{obs} = \{I_{t-3}, \ldots, I_t\}$ 和子目标图像编码的特征）生成 8 个空间路径点 $\mathcal{W}_t = \{\mathbf{w}_0, \ldots, \mathbf{w}_7\}$，再通过 MPPI 控制器（考虑差速两轮动力学）跟踪。值得注意的是，GSplatVNM 展示了一项独特能力：\textbf{导航至被观测到但未被物理访问过的位置}——3DGS 合成的视角图像可以覆盖图像采集路径未经过但被相机观察到的区域，这是 ITG 方法无法实现的。

实验在 AI Habitat 仿真器的三个 Gibson 室内环境（Greigsville 43.6 m$^2$、Ribera 50.6 m$^2$、skokloster-castle 约 190 m$^2$）中进行，每环境测试 20 次试验。GSplatVNM 在 SR 上一致优于 ITG 基线：在最大的 skokloster-castle 环境中，SR 达 0.59（ITG 基线 0.39），提升 20 个百分点。存储效率方面，以 skokloster-castle 为例，3000 张图像的 ITG 需 3400 MB，而对应 3DGS 仅 223 MB——图像越多，3DGS 反而越小（因高斯数量取决于场景复杂度而非图像数量），而 ITG 线性增长。在不同图像数据库稀疏度（300--3000 张）下，GSplatVNM 的 SPL 均保持稳定，仅在最稀疏条件（300 张）下因 3DGS 渲染质量下降而略有降低。

\paragraph{层次化评分导航}

SplatSearch 和 GSplatVNM 虽然利用 3DGS 的可渲染性合成候选视角，但仍依赖在连续 3D 空间中随机或沿轨迹采样大量视角来保证覆盖，重叠渲染图像带来显著的计算冗余。同属 3DGS 地图方法的 GaussNav 也面临类似问题——围绕预测物体实例穷举采样视角进行匹配，缺乏策略性引导。GauScoreMap（2025.06）~\cite{2506_07338} 从根本上规避了穷举采样，转而提出\textbf{全局语义-局部几何两阶段评分}范式，用仅两次前向推理即可定位目标（全局语义评分一次、局部几何评分一次），无需穷举渲染与比较。首先在地图构建阶段，GauScoreMap 采用层次化策略构建 3DGS 场景表示——先将局部区域的观测帧反投影为子地图，再合并为全局高斯场；随后将 CLIP 视觉特征通过渲染权重的加权平均提升（lifting）到每个 3D 高斯上：

\begin{equation}
f_g = \sum_{(i,(u,v)) \in S_g} \bar{w}_g(i,(u,v)) \, F_{i,(u,v)}, \quad \bar{w}_g(i,(u,v)) = \frac{w_g(i,(u,v))}{\sum_{(i,(u,v)) \in S_g} w_g(i,(u,v))}
\end{equation}

其中 $F_{i,(u,v)}$ 是第 $i$ 帧像素 $(u,v)$ 处的 2D CLIP 特征，$w_g$ 是高斯 $g$ 在该像素处的渲染权重（即 $\alpha$-混合 中的不透明度贡献），$S_g$ 为所有对高斯 $g$ 有贡献的帧-像素对集合。这一加权聚合方式确保多帧观测的融合忠实于高斯在每帧中的实际视觉贡献，形成支持开放词汇查询的语义特征场。在\textbf{全局语义评分}阶段，系统先用 Mask-RCNN 从目标图像提取物体类别标签（而非直接使用完整目标图像，因为图像中除目标实例外还包含大量背景信息，直接匹配会产生噪声分割），再以该类别的 CLIP 文本嵌入 $E_T$ 与每个高斯的 CLIP 特征 $f_g$ 计算余弦相似度 $S_g = E_T^\top f_g / (\|E_T\| \|f_g\|)$，生成覆盖整个环境的连续相关性得分场；经阈值 $\tau$ 分割后，遵循 LUDVIG 的策略结合场景几何对分割结果进行基于特征相似性的空间扩散，连接同一实例的碎片化部分，形成空间连贯的候选实例区域（连通分量），每个连通分量即一个候选目标。在\textbf{局部几何评分}阶段，该阶段进一步分为区域选择和精细位姿估计两个子步骤，且两个子步骤共用同一个预训练的射线-图像交叉注意力网络。

首先对每个候选区域沿高斯表面法线半球采样射线（由邻近高斯估计法线），每条射线由原点 $o_i$、方向 $d_i$ 和一阶球谐系数（颜色）$c_i$ 组成，经位置编码 MLP 编码为特征 $r_i = \text{MLP}(\gamma(o_i), \gamma(d_i), \gamma(c_i))$，与 DINOv2（自监督视觉Transformer，通过自蒸馏训练获得强大视觉特征）提取的目标图像特征 $F_g$ 通过\textbf{射线-图像交叉注意力}（以射线为查询、图像特征为键值）计算交叉注意力矩阵 $A = \text{CrossAttention}(r, F_g) \in \mathbb{R}^{K \times l}$，沿图像维度求和得到区域级得分 $S_{r_i} = \sum_{j=1}^{l} A_{i,j}$，选出最高得分区域。随后在该区域内进行第二轮更密集的高斯-射线对采样，选取 top-$k$ 高分射线对通过三角化估计 6D 位姿（位置与朝向）。这种先粗后细的两阶段设计将搜索空间从整个环境逐级收缩至单个区域再到精确位姿，显著提高了效率。路径规划阶段将高斯投影到鸟瞰图（BEV）栅格并使用 Fast Marching Method 计算无碰撞最优轨迹。

消融实验清晰验证了两阶段的各自贡献：移除全局语义评分后成功率下降 17.6\%（SR 从 0.784 降至 0.608，因室内遮挡严重，仅靠局部几何难以在全局范围内高效定位目标），移除局部几何评分后成功率骤降 36.3\%（SR 从 0.784 降至 0.421，因语义评分只能识别类别，无法区分同类别不同实例）。在 HM3D 基准的 1000 个验证回合上（涵盖 chair、couch、bed、toilet、television 五类共 795 个唯一实例），GauScoreMap 以 SR 0.784、SPL 0.605 超越同基于 3DGS 的 GaussNav（SR 0.725、SPL 0.578），SR 提升 5.9\%、SPL 提升 2.7\%，同时通过仅两次前向推理替代穷举视角采样，大幅降低了渲染与比较开销。论文还提供了真实室内环境的部署演示，验证了方法的实际可迁移性。

\paragraph{3DGS 增强的 VLM 导航}

上述方法均以 3DGS 作为几何定位地图或离线渲染器——无论前馈重建（IGL-Nav）、评分搜索（GauScoreMap）还是离线视角合成（GSplatVNM），导航决策最终都由专用几何模块或预训练策略网络完成，3DGS 仅充当地图或渲染器。3DGSNav（2026.02）~\cite{2602_12159} 则提出一个截然不同的思路：\textbf{将 3DGS 直接作为 VLM 的持久空间记忆}，让 VLM 本身成为导航的决策核心。这一设计的动机源自一个关键观察：现有零样本物体导航（ZSON）方法普遍依赖语义地图或文本描述等场景抽象（scene abstraction），将连续空间离散化为符号表示——这一抽象过程虽简化了规划，却丢失了物体间的空间结构和几何细节，低层感知精度掣肘了高层决策，阻碍了 VLM 固有视觉-空间推理能力的充分发挥。值得注意的是，GSplatVNM 虽也利用 3DGS 渲染视角图像供给导航策略（NoMaD），但其策略本质上是扩散模型生成的固定动作序列，VLM 从未参与路径选择；3DGSNav 则将决策权完全交给 VLM——渲染的视角图像不再驱动底层控制器，而是作为 VLM 的视觉输入激发其空间语义推理。3DGSNav 的核心假设是：如果能用 3DGS 保留环境的完整视觉-几何信息，并通过自由视角渲染为 VLM 提供可分析的视觉输入，VLM 就能直接利用其强大的空间推理能力进行高质量决策，无需经过任何场景抽象。具体而言，3DGSNav 围绕三个模块展开：

\textbf{主动感知}模块通过增量构建 3DGS 替代被动旋转扫描。系统渲染全景不透明度场检测低不透明度区域（阈值 $\tau=0.3$），用 DBSCAN 聚类在连续视点空间中确定需要补充观测的代表性离散视点，驱动机器人主动转向覆盖视觉缺失区域，避免了机械式旋转带来的定位误差和冗余观测。

\textbf{轨迹引导的自由视角渲染}是 3DGSNav 相对于 IGL-Nav 前馈定位范式的关键差异：IGL-Nav 用前馈网络直接预测高斯参数并经 3D 卷积在体素空间中搜索匹配位姿——整个定位过程由几何模块驱动；而 3DGSNav 利用 3DGS 的可渲染性，沿引导轨迹优化虚拟相机位姿以生成信息量最大的前沿视角图像，再将渲染结果交给 VLM 进行空间语义推理。具体流程为：前沿点先经距离场骨架分割自适应聚类（同一区域的前沿点共享环境上下文，避免逐点分析的冗余），对每个代表前沿点以 Dijkstra 算法构建包含障碍物指数惩罚 $\Omega(n)=\omega\cdot\exp(2\cdot(d_s-d(n)))$ 的引导轨迹 $\mathcal{T}_i$，然后在轨迹附近以曲率-距离联合评分 $p^* = \arg\max_{p_i \in P_v} \alpha \kappa_i + (1-\alpha) d_i$（其中 $\alpha=0.7$ 平衡曲率 $\kappa_i$ 与前沿距离 $d_i$ 的贡献，候选集约束为 $r_\text{min}=1.0$ m 至 $r_\text{max}=4.0$ m 的安全范围）初始化虚拟相机位姿，通过\textbf{多约束优化}（迭代 40 步，权重 $\lambda_o=0.01, \lambda_v=1.0, \lambda_c=0.01, \lambda_t=0.1, \beta=5$）求解最优视点：

\begin{equation}
\mathcal{L} = \lambda_o \mathcal{L}_\text{opa} + \lambda_v \mathcal{L}_\text{vis} + \lambda_c \mathcal{L}_\text{cos} + \lambda_t \mathcal{L}_\text{traj}
\end{equation}

其中 $\mathcal{L}_\text{opa}$ 衡量已观测/未观测区域比例（防止过多未观测区域阻碍前沿环境分析），$\mathcal{L}_\text{vis}$ 以深度一致性约束确保前沿点不被遮挡（沿射线采样并以 sigmoid 判定遮挡概率），$\mathcal{L}_\text{cos}$ 以余弦相似度保证相机朝向前沿点（引入遮挡概率 $Pr_v$ 自动降低被遮挡目标的约束权重），$\mathcal{L}_\text{traj}$ 以平滑最小距离约束相机贴近引导轨迹（采用 LogSumExp 平滑近似，权重由不透明度损失自适应调节）。优化后的视角渲染后叠加注视点、历史轨迹和未观测区域标记等\textbf{结构化视觉标注}——第一人称视角（FPV）提供局部视觉细节与注视区域，鸟瞰图（BEV）提供全局空间布局与历史运动上下文——与思维链（Chain-of-Thought）提示一起输入 Gemini3 规划 VLM，由 VLM 自主推理并选择最优前沿点。

\textbf{虚拟视角重验证}机制则进一步展示了 3DGS 持久记忆的独特优势。导航过程中，YOLO-E 开放词汇检测器实时扫描潜在目标；当检测到疑似目标时，系统不驱动机器人物理移动（避免额外的探索成本），而是调用 GLM-4.5v 作为动作决策 VLM，推理环境结构和物体遮挡关系后选择下一步动作（前进/后退各 0.25 m、左右转向各 $10^\circ$、或离开），将该动作投影到 3DGS 空间渲染新视角进行二次确认。这一"在记忆中移动而非在现实中移动"的策略，将物理探索成本转化为虚拟渲染开销，大幅提高了目标确认的效率与可靠性。

在 HM3Dv1、HM3Dv2 和 MP3D 三个基准的零样本物体导航任务中，3DGSNav 相比现有最优方法平均提升 SR 13.5\% 和 SPL 32.08\%；相比同样不依赖场景抽象的 RL 方法 ZSON，SR 和 SPL 分别提升 203\% 和 320\%；相比采用体素表示构建 3D 语义地图的 BeliefMapNav，SR 和 SPL 分别提升 25.25\% 和 51.65\%，充分验证了让 VLM 直接在 3DGS 视觉记忆上推理优于传统场景抽象方法。VLM 消融实验进一步揭示了这一范式的关键：相比随机前沿选择（SR 提升 42.07\%、SPL 提升 115.48\%），VLM 的空间语义推理是性能提升的主要来源；相比开源 Qwen3-VL-235b（SR 提升 18.24\%），Gemini3-Pro 在空间理解上的显著优势（SR 再提升 20.14\%）表明 VLM 推理质量直接决定了导航性能上限。在四足机器人的真实世界部署中，3DGSNav 同样展现了鲁棒的零样本迁移能力。

\paragraph{多目标开放词汇导航}

3DGSNav 的虚拟视角重验证证明了 3DGS 渲染对导航的价值，但其记忆服务于单次 VLM 决策——当任务从单个目标扩展为多房间、多模态、多目标的连续序列时，需要一个更根本的记忆重构。LagMemo（2025.10）~\cite{2510_24118} 正是从这一需求出发，将 3DGS 提升为\textbf{统一的多任务语言记忆}：一次探索构建的空间-语义关联，由所有后续任务共享，无需重复建图。其设计动机源于对现有模块化导航记忆两类瓶颈的分析——基于对象的方法（如 Modular GOAT）依赖固定词汇表的上游检测器填充记忆，无法捕获开放类别目标（如"Mickey Mouse 玩偶"在探索时被忽略则永远无法检索）；基于 2D 栅格的方法将语义特征投影至平面，丢失垂直空间细节且难以跨视角保持全局一致性。

与 GSMem 通过反向聚合为每个高斯赋予连续语言嵌入不同，LagMemo 采用\textbf{两级码本量化}策略构建离散化语言记忆：首先通过 SAM 提取实例掩码并用 CLIP 编码 2D 特征，经特征泼溅渲染和掩码级聚合优化高斯嵌入，确保同类实例内一致、异类实例间可分；然后以两级聚类将高斯特征离散化——粗粒度分区联合考虑 3D 位置和语言特征（$k_1 = 32$ 个簇），细粒度分区仅基于语言特征在每个粗簇内细分（$k_2 = 5$，共 160 个码本条目），确保跨视角实例对齐的鲁棒性。LagMemo 特别指出，选择 3DGS 而非离散点云的理由正在于其连续可微结构：几何上高斯能自然插值稀疏视角的孔洞，语义上可微渲染损失能够通过反向传播主动解决跨视角冲突，而非启发式地多视角平均。

为解决大规模多房间环境中帧间重叠有限导致的几何遗忘和表面孔洞问题，LagMemo 引入\textbf{关键帧检索机制}——每步先优化当前帧 $p_1 = 30$ 轮，再对 PSNR 较低的历史帧加权采样（PSNR 越低被采样概率越高）优化 $p_2 = 60$ 轮，使低保真区域更容易被回访修正。消融实验表明，该机制将重建 PSNR 从 21.15 提升至 27.20，同时使目标定位准确率从 66.3\% 提升至 70.8\%。面对多模态任务目标（文本、图像或类别），系统将目标编码为 CLIP 嵌入，通过余弦相似度在码本中检索候选实例，计算关联高斯的几何质心并投影至 2D 障碍图生成航点，再用 FMM 规划无碰撞路径。

在验证环节，LagMemo 做出了与 3DGSNav 的渲染验证截然不同的策略选择。3DGSNav 通过虚拟相机渲染新视角进行二次确认，但 LagMemo 指出快速探索中收集的稀疏视角会导致 3DGS 渲染常含模糊和伪影，严重干扰 2D 匹配模型的精度。因此 LagMemo 采用\textbf{现场验证}策略——3DGS 记忆仅作为粗粒度全局先验生成候选航点，到达候选航点 1.2m 范围内后执行全景扫描进行像素级二次验证：文本/类别目标采用 CLIP 相似度与 SEEM 开放词汇分割的像素掩码双阈值判定（$\tau_{\text{CLIP}} \geq 0.23$ 或 $\tau_{\text{SEEM}} \geq 1.1$ 即视为通过），图像目标用 LightGlue 特征匹配确认（内点匹配比例 $r_{\text{match}} \geq 5\%$）。消融实验表明，LightGlue + SEEM+CLIP 组合验证策略将整体 SR 从无验证的 41.3\% 提升至 56.3\%，远超仅使用 CLIP 验证的 46.7\%。验证通过后，系统利用 SEEM 掩码结合深度信息将目标足迹投影至障碍图，定位最近可通行栅格作为最终目标点，切换至局部 FMM 导航并触发停止。整个"码本查询—航点导航—现场验证"循环迭代执行，直至目标被发现或步数耗尽。

在 GOAT-Core 基准（每 episode 含 20 个多模态子任务，远高于 GOAT-Bench 的 7.88 个，平均场景面积 218.63 $\text{m}^2$）上，LagMemo 目标定位准确率达到 70.8\%（VLMaps 为 58.8\%），多目标导航 SR 达 56.3\%（Modular GOAT 为 38.3\%），SPL 达 35.3\%（Modular GOAT 为 29.7\%）。相比使用真实语义分割的 GT-Sem 上界（SR 75.0\%），LagMemo 仅差 18.7 个百分点，展现了接近理论极限的性能。在完整 GOAT-Bench 的 seen/synonyms/unseen 三个划分上，LagMemo 分别取得 SR 36.8\%/44.8\%/37.9\%，全面超越所有基线方法。相比 Modular GOAT，LagMemo 将记忆索引错误率从 37.1\% 降至 22.1\%，验证错误率从 22.1\% 降至 18.3\%，总失败率从 63.7\% 降至 43.7\%，充分验证了码本量化 3DGS 记忆在保留三维空间上下文和过滤多视角噪声方面的优势。

\subsection{语义探索与安全规划}
\label{ss:17_1b}

本节讨论两条技术路线：（1）语义搜索与主动探索；（2）安全规划与碰撞避免。

\subsubsection{语义搜索与主动探索}

本路线介绍利用 3DGS 可渲染性进行语义搜索与主动探索的方法。语义搜索要求机器人在未知环境中主动探索以寻找特定对象或回答问题。3DGS 的可渲染性使机器人能够"回顾"已探索区域，从最优视角重新观察可能遗漏的目标。

\paragraph{持久空间记忆}

GSMem（2026.03）~\cite{2603_19137} 修改了 3DGS 的\textbf{语言场构建环节}——用\textbf{权重一致的反向聚合}（weight-consistent reverse aggregation）替代传统的特征优化或后分配策略——3DGS 的几何与外观建模保持不变。

\textbf{核心问题与动机。}现有的具身导航记忆分为两类：基于对象的方法（如场景图）将视觉数据压缩为离散标签，单次检测失败即造成不可恢复的记忆缺失；基于视角的方法（如关键帧快照）仅保留特定视角的稀疏观察，若目标在初始探索时被遮挡或以次优角度捕获，则无法事后补救。GSMem 识别出这一根本缺陷——\textbf{缺乏事后可再观察性}（post-hoc re-observability）——并提出以 3DGS 的稠密连续辐射场作为持久空间记忆，使智能体获得\textbf{空间回忆}（Spatial Recollection）能力：从任意最优视角渲染已探索区域的照片级真实感新视角，即使该视角从未被物理占据。

\textbf{语言场的零优化构建。}为了让空间记忆支持语义级查询，GSMem 需要为每个 3D 高斯赋予语言嵌入。直觉上，正向渲染时 2D 像素的颜色是各 3D 高斯按深度排序加权混合的结果；GSMem 反转这一过程——将 2D 像素上的语义特征按同样的权重反向分配回 3D 高斯。具体而言，首先通过像素级 CLIP 编码器提取稠密 2D 特征 $f_{p,k}^{\text{2D}}$，再经预训练自编码器将 768 维压缩至 32 维。设 $w_{i,p,k}^t$ 为时刻 $t$ 关键帧 $k$ 中高斯 $g_i$ 在像素 $p$ 处的混合权重（与正向 RGB-D 渲染中完全相同），则高斯语言特征的在线更新公式为：

\begin{equation}
f_i^t = \frac{W_i^{t-1} f_i^{t-1} + \sum_{k \in \mathcal{T}_t} \sum_p w_{i,p,k}^t \, f_{p,k}^{\text{2D}}}{W_i^t}, \quad W_i^t = W_i^{t-1} + \sum_{k \in \mathcal{T}_t} \sum_p w_{i,p,k}^t
\end{equation}

其中 $\mathcal{T}_t$ 为当前优化帧集。该公式本质上是跨帧的加权移动平均：每次新关键帧到来时，用渲染过程中已计算的混合权重 $w_{i,p,k}^t$ 将 2D 像素特征累加回对应高斯，$W_i^t$ 记录累积权重总量以实现归一化。由于混合权重在正向 RGB-D 渲染中已自然计算，此策略无需引入任何额外优化步骤，即可零额外开销地生成稠密 3D 语言场。

\textbf{多级检索与最优视点选择。}GSMem 的空间回忆通过\textbf{多级检索-渲染}（multi-level retrieval-rendering）机制实现。检索阶段，系统同时维护两个互补表示：对象级场景图提供离散目标定位，语义级语言场支持开放词汇查询（通过余弦相似度阈值 $\tau_{\text{clip}}$ 检索相关高斯并聚类为空间连贯区域）。这种互补设计确保即使对象检测器遗漏了特定目标，语言场的语义嵌入仍可定位候选区域。视点选择采用\textbf{采样-评分}范式：在目标区域 3D 包围盒周围以 $10^\circ$ 间隔均匀采样 36 个方位角和 3 个仰角（$-10^\circ, 0^\circ, 15^\circ$），共生成 108 个候选位姿。两阶段评分筛选最优视点：第一阶段粗筛通过射线可见性评分 $S_{\text{vis}}$（3D 包围盒 14 个代表点的可见比例）和高斯惩罚面积评分 $S_A$ 过滤不可行位姿：

\begin{equation}
S_{\text{vis}} = \frac{|\text{visible}|}{14}, \quad S_A = \exp\!\left(-\frac{(A - A^*)^2}{2\sigma_a^2}\right)
\end{equation}

其中 $A^*$ 为目标投影面积，$\sigma_a$ 控制尺度偏差容忍度。此设计鼓励中等距离观察——太近则缺乏上下文，太远则丢失细节。第二阶段精选对粗筛保留的 top-10 位姿渲染透明度图，计算平均不透明度评分 $S_{\text{opa}} = \frac{1}{|\Omega|}\sum_{p \in \Omega} \alpha_p$，其中 $\Omega$ 为投影 ROI 像素集，$\alpha_p$ 为累积不透明度。由于 3DGS 的可组合性，更高的累积不透明度意味着更强的表面存在和更好的视觉可观测性。最终选择综合评分 $S_{\text{final}} = S_{\text{vis}} + S_A + S_{\text{opa}}$ 最高的位姿作为渲染视点。

\textbf{混合探索策略。}探索模块在经典的\textbf{前沿探索}（frontier-based exploration）框架上设计混合准则，同时平衡语义相关性和几何覆盖。语义评分由 VLM 对前沿观测输出归一化相关性分数 $s_i^{\text{sem}} \in [0,1]$；几何覆盖评分基于 3DGS 参数的 Fisher 信息矩阵，度量新观测对高斯场确定性的增益。由于高维全局 FIM 行列式计算不可行，GSMem 采用高效的 T-最优性代理——增量 FIM 的迹 $\text{Tr}(I_i)$ 作为几何不确定性代理，可直接从 3DGS 渲染函数梯度计算而无需真值监督。混合策略为：当存在前沿的语义评分超过阈值 $\tau_s$ 时选择最高语义评分前沿，否则选择最大化几何评分的前沿。实验表明，阈值 $\tau_s$ 过大将使探索偏向几何覆盖而延迟目标发现，过小则过早聚焦语义相关区域而降低整体成功率。

在 Active EQA 基准（HM3D 63 场景，184 个问题）上，GSMem 达到当时最优的 LLM-Match 评分；在 GOAT-Bench 终身导航基准（36 场景，2600+ 子任务）上取得更大的性能增益，验证了持久空间记忆在长视野场景中的特殊价值。消融实验显示，移除语言场导致成功率显著下降，移除混合探索策略则大幅降低 SPL，证实了两个组件的互补重要性。系统在单块 RTX 4090 上以约 1.2 秒/步的速率运行。

\paragraph{问题驱动的空间记忆探索}

GSMem 通过反向聚合构建语言场并以几何评分（可见性、面积、不透明度）筛选视点，其视角选择完全由可计算的几何指标驱动——这在视角质量容易量化的导航场景中效果良好，但在需要\textbf{理解遮挡和复杂空间关系}的问答推理中可能力有不逮：预录的训练相机视角未必能提供回答问题所需的最佳证据视角。此外，GSMem 的语言场基于 CLIP 嵌入的余弦相似度检索，本质上是为\textbf{类别级查询}（如"椅子""桌子"）优化的静态嵌入匹配——面对组合语义查询（如"厨房里的枕头是什么颜色？"），仅靠单个高斯级别的嵌入相似度无法捕获跨对象的空间关系。GaussExplorer（2026.01）~\cite{2601_13132} 正是针对这两重局限（固定视角约束与静态嵌入推理），将 VLM 引入视角选择循环，实现\textbf{问题驱动的探索和推理}。其核心洞察是：3DGS 的新视角合成能力天然充当 VLM 的"可调节相机"——不同于 GSMem 仅在已训练相机中搜索最优位姿，GaussExplorer 可在初始视角邻域内\textbf{合成全新的扰动视角}，由 VLM 直接评判哪个视角对回答问题最有信息量。这种设计将视角质量评估从 GSMem 的纯几何评分（$S_{\text{vis}} + S_A + S_{\text{opa}}$）转移至 VLM 的语义理解能力——后者天然擅长判断"此视角是否为回答该问题提供了充分证据"，而非仅仅度量"此视角的几何覆盖是否良好"。

\textbf{语义注入策略的差异。}GaussExplorer 与 GSMem 的第一处根本差异在于语义注入方式。GSMem 通过反向聚合将 2D CLIP 特征零优化地分配给 3D 高斯，生成开放词汇语言场；GaussExplorer 则采用\textbf{任务自适应的双轨语义构建}——对于具身问答（EQA，Embodied Question Answering）任务，使用闭集目标检测器（YOLOv8）为每个高斯赋予类别标签，直接按类别匹配激活高斯 $G^{\text{act}} = \{G_j \mid C_j \in C^{\text{selected}}\}$；对于 3D 定位任务，则使用 CLIP（对比语言-图像预训练）+SAM（Segment Anything Model，通用分割模型）开放词汇分割管线，通过文本-特征相似度阈值 $\tau$ 激活高斯 $G^{\text{act}} = \{G_j \mid \max_{c \in C^{\text{selected}}} \text{sim}(f_j, t_c) > \tau\}$。两种策略均将 2D 语义信息直接注册到 3DGS 表示中。

\textbf{问题驱动的视角选择管线。}GaussExplorer 的视角选择遵循"查询解析 $\to$ 粗选 $\to$ 精调 $\to$ 验证"四阶段管线。首先，LLM 从用户问题中提取"证据类别"（evidence categories）——例如对"在哪里可以小睡？"提取"枕头、垫子"等关键词。然后在语义 3DGS 中检索匹配高斯，用 HDBSCAN 聚类为实例并合并物理上邻近的簇，动态确定实例数 $L$。对每个实例，GaussExplorer 选择\textbf{可见性最高}的训练相机视角作为初始视角：

\begin{equation}
V_{l,i} = \frac{|G_l^{\text{cluster,vis}}|_i}{|G_l^{\text{cluster}}|}
\end{equation}

其中可见高斯 $G_l^{\text{cluster,vis}}$ 通过 splatting 渲染的激活图阈值化获得——该可见性评分实质上是对象表面覆盖率的代理指标，消融实验表明其优于体积评分（大尺度离群高斯会扭曲体积评分）。这一粗选阶段已使 GaussExplorer 超越仅使用预录快照的 3D-Mem 基线。

\textbf{VLM 引导的新视角精调。}GaussExplorer 的核心创新在于精调阶段：对 $L$ 个初始视角，通过平移（左右）和缩放（放大/缩小）各生成 4 个扰动新视角（共 $K = 4L$ 个），由 VLM 逐一评估每个渲染视角与问题的相关性。由于 VLM 在处理长视觉序列时性能下降，GaussExplorer 独立评估每个候选视角而非批量输入，再由 LLM 从各视角的独立预测中选择最佳答案——最佳答案隐式决定了最佳最终视角。最后，一个\textbf{验证阶段}（verification stage）通过跨视角响应一致性检验确认所选视角是否真正必要，消融实验表明验证阶段在无额外渲染成本的前提下进一步提升准确率。在 OpenEQA 基准（1,623 个问题）上，GaussExplorer 以相近的平均帧数达到最高 LLM-Match 评分，在 question-based 3D referring segmentation 任务上则展现出精确的细粒度 3D 定位能力。

\paragraph{LLM 引导主动建图}

如果说 GaussExplorer 将 VLM 引入视角选择是为了解决\textbf{问答推理中的证据收集}问题，那么多模态 LLM 引导的主动建图（2024.10）~\cite{2410_17422} 则将其用于解决\textbf{探索效率中的长期目标选择}问题——两者都借助基础模型的能力，但服务于截然不同的任务目标。与 GSMem 的混合探索策略（Fisher 信息用于\textbf{前沿级}评分，在语义评分与几何评分之间阈值切换）相比，该方法采用了更彻底的\textbf{分工策略}——将探索分为长视野规划和细节路径选择两个阶段，分别由不同性质的模型负责：多模态 LLM（GPT-4o）负责语义层面的"去哪里"，Fisher 信息驱动的信息论目标负责几何层面的"怎么走"。值得注意的是，该方法是首个\textbf{不依赖真值位姿}的 3DGS 主动建图系统——其定位依赖于 MonoGS 的在线位姿估计，需要在主动探索的同时自行完成定位，这使其面临的挑战比假定已知位姿的方法（如 UPEN、Active-INR）更为严峻：探索未重建区域意味着定位参考点稀缺，而定位精度下降又会反噬建图质量，形成潜在的恶性循环。

\textbf{长视野规划：LLM 零样本规划器。}在长视野规划阶段，利用 3DGS 的高质量鸟瞰图渲染能力将多模态 LLM（GPT-4o）作为\textbf{零样本规划器}：利用 3DGS 表示构建占据栅格（occupancy grid）识别前沿（自由空间与未观测空间的边界），渲染标注了当前位置、历史轨迹和候选前沿的鸟瞰图，通过思维链（Chain-of-Thought）提示让 LLM 先分析各前沿的语义含义再选择最优长期目标。一个精巧的设计是：LLM 不仅可以选择前往前沿，还被允许决定\textbf{不前往任何前沿}，转而回到已探索区域进行补充观测——这模拟了人类"先粗览再精查"的探索策略。通过将候选目标限制在已检测到的前沿集合中，系统还避免了 LLM 输出不可达目标的问题（若 LLM 输出无效答案，则回退到选择最大前沿）。

\textbf{细节路径选择：Fisher 信息驱动的路径评估。}在细节路径选择阶段，GSMem 在前沿粒度上用 Fisher 信息矩阵的迹 $\text{Tr}(I_i)$ 评分，而此方法将 Fisher 信息细化到\textbf{路径级}评估——对候选路径上每个位姿逐一计算期望信息增益（EIG）。具体而言，确定 LLM 指定的长期目标后，系统在目标区域均匀采样候选位姿集合 $\mathcal{T}_I$，用 EIG 粗筛为 $\mathcal{T}_F$，再通过 A* 算法构建从当前位置到各候选位姿的路径 $\mathcal{P}_j$。与 FisherRF 使用经验 Fisher 信息（在训练集上计算 Hessian 近似）不同，该方法利用一个关键观察：在主动建图场景中，相机位姿的分布 $p(\bx)$ 就是地图自由空间的均匀分布，因而是\textbf{可计算的}。由此通过 Monte Carlo 采样估计真实 Fisher 信息：

\begin{equation}
\mathcal{I}(\bw) = \mathbb{E}_{\bx \sim p(\bx)} \left[ \bH''[\by| \bx, \bw] \right] \simeq \sum_{k=1}^{N} \bH''[\by_k| \bx_k, \bw]
\end{equation}

其中 $\bx_k$ 从地图自由空间中均匀采样。路径级 EIG 还考虑了路径上各点间的互信息累积：

\begin{equation}
\mathtt{EIG}_{\mathcal{P}_j, i}(\bx_i) = \tr\left( \bH''[\by_i | \bx_i, \bw] \; \mathcal{I}_{\mathcal{P}_j, i}(\bw)^{-1}\right), \quad \mathcal{I}_{\mathcal{P}_j, i}(\bw) = \mathcal{I}(\bw) + \sum_{\bx_t \in \mathcal{P}_j, t < i} \bH''[\by_t | \bx_t, \bw]
\end{equation}

\textbf{定位不确定性代价。}如果仅最大化 EIG，机器人将倾向于探索未重建区域——但未重建区域也意味着定位精度的下降风险。该方法首次在主动建图中引入\textbf{定位不确定性}作为路径选择的约束。利用 Cram\'{e}r--Rao 下界，相机位姿对数映射 $\tau_i = \Log(\bx_i)$ 的协方差下界为 Fisher 信息矩阵的逆：$\Cov(T(\hat{\tau}_i)) \geq \mathcal{I}(\tau_i)^{-1}$。由此定义定位代价 $C_{\textit{local}}(\tau_i) = -\log\det(\nabla_{\tau_i} f(\tau_i, \bw)^T \nabla_{\tau_i} f(\tau_i, \bw))$。与 3DGS 参数的 Fisher 信息不同（参数量超过 2000 万，需要 Laplace 对角近似），定位代价的计算只涉及 6 维位姿参数的完整 Jacobian，可以直接在\textbf{完整的} Fisher 信息矩阵上计算而无需近似。最终路径选择目标为：

\begin{equation}
\arg\min_{\mathcal{P}_j} \sum_{\bx_i \in \mathcal{P}_j} C_{\textit{local}}(\Log(\bx_i)) - \eta \sum_i \log(\mathtt{EIG}_{\mathcal{P}_j,i}(\bx_i))
\end{equation}

其中 $\eta$ 控制信息增益与定位精度之间的权衡。

在 Gibson 数据集上，该方法取得 PSNR 23.28、SSIM 0.8067、深度 MAE（Mean Absolute Error，平均绝对误差）0.0696，全面超越最优基线 FBE（PSNR 21.45、SSIM 0.7618、深度 MAE 0.1028），同时定位精度 RMSE ATE 仅为 0.0226（FBE 为 0.1680，ANS 为 0.1105），场景覆盖率达 84.38\%（FBE 为 55.87\%）。在更大规模的 HM3D 数据集上，PSNR 达到 19.86（FBE 仅 15.80），RMSE ATE 为 0.0336（FBE 为 1.2004）。消融实验（HM3D）揭示了三个关键发现：（1）引入定位不确定性使 RMSE ATE 从 0.2478 降至 0.0226 量级，同时 PSNR 从 16.15 提升至 19.86；（2）引入 LLM 不仅提升了覆盖率和重建质量，还反过来降低了定位误差——更好的场景重建有助于提升跟踪精度（使用 GPT-4o 时 RMSE ATE 为 0.0336，使用 LLaVA-7B 时为 0.0159，而不使用 LLM 时为 0.1694）；（3）正则化参数 $\lambda$ 对 Laplace 近似的精度敏感——$\lambda = 2 \times 10^{-6}$ 时 RMSE ATE 为 0.0479，$\lambda = 5 \times 10^{-6}$ 时为 0.0551，显示过大的正则化会压制有效信息。系统每步执行约 2 秒（含 3DGS 渲染、Fisher 信息计算和位姿优化），在 2000 步内完成场景探索，动作空间为 5cm 前进和 $5^\circ$ 旋转。

\paragraph{多尺度语义地图}

与前述 GSMem 和 GaussExplorer 均将 3DGS 用作\textbf{底层表示}（高斯仅承载视觉/几何信息，语义通过 CLIP 反向聚合或闭集检测器注入高斯属性）的设计不同，GLMap（2026.05）~\cite{2605_01736} 提出 3DGS 的一个全新角色——将高斯本身作为语义地图的\textbf{原生视觉模态}，直接参与多尺度语义的组织与消费。这一设计的出发点是观察到现有语义地图在\textbf{显式几何、多尺度语义与大模型友好接口}三方面的固有矛盾——类别级地图（如 GOAT 的语义栅格，ObjectNav HM3D SR 50.6\%）语义过粗仅存物体标签，InstNav 上性能骤降至 17.0\%；特征级地图（如 g3D-LF 的 CLIP 栅格、Scene-LLM 的点云特征）虽能编码多尺度语义但需额外的特征-投影对齐训练才能被 LLM/VLM 消费，且基于 2D patch 提取的特征无法描绘清晰的物体轮廓，导致 InstNav SR 仅 11.5\%；稠密几何地图（如 Chat-Scene）依赖完整点云而不支持增量构建，无法用于导航任务。

拓扑地图（如 UniGoal 的 scene graph）虽能显式建模实例关系（InstNav SR 20.2\%），但仅靠文本边描述（如"next to"）缺乏视觉细节，在需要精细场景理解的 SQA 上受限（EM-1 34.2\%）。为此，GLMap 提出\textbf{多尺度高斯-语言地图} $\mathcal{M} = \{m, \mathcal{S}_o, \mathcal{S}_r\}$，以 2D 索引网格 $m$ 精确定位语义单元的世界坐标（比拓扑地图的 scene graph 边缘关系更精确），同时维护实例语义单元集 $\mathcal{S}_o = \{(n_o, o)\}$ 和区域语义单元集 $\mathcal{S}_r = \{(n_r, r)\}$ 两层粒度。每个语义单元\textbf{双模态}存储——自然语言描述可直接输入 LLM，3D 高斯表示可通过 Gaussian Splatting 快速渲染任务相关图像供 VLM 消费，无需任何特征投影对齐训练。实例单元 $o = (\mathcal{G}, T_o)$ 存储一组 3D 高斯 $\mathcal{G}$ 和开放词汇文本描述 $T_o$（如对象类别、属性、颜色）；区域单元 $r = (\mathcal{I}_r, T_r)$ 记录功能区域描述及其包含的实例 ID 集 $\mathcal{I}_r$——不直接存储高斯以节省内存，需要时可从构成实例的高斯联合渲染。

值得强调的是，这种双模态接口使 GLMap 成为首个\textbf{显式提供参考}（explicit reference）的语义地图——GPT4Scene 虽也提供显式参考（鸟瞰图+区域图），但 GLMap 在实例和区域两层均提供文本+渲染图像双通道，覆盖粒度更细。

GLMap 的另一关键创新是\textbf{解析高斯估计器} $f_{GE}$。在具身场景中深度图像和相机位姿均可获得，可直接重建高质量稠密点云，因此无需像传统 3DGS 那样依赖可微优化——估计器从稠密点云通过体素化网格解析推导高斯参数（均值 $\boldsymbol{\mu}_v$、协方差 $\Sigma_v$、颜色 $\mathbf{c}_v$、不透明度 $\alpha_v$）：

\begin{equation}
\boldsymbol{\mu}_v = \frac{1}{|\tilde{\mathcal{P}}_v|}\sum_{\mathbf{p}_i \in \tilde{\mathcal{P}}_v} \mathbf{p}_i, \quad \Sigma_v = \frac{1}{|\tilde{\mathcal{P}}_v|}\sum_{\mathbf{p}_i \in \tilde{\mathcal{P}}_v}(\mathbf{p}_i - \boldsymbol{\mu}_v)(\mathbf{p}_i - \boldsymbol{\mu}_v)^T + \epsilon I
\end{equation}

其中 $\tilde{\mathcal{P}}_v = \bigcup_{\tilde{v} \in \mathcal{N}(v)} \mathcal{P}_{\tilde{v}}$ 聚合了 Chebyshev 距离 $\|\tilde{v} - v\|_\infty \leq 1$ 的邻域体素点集以防止孤立高斯并确保空间重叠，$\epsilon I$ 为正则化项；颜色 $\mathbf{c}_v$ 取逐点均值，不透明度 $\alpha_v$ 固定为 0.8。体素分辨率设为 1\,cm，比传统 3DGS 的 SfM 初始点云密度更高，充分利用了具身场景中 RGB-D 传感器提供的稠密深度。估计后通过曲率感知合并策略去除冗余高斯：两个候选高斯的相似度由均值距离、协方差 Frobenius 范数距离和颜色距离的加权和 $D(G_i, G_j) = \|\boldsymbol{\mu}_i - \boldsymbol{\mu}_j\|_2 + \lambda_\Sigma\|\Sigma_i - \Sigma_j\|_F + \lambda_c\|\mathbf{c}_i - \mathbf{c}_j\|_2$ 度量（论文设 $\lambda_\Sigma=0.6, \lambda_c=0.4$，赋予形状相似度略高于颜色相似度的权重），当 $D(G_i, G_j) < 1 + \tau(\kappa(\Sigma_i) + \kappa(\Sigma_j))$ 时合并——其中曲率代理 $\kappa(\Sigma) = \lambda_{\min}(\Sigma)/\text{tr}(\Sigma)$ 由最小特征值（对应表面法线方向）与迹的比值定义，确保高曲率区域保留更多高斯以维持边缘清晰，平滑区域则加强合并以降低内存开销。合并后新高斯的参数由两个被合并高斯的参数均值获得。这一解析估计管线的关键优势在于\textbf{完全避免了梯度优化}——传统 3DGS 需要数千到数万次可微渲染迭代来精调高斯参数，而 GLMap 的 $f_{GE}$ 从稠密点云到高斯表示仅需一次前向计算，使实时增量构建成为可能。

增量更新时，给定新 RGB-D 观测，MLLM（Gemma3-27B）首先从 RGB 图像生成结构化语义表示 $\mathcal{S}_t = \{(\tilde{n}_o, T_o)\} \cup \{(\tilde{n}_r, T_r, \tilde{\mathcal{I}}_r)\}$，其中 $\tilde{n}$ 为帧内局部 ID，GLMap 维护全局 ID $n$，两者之间的映射在实例匹配过程中动态确定。对每个实例描述，GroundingDINO + MobileSAM 提取分割 mask 并投影到 3D 点云，再经高斯估计器转化为 3D 高斯表示。新观测实例通过\textbf{语义余弦相似度 + 高斯可合并性}的双重准则匹配已有实例：先检查文本嵌入（nomic-embed-text 模型）余弦相似度是否超过阈值 $\tau_s$（论文设为 0.8），再判断高斯集是否满足合并条件；匹配成功则合并高斯集并拼接文本描述（超出长度阈值 300 词时由轻量级 Qwen3-8B 压缩合并），同时更新 2D 索引网格并记录局部-全局 ID 映射 $\mathcal{T}_t$。若无可匹配实例，则注册新实例并分配新的全局 ID。区域单元的增量更新则额外要求\textbf{实例集一致性}——先将区域的局部实例 ID 通过映射 $\mathcal{T}_t$ 转为全局 ID，然后检查两个区域是否共享至少一个公共实例且语义相似度超过阈值才可合并，确保区域-实例层级的一致性。

\textbf{导航推理与语义价值地图。}GLMap 在下游导航任务中的核心机制是基于多尺度语义计算\textbf{语义价值地图}（value map），将目标描述与地图中各语义单元的匹配得分投影到 2D 栅格上。给定目标描述 $T_g$，计算每个语义单元 $u \in \mathcal{S}_o \cup \mathcal{S}_r$ 与目标的相似度 $s_u = f_S(u, T_g)$，其中 $f_S$ 可由 LLM（文本-文本匹配）或 VLM（渲染图像-目标匹配）实现。实例和区域的渲染图像通过 Gaussian Splatting 从智能体高度的可导航位置获得，视角选择为最大化 3D 高斯可见投影（非遮挡）的平面角度。价值地图通过 2D 高斯核平滑生成：

\begin{equation}
H(l) = \frac{1}{Z}\sum_{v \in m_t}\Big(\sum_{u \in \mathcal{S}_o \cup \mathcal{S}_r} s_u\, \delta(v - p_u)\Big)\mathcal{K}_\sigma(l - v)
\end{equation}

其中 $\delta(\cdot)$ 将相关得分分配到对应栅格单元，$\mathcal{K}_\sigma(\cdot)$ 为带宽 $\sigma$ 的 2D 高斯核实施空间平滑，$p_u$ 为语义单元的 2D 栅格坐标。下一航路点选为距价值地图最大值位置最近的前沿，再由局部导航策略（如 FMM）计算到航路点的路径。

在零样本 ObjectNav（MP3D SR 42.5\%/SPL 18.3\%，HM3D SR 62.7\%/SPL 33.7\%）、InstNav（HM3D SR 22.5\%/SPL 13.7\%）和 SQA（SQA3D EM-1 58.5\%/EM-R1 61.3\%）三个任务上，GLMap 均达到当时最优水平——特别是相比仅存物体标签的 GOAT（ObjectNav 50.6\% vs.\ GLMap 62.7\%）和仅存区域特征的 g3D-LF（SQA 47.7\% vs.\ GLMap 58.5\%），双模态多尺度语义带来显著提升。消融实验显示实例单元和区域单元各自带来独立增益且互补——仅加实例单元 SR 从 52.5\% 提升至 57.4\%（+4.9\%），仅加区域单元提升至 56.2\%（+3.7\%），两者结合进一步提升至 59.1\%（+6.6\%），表明实例级（物体类别、属性、颜色）和区域级（功能区域、空间关系、场景描述）语义提供了互补线索。此外，GLMap 作为即插即用模块无缝集成到 ESC（LLM 方法，SR +9.6\%/SPL +2.9\%）、VLFM（VLM 方法，SR +6.6\%/SPL +1.8\%）、ApexNAV（SR +3.1\%/SPL +0.7\%）和 GPT4Scene（MLLM 方法，EM-1 +1.3\%/EM-R1 +0.9\%）等基线中，验证了双模态接口对 LLM、VLM、MLLM 三类大模型架构的零样本兼容性——GLMap 不改变基线方法的推理框架，仅替换其语义地图输入即可获得一致增益。

\paragraph{个性化导航训练}

上述方法（GSMem、GLMap 等）均假设策略已在仿真中训练完成，3DGS 的角色局限于\textbf{在线建图与语义推理}——即 3DGS 在导航过程中作为场景表示被构建和消费，但不参与策略本身的训练。EmbodiedSplat（2025.09）~\cite{2509_17430} 则将 3DGS 的角色从导航阶段的\textbf{运行时表示}前置为训练阶段的\textbf{环境生成器}，提出一个互补视角：\textbf{能否利用 3DGS 的低成本场景重建能力，直接弥合导航策略的仿真-真实鸿沟？}该问题的核心瓶颈在于，传统高质量仿真场景（如 HM3D 的 800 训练场景和 HSSD 的 134 训练场景）依赖 Matterport 扫描仪采集，设备昂贵且耗时数小时，且场景风格偏重公寓布局——当部署环境为大学教室、会议室等分布外场景时，预训练策略的零样本性能急剧下降。EmbodiedSplat 正是针对这一分布外部署难题，提出用消费级手机快速重建部署环境并在其中个性化微调导航策略。

\textbf{Real-to-Sim-to-Real 管线。}EmbodiedSplat 的核心思路是用消费级手机快速重建部署环境，在重建场景中微调导航策略后部署到真实机器人。

具体管线分四个阶段：(1)~用 iPhone 13 Pro Max 通过 Polycam 应用采集 RGB-D 数据（约 20--30 分钟，手持无需稳定器），再用 Nerfstudio 处理并采样 1000 帧低模糊度对齐 RGB-D 图像及位姿；(2)~以 DN-Splatter（深度-法线正则化的 3DGS）训练 30{,}000 迭代，传感器深度损失权重 $\lambda_d = 0.2$，启用深度平滑和法线一致性损失，采用 Metric3D-V2 作为法线编码器（实证优于 Omnidata——在 MuSHRoom 数据集 10 个场景上，Metric3D-V2 在深度误差和法线一致性指标上均胜出），通过 Poisson 重建将高斯场转为网格（\texttt{.ply}），再经 Blender 转换为 \texttt{.glb} 格式；(3)~导入 Habitat 仿真器，仅选取每个场景最大导航网格岛屿（navmesh island）以确保可通行区域完整性（因采集场景尺度较小，通常仅 1--3 个房间，选取最大岛屿可避免将非地板实体误认为可通行区域），生成 1000 训练回合和 100 评估回合的 ImageNav 任务；(4)~在重建场景中以仅 20M 步微调预训练策略（学习率 LSTM $2.5\times10^{-6}$、视觉编码器 $6\times10^{-7}$），相比从零训练的 100M+ 步大幅降低计算成本。

与 GaussNav（非端到端语义高斯地图、仅仿真中的 Instance ImageNav）和 SplatNav（模块化规划与定位策略、无 sim-to-real 验证）不同，EmbodiedSplat 是首个在 3DGS 重建场景上实现端到端 ImageNav 策略训练并完成真实世界部署验证的方法。与 Phone2Proc（同样使用 iPhone 采集但仅生成房间布局、通过后处理生成多种场景变体用于 ObjectNav）相比，EmbodiedSplat 直接捕获完整房间的 3DGS 重建而非抽象布局，专注于 ImageNav 而非 ObjectNav，且不做场景后处理——更简单但更依赖重建的视觉保真度。整个采集-训练管线总时间约 1.5--2.5 小时（采集 20--30 分钟 + DN-Splatter 训练 1--2 小时），远低于 Matterport 方案。

\textbf{关键实验发现。}EmbodiedSplat 的实验分析围绕三个维度展开——重建质量基准、策略微调收益和仿真-真实可预测性。

\textbf{重建质量基准。}在 MuSHRoom 数据集的 10 个场景上，DN-Splatter（Metric3D-V2 法线编码器）生成的 DN mesh 在视觉保真度上普遍优于 Polycam 直接输出的网格，尤其在细粒度几何（如家具边缘、桌面物体）处更为明显。然而两种管线各有局限：DN mesh 在大面积平坦区域（如白墙、天花板）可能出现空洞或伪影，而 Polycam mesh 在复杂几何处更容易出现过度平滑。论文指出，选取 Metric3D-V2 而非 Omnidata 作为法线编码器的关键依据正是在 MuSHRoom 场景上的系统对比——前者在深度 L1 误差和法线角度误差上均更低。

\textbf{策略微调收益。}在大学环境（教室、休息室、会议室等四个场景——lounge、classroom、conf\_a、conf\_b，尺度与 MuSHRoom 场景类似，均为 1--3 个房间）的 Stretch 机器人真实部署中，EmbodiedSplat 揭示了微调个性化策略的显著收益。HM3D 预训练策略（基于 Silwal et al. 的 VC-1 端到端架构）在真实世界零样本成功率仅 50\%（目标场景为休息室，与 HM3D 的公寓风格差异显著），HSSD 预训练策略更降至 10\%；而在 DN-Splatter 重建场景中微调仅 20M 步后，HM3D 基线提升 20 个百分点至 70\%，HSSD 基线提升 30--40 个百分点至 40--50\%。值得注意的是，微调在 DN mesh 上的效果持续优于 Polycam mesh，进一步验证了 DN-Splatter 深度-法线正则化对策略迁移的正面影响。

\textbf{仿真-真实可预测性。}EmbodiedSplat 借鉴 Kadian et al. 提出的 Sim-vs-Real 相关系数（SRCC）框架，系统测量了仿真中策略性能与真实世界表现的 Spearman 相关性。结果表明，重建网格的仿真-真实相关性高达 0.87--0.97，验证了 3DGS 重建网格作为策略训练仿真环境的保真度——仿真中的性能排序可靠预测真实世界表现。一个令人意外的发现是：即使不使用大规模数据集预训练，仅对 Polycam 重建的\textbf{单个场景}过拟合训练也能在真实世界达到 50\% 成功率，暗示 3DGS 重建场景的视觉保真度足以支撑策略的部分迁移。消融分析还揭示了若干实用洞察：场景尺度与成功率呈负相关（越大的场景越难导航），验证 PSNR 与成功率呈正相关（重建质量直接影响策略表现），且 HM3D 预训练策略在约 400M 步后出现零样本性能退化——持续优化仿真数据集并不自动转化为对新场景的泛化提升，这为数据规模与策略性能之间的边际收益递减提供了实证依据。

\paragraph{场景重排与探索}

与前述 GSMem、GaussExplorer 将 3DGS 用作导航地图、EmbodiedSplat 将 3DGS 用作策略训练环境的思路截然不同，SplatR（2024.11）~\cite{2411_14322} 将 3DGS 引入了一个全新的具身任务——\textbf{经验目标视觉重排}（experience goal visual rearrangement）：智能体先在目标配置下的场景中漫游记忆，再在打乱后的同一场景中将物体恢复原位。这一任务的核心挑战在于构建一个精确的世界模型以编码目标状态，并与当前打乱状态进行高效对比——而 SplatR 恰恰将 3DGS 置于这一对比的枢纽位置。

SplatR 是首个将 Gaussian Splatting 用于重排任务的方法，其选择 3DGS 的动机在于：与仅存储离散视角的图像方法（如 Weihs et al. 的端到端 CNN+RNN）相比，3DGS 支持\textbf{从任意位姿渲染目标视角}，智能体无需精确回到初始拍摄位置即可获得对应的目标配置图像；与点云方法（如 Liu et al. 的类别无关点云匹配）相比，3DGS 渲染的高质量图像携带丰富的外观语义信息，可直接进行视觉对比；与场景图方法（如 Gadre et al. 的连续场景表示）相比，3DGS 不需要预定义物体间关系的图结构，而是通过稠密渲染隐式编码所有视觉细节。

具体地，\textbf{第一阶段}（walkthrough，最多 700 步）中智能体以基于 2D 障碍物地图的增量探索策略覆盖场景——随机采样未探索区域并通过 fast-marching 方法导航，仅在接近障碍物时才更新地图以确保获取物体的近距离观测——每帧保存 RGB 图像、智能体位姿和点云数据，探索结束后点云通过体素网格降采样以限制高斯初始化数量并降低计算开销，采集的多视角数据用于训练目标配置的完整 3DGS 表示（每个场景训练约 20--30 分钟，场景越大所需高斯越多，训练时间相应增长）。

\textbf{第二阶段}（unshuffle，最多 1200 步）中智能体在打乱场景中探索时（每个 episode 有 1--5 个物体状态变化，变化类型包括位置移动或开放程度改变），\textbf{虚拟相机}在训练好的 3DGS 中同步跟踪智能体位姿并渲染目标配置视角，使智能体在任意时刻同时拥有当前观测图像与对应的目标渲染图像——这是 SplatR 管线的核心设计：通过 3DGS 实现任意位姿的"目标状态回放"，将变化检测问题转化为同一视角下两帧图像的直接对比。变化检测通过\textbf{DINOv2 补丁级稠密特征匹配}实现——选择 DINOv2（自监督视觉Transformer，通过自蒸馏训练学习强大的通用视觉特征，而非任务专用模型）作为特征提取器的动机是其在大规模多样化语料上的任务无关训练赋予了跨数据集的泛化能力——计算当前帧与目标帧对应补丁（patch size $14\times14$）的余弦相似度
\begin{equation}
S_{ij} = \frac{f_{ij}^C \cdot f_{ij}^G}{\|f_{ij}^C\| \|f_{ij}^G\|}
\end{equation}

将相邻差异补丁聚类为物体候选。每个物体节点存储五元组信息：对应图像 $I$、掩码 $m$、CLIP 语义嵌入 $g$（沿物体 mask 裁剪的图像区域提取）、中心坐标 $p$ 和世界坐标系下的点云 $P$。为兼顾计算效率与分割精度，探索阶段仅存储 DINOv2 补丁级粗掩码，探索结束后才统一以 SAM（边界框提示）精化物体边界——准确 mask 对后续物体拾取操作至关重要，但每帧同时运行 SAM 和 DINOv2 在计算上不可行，因此将 SAM 推迟到探索结束后批量执行。针对 3DGS 渲染伪影（尤其出现在墙壁、镜面等平坦/反射表面上的假阳性检测），SplatR 计算 CLIP 特征与条件词（如"墙壁""镜子"）的相似度进行过滤——若检测区域与条件词的相似度最高则判定为假阳性并丢弃。

探索过程中，每帧新检测到的物体节点通过视觉相似度（CLIP 嵌入余弦相似度 $\sigma^{\text{vis}}$）和空间相似度（点云最近邻比率 $\sigma^{\text{sp}}$，即点云 $P$ 中在点云 $P_i$ 中存在最近邻的点占比）的加权融合 $\sigma_i^{\text{node}} = \lambda \sigma_i^{\text{vis}} + (1-\lambda)\sigma_i^{\text{sp}}$ 与已有节点匹配合并，合并后特征向量以递增方式更新 $g_j \leftarrow (n_o^j g_j + g)/(n_o^j + 1)$，其中 $n_o^j$ 为历史合并次数。探索结束后，当前配置物体节点集 $\mathcal{O}_s$ 与目标配置节点集 $\mathcal{O}_g$ 通过 CLIP 嵌入的余弦相似度构建二分图，以\textbf{类别无关}方式（不依赖物体分类，直接比较视觉特征，从而绕过物体检测器可能引入的分类错误）通过匈牙利算法求解最优匹配（消融实验表明匈牙利算法优于贪心匹配策略，尽管时间复杂度为 $O(n^3)$，但 RoomR 数据集为单房间环境，物体节点数足够小，不引入显著延迟），匹配对按相似度排序依次执行重排动作。在 AI2-THOR RoomR 数据集（120 个房间、涵盖厨房/客厅/浴室/卧室等 70+ 物体类别）的 26 个随机采样 episode 上，SplatR 的所有检测和分割模型均未在 RoomR 场景上训练，结果为零样本性能。

SplatR 在 \%Fixed Strict（36.35\%）和 \%Misplaced 指标上达到 SOTA，同时 \%Energy Remaining 为所有方法中最低，表明重排后的物体最接近目标位置。\%Fixed Strict 与 \%Fixed 的紧密对齐进一步表明智能体的假阳性检测率极低——极少干扰未被移位的物体。然而整体成功率仍较低，根本原因在于 SplatR 的管线仅处理物体位置重排，不区分拾取/放置与开放/关闭操作（如抽屉的开合）——对于需要调整开放程度的 episode，即使所有物体位置正确，成功率也为零。其主要局限在于 DINOv2 的 $14\times14$ 补丁分辨率对小物体（如铅笔）的检测精度不足，以及 3DGS 表示的内存开销限制了场景规模的可扩展性。

与 SplatR 在已完成重建的静态场景中进行物体重排的任务设定不同，AIR-Embodied（2024.09）~\cite{2409_16019} 将 3DGS 的角色从被动场景表示前置为\textbf{主动重建过程的优化目标}——其核心问题是：如何利用多模态大语言模型（MLLM）的高层推理能力，规划视点选择与交互动作以最大化 3DGS 的重建质量？该问题的瓶颈在于，传统 Next Best View（NBV，下一最优视角）方法（如基于不确定性的和基于 Fisher 信息的方法）仅从球面候选集中选择视点，依赖低层不确定性信号而缺乏对重建状态的高层理解，且无法处理物体自身遮挡（如物体底部）——即使选定最优视点，接触面仍不可见。AIR-Embodied 的设计动机正是突破这两层限制：(1)~用 MLLM 将像素级不确定性提升为语义级理解，实现自由空间中的视点规划而非限于球面候选集；(2)~引入机械臂交互操作暴露遮挡区域，使重建不再受限于被动观测。

\textbf{三阶段主动重建框架。}AIR-Embodied 采用 PGSR（平面正则化 3DGS 变体）作为底层面表示，在其上构建三阶段闭环管线。

\textbf{阶段一：重建状态理解。}系统首先定义\textbf{像素级不确定性}指标——基于 3DGS $\alpha$-混合 过程中沿深度的权重变化 $w_i = T(i) - T(i-1)$ 的分布，计算每个高斯基元的深度分布熵 $\xi_j = -\sum_{i=1}^{n} P(w_i, d_i) \log P(w_i, d_i)$：对于高质量重建表面，像素颜色由单个高斯主导，权重集中在表面附近（低熵）；而未充分重建区域则呈现多个分散峰值（高熵）。在以目标物体为中心构建的 3D 体素地图 $\gamma: \mathbb{Z}^3 \to \mathbb{R}^n$ 中，系统沿稀疏采样点 $p_i \in \mathcal{P} \subseteq \mathbb{Z}^3$ 向物体质心方向渲染，逐视点统计高熵像素占比 $\Omega(\mathbb{S}(R,t)) = \sum_1^N \mathbb{I}(\xi_j > \xi_t) / M$（$M$ 为总像素数，$\xi_t$ 为阈值），从而得到覆盖整个任务空间的不确定性体素地图 $\gamma = \{\varphi_\Omega\}$。随后，MLLM（GPT-4o）将不确定性分布与对应渲染图像结合进行多模态推理，生成高层理解 $\mathcal{L}_H = \{l_\text{t}, l_\Omega, l_\text{o}, l_d\}$——其中 $l_\text{t}$ 标识欠重建区域位置，$l_\Omega$ 量化不确定性程度，$l_\text{o}$ 判断可观测性（遮挡原因），$l_d$ 为决策因子（选择观察还是交互）。这一步的关键设计是将 MLLM 的自然语言推理能力桥接到 3D 空间：高层理解被映射回体素地图 $\gamma \to \gamma_q$，作为后续视点规划的解空间。

\textbf{阶段二：视点与交互规划。}基于 $\mathcal{L}_H$ 和质量地图 $\gamma_q$，系统通过三类约束对候选视点评分：(1)~任务加权约束 $W_\Phi(p) = \mathbb{I}(p \in \varsigma)$（$\varsigma$ 为 MLLM 推理生成的候选区域集，确保视点指向欠重建区域）；(2)~距离约束 $W_\kappa(p) = \exp(-\lambda_\kappa(\kappa_p - \kappa_r)^2)$（高斯衰减，惩罚过远或过近的视点，最优距离阈值 $\kappa_r$ 使目标在图像中占据合适尺寸）；(3)~密度约束 $W_\varkappa(p) = \exp(-\lambda_\varkappa(\kappa_p - \kappa_m)^2)$（抑制视点过于密集导致的冗余采集）。综合评分 $\varphi(p) = \varphi_\Omega(p) \cdot W_\Phi(p) \cdot W_\kappa(p) \cdot W_\varkappa(p)$ 最高者被选为下一视点。对于\textbf{不可观测区域}（如物体与桌面的接触面），MLLM 规划交互动作——例如将任务分解为"推倒物体"和"补充底面视点"两个子任务，通过贪心算法在体素地图的 $\varphi_a$ 属性中搜索高分离数点作为轨迹点，再由 MoveIt 生成平滑动作轨迹。交互完成后系统重新感知并更新体素地图 $\gamma \to \gamma_\text{new}$，重新规划视点。

\textbf{阶段三：闭环推理。}由于真实世界中感知误差和控制偏差导致开环指导偏离期望位姿，AIR-Embodied 设计了闭环验证模块——每步动作 $\tau$ 完成后，比较实际场景状态 $\mathcal{L}_i^S$ 与期望状态 $\mathcal{L}_\text{desired}^S$，计算偏差 $\Delta\mathcal{L}^S = \mathcal{L}_i^S - \mathcal{L}_\text{desired}^S$ 并据此微调修正，确保精确执行。

\textbf{实验验证。}在 OmniObject3D 数据集的 50 个物体（零样本泛化）上，AIR-Embodied 在 20 视点限制下达到 PSNR 30.846、Chamfer 距离 0.0058、F-score 0.9237，均显著优于 FisherRF（PSNR 28.687、Chamfer 0.0096、F-score 0.7648）和基于不确定性的 NBV 方法（PSNR 28.481、Chamfer 0.0108）。平均贡献率（ACR，衡量每个新视点对重建质量的平均贡献）为 5.08\%，高于 FisherRF 的 4.32\%，证实自由空间视点规划的效率优势。消融实验显示交互操作带来显著增益：去除操作模块后 PSNR 从 30.846 降至 29.113（$-$1.743），Chamfer 从 0.0058 恶化至 0.0089（$+$0.0031），表明机械臂交互暴露遮挡区域对重建完整性的关键作用。在真实世界实验中（UR5e 机械臂 + RealSense D455，覆盖简单/中等/复杂三类物体），闭环推理模块的贡献在复杂物体上尤为突出——有闭环时 Chamfer 为 0.0299，无闭环时恶化至 0.0372（$+$24\%），验证了闭环纠偏对感知敏感任务的必要性。此外，与 GPT-4o 直接规划（重复率 73.3\%、ACR 3.21\%）和 VoxPoser 风格规划（重复率 66.7\%、ACR 3.27\%）的对比表明，基于重建质量反馈的引导机制将重复率压至 13.3\%，有效克服了 LLM 对视点选择任务的程序化重复倾向。

\subsubsection{安全规划与碰撞避免}

本路线介绍在 3DGS 环境中进行安全规划与碰撞避免的方法。在 3DGS 表示的环境中进行安全导航，需要将高斯基元的几何信息转化为可计算的碰撞约束。直觉上，每个 3DGS 高斯本质上是一个 3D 椭球体，而椭球体与机器人之间的碰撞检测在数学上有成熟的理论基础——这使得 3DGS 场景天然适合形式化安全导航问题，而不需要先将高斯转换为网格或体素等中间表示。3DGS 的显式椭球表示为碰撞检测提供了数学上优雅的形式化。

\paragraph{控制障碍函数方法}

SAFER-Splat（2024.09）~\cite{2409_09868}（Simultaneous Action Filtering and Environment Reconstruction）修改了 3DGS 的碰撞检测与安全控制环节——用基于欧氏距离的控制障碍函数（Control Barrier Function, CBF）替代传统的 Mahalanobis 距离度量或 NeRF 体积查询——3DGS 的渲染管线和训练过程保持不变。该方法是首个在 3DGS 地图上实现实时安全过滤的系统，核心思路是直接利用 3DGS 高斯基元的椭球几何，将安全导航问题转化为可快速求解的二次规划（Quadratic Programming, QP）。

\textbf{场景表示与 Baseline 做法。}在标准 3DGS 中，每个高斯基元由均值 $\mu_i \in \mathbb{R}^3$、协方差矩阵 $\Sigma_i = R S S^T R^T$（$R$ 为旋转矩阵，$S$ 为对角缩放矩阵）、不透明度 $\alpha$ 和球谐系数参数化。对于碰撞检测而言，3DGS 的一个关键属性是：每个高斯的几何外壳天然是一个 3D 椭球体。具体地，SAFER-Splat 将每个高斯基元视为 $1\sigma$ 置信椭球体 $\mathcal{E}_i = \{x \in \mathbb{R}^3 | (x-\mu_i)^T \Sigma_i^{-1}(x-\mu_i) \leq 1\}$，整个场景表示为 $K$ 个椭球体的并集 $\mathcal{G} = \{\mathcal{E}_i\}_{i=1}^K$。这意味着 3DGS 场景可以直接作为碰撞地图使用，无需先将高斯转换为网格或体素等中间表示。

\textbf{设计动机。}现有基于 NeRF 的安全导航方法（如 NeRF-CBF）需要通过体积渲染查询碰撞信息，计算成本高且难以实时运行；而直接使用 Mahalanobis 距离作为 CBF 度量虽然可行，但 Mahalanobis 距离是一种非欧氏度量——对于小型椭球体，即使实际欧氏距离很近，Mahalanobis 距离也可能剧烈变化，导致安全控制器行为不平滑。此外，场景中同时存在大小差异巨大的椭球体，Mahalanobis 距离在不同椭球上的变化率差异极大，给控制器设计带来困难。因此，SAFER-Splat 选择直接计算球体机器人与椭球之间的最小欧氏距离作为 CBF 的基础。

\textbf{改了什么：欧氏距离 CBF。}SAFER-Splat 的核心创新是将球体-椭球最小距离问题形式化为一个二次规划，并通过解析求解其对偶问题。直觉上，给定一个以位置 $\mathbf{p}$ 为中心、半径 $r$ 的球形机器人，寻找椭球 $\mathcal{E}_i$ 表面上离 $\mathbf{p}$ 最近的点 $\mathbf{y}^*$，即为球心到椭球的最短距离。由于距离关于平移和旋转不变，问题可以转换到椭球对齐坐标系中求解：

\begin{equation}
    d_i(\hat{\mathbf{p}}) = \min_{\hat{\mathbf{y}}} \|\hat{\mathbf{p}} - \hat{\mathbf{y}}\|_2^2 \quad \text{s.t.} \quad \hat{\mathbf{y}}^T S^{-2} \hat{\mathbf{y}} \leq 1,
\end{equation}

其中 $\hat{\mathbf{p}} = FR^T(\mathbf{p} - \mu_i)$ 是转换到椭球对齐坐标系后的机器人位置，$S$ 为缩放矩阵，$F$ 为反射矩阵（利用椭球对称性将问题限制在正象限）。该二次规划只有单个约束，始终可行。

\textbf{实现原理。}作者不使用通用 QP 求解器，而是通过拉格朗日对偶解析求解。对偶问题的一阶最优性条件给出椭球表面最近点的闭式表达 $\hat{\mathbf{y}}^*(\lambda) = \left(\ldots, \frac{S_{ii}^2 \hat{p}_i}{\lambda + S_{ii}^2}, \ldots\right)^T$，其中对偶变量 $\lambda$ 通过求解单调递减函数 $q(\lambda) = \sum_{i=1}^n \frac{S_{ii}^2 \hat{p}_i^2}{(\lambda + S_{ii}^2)^2} - 1 = 0$ 确定。由于 $q(\lambda)$ 在可行域内严格单调且从 $+\infty$ 趋向 0，方程有且仅有一个解，可用 GPU 并行二分搜索高效求解到任意精度。

在此基础上，SAFER-Splat 定义控制障碍函数为：

\begin{equation}
    h_i(\mathbf{p}) = \phi_i(\hat{\mathbf{p}}) \cdot d_i(\hat{\mathbf{p}}) - (r + \epsilon)^2, \quad \phi_i = \text{sign}(\hat{\mathbf{p}}^T S^{-2} \hat{\mathbf{p}} - 1),
\end{equation}

其中 $\phi_i$ 是符号函数（机器人在椭球外时为 $+1$，穿透椭球时为 $-1$），$r$ 为机器人半径，$\epsilon$ 为安全缓冲距离（补偿 3DGS 重建误差和机器人定位不确定性）。当 $\epsilon = 0$ 时，$h_i > 0$ 表示安全（机器人与椭球无交叉），$h_i = 0$ 表示恰好接触，$h_i < 0$ 表示碰撞。CBF 的梯度通过包络定理（Envelope Theorem）高效计算：$\nabla_{\mathbf{p}} h_i = 2\phi_i(\mathbf{p} - \mathbf{y}^*)^T$，Hessian 矩阵通过隐函数定理（Implicit Function Theorem）获得闭式表达。

系统采用双积分器动力学模型 $\dot{x} = [\dot{\mathbf{p}}; \dot{\mathbf{v}}] = [\mathbf{v}; \mathbf{u}]$（无人机底层姿态稳定控制器使其近似双积分器行为），将 CBF 约束嵌入最小侵入性二次规划：

\begin{equation}
    \min_{\mathbf{u}} \|\mathbf{u} - \bar{\mathbf{u}}\|_2^2 \quad \text{s.t.} \quad 2\phi_i(\mathbf{p} - \mathbf{y}^*)^T(\mathbf{u} + (\alpha + \beta)\mathbf{v}) \geq -\mathbf{v}^T \mathbf{H}_i \mathbf{v} - \alpha\beta h_i,
\end{equation}

其中 $\bar{\mathbf{u}}$ 为期望控制输入（来自操作员或规划器），$\alpha, \beta > 0$ 为 CBF 增益参数。该 QP 的目标是在满足安全约束的前提下，使安全控制 $\mathbf{u}$ 尽可能接近期望控制 $\bar{\mathbf{u}}$，即仅当动作不安全时才进行修正。作者证明该 QP 始终可行（$u_0 = -(\alpha + \beta)\mathbf{v}$ 是一个闭式可行解），且安全集是前向不变的。此外，速度和加速度限制也可纳入同一 QP 框架，参数约束为 $\alpha + \beta \leq a_{\max}/v_{\max}$。

\textbf{约束剪枝与计算效率。}场景中可能有数十万高斯基元，直接为每个椭球生成 CBF 约束会导致 QP 规模过大。SAFER-Splat 利用动力学可行性进行约束剪枝：给定最大速度 $v_{\max}$ 和安全层推理时间 $\delta t$，机器人在两次 CBF 求解之间只能到达以 $\mathbf{p}$ 为中心、半径 $v_{\max}\delta t + r$ 的球内。因此，与机器人距离超过此半径的椭球对应的约束在当前时刻必然冗余，可以安全移除。这种剪枝在稠密 GSplat 中收益最大，可剪除约 $99.9\%$ 的高斯约束。该策略不会影响安全性，因为被移除的约束在当前时刻不可能被激活。

\textbf{性能表现。}在仿真实验中，SAFER-Splat 在 100K--500K 高斯组成的 5 个场景中，比 NeRF-CBF 快 20--50 倍、更安全且更少保守性。在真实无人机硬件实验中，系统在 150K--300K 高斯的场景上以 15 Hz 运行安全过滤，同时进行在线 GSplat 训练，GPU 计算仅占总时间的一小部分，不干扰训练。即使在操作员反复尝试驱使无人机撞击障碍物的情况下，安全过滤器也能可靠防止碰撞（CBF 值 $h_{\min} > 0$），仅在接近安全边界时才激活修正，实现了最小侵入性控制。该方法还可扩展至椭球形机器人，通过线性变换将椭球-椭球距离问题转化为球-椭球问题。

\paragraph{速度空间碰撞锥安全过滤}

SAFER-Splat 的安全过滤器虽然在位置空间中实现了高效的欧氏距离 CBF，但其约束激活方式本质上是\textbf{反应式}的——只有当机器人已经接近障碍物时才触发规避修正。碰撞锥 CBF 方法（2025.09）~\cite{2509_14421}（Tscholl, Nakka \& Gunter, Georgia Tech, IEEE L-CSS/ACC）保留了 SAFER-Splat 的整体框架（双积分器动力学模型 $\dot{x} = [0, I; 0, 0]x + [0; I]u$ + 最小侵入性安全 QP + PD 期望控制器 $\bar{u}$），但将安全约束的定义域从\textbf{位置空间}切换到\textbf{速度空间}——SAFER-Splat 的 CBF $h_i(\mathbf{p})$ 仅依赖机器人位置 $\mathbf{p}$ 与椭球表面的欧氏距离，而碰撞锥 CBF 的 $h_i(p,v)$ 同时编码位置和速度，将"距离多近"的几何查询替换为"当前速度方向是否将导致碰撞"的运动学预判——从而实现\textbf{主动式}约束激活：安全过滤器在机器人尚远离障碍时即开始预调整速度方向，产生更平滑的规避机动。作者的核心观察是：3DGS 高斯基元的二次型椭球几何使得射线-椭球交点问题具有闭式解，由此可直接推导出前向碰撞锥（forward collision cone）的解析条件，而无需对障碍物进行保守的几何近似或膨胀——这是传统碰撞锥方法在处理任意形状障碍物时无法做到的。

\textbf{差异的核心：碰撞锥 vs.\ 欧氏距离。}SAFER-Splat 回答的问题是"机器人当前位置离椭球有多近"——这是一个纯几何的距离查询，与速度无关。碰撞锥方法回答的问题是"给定机器人当前速度，是否\emph{最终}会撞上高斯基元"——这是一个将位置与速度耦合在一起的运动学预判。具体而言，给定高斯基元的 $99\%$ 置信椭球 $\mathcal{E} = \{x \in \mathbb{R}^3 \mid (x-\mu)^T A (x-\mu) \leq c^2\}$（$A=\Sigma^{-1}$，$c^2 = \chi^2_{3,0.99}$ 为自由度 3 的卡方分布第 99 百分位），视线向量 $r = \mu - p$ 和速度 $v$，作者将射线 $x(t) = p + tv$ 代入椭球不等式得到二次函数 $\phi(t) = at^2 - 2bt + d$（$a = v^T\!Av$，$b = r^T\!Av$，$d = r^T\!Ar - c^2$），则碰撞发生的充要条件为最小值非正且最小值时刻非负，即
\begin{equation}
\underbrace{(v^T\!Av)}_{\beta}\underbrace{(r^T\!Ar - c^2)}_{\gamma} - \underbrace{(r^T\!Av)^2}_{\delta^2} \leq 0, \quad r^T\!Av \geq 0.
\end{equation}

直觉上，$\beta\gamma - \delta^2$ 量度速度向量 $v$ 相对于椭球方向的"对齐程度"——当该值小于零且机器人朝向椭球运动时（$r^T\!Av \geq 0$），速度向量落入了由椭球张成的正向碰撞锥内部，意味着沿当前方向不可避免地将与高斯基元相交。作者通过白化变换 $L = S^{-1}R^T$ 将椭球映射为单位球，证明该条件等价于经典的 Chakravarthy \& Ghose 三维碰撞锥条件（Lorentz 锥不等式）：在白化坐标中将速度 $\tilde{v}$ 分解为球坐标分量 $(V_r, V_\theta, V_\phi)$，碰撞锥条件等价于 $\|\tilde{r}\|^2(V_\theta^2 + V_\phi^2) \leq c^2(V_\theta^2 + V_\phi^2 + V_r^2)$ 且 $V_r \geq 0$——即速度的切向分量不超过某一角度阈值。由于白化矩阵 $L$ 非奇异，原始坐标中的碰撞条件与白化坐标中的 Lorentz 锥条件严格等价，从而保证了与已有碰撞锥理论的一致性。

\textbf{一阶 CBF 构造。}CBF 候选函数取碰撞锥条件的补集 $h_i(p,v) = \beta\gamma - \delta^2 \geq 0$，定义安全集 $\mathcal{C}_i = \{(p,v) : h_i \geq 0\}$。与 SAFER-Splat 的 CBF $h_i(\mathbf{p})$（仅依赖位置）不同，碰撞锥 CBF 同时依赖于位置和速度。关键优势在于：由于控制输入仅影响加速度项（双积分器模型中 $\dot{v} = u$），位置项对时间的导数 $\nabla_p h^\top \dot{p} = 2(\beta Ar - \delta Av)v = 0$ 恰好为零，因此 $h$ 的时间导数中不出现位置梯度——这使得标准一阶 CBF 即可满足要求，无需为处理高阶相对阶而引入高阶 CBF 扩展（HOCBF）。Lie 导数解析可得 $L_f h = 2w(x)^\top f_v(x)$，$L_g h \cdot u = 2w(x)^\top g_v(x) u$，其中 $w(x) = \gamma Av - \delta Ar$。对双积分器系统（$f_v = 0$，$g_v = I$），CBF 约束简化为 $w(x)^\top u \geq -\frac{p_k}{2}h(p,v)$（$p_k > 0$ 为类-$\mathcal{K}$ 函数增益），这是一个关于控制输入 $u$ 的线性约束，可直接嵌入与 SAFER-Splat 相同形式的最小侵入性 QP 中。

\textbf{球形机器人扩展。}对于物理尺寸不可忽略的机器人，SAFER-Splat 通过将机器人半径叠加到安全缓冲 $\epsilon$ 中处理。碰撞锥方法则通过 Minkowski 和将机器人球体半径编码到有效碰撞参数中：在白化空间中，单位球（来自椭球）与机器人球（在白化空间中变为椭球）的 Minkowski 和给出精确的方向依赖碰撞半径 $c_M = c + \rho/\sqrt{n^\top(S^\top S)n}$（$\rho$ 为机器人半径，$n$ 为射线到椭球中心的正交方向）。一种更保守但更简单的近似取 $c_M = c + \rho/s_{\min}$（$s_{\min} = \min(s_1,s_2,s_3)$），此时只需将碰撞锥条件中的 $c^2$ 替换为 $c_M^2$，CBF 结构无需任何修改。

\textbf{主动式约束激活的平滑性优势。}碰撞锥 CBF 的核心设计动机是克服 SAFER-Splat 的反应式约束激活。在 SAFER-Splat 中，欧氏距离 CBF 仅当机器人接近椭球表面时才产生显著的约束力，迫使控制器进行急剧的速度修正；此外 SAFER-Splat 需要通过包络定理和隐函数定理计算 Hessian 矩阵，计算开销较大。碰撞锥方法在两个维度上实现了改进。第一，\textbf{约束激活方式}：即使机器人距椭球尚远，只要速度方向指向椭球，碰撞锥约束即被激活——安全过滤器在远距离就开始微调速度方向，而非等到近距离才急剧转向。作者指出这类似于驾车时提前变道与紧急变道的区别。第二，\textbf{计算效率}：碰撞锥 CBF 为一阶（仅需 Lie 导数 $L_g h$，无需 Hessian），且约束关于控制输入 $u$ 是线性的，QP 规模更小、求解更快。在实现层面，约束在 GPU 上通过 PyTorch 并行计算，QP 使用内点法求解器 Clarabel 求解；训练阶段对高斯基元施加缩放正则化以避免高度各向异性的椭球导致 $\Sigma^{-1}$ 特征值达到 $10^{16}$ 量级的数值问题。

在约 170K 高斯的 Stonehenge 合成场景中（通过 Nerfstudio/Splatfacto 训练，50 条起点均匀分布、目标点在对面 $180^\circ$ 的轨迹），碰撞锥方法的规划时间比 SAFER-Splat 缩短约 3 倍，轨迹 jerk（加加速度）在三种度量（归一化 jerk、RMS jerk、积分平方 jerk ISJ）下均显著降低，轨迹也更短（受益于碰撞锥的主动激活：约束更早介入、更早终止不可行的候选方向），成功率与 SAFER-Splat 几乎一致。实验中也观察到 ISJ 分布存在较重的尾部，这对应于目标点附近的急剧修正。一个已观察到的局限是：当起始位置非常靠近障碍物时，机器人可能缺乏足够的角速度在一个时间步内逃出碰撞锥，导致 $h < 0$ 使 QP 不可行——这是因为碰撞锥约束作为硬约束实现，引入松弛变量或在 MPC 框架中处理可缓解此问题（作者将此列为未来工作）。

\paragraph{HJB 最优控制图像导航}

以上方法（SAFER-Splat、碰撞锥 CBF（2025.09）~\cite{2509_14421}）均围绕 CBF 安全过滤范式展开——将 3DGS 中的高斯基元视为静态障碍，在已有控制器之上叠加位置空间或速度空间的碰撞约束，任务目标是指定目标的点对点导航。BEINGS（2024.09）~\cite{2409_10216} 从任务目标到 3DGS 的使用方式都与之根本不同：它不处理碰撞避免，而是将\textbf{图像目标导航}（ImageNav：给定目标图像 $I_g$，导航至该图像的拍摄位置）形式化为模型预测控制（MPC）框架下的 Hamilton-Jacobi-Bellman（HJB）最优控制问题。在上述 CBF 方法中，3DGS 场景仅作为障碍物几何的表示（椭球位置、协方差），安全过滤器被动地查询场景以计算距离或碰撞锥——3DGS 不参与控制策略的生成。BEINGS 则将 3DGS 的\textbf{视角合成}能力作为核心决策引擎：从任意候选位姿渲染未来观测图像，在不依赖外部经验数据的前提下评估每条候选轨迹的代价，并通过贝叶斯更新随观测动态精化搜索策略。这一设计的动机是弥合 ImageNav 中两类方法的鸿沟：学习型方法（如 OVRL、PoliFormer）依赖数百万训练回合和高度逼真的仿真器（如 Habitat），泛化能力有限；探索型方法（如 Frontier、Stubborn）无需大量数据但缺乏策略精化能力，在复杂环境中难以收敛。BEINGS 的核心洞察正是利用 3DGS 的渲染能力赋予机器人类似人类的"想象力"——\emph{先看见未来，再决定怎么走}。

具体而言，BEINGS 将搜索空间 $\mathbb{S}$ 离散为等体积 Borel 子集 $\{\mathbb{S}_1, \dots, \mathbb{S}_M\}$，全局代价定义为离散 HJB 方程：

\begin{equation}
    J_\pi(s[0], U_T) = \mathcal{D}(I_r[T\!+\!1], I_g) + \sum_{k=0}^{T} \mathcal{L}_\pi(s[k], u[k]),
\end{equation}

其中终端代价 $\mathcal{D}(I_r, I_g)$ 为渲染图像与目标图像的 VLAD 描述符距离，阶段代价 $\mathcal{L}_\pi$ 定义为移动成本除以目标概率与检测成功率之积：

\begin{equation}
    \mathcal{L}_\pi(s[k], u[k]) = \frac{\mathcal{C}(s[k], u[k])}{\mathcal{P}_\pi(s[k\!+\!1]) \cdot \mathcal{Q}(s[k\!+\!1])}.
\end{equation}

直觉上，机器人应优先前往"移动成本低、目标概率高、检测成功率高"的区域——这体现了最优搜索理论的核心思想。

该方法包含三个耦合组件。第一，\textbf{过程感知贝叶斯更新}利用 AnyLoc-VLAD-DINO（一种基于自监督视觉特征和 VLAD 聚类的通用位置描述符）图像描述符定义子集 $\mathbb{S}_i$ 内的检测成功率 $q_i = 1 - \mathcal{D}(I_g, \mathcal{H}(s[k]))$。当机器人在 $\mathbb{S}_i$ 中未发现目标时，该子集的后验概率质量按贝叶斯搜索理论更新为 $p_i[k\!+\!1] = p_i[k] \cdot (1 - q_i[k]) / (1 - p_i[k] q_i[k])$，其余子集的后验按 $r[k\!+\!1] = r[k] / (1 - p_i[k] q_i[k])$ 归一化——直觉上，每步观测要么缩小目标可能区域的范围，要么排除当前子集，概率质量从已搜索区域流向未搜索区域。第二，\textbf{3DGS 模型预测}利用高斯泼溅的渲染能力，对任意候选控制序列前向积分机器人动力学 $\mathcal{F}$（BEINGS 使用小型飞艇机器人的 4-DoF 离散运动模型），从 3DGS 渲染预测长度为 $K$ 的图像序列 $I_{1:K}$，计算上述 HJB 代价的滚动时域近似 $\mathcal{J}_\pi(s[k], k\!+\!K)$。第三，\textbf{Monte Carlo MPC} 不直接求解非凸的 HJB 方程，而是将优化目标转化为分布匹配问题：寻找采样分布 $\hat{\mathbf{U}}$ 使之最小化与最优分布之间的 KL 散度 $D_{KL}(\mathbf{U}_T^* \| \hat{\mathbf{U}})$。具体地，从 $\hat{\mathbf{U}}$ 采样 $N$ 条 $K$ 步控制序列，按代价计算非归一化权重 $w_i^k = e^{-\mathcal{J}_\pi}$ 和归一化权重 $W_i^k = w_i^k / \sum_j w_j^k$，以 $W_i^k$ 为概率对控制序列进行重要性重采样，逐步将 $\hat{\mathbf{U}}$ 逼近最优分布 $\mathbf{U}^*$。

每次执行权重最高的控制序列的首步输入，接收新观测后重新启动贝叶斯更新和 MPC 循环。BEINGS 证明了全局价值函数满足 Bellman 最优性分解 $\mathcal{V}(s[0], T) = \mathcal{V}(s[0], k) + \mathcal{V}(s[k\!+\!1], T)$，确保滚动时域优化在理论上与全局最优一致。

\textbf{性能表现。}BEINGS 在室内飞艇平台（$10 \times 10 \times 2$ m 测试场，飞艇配备单目相机，地面站配备 i7-13700KF CPU 和 RTX 4080 GPU）上验证。3DGS 场景由 534 幅 $1920 \times 1440$ RGB 图像（使用 iPhone 在 OptiTrack 系统中标定拍摄）构建。实验设置三个难度级别：Easy（无障碍，仅调方向即可看到目标）、Medium（1 个 $1.5$ m 高障碍物）、Hard（3 个 $1.5$ m 高障碍物，需长距离绕行后搜索目标）。BEINGS 在三个难度上的成功率（SR）分别为 $100\%$、$100\%$、$78\%$，SPC 分别为 $0.98$、$0.73$、$0.53$，步数（NS）分别为 2、7、12——均在 50 步上限内完成。

作为对比，纯 3DGS + 直接策略在 Hard 任务上成功率仅 $5.11\%$，语义栅格地图 + Frontier 在 Hard 任务上甚至无法在 50 步内完成，VPR 数据库 + 直接策略在 Hard 任务上 SR 为 $64\%$、SPC 降至 $0.00$（碰撞惩罚过高所致）。消融实验表明，贝叶斯更新和 MC-MPC 两个模块缺一不可：去除贝叶斯更新后 MC-MPC 仅依赖图像相似度重采样，机器人无法收敛到特定位置；去除 MC-MPC 后仅使用贝叶斯估计器，机器人虽可收敛但收敛至错误目标；两者同时禁用则退化为随机策略，导航误差最高。在真实硬件实验中，飞艇在预训练的 3DGS 场景中成功执行 ImageNav，终端图像与目标图像视觉上高度一致，验证了仿真到真实的迁移可行性。

\paragraph{风险感知轨迹优化}

SAFER-Splat 和碰撞锥 CBF 代表了一种确定性安全范式：计算机器人到高斯基元的距离或速度锥，以二值判断（安全/不安全）驱动约束。这种范式的局限在于，它将高斯基元的置信椭球视为刚性障碍边界，而 3DGS 场景本质上是连续概率密度场——高斯基元并非硬边界，而是模糊的概率密度分布。SPLANNING（2024.09）~\cite{2409_16915} 的核心洞察正是利用这一概率结构：不再追求确定性安全（绝对不碰障碍），而是保证碰撞概率低于用户指定阈值 $\alpha \cdot \beta$。这一\textbf{概率化安全约束}的直接好处是在安全性与任务可达性之间提供连续可调的折中——用户可通过风险阈值灵活控制风险容忍度，而非被硬性安全边界锁定在过度保守的轨迹上。此外，SAFER-Splat 和碰撞锥 CBF 均面向自由飞行无人机的安全过滤（在已有控制器之上叠加碰撞约束），而 SPLANNING 面向操控器的轨迹级运动规划（从零生成完整轨迹），两者的规划层次和机器人形态均不同。

\textbf{归一化高斯泼溅（Normalized 3DGS）。}概率碰撞约束的前提是将 3DGS 场景解释为有效的概率密度场。然而标准 3DGS 使用未归一化高斯函数表示密度，其 3D 到 2D 变换是简单的正交投影（project）而非沿射线的积分（integrate），因此 3D 高斯不能解释为体积密度函数的基底，碰撞概率也无从计算。为解决这一根本问题，SPLANNING 将标准 3DGS 重新表述为 $n_G$ 个归一化 3D 高斯的基底：$\sigma(x) = \sum_{n=1}^{n_G} w_n G_n(x)$，其中每个 $G_n(x) = \frac{1}{(2\pi)^{3/2}|\Sigma_n|^{1/2}} \exp\!\bigl(-\frac{1}{2}(x - \mu_n)^T \Sigma_n^{-1}(x - \mu_n)\bigr)$ 是严格归一化的高斯密度函数（积分为 1），$w_n \in \mathbb{R}^+$ 为权重参数。在此表述下，透射率函数 $T_a^b[r] = \exp\!\bigl(-\int_a^b \sigma(r(t))\, dt\bigr)$ 可精确计算沿射线的碰撞概率 $P(C_a^b[r]) = 1 - T_a^b[r]$。SPLANNING 同时省略了标准 3DGS 中用于减少渲染伪影的屏幕空间低通滤波步骤（因为该滤波的效果仅存在于 2D 像素空间，无法在 3D 中进行碰撞推理），代之以在世界坐标的 3D 高斯上施加协方差为 $10^{-6}$ 的各向同性平滑——这一操作在 3D 空间中进行，不破坏概率密度场的可解释性。实验表明，归一化 3DGS 的重建质量（SSIM、PSNR、深度 RMS）接近标准 3DGS，同时是唯一支持风险感知规划的 3DGS 变体。

\textbf{碰撞概率上界的推导。}给定归一化 3DGS 表示的场景，机器人被建模为球的并集——球形前向占据（Spherical Forward Occupancy, SFO），该占据通过多项式 Zonotope 算术对机器人正向运动学进行过近似，确保覆盖机器人沿参数化轨迹运动时的全部可能位形。碰撞概率的核心推导分为两步。第一步，利用 Markov 不等式和 Jensen 不等式，将球 $S$ 与场景碰撞的概率上界化为场景密度在球上的体积积分：$P(C(S)) \leq \frac{1}{\alpha}\bigl(1 - \exp\!\bigl(-\frac{1}{4\pi}\int_S \sigma(x)\, dV\bigr)\bigr)$，其中 $\alpha \in \mathbb{R}^+$ 为射线级碰撞风险阈值，$1/(4\pi)$ 因子来自将球面均匀分布的射线积分转换为体积积分的归一化。第二步，将球近似为 $L_\infty$ 球（轴对齐立方体），使归一化高斯密度的体积积分可分解为三个误差函数（error function）的乘积，从而获得闭式上界 $H(S) = \sum_{n=1}^{n_G} \eta'_n w_n \prod_{\ell=1}^{3} \sqrt{\frac{\pi \lambda'_{n,\ell}}{2}} \bigl[\mathrm{erf}\!\bigl(\frac{\rho - \mu'_{n,\ell}}{\sqrt{2\lambda'_{n,\ell}}}\bigr) - \mathrm{erf}\!\bigl(\frac{-\rho - \mu'_{n,\ell}}{\sqrt{2\lambda'_{n,\ell}}}\bigr)\bigr]$（$\mu'_n = R_n^T \mu_n$、$\lambda'_n$ 为旋转至轴对齐后的均值和特征值，$\eta'_n$ 为归一化常数）。最终碰撞概率约束为 $1 - \exp\!\bigl(-\frac{1}{4\pi} H(S_{j,i,m})\bigr) < \alpha \cdot \beta$（$\beta \in (0,1]$ 为配置级风险阈值），对 SFO 中每个球 $S_{j,i,m}$ 均需满足。实践中作者取 $\alpha = \beta$ 以简化双参数的相互依赖。

\textbf{滚动时域轨迹优化。}SPLANNING 以滚动时域方式在线求解轨迹优化：在离线阶段预定义参数化轨迹族 $q(t;k)$（满足初始条件匹配和末端制动到零速度的性质），运行时在每个规划周期（不超过 $t_p \leq 0.5$ 秒）选取参数 $k$，最小化目标代价（如接近目标位形）并满足关节位置/速度限制和上述概率碰撞约束。由于 $H(S_{j,i,m})$ 关于轨迹参数 $k$ 可微（通过前向运动学的多项式 Zonotope 导数传播），碰撞约束及其梯度可直接提供给数值优化器，无需自动微分。与 SAFER-Splat 和碰撞锥 CBF 面向自由飞行无人机（6-DoF 或双积分器动力学）不同，SPLANNING 面向串联操控器的关节空间规划，机器人几何通过多项式 Zonotope 过近似而非简单球体建模。

实验中，SPLANNING 在含 10--40 个随机立方体障碍物的 300 个场景中以 Kinova Gen3 7-DOF 操控器为平台，碰撞分类精度（precision-recall）在所有辐射场方法中最高，且通过调节 $\alpha,\beta$ 可在激进与保守之间精细控制约束行为；CATNIPS 精度最低（假阳性多导致过度保守），SplatNav 居中但缺乏连续风险调节能力。在 $\alpha = \beta = 0.025$--$0.05$ 时，SPLANNING 的成功率（60--69/100，10 障碍物；47--52/100，20 障碍物）超过所有使用真实障碍几何的基线（cuRobo、MPOT、TrajOpt 等），碰撞率接近零。平均每步规划时间约 0.20--0.32 秒（10--40 障碍物），约束评估时间 6.6--8.1 毫秒。在真实硬件实验中，系统在杂乱场景中实时生成无碰撞轨迹并成功执行操控任务，RRT 高层规划与 SPLANNING 轨迹优化在不同 GPU 上并行运行。

上述方法（SAFER-Splat、碰撞锥 CBF、SPLANNING）分别从欧氏距离、速度空间碰撞锥和概率上界三个角度解决了安全导航问题，但都隐含一个假设：感知与安全是可分离的两个阶段。冲突感知框架（2026.05）~\cite{2605_20566} 指出这一假设在主动感知场景中不成立，并系统性地解决了安全与感知之间的根本冲突。

\textbf{设计动机：安全--感知死锁。}在主动感知任务中，信息增益最大的视角通常位于未观测或不确定区域——而这些区域往往靠近障碍物。因此，安全策略要求机器人减速并远离障碍，而感知策略要求机器人加速朝向信息增益最大的方向（即靠近障碍），两者形成\textbf{安全--感知死锁（safety--perception deadlock）}。SAFER-Splat 仅处理安全过滤，SPLANNING 引入概率碰撞约束但未考虑感知目标，碰撞锥 CBF 关注速度空间预判但同样未整合主动感知——这些方法在纯导航中有效，但在需要同时建图和导航的场景中会遇到上述死锁。

\textbf{安全屏障函数：基于 AV@R 的碰撞风险。}每个高斯基元被建模为局部不确定障碍，具有均值 $\mu_i$ 和协方差 $\Sigma_i$。机器人到高斯基元 $g_i$ 的有符号距离 $d_i(x)$ 本身为高斯随机变量（因为高斯基元位置的不确定性），因此碰撞风险需要用概率度量而非确定论距离。作者选择 Average Value-at-Risk（AV@R$_\alpha$）——在置信水平 $\alpha \in (0,1)$ 下对距离分布的尾部期望——来量化每个高斯基元的碰撞风险 $\rho_i(x) := \text{AV@R}_\alpha(d_i(x))$。直觉上，$\alpha$ 越高，风险度量越保守，等效于放大高斯基元的感知障碍体积。安全要求机器人对所有高斯基元保持非负风险裕度 $\min_{g_i \in G} \rho_i(x) \geq 0$。由于 $\min$ 算子不可微，作者用 log-sum-exp 平滑近似得到全局可微的安全屏障 $h_s(x) = -\frac{1}{\beta}\log\bigl(\sum_{g_i \in G}\exp(-\beta \rho_i(x))\bigr)$，其中 $\beta > 0$ 控制近似精度——$\beta \to \infty$ 时 $h_s \to \min_{g_i} \rho_i$。安全类-$\mathcal{K}$ 函数设计为 $\kappa_s(h_s) = 2 a_{\max} \sqrt{h_s}$，将安全约束直接关联到机器人的物理制动能力：当安全裕度减小时，允许的减速时间也相应缩短。

\textbf{感知屏障函数：空间与角度信息渐升。}感知目标包含两个互补的几何约束。

\textbf{空间感知屏障} $h_p$ 要求相机朝向与局部空间信息渐升方向 $\nabla I|_p$ 对齐——即强制相机朝向信息增益梯度最大的方向，允许偏差不超过半角 $\delta_p \in (0, \pi/2)$。其类-$\mathcal{K}$ 函数 $\kappa_p(h_p) = (1 + e^{-\gamma \|\nabla I\|}) h_p$ 自适应地根据信息梯度幅值调整约束强度，信息丰富的区域约束更强。

\textbf{角度感知屏障} $h_\omega$ 引导机器人转向角度信息渐升方向 $d_\omega$，其类-$\mathcal{K}$ 函数参数 $k_\omega$ 根据机器人角速度限制选取。两类感知屏障的 Lie 导数可解析计算，嵌入统一的 CBF-QP 框架。

\textbf{冲突感知 QP：安全硬约束 + 感知松弛。}核心优化形式为最小化控制偏差 $\|u - u_\text{ref}\|^2$ 加感知松弛惩罚 $w_p \xi_p^q + w_\omega \xi_\omega^q$，满足安全 CBF 约束（无松弛）和感知 CBF 约束（通过 $\xi_p, \xi_\omega \geq 0$ 松弛）。作者通过 KKT 条件严格分析了冲突解决机制：由于安全约束无松弛变量，其拉格朗日乘子 $\lambda_s$ 可以无界增长，在冲突时压制感知需求。惩罚指数 $q$ 的选择决定了松弛行为模式：$q=1$（线性惩罚）产生稀疏激活——松弛变量一旦激活即被固定在饱和值 $w_p$，不再随违反幅度增长；$q=2$（二次惩罚）则使惩罚与违反幅度成正比——小冲突时感知约束轻微放松，严重冲突时感知约束被强力压制，形成比例反馈式的自适应行为。实验表明 $q=2$ 在轨迹全程产生更平滑的松弛，并带来略大的最小安全裕度。

\textbf{风险感知 EIG 与轨迹依赖掩码。}标准 NBV 选择在整个地图上最大化信息增益，但安全导航只关心机器人未来路径附近的区域。作者定义轨迹依赖掩码半径 $r_\text{mask}(p_k) = \gamma_1 \exp(-\gamma_2\, r_\alpha(p_k))$——碰撞风险越高的路径点，掩码半径越大——将信息获取聚焦于未来轨迹附近的球形区域并集 $\mathcal{M}_\tau = \bigcup_{p_k} B(p_k, r_\text{mask}(p_k))$ 内的高斯基元。这使 Fisher 信息矩阵只在掩码区域内计算，将感知计算复杂度从全地图降低到局部区域。

\textbf{性能表现。}在双积分器模型实验中，AV@R 安全屏障（$\alpha=0.75, 0.90$）与 SAFER-Splat 达到可比的安全率，同时保持更大的最小安全距离，计算速度提升 30--50 倍。在地面机器人（unicycle 动力学）实验中，冲突感知框架在安全区域主动朝向信息渐升方向，在狭窄区域自动松弛感知约束以保证安全，$q=2$ 的松弛模式在整个轨迹上表现出更平滑的安全--感知切换。该方法将安全与感知统一在同一 CBF-QP 中求解，无需外部调度或分层架构。

\paragraph{高斯-高斯碰撞检测}

FOCI（2025.05）~\cite{2505_08510} 提出基于高斯重叠积分（Field Overlap Collision Integral）的轨迹优化算法。该方法从 SAFER-Splat 的一个基本限制出发：SAFER-Splat 将机器人建模为球体，通过计算球心到高斯椭球表面的最小欧氏距离来判断碰撞，这种点-椭球距离度量无法捕捉机器人的完整几何形状（特别是非球形机器人的朝向信息），因此需要额外的安全余量补偿几何近似误差，限制了规划器在狭窄空间中的通过能力。FOCI 的核心思路是消除这种近似——将机器人自身也建模为一组 3D 高斯椭球（$M$ 个），与环境中的 $N$ 个高斯基元处于同一表示空间，使碰撞检测从"点 vs 椭球"提升为"高斯场 vs 高斯场"。这一设计选择的动机是：当环境和机器人都用高斯表示时，两者之间的碰撞度量可以自然地定义为重叠体积积分，而两个高斯函数的乘积在 $\mathbb{R}^3$ 上的积分具有闭式解析解。

具体实现上，环境密度场 $\rho_e(\mathbf{r}) = \sum_{i=1}^{N} \mathcal{N}(\mathbf{r}; \mu_i, \Sigma_i)$ 和机器人密度场 $\rho_r(\mathbf{r}) = \sum_{j=1}^{M} \mathcal{N}(\mathbf{r}; \mu_j^r, \Sigma_j^r)$ 之间的重叠体积积分为：

\begin{equation}
    \int_{\mathbb{R}^3} \rho_e(\mathbf{r}) \rho_r(\mathbf{r}) dV = \sum_{i=1}^{N} \sum_{j=1}^{M} \mathcal{N}(\mu_j^r; \mu_i, \Sigma_i + \Sigma_j^r),
\end{equation}

其中每一对 $(i,j)$ 的贡献等价于在联合协方差 $\Sigma_i + \Sigma_j^r$ 下评估的 Mahalanobis 距离的指数衰减 $\exp(-\frac{1}{2} d_M^2(\mu_j^r, \mu_i))$。该碰撞度量具有三个关键性质：（1）当机器人与环境无重叠时，由于高斯的指数衰减特性，度量趋近于零；（2）重叠程度越大，度量值越大；（3）关于机器人位姿完全可微，可直接提供优化梯度。由于机器人高斯的均值和协方差均可由基座位姿 $(p, \theta)$ 通过运动学链参数化，梯度可通过链式法则反向传播至轨迹参数。这使得 FOCI 能够实现\textbf{朝向感知规划}——优化器主动调整非球形机器人（如 ANYmal 足式机器人）的朝向以通过狭窄通道，而非简单增大安全余量。

\textbf{轨迹参数化与约束。}轨迹采用三次 B 样条参数化，优化变量为控制点 $q_i = [p_i^T, \theta_i]^T \in \mathbb{R}^4$（位置加航向角）。速度和加速度约束利用 B 样条的凸包性质：由于每个样条段位于其控制点的凸包内，只需约束控制点的速度/加速度范数即可保证整段轨迹满足约束。但标准 B 样条基底的凸包体积过大，导致约束过度保守。FOCI 采用 MINVO 基底替换标准基底——MINVO 基底在所有同阶样条基底中产生最小体积的凸包，从而以最小保守性约束速度 $v_{\max}$、加速度 $a_{\max}$ 以及角速度/角加速度界限。障碍物代价作为目标函数项（而非硬约束），因为重叠积分非归一化，难以确定统一的安全阈值——作者指出这与 NeRF 轨迹优化中面临的类似问题一致。初始轨迹通过在 3DGS 均值点构建的占据栅格上运行 A* 获得。

\textbf{GPU 并行求解。}由于碰撞度量的计算结构为 $N \times M \times K$ 的三重独立求和（$K$ 为轨迹离散化点数），FOCI 利用 NVIDIA Warp 框架在 GPU 上并行计算所有高斯对的重叠积分及其梯度，再由 CasADi 构建优化问题并通过 IPOPT 内点法求解。与多核 CPU（OpenMP）相比，GPU 实现提速 320 倍，使得在包含数十万环境高斯的场景中，完整规划（A* 初始猜测 + 轨迹优化）可在数秒内完成。在 ANYmal 硬件实验中，该方法在无人机航拍重建的真实 3DGS 场景中通过 ICP 定位后，成功规划出利用机器人非对称外形旋转穿过狭窄开口的轨迹——这正是因为高斯-高斯碰撞度量能精确评估全身几何与环境的交互。

\paragraph{可通行性感知导航}

上述方法均将 3DGS 场景中的高斯基元视为刚性障碍物（SAFER-Splat）或概率密度分布（SPLANNING），通过距离或碰撞概率度量来规避。然而在户外非结构化环境中，障碍物具有不同的几何形态、密度和不透明度——高草可以压过通行，树木则不可——这种二值化的障碍物建模将导致机器人在可通行植被前冻结。Splatblox（2025.11）~\cite{2511_18525} 针对此问题，将 SAFER-Splat 的二值碰撞检测范式增量扩展为可通行性感知建图与规划系统，核心思路是为每个高斯基元附加连续可通行性代价 $c_i \in [0,1]$，使 ESDF 不仅编码几何距离，还编码地形语义。

\textbf{可通行性估计与语义嵌入。}系统利用 CLIPSeg 提示式分割模型（以"混凝土""沙地""岩石""灌木"等地形描述词为提示）对 RGB 图像进行实时语义分割，按照 Elnoor et al.\ 的四级方案将分割结果映射为递增的可通行性代价：稳定地形（如混凝土，代价最低）、颗粒地形（如沙地）、不良支撑/岩石地形、以及高阻力/植被（代价最高）。分割输出转换为 2D 可通行性代价图后，每个 LiDAR 点通过外参标定 $T_{\text{lidar}} \in SE(3)$ 投影到相机坐标系 $p_i^{\text{cam}} = T_{\text{lidar}} \, p_i^{\text{lidar}}$，从代价图继承可通行性代价 $c_i$，然后被建模为各向异性高斯基元 $G_i(x) = \alpha_i \exp\!\bigl(-\frac{1}{2}(x - p_i^{\text{cam}})^T \Sigma_i^{-1}(x - p_i^{\text{cam}})\bigr)$，其中协方差 $\Sigma_i$ 由缩放 $S_i$ 和旋转 $R_i$ 因子分解为 $\Sigma_i = R_i S_i S_i^T R_i^T$，不透明度 $\alpha_i$ 受可通行性代价 $c_i$ 调制——代价越高（越不可通行）的高斯对体代价场的贡献越大。重要的是，这里的渲染目标不是颜色而是可通行性代价图：通过将高斯基元投影到图像平面并执行 $\alpha$-混合 加权合成 $C(x) = \sum_{i=1}^K c_i \alpha_i \prod_{j=1}^{i-1}(1 - \alpha_j)$，其中像素强度编码语义代价而非颜色。为确保渲染的可通行性图与真值对齐，高斯参数通过 L1 损失加 SSIM 结构相似性损失联合优化：$\mathcal{L} = \|I - I^{\text{gt}}\|_1 + (1 - \text{SSIM}(I, I^{\text{gt}}))$。系统以增量方式在线运行——随新传感器数据到达不断插入和更新高斯基元，维护一个持续演化的可通行性感知体代价场。

\textbf{从体代价场到 ESDF 的转换与融合。}聚合的可通行性场 $C(x) = \sum_i G_i(x)$ 被转换为距离场：$d_{\text{ESDF}}(x) = (1 - c) \cdot d_{\max}$，使完全可通行区域对应大的正距离值，障碍物区域趋近零值。为扩展规划视野，系统设计了一种分区融合策略：在前视相机覆盖区域 $F \subset \mathbb{R}^2$ 使用 GSplat-ESDF（包含语义推理），在侧视和后视区域回退至基于 Nvblox 的 LiDAR-ESDF（纯几何），即
\begin{equation}
    d_{\text{fused}}(x) = \begin{cases} d_{\text{GSplat}}(x), & x \in F \\ d_{\text{LiDAR}}(x), & x \notin F \end{cases}
\end{equation}

这一设计利用前视相机的语义推理能力处理最关键的前方区域（可通行性线索最为丰富），同时依靠 LiDAR 提供 360$^\circ$ 几何覆盖以支持远程规划。对障碍物边界统一施加 0.8 m 的膨胀核以增加安全余量。融合后的可通行性感知 ESDF 以 5 cm 分辨率光栅化到地面平面，生成供 Nav2 框架中 MPPI 控制器使用的连续代价图。

\textbf{边缘部署与性能。}该系统能够区分可通行植被（如高草）和刚性障碍物（如树木），在 Ghost Vision 60 四足机器人和 Clearpath Jackal UGV 上验证。在植被密集的户外场景中，Splatblox 比当时最优方法成功率高 50\% 以上、冻结率降低 40\%、路径缩短约 5\%、到达目标时间减少 13\%，在 100 米远程任务中达到 100\% 成功率（所有基线方法至少在一个场景中失败）。为适应机载边缘 GPU（RTX 3070，$<$6 GB VRAM）的内存约束，活跃高斯数量被限制在 100k 以内，优先保留视野内及附近区域的高斯，视野外的予以丢弃，整体系统以约 2 Hz 运行。消融实验表明，使用高斯基元的完整点云（而非仅使用高斯均值点）构建 ESDF 能一致提升成功率，更稠密的表示提供更精细的几何细节，使可通行性估计更准确；ESDF 网格深度从 1 m 增至 2 m 时成功率从 60\% 提升至 100\%，因为 2 Hz 的重建速率下更大的前视深度可补偿机器人运动产生的盲区。

\paragraph{统一表示导航}

Gaussian Splatting 作为统一表示（2025.05）~\cite{2505_11794} 的研究表明，高斯泼溅的稠密几何和光度信息对非结构化户外环境导航具有重要价值，语义信息可嵌入高斯地图以支持大规模任务驱动导航。与上述方法将 3DGS 用于单一功能（碰撞检测、轨迹优化或可通行性建图）不同，该方法展示了一个完整的自主导航系统：采用子图策略处理大规模场景（随机器人移动动态加载/卸载局部子图），通过 GPU 加速的穷举式碰撞检查替代传统光栅化方法，利用语言嵌入支持任务驱动规划。系统在 Clearpath Jackal 平台上验证了从地图构建到层次化规划的完整管线。

SAGE-3D（2025.10）~\cite{2510_21307} 提出将 3DGS 升级为可执行、语义和物理对齐的导航环境，引入 InteriorGS（1K 物体标注 3DGS 室内场景数据）和 SAGE-Bench（首个 3DGS 视觉-语言导航基准，2M VLN 数据），基线性能提升 31\%。与上述方法将 3DGS 用作被动地图表示不同，SAGE-3D 的核心范式转换是将 3DGS 从渲染资产提升为可执行仿真环境：$G + M + \Phi \rightarrow E_{exec}$，其中语义层 $M$ 提供物体级实例分割（通过 InteriorGS 的双校验标注覆盖 554k 物体实例、755 类别），物理层 $\Phi$ 提供碰撞检测和动力学交互。为实现物理可执行性，SAGE-3D 采用 3DGS-网格混合表示：3DGS 提供照片级渲染外观，CoACD 凸分解生成的三角网格提供碰撞几何和接触动力学——两者通过 USDA 场景组装为统一资产，3DGS 文件保持可见性，碰撞体设为不可见的刚体。这种设计避免了在 3DGS 上直接进行射线追踪碰撞检测的困难。仿真器同时暴露离散指令（转向/前进/停止）和连续控制（地面机器人速度 $(v,\omega)$，无人机 6-DoF 速度/姿态指令），支持 Unitree G1/Go2/H1 等多种平台。InteriorGS 数据集的场景通过艺术家创建网格渲染平均 3000 个相机视角，再用 GSplat 估计 3DGS 参数以确保遮挡密集的室内场景的重建质量。基线性能提升 31\%。

\subsection{无人机导航}
\label{ss:17_1c}

本节讨论基于 3DGS 的无人机导航方法，涵盖仿真训练到真实部署的全流程。

\subsubsection{无人机导航}

本路线介绍基于 3DGS 的无人机导航方法，涵盖仿真训练、域迁移与自主飞行等关键技术。无人机导航是 3DGS 在机器人领域最活跃的应用方向之一。无人机的高速运动、三维自由度和严格机载计算约束，使其对感知系统的效率和鲁棒性要求极高。3DGS 的实时渲染能力使构建照片级真实感仿真环境成为可能，极大降低了无人机策略学习的成本和风险。

\paragraph{零样本迁移}

SOUS VIDE（2024.12）~\cite{2412_16346} 修改了 3DGS 的渲染-策略学习管线——用 3DGS 场景重建替代手工仿真环境作为视觉数据源，用领域随机化的 MPC 蒸馏替代人类示范作为策略训练数据——3DGS 的标准渲染管线保持不变。该方法的核心思路是：从手持摄像机录制 2--3 分钟场景视频，用 Nerfstudio 训练 3DGS 模型 $\mathcal{GS}_\phi$，然后在该模型中嵌入无人机动力学，构建名为 FiGS（Flying in Gaussian Splats）的轻量仿真器，以 130 fps 渲染第一人称视角照片级真实感飞行图像。与手工构建数字孪生环境（通常耗时数周）相比，FiGS 从视频采集到可仿真只需约一小时。

\textbf{无人机动力学与渲染耦合。}
SOUS VIDE 采用 10 维半运动学状态向量描述无人机运动。直觉上，这个模型只关心无人机"在哪里、飞多快、朝哪个方向"，而不追踪旋翼转速等细节：

\begin{equation}
\bm{x} = [\bm{p}_\mathcal{W},\; \bm{v}_\mathcal{W},\; \bm{q}_{\mathcal{BW}}]^T, \quad
\bm{u} = [f_{th},\; \bm{\omega}_\mathcal{B}]^T
\end{equation}

其中 $\bm{p}_\mathcal{W}=(p_x, p_y, p_z)$ 为世界坐标系下的位置，$\bm{v}_\mathcal{W}$ 为速度，$\bm{q}_{\mathcal{BW}}$ 为四元数表示的机体姿态，$f_{th}$ 为归一化推力，$\bm{\omega}_\mathcal{B}=(\omega_x, \omega_y, \omega_z)$ 为机体角速率。控制输入 $f_{th}$ 和 $\bm{\omega}_\mathcal{B}$ 等效于专业飞手使用的 Acro 模式，是大多数飞控的标准低级输入——这意味着策略的输出可直接映射到真实飞控，无需额外的姿态控制层。

动力学方程为：

\begin{equation}
\dot{\bm{p}}_\mathcal{W} = \bm{v}_\mathcal{W}, \quad
\dot{\bm{v}}_\mathcal{W} = g\bm{z}_\mathcal{W} - k_{th}\frac{f_{th}}{m_{dr}}\bm{z}_\mathcal{B}, \quad
\dot{\bm{q}}_{\mathcal{BW}} = \tfrac{1}{2}\bm{W}(\bm{\omega}_\mathcal{B})\bm{q}_{\mathcal{BW}}
\end{equation}

其中 $g$ 为重力加速度，$\bm{z}_\mathcal{W}$ 和 $\bm{z}_\mathcal{B}$ 分别为世界系和机体系的 $z$ 轴单位向量，$\bm{W}(\bm{\omega}_\mathcal{B})$ 为四元数乘法矩阵。参数向量 $\bm{\theta}=(k_{th}, m_{dr})$ 存储推力系数和无人机质量。直觉上，速度变化由两个力决定：重力 $g\bm{z}_\mathcal{W}$ 向下拉、推力 $k_{th}f_{th}/m_{dr}$ 沿机体垂直轴推；姿态变化则由机体角速率通过四元数动力学驱动。由于控制输入是推力和角速率（而非力矩），模型省略了欧拉旋转方程，极大简化了动力学建模。

前向积分上述方程得到状态轨迹 $\mathbf{X}=\{\bm{x}_0, \dots, \bm{x}_K\}$ 后，通过机体-相机变换 $T_\mathcal{C}^\mathcal{B}$ 将位姿映射到相机坐标系，再从 3DGS 模型渲染图像序列 $\bm{\mathcal{I}}=\{\bm{I}_0, \dots, \bm{I}_K\}$。这一步就是标准 3DGS 的渲染调用 $\bm{I} = \mathcal{GS}_\phi(\bm{p}, \bm{q})$，只不过相机位姿来自无人机动力学积分而非人工指定。仿真器使用 ACADOS 轨迹优化器进行高效前向积分。

\textbf{MPC 专家与领域随机化数据合成。}
传统无人机策略学习的瓶颈在于数据：人类飞手示范难以大规模获取，而仿真数据又存在域间隙。SOUS VIDE 的解决方案是让一个拥有特权信息（地面真值状态和动力学参数）的 MPC 专家在 FiGS 中生成大规模示范数据。

具体地，给定期望轨迹 $(\mathbf{X}^d, \mathbf{U}^d)$，MPC 在 $N$ 步滚动时域内求解：

\begin{equation}
\min_{\bm{u}} \sum_{k=0}^{N-1}(\delta\bm{x}_k^T Q_k \delta\bm{x}_k + \delta\bm{u}_k^T R_k \delta\bm{u}_k) + \delta\bm{x}_N^T Q_N \delta\bm{x}_N
\end{equation}
\begin{equation*}
\text{s.t.} \quad \bm{x}^s_{k+1} = f_d(\bm{x}^s_k, \bm{u}^s_k), \quad g_c(\bm{u}^s_k) \leq \bm{0}
\end{equation*}

其中 $\delta\bm{x}_k = \bm{x}^s_k - \bm{x}^d_k$ 为当前状态与期望轨迹的偏差，$(Q_k, R_k)$ 为阶段权重，$Q_N$ 为终端权重。每次 rollout 开始时，从均匀分布 $\bm{\theta}_s \sim U(\bm{\theta}_{\min}, \bm{\theta}_{\max})$ 采样无人机参数，并从期望轨迹附近扰动初始状态 $\bm{x}^s_0 \sim (\bm{x}^d_i - \Delta\bm{x},\; \bm{x}^d_i + \Delta\bm{x})$。这种领域随机化迫使策略学会适应不同的质量、推力系数和初始条件。最终生成 100k--300k 图像/状态-动作对，涵盖视觉观察、可观测状态和控制指令。

\textbf{SV-Net 策略架构。}
SV-Net（Sous Vide Network）是运行在机载 Orin Nano 上、以 20 Hz 输出推力和机体角速率指令的轻量策略网络。其输入仅为机载传感器数据：机载相机图像，以及由 EKF（扩展卡尔曼滤波器）融合 IMU（惯性测量单元，Inertial Measurement Unit）、下视 ToF（Time of Flight，飞行时间）传感器和光流传感器估计的高度 $p_z$、速度 $\bm{v}_\mathcal{W}$ 和姿态 $\bm{q}_{\mathcal{BW}}$。值得注意的是，策略不观测水平位置 $(p_x, p_y)$，导航完全依赖视觉信息。

SV-Net 由三个组件构成：\textbf{特征提取器}使用 SqueezeNet 处理彩色图像，将输出与估计高度和姿态拼接后通过 MLP 生成位姿感知特征向量；\textbf{历史网络}采用 RMA（Rapid Motor Adaptation）技术，利用滑动时间窗口内的可观测状态历史数据生成潜在编码，实时适应电池衰减、旋翼下洗和阵风等变化条件；\textbf{指令网络}将特征提取器和历史网络的输出与可观测状态及目标向量 $\mathcal{O}^k$（编码位置变化、初末速度、初末姿态和轨迹总时间）融合，通过 MLP 输出机体角速率指令。

\textbf{RMA 自适应模块的设计动机与原理。}
标准 3DGS 仿真管线的一个关键问题是：仿真中的动力学参数与真实世界存在差异（电池衰减导致质量/推力变化、阵风引入外力、旋翼下洗产生气动力），固定参数的策略无法应对这些运行时变化。RMA 模块的核心思想是：如果将推力系数和质量视为可在线调整的变量 $c = k_{th}/m$，那么它们可以补偿广泛的不确定性。

定义等价推力-质量比 $c = k_{th}/m$ 和外部力 $\bm{f}_{add}$（包括气动力、地面效应等），则等价参数 $\hat{c}$ 满足：

\begin{equation}
g\bm{z}_\mathcal{W} - \hat{c} f_{th} \bm{z}_\mathcal{B} = g\bm{z}_\mathcal{W} - c f_{th} \bm{z}_\mathcal{B} + \bm{f}_{add}
\end{equation}

直觉上，这个公式说的是：如果给 $c$ 换一个"等价值" $\hat{c}$，使得修改后的推力项恰好等于原始推力加所有未建模外力，那么简化的动力学模型就能准确描述真实飞行。对 $\hat{c}$ 取最小二乘估计：

\begin{equation}
\hat{c} \approx c - \frac{\bm{z}_\mathcal{B}^T \bm{f}_{add}}{f_{th}}
\end{equation}

当外力 $\bm{f}_{add}$ 与机体 $z$ 轴 $\bm{z}_\mathcal{B}$ 近似共线时（悬停附近的重力变化、高速飞行时的气动阻力均满足此条件），$\hat{c}$ 能准确捕获外部力的影响。历史网络通过滑动窗口内的可观测状态历史数据在线估计这个潜在编码，无需显式知道 $\bm{f}_{add}$ 是什么。

\textbf{两阶段训练。}
SV-Net 的训练分两个阶段：第一阶段，用示范数据中的 $(\bm{X}^s, \bm{U}^s)$ 计算历史数据，训练历史网络估计无人机参数 $\bm{\theta}_s$；第二阶段，冻结历史网络，端到端训练剩余组件（包括 SqueezeNet 图像编码器），以可观测状态和图像预测机体角速率指令。实验发现，使用历史网络的倒数第二层（而非显式参数估计 $\hat{\bm{\theta}}$）作为指令网络的输入效果最佳。

\textbf{实验验证。}
在 105 次硬件实验中，SOUS VIDE 策略展现出强鲁棒性：对 30\% 质量变化、40 m/s 阵风、60\% 亮度变化、场景物体移动或移除、以及人员穿越视野等干扰均能稳定飞行。消融实验中，完整 SV-Net 的轨迹跟踪误差为 0.17 m、路径接近度（30 cm 内）96\%、碰撞率 0，相比无 RMA 模块的变体性能提升超过一倍。在消融中，真实世界的等价参数估计 $\hat{c}$ 的均值为 $\mu_{\hat{c}}=6.05$（真值 $c=6.03$），标准差 $\sigma_{\hat{c}}=0.25$，验证了 RMA 模块在线适应的有效性。

SOUS VIDE 的模仿学习范式虽然有效，但面临三个局限：需要大量高质量专家飞行数据、训练时间长、且难以泛化至训练中未见过的新任务实例。GRaD-Nav（2025.03）~\cite{2503_03984} 从策略训练范式层面改变了这一思路——不再依赖专家示范，而是用可微深度强化学习（Differentiable Deep Reinforcement Learning, DDRL）直接在 3DGS 环境中端到端训练策略，通过可微分动力学仿真器与 3DGS 渲染的梯度回传优化策略参数，从而同时提升采样效率和任务泛化能力。

\textbf{混合可微仿真器。}
GRaD-Nav 构建了并行化可微四旋翼动力学仿真器（PyTorch 实现），接收机体角速率 $\omega^d_t \in \mathbb{R}^3$ 和归一化推力 $c_t \in [0,1]$ 为控制输入，输出下一状态 $s_{t+1} = (p, q, v, \omega, a) \in \mathbb{R}^{16}$（位置、四元数姿态、线速度、角速度和加速度）。动力学完整建模了 PD 姿态控制器、四元数积分、电机延迟和空气阻力，且 $\partial s_{t+1}/\partial s_t$ 和 $\partial s_{t+1}/\partial(\omega^d_t, c_t)$ 通过 PyTorch 自动微分计算。每步仿真器向预训练 3DGS 模型发送批量化无人机位姿 $T = (p, q) \in \mathbb{R}^7$，后者渲染第一人称 RGB 图像作为策略视觉输入。同时，3DGS 点云用于设置奖励航点、规划参考轨迹和碰撞检测。在 128 架无人机并行仿真（房间约 100 m$^2$，3DGS 模型约 1.5 GB）下，仿真时间步 0.05 s（20 Hz），每步墙钟时间约 0.7 s（i9-13900K + RTX 4090），其中可微动力学占 33.5\%、3DGS 渲染占 55.7\%、碰撞检测占 10.8\%。

\textbf{策略架构与 CENet 上下文估计。}
策略架构包含四个网络。

\textbf{感知网络}使用预训练 SqueezeNet 处理 RGB 图像，通过全连接层将 512 维特征压缩至 24 维视觉嵌入 $e_t \in \mathbb{R}^{24}$。

\textbf{策略网络} $\pi_\theta(a_{t+1}|o_t, z_t, e_t)$ 为三层 MLP（512/256/128 神经元），输入为观测 $o_t = [h_t, q_t, v_t, a_t, a_{t-1}]^T \in \mathbb{R}^{16}$（高度、四元数、线速度、当前/上一步动作——注意不需要水平位置，导航完全依赖视觉）、CENet 潜在编码 $z_t \in \mathbb{R}^{16}$ 和视觉嵌入 $e_t$，输出机体角速率和归一化推力。

\textbf{价值网络} $V_\phi(s_t, z_t)$ 为三层 MLP，使用特权状态 $s_t \in \mathbb{R}^{43}$（在观测基础上额外包含位置 $p_t \in \mathbb{R}^3$ 和 3DGS 深度图池化后的深度先验 $d_t \in \mathbb{R}^{24}$），仅在训练时使用。

\textbf{CENet} 基于 $\beta$-VAE 架构，以最近 5 步历史观测 $o_t^5$ 为输入，编码周围环境的空间-障碍物关系为潜在向量 $z_t$，训练目标包含观测重构损失和 KL 正则化项，使策略在运行时适应不同环境配置——这是 GRaD-Nav 实现任务内泛化的关键模块。

\textbf{训练算法与奖励设计。}
GRaD-Nav 采用 Short-Horizon Actor-Critic（SHAC）算法：将完整任务时间线划分为若干短子窗口，在每个学习回合中并行采样 $N$ 条短视界 $h \ll H$ 的轨迹，利用可微仿真器在子窗口内回传策略梯度，同时训练策略网络和价值网络。这种设计既利用了可微仿真的梯度信息提升采样效率（相比 PPO（Proximal Policy Optimization，近端策略优化）和 BPTT），又通过截断视界避免了梯度消失/爆炸问题。奖励函数分为安全控制（7 项：存活 $+8.0$、线速度 $-0.5$、姿态偏差 $-0.5$、高度偏差 $-2.0$、动作惩罚 $-1.0$、动作变化率 $-1.0$、平滑度 $-1.0$）和高效导航（5 项：偏航对齐 $+0.25$、航点接近 $r_{wp} = \exp(-\|p - w_{next}\|^2)$，权重 $+2.0$、障碍物保持 $+1.0$、越界惩罚 $-2.0$、参考轨迹跟踪 $-2.0$），整体设计遵循简洁原则以增强跨环境迁移性。领域随机化对质量、最大推力、惯性、电机延迟和阻力系数进行随机采样，增强策略鲁棒性。

\textbf{课程训练与任务内泛化。}
为训练可泛化至未见门位置的导航策略，GRaD-Nav 采用课程训练：在 3 个不同 3DGS 环境间轮换训练（各 100 轮 $\times$ 5 次），每个环境设置独特的门位置，每次环境切换时重置学习率。训练完成后，策略无需任何环境相关信息（地图名、门位置、参考轨迹），仅凭第一人称 RGB 图像即可``回忆''在哪个环境中并引导无人机飞过正确位置的门。

\textbf{实验验证。}
与 PPO 和 BPTT 的对比实验表明，非可微 RL（PPO）难以训练出满意的端到端视觉导航策略，而 GRaD-Nav 相比 BPTT 实现超过 300\% 的采样效率提升，仅用 20\% 训练时间即获得更优策略。消融实验在两条轨迹（中门和右门）上验证各模块必要性：完整方法在两条轨迹上分别达到 8/10 和 7/10 仿真成功率；去掉 CENet 后成功率降至 4/10 和 5/10，且所有失败均为障碍物碰撞，说明 CENet 对环境理解至关重要；去掉视觉输入或替换为深度图则成功率为 0。在真实硬件实验中（Pixracer 飞控 + Jetson Orin Nano + Realsense D435，策略推理频率 30 Hz），课程训练后的泛化策略在三个未见门位置（左/中/右）的仿真成功率为 10/10、10/10、9/10，真实世界零样本迁移成功率为 7/10、7/10、6/10。CENet 的 PCA 分析显示，潜在向量在穿越门阶段呈现更紧凑的分布，验证了其环境理解能力。去掉 CENet 后真实世界成功率骤降至 0/10、2/10、1/10，进一步证明 CENet 对 sim-to-real 迁移的核心作用。

\paragraph{视觉-语言导航}

SINGER（2025.09）~\cite{2509_18610}（Semantic In-situ Navigation and Guidance for Embodied Robots）直接构建在 SOUS VIDE 的 FiGS 仿真器之上，解决后者的一个核心局限：FiGS 缺乏语义知识，策略只能在训练时见过的环境中沿固定轨迹飞行，无法响应"飞向红色沙发"这类开放世界指令。SINGER 的核心增量在于将自然语言语义引入整个仿真-训练-部署管线，使策略以语义而非像素外观为纽带实现跨环境泛化。

具体而言，SINGER 对 FiGS 的改造分三个层面。

\textbf{语义 3DGS。}
SINGER 将 CLIP 图像嵌入蒸馏到 3DGS 场景中，通过多分辨率哈希网格加 MLP 参数化语义场 $f: \mathbb{R}^3 \to \mathbb{R}^l$，使场景中每个 3D 点可查询其语义嵌入。用户输入的语言查询经 CLIP 文本编码后，在 3DGS 稀疏点云中检索语义匹配区域，定位目标物体的 3D 质心。这一语义场与 3DGS 联合训练，直接为后续的轨迹规划和碰撞检测提供语义锚点。

\textbf{语言条件轨迹生成。}
以每个语义目标质心为根构建 RRT* 树，在目标高度的水平面上向环境边界扩展，生成空间覆盖广泛的无碰撞轨迹集合。稀疏点云的包围气泡充当碰撞禁区，目标质心附近额外设置目标区域（goal-region）以引导轨迹从环境中心方向接近目标而非从上方飞越。每个 RRT* 节点对应的粗轨迹经三次样条平滑后，由 SOUS VIDE 同款的鲁棒 MPC 专家在 FiGS 中倒序（叶到根）飞行执行，同时进行 30\% 范围内的质量/推力随机化和位姿扰动，利用管式 MPC 思想构建鲁棒数据集。

\textbf{语义热图策略训练。}
与 SOUS VIDE 训练策略处理原始 RGB 图像不同，SINGER 用 CLIPSeg 将图像转换为 3 通道语义热图（目标物体区域呈亮红色，语义无关区域呈深蓝色），策略网络在这个语义-空间对齐的图像空间上训练。架构沿用 SV-Net（历史网络 + 特征提取器 + 指令网络）和两阶段训练流程，但特征提取器接收 CLIPSeg 处理后的语义热图而非原始 RGB。部署时加入语义归一化技巧：将当前 CLIPSeg 相似度分数对飞行中历史最高分归一化，赋予策略"记忆"——只引导自己飞向比已见区域更语义相似的方向。CLIPSeg（基于 ViT-B/16，86M 参数）通过 ONNX 追踪优化，在 Jetson Orin Nano 上从 3 Hz 提升至 12 Hz，以异步方式与 20 Hz 策略推理并行运行。在 5 个 3DGS 环境（室内+室外）中生成约 1650 条轨迹、15 种语义查询，共计 907,440 个观测-动作样本。

硬件实验中，SINGER 在训练时从未见过的环境和语义查询上实现零样本迁移：与基于 CLIPSeg + PD 控制的基线相比，到达目标 1 m 内的成功率提升 23.33\%，碰撞率降低 10\%，目标保持在视野内的比例提升 16.67\%，且全程仅依赖机载传感和计算（无需运动捕捉或外部位姿估计）。

GRaD-Nav 虽然通过可微强化学习实现了高效的端到端策略训练，但策略只能在课程训练中见过的固定门位置执行导航——它无法理解"从左侧通过门然后飞向梯子"这类自然语言指令，也无法在飞行中切换任务目标。GRaD-Nav++（2025.06）~\cite{2506_14009} 正是从这一局限出发，将 GRaD-Nav 的可微强化学习框架扩展为\textbf{轻量级 VLA（Vision-Language-Action）框架}，使策略能够直接从自然语言指令推理并执行导航动作。与 GRaD-Nav 使用 SqueezeNet（0.5M 参数）编码原始 RGB 图像不同，GRaD-Nav++ 的核心差异在于引入\textbf{视觉-语言模型（VLM）}替代独立 CNN 视觉感知：预训练 CLIP-ViT-Base-Patch32（150M 参数）同时编码自然语言指令和第一人称 RGB 图像，通过可训练线性层融合并下采样为 512 维特征向量 $e_t$，为策略提供语义决策能力。

这意味着策略不再使用独立卷积网络处理 RGBD，而是直接将 VLM 特征向量 $e_t$ 作为高层视觉语义表示，仅用最小池化深度特征 $d_t \in \mathbb{R}^{32}$ 进行低层避障——这一设计选择的主要动机是优化机载推理速度，避免运行独立的 CNN 视觉感知模块。与 GRaD-Nav 的单策略网络相比，GRaD-Nav++ 采用\textbf{混合专家（MoE）策略网络}（源自 MENTOR 架构）：2 个专家子网络，路由器每次激活 top-$k=2$ 个专家，在保持稀疏计算的同时确保所有专家被利用，并通过动态路由机制避免跨环境/跨任务训练时的灾难性遗忘——这是 GRaD-Nav 多环境课程训练中未曾解决的问题。每个专家是多头部网络：每个头是三层 MLP（512/256/128），将处理后的 VLM 特征与 72 维观测 $o_t = [d_t, z_t, r_t]^T$ 融合后，经另一个三层 MLP（256/128/64）输出 4 维动作（机体角速率 $\omega^d_t \in \mathbb{R}^3$ + 归一化推力 $c_t \in \mathbb{R}$）。

观测向量 $o_t$ 的 72 维构成包括：最小池化深度图 $d_t \in \mathbb{R}^{32}$（用于碰撞避免）、CENet 潜在编码 $z_t \in \mathbb{R}^{24}$（与 GRaD-Nav 相同的 $\beta$-VAE 上下文估计器，编码周围环境的空间-障碍物关系）、以及原始动力学观测 $r_t$（高度、四元数、线速度、当前/前一步动作）。与 GRaD-Nav 的观测维度（$\mathbb{R}^{16}$，纯动力学量）相比，GRaD-Nav++ 的 72 维观测大幅扩展了感知带宽，这主要得益于深度特征（32 维）和 CENet 编码（24 维，较 GRaD-Nav 的 16 维也有所扩充）的加入。价值网络 $V_\phi(s_t, e_t)$ 共享相同的 MoE 结构但访问特权观测 $s_t \in \mathbb{R}^{74}$（额外包含机体 $x$-$y$ 位置）。为适配机载 Orin Nano 的计算约束，CLIP 以异步方式运行（每 10 步更新一次特征向量，运行时内存占用约 6 GB），策略以约 30 ms/步的延迟运行，整体控制频率 25 Hz。

奖励函数沿用 GRaD-Nav 的设计但引入任务随机化。奖励项包括：航点奖励 $r_{wp} = \exp(-\|p - w_{next}\|^2)$（权重 2.0）、参考轨迹跟踪惩罚（权重 $-2.0$）、障碍物惩罚 $d_{obs}$（最近障碍物距离 $< 0.5$ m 时激活，权重 1.0）、以及安全控制奖励（生存、速度、姿态、高度、动作平滑性等，总权重约 $-3.5$）。每次环境重置时从 12 个两阶段任务中均匀采样——每个任务要求无人机先从 \{through, left, right, above\} 中选择正确方向通过门，再飞向 \{ladder, cart, monitor\} 中的指定目标物体。

实验分两组。

\textbf{多任务泛化实验}在单一 3DGS 环境中训练 1000 epochs（约 8 小时，RTX 4090），训练 8 个任务、零样本评估 4 个未见任务：仿真训练任务成功率 83\%、未见任务 75\%；真实硬件训练任务 67\%、未见任务 50\%。消融实验表明，去除 VLM（替换为 one-hot 编码）后仿真成功率从 83\% 骤降至约 4\%，证明 VLM 不仅充当文本编码器，更实现了指令到视觉上下文的语义接地；标准 PPO 在 $10^7$ 步内完全无法学到有效策略，验证了 DiffRL 的样本效率优势。还测试了退化 3DGS 模型（仅用 $1/5$ 训练步数，渲染质量显著下降）下的鲁棒性，仿真成功率降至约 58\%，但策略仍能维持合理性能。

\textbf{多环境适应实验}在两个不同 3DGS 环境间交替训练（各 1500 epochs，每 100 epochs 切换），每个环境有不同门位置和干扰物体。4 个方向任务的平均仿真成功率 81\%、真实硬件 67\%。

与单专家（SE）和大网络（LN，参数量与 MoE 匹配但无路由）的消融对比表明：在第二个环境（mid-gate）中，SE 和 LN 因灾难性遗忘成功率分别骤降至 25\%/33\%（仿真）和 8\%/25\%（真实），而 MoE 保持 70\%/67\%——路由器通过动态分配专家避免了跨环境知识的覆盖。MoE 可扩展性实验进一步测试了 4/8 专家变体：仿真性能随专家数微增，但真实硬件上 8 专家因控制频率降至约 10 Hz 而性能骤降，表明 MoE 扩展受机载计算约束。此外，论文设计了\textbf{任务移位实验}验证 VLM 的在线适应能力：在无人机执行长视界任务的第 100 步，将指令从"从左侧通过门然后飞向梯子"突然切换为"从右侧通过门然后停在显示器上方"——VLM 文本-视觉余弦相似度在指令切换后立即显著下降，随后随无人机调整行为逐步回升，且策略成功完成了新任务。

这一实验直接证明了 GRaD-Nav++ 的 VLM 不仅编码静态指令，更能动态地将新指令与当前视觉上下文重新接地，实现飞行中的任务切换——这是 GRaD-Nav 完全不具备的能力。

\paragraph{空中导航基准}

VLA-AN（2025.12）~\cite{2512_15258} 是首个面向无人机自主导航的大规模 VLA 框架。GRaD-Nav++ 以 CLIP-ViT-Base（150M 参数）+ 2 专家 MoE 策略网络构建轻量 VLA，其语义理解能力受限于 CLIP 的对比学习范式——仅能对齐指令与视觉的全局嵌入，无法进行链式推理或场景级空间推断。VLA-AN 的核心差异在于将策略引擎从\textbf{轻量 VLM}全面升级为\textbf{完整多模态大语言模型（MLLM）}——基于 Qwen2-VL 骨干（7B/3B/2B 三种规模，最大模型参数量约为 GRaD-Nav++ CLIP 编码器的 47 倍），将 Vision Transformer 编码的视觉特征经 MLP 投影对齐到 LLM 潜在空间，由 LLM 自回归解码器直接从视觉-语言输入进行多步推理并输出语义控制动作。这一架构升级的动机在于：无人机在复杂环境中需要理解``找到从近到远第三棵树''这类需要距离排序+计数组合的复合指令，CLIP 的单次嵌入对比无法完成此类链式空间推理，而 MLLM 的自回归推理链可以逐步分解子目标。同时，与 SOUS VIDE 和 GRaD-Nav 输出机体角速率+推力等低级控制指令（4 维连续动作）不同，VLA-AN 输出\textbf{3D 航点和偏航角}（4 维导航目标），将策略输出从电机级命令抽象为语义级导航目标——这一抽象使模型能从异构机器人数据集（不同配置无人机乃至人形/四足机器人）学习导航知识，大幅提升训练效率并避免 OpenVLA 和 $\pi_0$ 在任务精调时的灾难性遗忘问题。该方法包含四个核心创新。

第一，\textbf{高保真混合数据集}：GRaD-Nav++ 仅在 1--2 个 3DGS 环境中通过 RL 在线采样生成训练数据，数据规模受限于仿真吞吐量（每步约 0.7 s 墙钟时间）。VLA-AN 转向离线数据集构建范式，集成 3DGS 照片级真实感渲染（保留真实世界光照和几何细节）与 Unity 网格场景（支持光照/遮挡/目标物体的灵活编辑）和真实数据，按比例混合以平衡视觉一致性和环境多样性。3DGS 场景从室内外视频流（办公室、走廊、公园、街道）构建后导入 Unity，利用梯度规划器生成满足无人机运动学约束的全空间碰撞自由轨迹，四视角同步（前/后/左/右）以 10 Hz 采集 RGB、深度和 4D 位姿数据。消融实验定量验证了 3DGS 数据的 sim-to-real 桥接效果：仅用 3DGS 数据训练的模型在已见数据集上物体导航成功率 98.6\%、空间定位 95.2\%，与仅用真实数据训练（98.8\%/95.0\%）几乎持平；而仅用网格数据的模型在未见数据集上物体导航成功率骤降至 70.3\%、长程导航仅 61.8\%（3DGS 数据对应 95.8\%/82.9\%），证明 3DGS 的连续几何表示和一致光照对真实世界泛化至关重要。最终构建的混合数据集包含超过 100K 条导航轨迹和 100 万+ 多模态样本，数据规模远超 GRaD-Nav++ 的在线 RL 采样量。

第二，\textbf{渐进式三阶段训练}：GRaD-Nav++ 使用单一可微 RL 算法（SHAC）在一个固定任务上端到端训练整个策略网络，无法赋予模型通用的视觉-语言推理能力。VLA-AN 针对大模型特性设计渐进式三阶段训练，每个阶段聚焦不同能力层次。阶段 I 进行全参数 SFT，在大规模混合数据集（VQA、空间定位、推理、STEM 任务）上增强视觉理解和空间推理，引入多帧多视角图像推理和跨动作序列的时间一致性建模，使模型从物体级感知扩展到场景级空间推断；消融实验显示仅阶段 I 即使场景推理成功率从基线的 77.4\% 提升至 87.2\%、空间定位从 68.2\% 跃升至 94.5\%，但物体导航仅 38.1\%，表明通用推理能力不足以支撑实际导航。阶段 II 注入高质量无人机导航数据并按比例混合精选 VQA 推理数据，显式教授 3D 航点规划、偏航预测和动态任务重规划；仅阶段 II 即可将物体导航提升至 95.3\%，但场景推理仍停滞在 75.6\%。阶段 III 用 GRPO 强化学习精调（RFT），奖励函数覆盖视觉推理（严格答案匹配）和空间定位（预测与真值的 IoU），对每条指令采样多个响应并在组内归一化计算优势估计，纠正前期失败模式以提升决策一致性。完整三阶段模型在四个维度上达到最优或次优：场景推理 85.5\%、空间定位 94.6\%、物体导航 98.1\%、长程导航 85.7\%，其中物体导航比无训练基线（34.2\%）提升 63.9 个百分点，验证了渐进式训练的必要性。

第三，\textbf{时间比较与闭环重规划}：GRaD-Nav++ 的 CLIP+MoE 策略是纯前馈的单帧推理——每步独立处理当前观测，不具备显式的任务进度追踪能力。VLA-AN 引入时间比较模块 $f_{\text{cmp}}$，通过对比初始帧 $I_0^{rgb}$ 和当前帧 $I_t^{rgb}$ 与语言指令 $L$ 判断任务完成状态 $s_t \in \{0,1\}$，赋予模型``记住起点、感知进度''的时序推理能力。无人机系统建模为状态转移过程 $\mathcal{S}=\{\text{IDLE},\text{NAVIGATING},\text{REPLANNING},\text{TASK\_COMPLETE}\}$，当目标尚未到达时自动触发全局重规划更新导航目标，实现闭环自主协调。这一机制使模型在目标跟踪样本仅占训练集不到 1\% 的情况下仍达到 82.0\% 的跟踪成功率——论文推测这是因为导航导向训练隐式促进了任务重规划能力，使模型能通过聚合多帧时序观测推断目标状态变化。

第四，\textbf{实时安全动作模块}：GRaD-Nav++ 直接输出机体角速率+推力等 4 维低级控制指令，碰撞避免完全依赖奖励函数中的障碍物惩罚项在训练中隐式学习。VLA-AN 输出的是 3D 航点+偏航角等语义级导航目标，碰撞避免由独立的安全动作模块显式保障：当 VLA 生成的动作轨迹与障碍物相交时，从深度图提取局部几何线索构建表面锚点和排斥梯度力，即时调整轨迹并重新估计时间分配，在安全约束下重建参考轨迹——这避免了大型生成模型（如 $\pi_0$ 的 0.3B flow-matching 动作专家）的推理延迟和随机生成噪声带来的碰撞风险。

部署在 NVIDIA Jetson Orin NX 16GB（约 100 TOPS 算力，总增重约 80g）上，采用 AWQ 量化将模型压缩至可部署规模，通过 Flash-Attention、FFN-Norm 算子融合、KV-cache 预加载和 CUDA graph 管线调度实现系统性加速：ViT 阶段推理时间从 2350ms 降至 120ms（贡献最大绝对加速），投影+LLM 预填充从 550ms 降至 130ms，解码从 1000ms 降至 230ms，总推理时间从约 4100ms 降至 494ms，实现 8.3 倍吞吐量提升（7.1--9.4 倍视模型规模），在 2--3 Hz 实时推理频率下支持敏捷自主飞行。三种模型规模的浮点运算量分别为 7B-AWQ 约 1815.8 GFLOPs（每 token 解码 0.110 s）、3B-AWQ 约 946.7 GFLOPs（0.051 s）、2B-AWQ 约 563.3 GFLOPs（0.032 s），用户可根据机载算力限制灵活选择。

部署在两种不同无人机平台上，配备 Intel RealSense 相机，异步图像传输与连续动作执行结合实现高频实时动作推理。在仿真基准测试中，VLA-AN 在物体导航任务上平均成功率超过 98\%，在空间定位、场景推理和长程导航等复杂任务上全面超越 OpenVLA、$\pi_0$ 和 Groot N1.5；最高单任务成功率达 98.1\%，平均成功率超过 90\%。值得注意的是，OpenVLA 和 $\pi_0$ 在任务精调时出现严重的灾难性遗忘——视觉-语言对齐能力显著退化，而 VLA-AN 的渐进式训练通过保留推理数据混合比例有效缓解了这一问题。

\paragraph{生成式场景导航基准}

FlyMirage（2026.05）~\cite{2605_19600} 利用 LLM 作为环境设计器、生成式世界模型实例化高保真 3DGS 场景，构建大规模多样化航拍视觉-语言导航数据集。与 VLA-AN 集成 3DGS + Unity + 真实数据的混合数据集（仍依赖人工场景构建和真实采集）不同，FlyMirage 的核心差异在于\textbf{全自动化场景生成——从场景设计到轨迹采集无需任何人工干预}。与 SOUS VIDE 和 SINGER 从真实视频重建 3DGS 的路径也不同，FlyMirage 利用生成式世界模型直接从文图提示创建 3DGS 世界，使场景类型可覆盖难以采集的危险或受限空间（如辐射区域、化学设施）。

\textbf{场景生成。}
FlyMirage 利用层次化场景类别分类法约束 GPT-5.4/Gemini 3.1 的采样分布，先生成文本场景描述再用 GPT Images 2.0 生成场景图像，最终通过 Marble 1.1 Plus 从文图对生成 3DGS 场景。层次化分类的设计动机是：实验发现直接 prompting LLM 生成随机场景描述时，模型存在显著的采样偏差（例如 GPT-5.4 过度倾向生成医疗相关描述），通过先从分类法中随机选择子类别再让 LLM 描述具体场景，将场景类型选择与 LLM 的生成偏好解耦，实现更均衡的场景多样性。生成的图像相比纯文本描述提供了额外的空间布局、材质外观和光照条件等视觉线索，作为文本的互补输入提升了 3DGS 场景的几何和纹理质量。

\textbf{迭代式场景标注。}
场景标注采用迭代式多视角探测策略。由于 Marble 仅输出 3DGS 表示（无训练图像或物体语义信息），FlyMirage 先以场景中心 $\mathbf{c}$ 为球心、$r_{\mathrm{orb}}=0.1$ m 为半径的轨道上放置相机（多仰角、完整偏航扫描），渲染 RGB 和深度图后用 Boxer 进行开放词汇 3D 边界框检测。核心问题在于 Boxer 对远处物体的检测精度急剧下降——因此引入距离感知的目标选择算法（Algorithm 1）：设定两个距离阈值 $d_{\mathrm{th1}}=3.0$ m 和 $d_{\mathrm{th2}}=4.0$ m，优先选择距离 $d_i > d_{\mathrm{th2}}$ 的远处候选目标（最多 $N_t=5$ 个），并按距离 $d_{\mathrm{prune}}=d_{\mathrm{th1}}$ 和方位角阈值 $\theta_{\mathrm{th}}=35^\circ$ 剪除空间上过近或方向重叠的候选以保证目标空间多样性。对选中的远处目标，利用体素 A* 搜索在膨胀占据栅格上规划碰撞自由相机路径，沿途四视角（前/左/右/后）同步采集近距观测，与初始轨道数据合并后重新运行 Boxer 并施加每帧 $d_{\mathrm{th1}}$ 距离剪枝以消除远距伪影，最终融合得到精确的全局 3D 边界框标注。

\textbf{导航数据采集。}
导航数据通过 EGO-Planner 生成碰撞自由轨迹。与 OpenFly 和 SAGE-3D 使用 A* 搜索器生成运动学不可行的路径不同，FlyMirage 的 EGO-Planner 是经过大量仿真和真实实验验证的动力学可行规划器，输出的轨迹满足无人机运动约束——这是首个全自动管线生成真正 6-DoF 动力学可行航拍轨迹的数据集（此前仅有人工遥控采集的 UAV-Flow 和 OpenUAV 提供 6-DoF）。每个场景指定目标集数 $N_s$ 和每集目标数 $N_o$，候选目标须通过位置验证（无人机不在物体内）、可见性检查（射线检测遮挡）、航点安全（目标周围无占据点）和距离约束（2.0--10.0 m）四重筛选。从 3DGS 渲染四视角 RGB 和深度数据，经 Qwen-3.5-Flash 视觉质量评估过滤渲染伪影后存档。每条有效轨迹生成最多三种指令变体（物体中心如"找到书架"、相对位置如"去书架旁的椅子"、外观描述如"飞向绿色沙发"），支持多样的指代表达风格。

\textbf{数据集规模与成本。}
最终生成的数据集包含 500 个 3DGS 场景（跨 6 大类别：交通空间、办公场所、商业零售、工业设施、休闲场所、住宅空间——相比 SAGE-3D 约 75\% 为住宅场景的分布更为均衡），自动识别 5,000+ 唯一物体标签（每个典型场景含 60--100 个物体实例），约 50,000 条导航轨迹（平均长度 4.33 m，中位数 4.06 m，可组合为约 20 m 的长程任务），每条轨迹关联 2--3 条语言指令。成本方面，每场景 API 费用约 \$2（含 Marble 和所有 LLM 调用），批量渲染仅需消费级 RTX 4070 GPU，每场景全流程约 1 小时且可并行——相比之下 HM3D 需价值约 \$3,000 的 Matterport Pro2 传感器扫描，SAGE-3D 需逐场景人工设计，UAV-Flow 需约 \$100/小时的人工遥控飞行。

\paragraph{异构引擎导航基准}

OpenFly（2025.02）~\cite{2502_18041} 从平台构建层面改变了上述方法各自依赖单一渲染引擎/场景生成技术的思路——与 FlyMirage 利用生成式模型全自动化创建 3DGS 场景（500 个纯合成场景）不同，OpenFly 的核心差异在于\textbf{异构渲染引擎集成}：不创造新场景，而是将四种互补的渲染引擎统一到同一个数据生成平台中，每个引擎覆盖不同维度——Unreal Engine 提供 8 个城市场景（超 100 km$^2$，含建筑、车辆、行人等丰富数字资产），GTA V 提供以洛杉矶为原型的高度真实城市景观，Google Earth 覆盖 4 个城市区域（Berkeley/大阪/华盛顿 D.C./圣路易斯，共 53.60 km$^2$），\textbf{层次化 3DGS}（hierarchical 3DGS）从无人机采集的真实图像重建 5 个校园场景（超 7 km$^2$），是实现 real-to-sim 渲染的唯一引擎——使用 COLMAP 从输入图像生成稀疏点云，在真实世界几何和纹理的基础上构建可自由导航的仿真环境。数据集按引擎分工：UE 贡献 75.7\% 的训练数据（场景最多），3DGS 贡献近 20\%（提供 real-to-sim 真实感），Google Earth 仅贡献 4.46\%（为保证视觉质量仅采集高空视角数据）。

\textbf{自动化数据生成工具链。}
OpenFly 设计了三个统一接口控制跨引擎的智能体运动和传感器数据获取，并在此基础上构建四阶段自动化管线。

\textbf{点云获取}分两种策略：UE/GTA V 场景使用光栅化采样重建（在适当分辨率的采样点获取局部点云并拼接），3DGS 场景直接使用 COLMAP 稀疏重建步骤生成的点云。

\textbf{场景语义分割}提供三种灵活选择的方法：基于 Octree-Graph 的 3D 场景理解（从俯视图序列提取语义 3D 实例）、点云投影 + 轮廓提取 + GPT-4o 语义标注、以及手动标注（低质量点云或需精细分割时的备选方案）。

\textbf{轨迹生成}从场景点云构建全局体素地图 $M_{global}$，随机选择地标作为目标并在一定距离选起点，基于 $M_{global}$ 和自定义动作空间使用 A* 算法搜索碰撞自由轨迹——与 FlyMirage 使用 EGO-Planner 生成满足无人机动力学约束的轨迹不同，OpenFly 的 A* 生成运动学可行但非动力学最优的路径，步长设为 3 m、6 m 和 9 m 三档以适应不同距离目标。迭代选择终点为新起点可生成复杂轨迹。

\textbf{指令生成}采用子轨迹分割策略：按动作转换点将完整轨迹切分为子轨迹（忽略细微角度调整），提取每段子轨迹的关键动作和最后 3 帧图像送入 GPT-4o 生成子指令，最后由 LLM 整合为完整导航指令——这一策略显著优于直接输入全部图像的方法。随机抽检 3K 样本达到 91\% 合格率。

\textbf{数据集规模。}
最终构建的 OpenFly 数据集包含 100K 条轨迹（是此前最大空中 VLN 数据集 CityNav 的 3 倍、AerialVLN 的 12 倍），覆盖 18 个场景，词汇量 15.6K，平均轨迹长度 99.1 m（刻意设计为短-中程，更符合自然人类使用习惯），平均指令长度 59 词，每条轨迹平均 35 个动作。6 种离散动作定义为 $\{$Forward, Turn Left, Turn Right, Move Up, Move Down, Stop$\}$（4-DoF）。数据集划分为 Train/Test Seen/Test Unseen，其中 Test Unseen 包含 3 个未见场景（UE-smallcity、3DGS-sjtu02、GTA V 洛杉矶城市），共 1200 条轨迹。

\textbf{OpenFly-Agent 关键帧感知模型。}
与 VLA-AN 使用完整 MLLM（Qwen2-VL 7B）作为策略引擎不同，OpenFly-Agent 基于 OpenVLA（在 1M 操作数据上预训练的 VLA 模型）构建，核心创新是解决空中 VLN 的视觉序列冗余问题。OpenVLA 原始设计处理单帧图像，而空中导航需要处理长视频序列——均匀帧采样会错过包含关键地标的帧。OpenFly-Agent 提出两个协同策略。

\textbf{关键帧选择（KS）}：启发式识别无人机运动变化点，提取变化点及前后各 2 帧构成候选集；地标定位模块（3 层交叉注意力，融合 LLM 隐状态中的文本和图像特征）预测指令相关地标的边界框 $\mathbf{b} \in \mathcal{R}^4$，候选帧中边界框面积大于阈值 $\theta = 0.25 \times$ 输入图像尺寸的保留为最终关键帧。

\textbf{视觉 token 融合（VTM）}：选取关键帧集中边界框最大的帧为参考帧，密集计算参考帧与其他帧各视觉 token 对的余弦相似度，高相似度 token 通过平均融合，未融合的对比帧 token 被丢弃；维护容量为 $K=2$ 的 FIFO 记忆库保留最新关键帧。关键帧内部通过 grid pooling 进一步压缩至 1 个 token（当前帧保持 256 个 token 未压缩以获取最新观测）。动作预测使用词汇表最后 256 个 token 作为 6 种无人机动作的特殊 token。

\textbf{实验结果。}
在 OpenFly 基准测试中，OpenFly-Agent 在 test-seen 上实现 NE 93 m、SR 34.3\%、OSR 64.3\%、SPL 24.9\%，分别超越次优方法 NaVila 14.0\% 和 7.9\% 的成功率优势（test-unseen SR 22.6\% vs NaVila 14.7\%）。消融实验揭示了关键帧选择与视觉 token 融合的协同效应：单独使用 KS 时 SR 仅 9.2\%，单独使用 History + VTM 为 16.6\%，两者组合后 SR 飙升至 34.3\%——原因是 VTM 解决了文本-图像 token 数量严重不平衡问题，避免背景噪声稀释对关键线索的注意力。在 23 个真实户外场景的硬件实验中（Q250 无人机 + Jetson Xavier NX，轨迹长度 50--500 m），OpenFly-Agent 达到 SR 26.09\%、OSR 34.78\%，且用 OpenFly 多引擎数据训练的模型显著优于用 AerialVLN 单一引擎数据训练的模型，证明多引擎异构数据有效弥合 sim-to-real gap。

\paragraph{单目视觉避障}

Flying in Clutter（2025.12）~\cite{2512_17349} 与 SOUS VIDE 出自同一团队（ZJU / Differential Robotics），共享"3DGS 仿真 + 真实部署"管线框架，但在三个关键维度上做出了根本性改变：\textbf{任务}从轨迹跟踪变为杂乱场景自主导航（仅给定目标点，无参考轨迹）；\textbf{训练范式}从 MPC 专家蒸馏变为端到端 PPO 强化学习；\textbf{域适应}从 RMA 在线动力学适应变为对抗域适应（Adversarial Domain Adaptation）解决视觉域间隙。

\textbf{渲染加速。}
SOUS VIDE 的 FiGS 仿真器以 130 fps 渲染，足以支撑模仿学习的数据需求。但 RL 训练需要 1,024 个并行环境以 rollout 长度 128 同步采样，对渲染吞吐量提出更高要求。本方法引入 Speedy Splat 的双重加速策略：其一，高斯-分块映射从原始 3DGS 的正方形包围盒（基于半径 $r = \lceil 3\sqrt{\lambda_{\max}} \rceil$，其中 $\lambda_{\max}$ 为投影二维协方差矩阵的最大特征值）改为严格按椭圆投影区域分配，减少无效分块计算；其二，为每个高斯分配重要性分数 $\tilde{U}_i = (\nabla_{g_i} I_g)^2$（其中 $I_g$ 为渲染图像，$g_i$ 为第 $i$ 个高斯的二维投影标量值）并据此裁剪低贡献基元。结合 gsplat 并行光栅化后端的向量化 GPU 批处理操作（涵盖投影、排序、合成全流程），将峰值渲染速度提升至约 30,000 fps，比 FiGS 提升约 230 倍。训练环境包含 10 个通过手持摄像机扫描重建的杂乱场景（9 个训练 + 1 个评估），每个场景包含障碍物的添加、移除或重新排列等变化，COLMAP 点云导入 Isaac Sim 用于碰撞检测和深度渲染，RGB 观测则由 3DGS 实时提供。

\textbf{Actor-Critic 架构与单目感知。}
策略输入为 $60 \times 80$ RGB 图像和 20 维状态向量 $\bm{s}_t = [\bm{v}_b, \bm{R}, \bm{d}_g, z_c, z_t, \bm{a}_{\text{last}}]$（机体系速度、旋转矩阵、归一化目标方向、当前/目标高度、上一步动作），经归一化模块标准化后，由三层 CNN 提取 4,480 维视觉潜在向量 $\bm{z}_{rgb}$，再与归一化状态向量拼接，送入 512 单元 GRU 捕捉时序动态，最终通过 MLP 和投影层输出归一化动作 $\bm{a}_t = [T, \phi, \theta, \psi]$ 并经级联 PID 转为力矩和推力指令。与 SOUS VIDE 的 SV-Net 相比，本方法的 Actor 完全不依赖水平位置 $(p_x, p_y)$ 的观测，导航方向仅由视觉信息和目标方向向量 $\bm{d}_g$ 隐式推断。Critic 采用非对称 Actor-Critic 范式：额外获取 $60 \times 80$ 深度图作为特权信息，为价值函数提供准确的空间理解，稳定训练过程——消融实验表明，去除特权深度后策略奖励降低约 20\%，无法有效收敛，证实深度特权信息对 Critic 形成稳定价值函数的关键作用。

\textbf{奖励函数与终止条件。}
奖励函数采用加权和设计 $r_t = \sum_i w_i r_i$，以距离奖励 $r_{\text{distance}} = \text{clamp}(d_{t-1} - d_t, -d_{\text{step}}, d_{\text{step}})$（权重 30.0，其中 $d_{\text{step}} = v_{\max} \cdot \Delta t$，$v_{\max} = 2$ m/s）为核心驱动，配合成功奖励（80.0）和碰撞惩罚（$-80.0$）构成强烈的正负反馈；正则项包括 Z 轴速度惩罚（$-0.5$）、动作幅度惩罚（$-0.3$）、动作变化惩罚（$-0.6$）和速度超限惩罚（指数衰减，$-1.0$，仅当 $\|\bm{v}_t\| > v_{\max}$ 时激活）。与 SOUS VIDE 依赖 MPC 专家轨迹的跟踪奖励不同，本方法的奖励完全不依赖参考轨迹，仅通过目标距离变化引导探索。终止条件为三项之一：高度超出安全范围 $[0.1, 1.7]$ m、碰撞检测（基于膨胀点云占据图的实时检查）、或成功到达目标。

\textbf{对抗域适应。}
SOUS VIDE 通过 RMA 模块在线适应动力学参数（推力系数、质量）的变化，本质上解决的是\emph{动力学}域间隙。本方法面对的是更棘手的\emph{视觉}域间隙：3DGS 重建场景用 DJI Osmo Action 5 Pro 拍摄，但部署时用的是海康威视工业相机——镜头畸变、色彩风格、曝光参数和 3DGS 重建伪影共同构成显著的视觉分布偏移。方案是在 Actor 的 CNN 编码器后接一个判别器，通过梯度反转层（Gradient Reversal Layer）在反向传播时反转梯度并缩放因子 $\lambda$，使编码器被迫学习域不变特征。判别器用二元交叉熵损失训练：$\mathcal{L}_{\text{DA}} = -\mathbb{E}_{\bm{z}_s \sim \mathcal{D}_s}[\log D(\bm{z}_s)] - \mathbb{E}_{\bm{z}_t \sim \mathcal{D}_t}[\log(1 - D(\bm{z}_t))]$，其中 $\mathcal{D}_s$ 和 $\mathcal{D}_t$ 分别为 3DGS（源域）和真实世界（目标域）的潜在特征。与将域适应与 RL 训练解耦的做法不同，本方法将 DA 损失直接集成到 PPO 训练目标中：$\mathcal{L}_{\text{total}} = \lambda_1 \cdot \mathcal{L}_{\text{PPO}} + \lambda_2 \cdot \mathcal{L}_{\text{DA}}$，避免分阶段训练造成的策略性能退化。DA 仅需数分钟的真实世界相机视频即可完成对齐。

\textbf{域随机化。}
动力学随机化（DY）模拟执行器缺陷：向控制指令注入高斯噪声和可变延迟（均匀采样 $D \sim \mathcal{U}(0, 80)$ ms），延迟建模为历史窗口上的移动平均，噪声和延迟在每轮内以随机间隔 $T \sim \mathcal{U}(10, 100)$ ms 重新采样，迫使策略适应时变扰动而非过拟合固定延迟模式。感知随机化（DR）包含两个层面：状态估计噪声采用均值回归过程 $n_{t+1} = (1 - \theta) n_t + \mathcal{N}(0, \sigma^2)$（$\theta = 0.1$），分别对速度（$\sigma_{vel} = 0.08$）、方向（$\sigma_{dir} = 0.05$）和高度（$\sigma_z = 0.03$）施加时序相关扰动，模拟真实传感器估计误差的非独立性；视觉增强采用 RoboSplat 方法直接变换 3DGS 参数 $C'_{dc} = \alpha \cdot C_{dc} + \beta + \bm{\epsilon}$（颜色缩放 $\alpha \sim \mathcal{U}(0.8, 1.3)$、偏移 $\beta \sim \mathcal{U}(-0.05, 0.05)$、加性噪声 $\sigma_{rgb} = 0.025$），并随机化视场角 FOV $\sim \mathcal{U}(67^\circ, 106^\circ)$ 以增强对不同障碍物尺度的泛化。

\textbf{三阶段训练与实验验证。}
训练在单张 NVIDIA L40 GPU 上约两天完成，采用三阶段递进：第一阶段训练基线策略，第二阶段引入视觉/动力学域随机化（DR+DY）微调，第三阶段加入对抗域适应（DR+DA+DY）对齐潜在特征。t-SNE 可视化从两个维度揭示 DR 与 DA 的互补机制：（1）域对齐层面，DA 将仿真与真实域的 GSI 从基线的 0.06 降至接近 0，而单纯 DR 反而使 GSI 升至 0.17——说明 DR 扩展了数据覆盖范围但真实世界观测对编码器仍属分布外；（2）特征区分度层面，用 Logistic Regression 对五个数据增强环境的潜在特征分类，基线准确率高达 91\%（模型依赖环境特定风格），DR 降至 71\%（通过覆盖更多风格实现鲁棒性），DA 降至 54\%（即使仅对齐原始仿真与真实数据，也泛化到未见增强环境），DA+DR 进一步降至 40\%——这表明 DR 靠扩展数据覆盖，而 DA 迫使网络学习与域偏移无关的通用几何与语义特征，二者的协同效果最优。

消融实验中 RGB Actor 策略的奖励与基于深度的基线策略相当，证实策略能从纯光度数据中隐式提取可通行性线索；去除 Critic 特权深度后奖励降低约 20\%。真实世界部署在 Jetson Orin NX 上以约 2 ms 推理延迟运行，实时状态估计（归一化方向向量、高度、速度）由 EKF 融合下视光流与 IMU 测量值提供。在 10 条测试轨迹上达到 80\% 成功率（仿真中为 90\%），失败主因为旋翼下洗导致枝叶摆动——这一动态因素未被仿真建模。消融对比表明：仅用 DY 的基线因动力学间隙而失控，DY+DR 因视觉域间隙仍频繁碰撞复杂障碍物（如枝叶），DY+DA 缺乏对传感器噪声和环境光变化的鲁棒性，只有 DR+DA+DY 的完整组合同时解决动力学间隙、视觉鲁棒性和视觉域迁移三个问题。连续往返导航任务中，在飞行段间随机重排障碍物位置（包括树木和柱子），策略仍能生成平滑无碰撞轨迹，验证了其泛化能力。

快速光照变化实验中（灯光反复开关），策略在极端光度扰动下维持稳定控制命令，证明 DA+DR 使编码器提取了超越低层像素强度的语义几何表示。

\paragraph{扩散规划无人机导航}

基于扩散规划的无人机规划器（2025.10）~\cite{2510_03545} 从 3DGS 场景构建 SketchPlan 系统。与 SOUS VIDE 的 MPC 专家轨迹不同，该方法的核心差异在于允许人类通过手绘草图指定 3D 飞行路径。SketchPlan 由两个独立模块组成：SketchAdapter 用偏最小二乘法（PLS）学习从手绘草图 $\mathcal{S}$ 到 2D 相机投影 $\mathcal{P}$ 的仿射变换（5 个主成分的潜空间），DiffPath 用 1D UNet 扩散模型从 2D 投影和深度图像重建 3D 路径——编码器将深度图经 6 层卷积提取全局特征，沿投影坐标采样局部嵌入，拼接后条件化去噪网络。DiffPath 用人工标注数据训练 SketchAdapter，用全量自动标注数据训练扩散模型，实现半监督学习。生成的路径通过凸多面体碰撞过滤器后处理：从累积深度点云构建无碰撞半空间约束，以二次规划实时求解最近碰撞自由航点（50 Hz 以上）。系统从 3DGS 场景构建 32k 飞行路径合成数据集，实现零样本仿真到真实迁移。

\paragraph{固定翼无人机导航}

FalconWing（2025.05）~\cite{2505_01383} 将上述方法的 GSplat 仿真范式从旋翼无人机扩展到\textbf{固定翼无人机}——这一扩展并非简单的平台替换，而是面临根本不同的物理约束：固定翼必须持续维持最低空速以产生升力（FalconWing 的 150 g 平台约需 7 m/s），转弯半径受协调转弯关系 $R = V^2 / (g\tan\phi)$ 约束（在 7 m/s 空速和 30$^\circ$ 坡度下约 8.7 m），且每增加 1 g 载荷均会提高最小空速与转弯半径（例如若像旋翼平台那样加装 174 g 的 Jetson Orin 机载计算机，总重增至约 300 g 后最小空速升至 10 m/s、最小转弯半径扩大至约 17.5 m，几乎等于典型室内场地 20 m 的宽度，使机动变得不可行）。这意味着旋翼平台上理所当然的机载计算方案在固定翼室内场景中不可行——室内空间对固定翼平台提出了比旋翼严苛得多的重量约束。

为此 FalconWing 采用与 SketchPlan 和 SOUS VIDE 截然不同的传感器极简策略：整个平台仅 150 g（137 g UMX Turbo Timber 泡沫机身 + 9 g RunCam Spotter 模拟 FPV 相机，通过定制 3D 打印支架 5 g 沿机身中心线安装以避免螺旋桨遮挡并保持重心），所有计算卸载至地面站（RTX 4090），图像通过 5.725 GHz 模拟链路传输至地面 USB 采集卡以 640$\times$480 分辨率、20 Hz 流入 ROS 话题，控制指令通过 Arduino Mega 2560 桥接 ROS 与 Spektrum NX-8 发射机训练端口以 20 Hz 下发四通道 PPM 信号（油门、副翼、升降舵、方向舵），人类飞行员可随时一键切换手动接管。为缓解模拟视频在飞行中因振动退化的风险，系统部署基于 SSIM 的帧质量监控：连续 5 帧低于阈值 0.7 即触发告警由飞行员接管。

设计理念是将消费级航模（70 cm 翼展，配备扰流板以兼顾机动性和额外升力）通过最小改装转化为 ROS 兼容的视觉自主研究平台——新用户按 15 页手册可在 10 小时内完成组装首飞，中轻度坠机的修复成本极低（ESC 约 \$80、螺旋桨 \$8，更换时间 $<$2 小时），这正是 FalconWing 被定位为类似 F1-TENTH 的教育"飞行套件"的出发点。

\textbf{GSplat 仿真与系统辨识。}
与 SOUS VIDE 和 FalconGym 直接使用 NeRF 渲染不同（FalconWing 的前身 FalconGym 1.0 使用 NeRF，渲染速度和质量受限），FalconWing 用 GSplat 替代 NeRF 实现更快速和更高质量的渲染。系统用机载 FPV 相机在 40$\times$20$\times$5 m 室内飞行场中采集约 2000 张图像（覆盖不同位置和视角），COLMAP 估计相机内参和位姿后，以放置在跑道旁的 80 cm ArUco 标定板为全局原点，经 Kabsch-Umeyama 算法统一 COLMAP 坐标系到世界系——这一标定策略与 SOUS VIDE 的 ArUco 世界坐标对齐方法一脉相承，但 FalconWing 消除了对动作捕捉系统的依赖。NeRFStudio Splatfacto 管线在 RTX 4090 上约 15 分钟收敛，单帧 960$\times$720 渲染仅 0.004 秒。除模拟相机版本外，还额外构建了数字 ArduCam RGB 相机的仿真变体，为重视图像质量胜过轻量化的研究者提供选择。

非线性系统辨识采用降阶运动学模型（9 维状态 $[p_x, p_y, p_z, \theta, \gamma, \phi, v_x, v_y, v_z]$；4 维输入 $[u_T, \delta_a, \theta_c, \gamma_c]$，其中前 6 维状态变量恰好就是相机位姿 $p$），通过混合视觉状态估计获取状态-动作对——近距离用 OpenCV ArUco 检测器直接估计，远距离用 ViT 骨干的 iGSplat 单帧位姿估计网络（冻结早期 Transformer 层保留 ImageNet-21k 预训练特征，回归头预测角度的 $\sin/\cos$ 值以避免角度回归不连续问题），平均位置误差 0.42 m、偏航误差 2.37$^\circ$——这一混合方案的设计动机在于避免 iNeRF 等迭代位姿优化方法的计算开销，以单次前向推理换取实时速度。过滤噪声帧（模拟视频闪烁和明显异常位姿估计）后以非线性最小二乘（SciPy）拟合动力学参数 $K^*$，优化目标为 $\min_K \sum_{(x_t, u_t) \in D_x} \|f_K(x_t, u_t) - x_{t+1}\|_2^2$。此外还构建了 UMX 飞机本身的 GSplat 资产，支持在仿真中放置和渲染逼真的长机飞机——这一资产利用了 FalconGym 的 Edit API 对高斯基元的操作能力，是后续编队跟踪任务中生成训练数据的基础。

\textbf{视觉策略训练与 sim-to-real 验证。}
FalconWing 在两个 case study 上验证，均采用模仿学习范式——这与 FalconGym 的成功经验一致，但 FalconWing 针对固定翼特性消除了对已知目标位置和 IMU 读数的依赖（FalconGym 需要这两项输入），仅用单 ViT 网络替代 FalconGym 的双 ViT 架构以减小模型体积和推理延迟。

在编队视觉跟踪任务中，设计了三种策略变体：端到端 RGB 控制器（单 ViT 网络以当前图像 $I_t$ + 过去 30 步控制历史 $u_{t-30:t-1}$ 为输入，通过轻量自注意力模块融合后直接输出动作 $u_t$，训练目标为最小化 $\mathcal{L}(\theta) = \frac{1}{|D_C|} \sum \| \pi_\theta(I_t, u_{t-30:t-1}) - u_t^* \|_2^2$，ART 约 0.02 s 最快但泛化差——论文分析其过拟合原因在于长机仅占图像像素的极小比例，网络过度依赖背景外观而非长机本身的视觉表示）、Mask 控制器（先以 UNet 从 RGB 预测长机二值掩码——掩码真值通过 UMX GSplat 资产中已知的翼尖和鼻尾端点三维位置近似为 3D 椭球体投影生成，收集 10K 对 RGB-掩码图像训练 UNet——再由轻量 ResNet 从掩码 + 控制历史输出动作，约 100K 掩码-动作对训练控制器，泛化更好但在一种未见机动中 UNet 将背景明亮窗户的白色条纹误检为长机导致跟踪失败）、以及 RGB+Mask 控制器（将掩码作为第四通道叠加在 RGB 上输入网络，兼具外观上下文与分割优势，整体 SR 和 ATE 最优但 ART 最高使硬件部署风险增大）。

训练数据由专家状态控制器在 GSplat 仿真中生成——专家控制器首先根据标准固定翼控制设计实现，配置长机使用专家控制器跟踪预定义航点，僚机专家控制器同时访问双方真值状态计算最优动作，并在专家输出上注入高斯噪声 $u_t = u_t^* + \mathcal{N}(0, \sigma^2)$ 以拓宽状态-动作分布覆盖。训练中对长机的 GSplat 高斯基元施加颜色/尺度域随机化增强鲁棒性。最佳 RGB+Mask 策略在 3 种未见机动（左转 S 形下降、右转 S 形爬升、右转急爬升）上实现 100\% 跟踪成功率（30 次试验），且对 0.5$\times$--2$\times$ 尺度和 30\% 椒盐噪声的领导机外观扰动保持稳健（仅在 50\% 缩放下因检测不可靠而性能退化）。由于缺乏动作捕捉系统，所有跟踪实验仅在仿真中进行，sim-to-real 验证在着陆任务中完成。

在自主着陆任务中，选用满足 20 Hz 实时约束的 RGB 策略（唯一 ART $<$0.05 s 的变体），复用编队跟踪的 RGB 模仿学习框架但移除长机资产，专家控制器的目标改为跑道对齐和指定垫位触地。着陆场景设定为 40$\times$20$\times$5 m 室内飞行场一端的蓝色着陆垫（13$\times$2$\times$0.1 m），每次试验由人类飞行员手动飞至距跑道约 20 m、高度 1.5 m 的初始位置（此时着陆垫刚进入模拟相机视野）后切换自主模式，飞行员在出现不安全行为时立即接管。着陆评估指标为成功率（触地点在着陆垫范围内）和绝对侧向偏差（ALD，触地点到跑道中心线的距离，跑道宽 2 m 故可接受偏差 $\leq$1 m）。由于无动作捕捉，真实触地位置通过在起落架上涂粉末测量触地痕迹确定。仿真复现方面，记录自主模式启动时的图像并用 iGSplat 估计初始位姿后在仿真中重放，与状态专家对比。仿真中状态控制器和视觉控制器均在 10 次试验中 100\% 成功着陆，平均 ALD 分别为 15 cm 和 37 cm。真实硬件 10 次着陆试验中 8 次成功（80\%），平均 ALD 41 cm，两次失败（Run 3 和 Run 10）源于飞机交接时俯仰角过陡和模拟视频闪烁噪声——真实 ALD 略大于仿真，论文归因于动力学估计误差和图像渲染质量差异。整个软件栈（双仿真环境 + 系统辨识参数）打包为 conda 环境开源分发，设计目标是作为可复用的视觉固定翼控制基准。

\paragraph{导航仿真器}

FalconGym 2.0（2025.10）~\cite{2510_02248} 沿袭 SOUS VIDE 的 GSplat 仿真范式（COLMAP 位姿估计 + NeRFStudio Splatfacto 训练场景 + ArUco 标定板做世界坐标系对齐），但识别出后者的核心局限：视觉策略在单一赛道上训练后严重过拟合——即使用 Edit API 删除赛道中的门，策略仍按记忆路径飞行，表明其依赖背景外观而非门相关特征。FalconGym 2.0 的设计动机正是打破这种单赛道特化，使一个视觉策略泛化到任意未见赛道。

\textbf{Edit API：可编程赛道生成。}
与 SOUS VIDE 的静态 GSplat 场景不同，FalconGym 2.0 在 GSplat 表示之上封装了七个可组合的张量操作——add、translate、rotate、scale、duplicate、delete、lighting——每个操作直接修改选中高斯基元的均值 $\mu_j$、协方差 $\Sigma_j$（通过旋转 $R_j$ 和尺度 $s_j$）、颜色 $c_j$ 和不透明度 $\alpha_j$，在 RTX 4090 上单次操作仅约 0.004 秒。用户在世界坐标系下指定编辑指令，后端自动处理内部坐标系转换和相机-世界变换。这一 API 不仅支持静态门的程序化放置，还支持 4D 时变仿真（门随时间运动），为后续的性能引导精化提供基础。直觉上，SOUS VIDE 用 GSplat 重建"一个固定场景"，而 FalconGym 2.0 用 Edit API 从一个场景中"变出无数个赛道"。

\textbf{模块化视觉策略。}
为解决端到端 ViT 策略（如 FalconGym 1.0 的双 ViT 架构）的过拟合问题，FalconGym 2.0 采用感知-控制分离的模块化设计：感知模块为 U-Net 门分割网络，将机载 RGB 图像预测为二值掩码（白色像素表示门，黑色像素表示背景）；轻量控制器以掩码和短期控制历史为输入输出下一步动作。训练数据仅用两门赛道布局——直觉是，初始状态加两扇连续门足以张成更长赛道的局部几何变异性，泛化到多门赛道由不变性和组合性保证。门分割的真值掩码由已知的门三维位置和几何参数通过标准 3D-to-2D 投影解析生成，无需人工标注。

\textbf{性能引导精化（PGR）算法。}
均匀采样两门赛道空间在大型高维空间中样本效率低。PGR 将赛道采样建模为 min-max 对抗优化问题 $\min_\theta \max_{g \in G} \mathbb{E}_{\tau \sim \pi_\theta}[\mathcal{L}(\tau; g)]$，其中 $g$ 参数化两门赛道，$G$ 为动力学可行且可观测的两门空间，损失 $\mathcal{L}$ 编码碰撞/超时和穿越误差。具体实现采用网格化迭代：将两门空间 $G \subset \mathbb{R}^8$（每门有 $(x, y, z, \text{yaw})$）划分为网格（固定翼案例约 20k 网格），每轮迭代先用状态专家在采样赛道上收集模仿学习数据训练视觉策略，再在保留验证集上评估每个网格的策略表现 $\ell_i$，然后以与 $\ell_i$ 成正比的概率重新采样下一轮赛道（混合 $\beta=0.05$ 的均匀采样防止模式坍塌）。PGR 同时应用于感知模块和控制器的训练。与 SOUS VIDE 的领域随机化（从参数分布中均匀采样）不同，PGR 是一种主动课程学习策略，将训练集中到策略表现最差的赛道区域。

\textbf{实验验证。}
在固定翼 UAV 和四旋翼两个案例研究中，PGR 训练的单一视觉策略在三条未见赛道上均实现 100\% 成功率，而基准方法（FalconGym 1.0 端到端 ViT、Geles 等的异步 actor-critic）在训练赛道上表现良好但在未见赛道上严重退化。在门位姿扰动鲁棒性实验中，PGR 策略的成功率曲线紧随状态专家，显著优于基准。在四旋翼零样本 sim-to-real 迁移中，策略部署到 280g ModalAI Starling 2 四旋翼（PX4 飞控 + VOXL 2 机载计算机 + Adreno 650 GPU，8 Hz 闭环频率），30 次硬件试验覆盖两条三门静态赛道和一条移动门赛道，总成功率 98.6\%（70 个门中通过 69 个），主要失败模式为 U-Net 将背景梯子误检为门和 8 Hz 闭环延迟限制。

\paragraph{人体感知导航仿真}

与 SOUS VIDE 自建轻量仿真器 FiGS 不同，Habitat-GS（2026.04）~\cite{2604_12626} 的出发点是将 3DGS 渲染集成到成熟的 Habitat 具身智能生态系统中——这带来了一个 SOUS VIDE 不需要面对的技术挑战：Habitat-Sim 的渲染管线基于 OpenGL，而高性能 3DGS 光栅化依赖 CUDA 分块渲染，两个异构后端之间需要高效桥接。同时，SOUS VIDE 只处理静态场景渲染（策略学习通过领域随机化补偿域间隙），而 Habitat-GS 进一步引入了\textbf{可驱动高斯化身}模块，使仿真环境中首次出现照片级真实感的动态人类，支持人感知导航训练。

\textbf{视觉-导航解耦设计。}
Habitat-GS 的核心设计原则是\emph{视觉-导航解耦}：3DGS 负责所有视觉渲染（RGB、深度），传统 NavMesh 继续负责导航逻辑。直觉上，这个设计源于 3DGS 的根本限制——高斯基元是各向异性椭球体的集合，没有显式几何表面，无法直接用于碰撞检测。将视觉和导航解耦后，3DGS 渲染提供照片级真实感的传感器观测，而 NavMesh 保证可靠的路径规划，两者各司其职。

\textbf{零拷贝 CUDA-OpenGL 互操作。}
为实现 3DGS 渲染与 Habitat 传感器管线的无缝集成，Habitat-GS 采用零拷贝 CUDA-OpenGL 互操作机制：初始化时将 OpenGL 颜色和深度纹理注册为 CUDA 可访问资源；每帧渲染时，CUDA 光栅化器执行前向泼溅，将 RGB 和深度结果直接写入映射的 GPU 缓冲区，无需 CPU 侧数据搬运。之后通过全屏深度合成着色器将 3DGS 深度缓冲与 OpenGL 深度缓冲合并，正确处理 3DGS 场景、网格对象和多个高斯化身之间的遮挡关系。

\textbf{可驱动高斯化身与代理胶囊。}
在视觉-导航解耦和零拷贝互操作解决了静态场景渲染问题之后，Habitat-GS 进一步引入化身模块，使仿真环境中首次出现照片级真实感的动态人类——这是对前置模块的递进扩展，因为动态化身的渲染、驱动和碰撞检测分别依赖 CUDA-OpenGL 互操作管线和 NavMesh 导航系统。现有高斯化身方法（GaussianAvatars、AnimatableGaussians）聚焦于独立重建与渲染，未与仿真器传感器管线或导航网格系统耦合。Habitat-GS 的化身模块设计目标是让每个动态人类同时充当两个角色：照片级真实感的视觉实体和有效的导航障碍物。

视觉渲染方面，系统离线预烘焙标准空间的高斯属性（位置、SH 系数、不透明度、尺度、旋转和线性混合蒙皮权重），运行时通过 CUDA 加速的 LBS 核心将标准高斯变形到任意 SMPL-X 姿态，避免了代价高昂的神经网络推理。为支持多化身场景下的高效渲染，所有活跃化身的变形高斯被拼接到单一 GPU 缓冲区，在一次 CUDA pass 中完成光栅化，随后与主帧缓冲合成——这确保了 3DGS 场景、网格对象和任意数量化身之间正确的深度排序。

运动轨迹由 GAMMA 模型在 NavMesh 上生成场景感知的物理可行行走路径。具体管线为：给定起点、终点和若干途经点（在场景 NavMesh 上采样），先计算连接所有航点的最短路径，再将其输入 GAMMA 生成时序一致的身体姿态和根平移轨迹；对每帧通过前向运动学计算 SMPL-X 关节变换矩阵。运行时，化身模块在当前仿真时间对关节矩阵和根轨迹进行时间插值，将结果传递给 CUDA LBS 核心完成变形。

碰撞方面，由于 3DGS 无法直接用于碰撞检测，Habitat-GS 引入\textbf{代理胶囊}机制将化身的视觉外观与碰撞几何解耦：在轨迹合成阶段离线生成 SMPL-X 骨骼段对应的胶囊原语，并预计算每帧的世界空间位置；运行时仅通过时间插值检索当前胶囊位置，计算开销可忽略。由于静态 NavMesh 无法表达移动化身随时间变化的占位，Habitat-GS 扩展了 Habitat-Sim 路径规划器以支持动态胶囊障碍物：每个仿真步骤，化身管理器收集所有活跃化身的代理胶囊并注入路径规划器作为动态障碍物；步裁剪例程检查智能体包围胶囊与每个化身胶囊的相交，在碰撞边界裁剪智能体运动，确保其不会穿越化身身体。此外，Habitat-GS 为 Habitat-Lab 提供两个化身感知查询 API：一个返回到最近化身胶囊的距离，另一个指示候选步是否被阻挡——这两个 API 为研究人员设计基于距离的接近奖励、碰撞惩罚和跟踪指标提供了接口。

实验中，智能体先在静态 PointNav 上训练 $5 \times 10^7$ 步建立基础导航能力，再在 20 个 GS 场景上精调 $5 \times 10^6$ 步（仅占预训练预算的 10\%）学习化身感知导航。在 3 个高斯化身同时行走的场景中，训练出的智能体在 GS 测试集上碰撞率仅 2.342\%，个人空间侵入度（PSI）为 0.068；且该能力零样本迁移至网格测试集（碰撞率从 2.521\% 降至 2.342\%，PSI 从 0.075 降至 0.068），验证了高保真化身感知线索（识别人形轮廓、预判行走方向、估计个人空间）的跨域泛化能力。

\textbf{混合域训练策略。}
在点导航实验中，Habitat-GS 揭示了一个 SOUS VIDE 未涉及的发现：纯网格训练快速收敛但跨域泛化差（GS 测试集 SR 61.70\%），纯 3DGS 训练构建更强能力但收敛过慢，而混合域训练（20\% 网格 + 80\% GS）为最优策略——少量网格场景建立基础导航能力，大量 GS 场景增强视觉鲁棒性。该策略在 GS 测试集上达到最高 SR 的同时保持网格测试集性能与纯网格训练持平。

\paragraph{动态障碍物导航仿真}

与 Habitat-GS 将 GS 渲染嵌入 NavMesh 导航框架不同，ReaDy-Go（2026.02）~\cite{2602_11575} 走了另一条路线：放弃 NavMesh 和物理引擎，让动态障碍物（人类）也完全用 GS 表示，从场景到障碍物统一在 GS 表示下完成渲染和规划。这一设计选择的动机是：基于 NavMesh 的导航依赖异构表示（GS 负责视觉、NavMesh 负责碰撞），需要额外的网格提取和后端桥接；而统一 GS 表示可以实现端到端的纯视觉仿真管线——输入只需一段环境视频，即可自动生成带动态障碍物的照片级真实感导航数据集。

\textbf{动态 GS 仿真器。}
ReaDy-Go 的仿真器由两个 GS 模块组成。静态场景用 PGSR 重建——相比标准 3DGS，PGSR 通过平面约束和几何正则化将高斯基元压缩为扁平平面，显著提升表面重建精度，这对于后续场景体素化和占用图提取至关重要。动态人类用 HUGS 从 NeuMan 数据集视频中提取可动画人类 GS 化身：在 SMPL 身体模型初始化的标准空间中参数化人类高斯属性（位置、SH 系数、不透明度、尺度、旋转），并学习线性混合蒙皮（LBS）权重，使人类 GS 可由 SMPL 关节参数驱动到任意姿态。人类运动由 PriorMDM 生成：给定一条用户指定的 2D 轨迹，将其转换为体根线速度和旋转速度，归一化到 HumanML3D 表示后输入 PriorMDM（运动扩散模型作为生成先验），预测 3D 身体关节位置，再通过 SMPLify 拟合为 SMPL 参数并变换到世界坐标系，实现沿指定轨迹的合理人体动画。

\textbf{面向动态 GS 的专家规划器。}
由于现有 GS 规划方法（如 Splat-Nav）仅适用于静态场景且针对全向飞行器，ReaDy-Go 设计了面向地面车辆的 Hybrid A$^*$ 专家规划器。其核心挑战是：动态人类也是 GS 表示，需要从高斯基元中提取准确的 2D 占用信息。解决方案是：首先将静态场景体素化，基于累积不透明度过滤地面附近的高斯噪声（这是 PGSR 虽有几何正则化但仍残留的问题），生成机器人可导航图和人类可行走图；然后对人类 GS 基元进行清洗——移除协方差迹超过 75\% 百分位的高不确定度高斯椭球，再基于均值位置进行体素网格下采样——投影到 2D 可导航图，并按常速预测在 2 秒时域内膨胀安全区域。

规划器使用运动原语库（三种速度比例 $\times$ 三种转向命令）扩展节点，代价函数为弧长、转向惩罚和启发式代价的加权和。由于初始规划依赖短时域常速预测，运行时引入反应式重规划机制：当重规划周期到达，或动态障碍物与当前执行轨迹的前瞻段之间距离低于安全阈值时触发重规划。人类规划器则故意采样与机器人起点-终点连线交叉或平行的人类起终点对，以生成人-机器人交互密集的场景。

\textbf{策略训练与实验。}
策略架构简洁：10 层残差卷积网络编码 RGB 图像（20 维），取连续三帧编码结果与前一步动作和相对目标位置拼接（64 维），通过 3 层 MLP 输出线速度和角速度。每个环境仅用约 6 分钟视频（1000--1500 帧）重建 GS 场景，生成 400 个训练回合（80k--120k 样本）。实验中有三个关键发现：（1）环境特异性策略（ReaDy-Go 和 Vid2Sim）在目标环境中一致优于参数量 15 倍的通用导航模型（GNM、ViNT、NoMaD）；（2）在静态环境中 ReaDy-Go 与 Vid2Sim（使用 Unity 引擎人类资产）性能相当，但在动态环境中差距显著——ReaDy-Go 保持高成功率而 Vid2Sim 大幅下降，表明照片级真实感动态障碍物对动态场景导航至关重要；（3）在三个环境联合训练的策略零样本部署到未见环境，在静态和动态任务中均达到 50\% 以上成功率。

\paragraph{无人机数字孪生}

与 SOUS VIDE 将 3DGS 仿真用于导航策略训练不同，UAVTwin（2025.04）~\cite{2504_02158} 将 3DGS 数字孪生用于\emph{感知算法的数据增强}——从无人机采集的多段视频中重建背景 3DGS 模型，再在背景中合成可控人体前景，生成带像素级标注的训练数据。该方法面对两个 SOUS VIDE 未涉及的挑战：(1) 多架无人机在不同时间、不同光照下采集的视频存在显著外观差异，标准 3DGS 的球谐颜色表示难以处理；(2) 场景中的行人、车辆等动态物体会在 3DGS 训练中引入瞬态伪影。

\textbf{多序列高斯泼溅（MsGS）。}
UAVTwin 提出 MsGS（Multi-sequence Gaussian Splatting），为每段视频序列分配一个 32 维可学习嵌入 $q_i$，并通过一个 3 层颜色 MLP $f$ 将序列嵌入、视角方向嵌入、逐高斯嵌入 $h_i$ 和基色 $\hat{c}_i$ 映射为仿射变换参数 $(\alpha_i, \beta_i)$，输出调色后的颜色 $c_i^a = \alpha_i \hat{c}_i + \beta_i$。直觉上，球谐系数编码的是"这个高斯是什么颜色"，而仿射变换编码的是"在这个序列的光照条件下，颜色应该偏移多少"——这比球谐的显式颜色表示更适合处理大幅度的外观差异。关键优势在于：由于嵌入是\emph{逐序列}而非逐图像的，渲染新视角时只需选择一个序列嵌入即可，无需像逐图像方法那样依赖测试图像提取外观特征，因此能自由生成任意无人机轨迹下的图像。

\textbf{掩码精化模块。}
UAVTwin 用 GroundingSAM 提取动态物体掩码，但无人机视角下动态物体通常很小且稀疏，SAM 的分割质量有限。掩码精化模块的核心思路是：用实体分割获取额外的候选掩码，再根据每个候选掩码区域的平均光度误差判断其是否为动态物体——高误差区域（说明 3DGS 无法重建该区域，暗示是瞬态物体）被加入掩码，低误差区域则被移除。这种基于重建误差的筛选机制比 SAM 的语义分割更可靠，因为判断标准不再是"这是什么物体"，而是"这个区域在 3DGS 训练中是否产生了不一致"。

\textbf{数据生成管线。}
MsGS 训练完成后，通过 TSDF 融合提取背景网格 $B$，并将其导入 Blender：在网格上放置 10--15 个 SynBody 合成人体模型（SMPL-XL 骨架 + AMASS 动作库），通过 DiffusionLight 生成 HDR 环境光照，用 Cycles 渲染器渲染前景人体并添加运动模糊，最后通过 $\alpha$-混合与 MsGS 渲染的背景合成最终图像 $I_{comp} = \alpha G(I_{fg}) + (1-\alpha)I_{bg}$。管线支持平移、偏航旋转、环绕和变高四种无人机轨迹模式，并为每帧自动生成分割掩码、深度图和动作标签等像素级标注。在 Okutama-Action 数据集上，MsGS 相比 2DGS 和 PGSR 分别提升 PSNR 1.23 dB 和 2.79 dB；用合成数据微调 YOLOv8 后，人体检测 mAP50 提升 2.5\%--13.7\%。

VR-Robo（2025.02）~\cite{2502_01536} 将 SOUS VIDE 的 Real-to-Sim 管线从无人机扩展到四足机器人导航与运动——保留 3DGS 渲染作为视觉数据源，但新增 GS-mesh 混合表示解决物理交互问题：Gaussian 提供第一人称照片级真实感图像，TSDF 融合提取的 Mesh 提供精确碰撞检测，两者通过四点配准的齐次变换矩阵 $T_\text{homo}$ 实现坐标对齐。为支持鲁棒 RL 训练，VR-Robo 引入遮挡感知场景组合策略：分离重建的环境和物体高斯在自我中心视角下通过 z-buffer 实现正确的遮挡关系，并通过机器人位姿和锥桶位姿的随机采样大幅增加训练多样性。策略采用两级异步架构——高层策略（5 Hz）以冻结 ViT 编码 RGB 图像并输出速度指令，低层策略（50 Hz）基于本体感知执行关节控制——在 Isaac Sim 中用 PPO 训练，结合相机位姿噪声、图像增强和帧延迟等域随机化，在 Unitree Go2 四足机器人上实现纯 RGB 零样本仿真到真实迁移。

\paragraph{森林环境导航}

可重光照 3DGS（2026.02）~\cite{2602_07101} 在 SOUS VIDE 的基础上识别到一个关键瓶颈：标准 3DGS 的光照与几何固有耦合——球谐系数 $\mathbf{f}_i$ 同时编码了表面反照率和采集时的光照条件——导致策略在训练时只见过一种固定光照，迁移到不同时段或天气的户外场景时视觉特征崩溃。该方法的增量贡献是在 3DGS 表示层面解耦光照，使仿真器能够合成物理一致的多光照训练数据，从而迫使策略学习光照不变的几何特征。

\textbf{可重光照着色方程。}
标准 3DGS 中，每个高斯的颜色 $c_i$ 直接由球谐系数 $\mathbf{f}_i$ 计算，光照与材质不可分。可重光照 3DGS 将 $c_i$ 分解为三个独立因子：

\begin{equation}
\mathbf{c}_i = \boldsymbol{\rho}_i \odot \left(\sum_{m=1}^{(l+1)^2} \mathbf{L}_{env}^{m} \cdot \mathbf{O}_i^{m} \cdot \mathbf{d}_i^{m}\right)
\end{equation}

其中 $\boldsymbol{\rho}_i \in \mathbb{R}^3$ 为可学习的漫反射反照率（与光照无关的固有材质颜色），$\mathbf{L}_{env}^{m}$ 为全局环境光照的第 $m$ 个球谐系数（场景级共享变量），$\mathbf{O}_i^{m}$ 为遮挡系数（编码高斯中心处的方向性光线遮挡），$\mathbf{d}_i^{m}$ 为可学习的传递系数（建模法线朝向和间接反射等局部几何交互）。直觉上，这个公式说的是：每个高斯的最终颜色等于其"材质本色" $\boldsymbol{\rho}_i$ 乘以"有多少光到达这里"（$\mathbf{L}_{env}$ 经过遮挡 $\mathbf{O}_i$ 调制后的结果）。关键设计在于 $\mathbf{L}_{env}$ 是全局变量——修改它即可改变整个场景的光照，而几何 $\mathbf{O}_i$ 和材质 $\boldsymbol{\rho}_i$ 保持不变，这正是零样本迁移的基石。

\textbf{遮挡感知可见性建模。}
户外仿真中准确的阴影合成至关重要，但标准 splatting 框架通常忽略此问题。该方法引入预计算遮挡场来建模全局光传播：将场景空间离散化为均匀体素网格，在每个网格节点沿六个主轴渲染正交深度图以捕获周围几何，将二值可见性投影到球谐基上得到遮挡系数 $\mathbf{B}_k$，再通过三线性插值将离散探针值分配到连续的高斯中心位置：

\begin{equation}
\mathbf{O}_i = \sum_{k \in \mathcal{N}(\boldsymbol{\mu}_i)} w_k \cdot \beta_k \cdot \mathbf{B}_k
\end{equation}

其中 $w_k$ 为空间插值权重，$\beta_k$ 为背面剔除掩码（确保探针方向与高斯法线一致）。这个预计算步骤在训练前一次性完成，运行时只需查表插值，不增加渲染开销。

\textbf{HDR 光照先验与稳定解耦。}
联合优化几何与光照容易产生歧义——暗纹理可能被错误建模为阴影，阴影可能被烘焙进反照率。该方法采用深度全景光照（DPL）管线从低动态范围训练图像恢复 HDR 照明：先用预训练自编码器将 LDR 提升到 HDR，再投影到球谐系数 $\mathbf{L}_{lm}$。关键在于，这些系数在 3DGS 优化开始前即被冻结为强先验，迫使网络将光度残差归因于反照率和遮挡更新，而非调整光照来"作弊"。

\textbf{光度量域适应策略。}
在每个训练回合开始时，对全局光照系数 $\mathbf{L}_{env}$ 施加三类扰动：绕垂直轴旋转（模拟不同时段和太阳方位角）、强度缩放（模拟不同曝光和云量）、色偏着色（模拟不同天气，如黄昏暖色调、阴天冷色调）。同时保持几何 $\mathbf{O}_i$ 和材质 $\boldsymbol{\rho}_i$ 严格不变，迫使策略只能从几何结构中提取导航线索，而非依赖特定光照下的外观模式。

\textbf{端到端策略网络与训练课程。}
策略网络处理两个输入流：CNN 从单目 RGB 图像 $I_t$ 提取空间几何特征，MLP 编码本体感受状态 $\mathbf{s}_t = [\mathbf{p}_{rel}^T, \mathbf{v}^T, \psi, \dot{\psi}]^T$，GRU 融合两者捕获时序上下文，最终通过 actor-critic 头输出连续偏航率指令 $a_t \in [-u_{\max}, u_{\max}]$。训练采用两阶段课程：第一阶段在标准 3DGS 场景（固定光照）下建立基本几何理解；第二阶段激活可重光照 3DGS 引入随机光照变化，策略先经历短暂性能下降，随后恢复并超越第一阶段水平，仿真成功率从约 90\% 提升至超过 90\% 的目标阈值。

\textbf{自适应速度调度。}
为确保高速机动时的动力学稳定性，前向速度 $v_{tar}$ 与偏航率 $u$ 耦合：

\begin{equation}
v_{tar} = v_{\min} + (v_{base} - v_{\min}) \cdot \left(1 - \frac{|u|}{u_{\max}}\right)^{0.5}
\end{equation}

直觉上，直线飞行时速度最大化为 $v_{base}$，转弯时速度平滑衰减至 $v_{\min}$，防止急转弯时的侧滑。

\textbf{实验验证。}
在 10 个真实森林场景（每个约 $60\times 60$ 米）的广泛实验中，该方法在仅使用单目 RGB 相机（无 LiDAR、无深度相机）的轻量四旋翼平台上实现了最高 10 m/s 的鲁棒无碰撞导航。核心消融实验表明：在强光、阴天和黄昏三种自然光照条件下，无可重光照域适应的策略成功率分别为 6/10、9/10、3/10（黄昏几乎完全失败），加入可重光照后提升至 8/10、10/10、8/10。注意力可视化显示，策略自动学会聚焦于树干和树枝等障碍物边界，而非光照变化的区域，验证了光照不变特征学习。在 Table III 的定量对比中，该方法是唯一使用单目 RGB 传感器在户外环境达到 10 m/s 飞行速度的方法，且显式支持随机光照训练。

Liquid 网络（2024.06）~\cite{2406_15149} 沿用 SOUS VIDE"3DGS 仿真器 + 策略学习"的全栈范式，但在三个关键环节做了不同选择：仿真器层面，该方法使用 PyBullet 物理引擎（240 Hz）驱动四旋翼动力学，并将 3DGS 相机绑定到仿真位姿渲染照片级真实感图像——与 SOUS VIDE 的 FiGS 相比，构建流程更轻量（仅需 iPhone 采集约 430 张图像、15 分钟训练 GS 模型），但不含 SOUS VIDE 的领域随机化参数采样；数据层面，该方法设计了"隐式闭环增强"策略，将初始位姿设为目标出现在图像边缘的偏离位置，迫使网络从训练数据本身学会恢复轨迹，而非依赖 SOUS VIDE 的 MPC 专家示范；策略架构层面，该方法使用 Closed-form Continuous-time（CfC）Liquid 神经网络——其隐藏状态演化由显式含时间的微分方程闭式解驱动，可在任意时间间隔求值，因此支持"不规则时间采样增强"（训练时随机化帧间时间间隔至 1--10 Hz），这是 SOUS VIDE 的离散时间 SV-Net 无法做到的。

网络仅包含约 12k 参数（为 LSTM 基线的 1/25），以 CNN 提取图像特征后送入 34 个 Liquid 神经元，输出机体系三轴速度和偏航角速率。实验中，该方法在单个室内场景训练后，不仅在 GS-to-Real 室内测试中保持 75\% 成功率（LSTM 从 72\% 骤降至 12.5\%），更零样本泛化到完全未见的户外校园草坪环境（50\% 成功率），并在户外连续通过 5 个检查点完成"徒步导航"任务——这是 SOUS VIDE 未涉及的跨域泛化场景。

\paragraph{GPS 拒止环境}

与 Liquid 网络沿 SOUS VIDE"3DGS 仿真器 + 策略学习"全栈范式改进策略架构不同，Pan 等人（2024.08）~\cite{2408_06030} 彻底改变了 3DGS 在系统中的角色定位——不再将 3DGS 作为策略训练的仿真渲染器，而是将其放在下游：先用分层感知与规划系统自动识别检查目标并生成扫描路径，MAV 执行完数据采集后再用 3DGS 重建高保真结构模型。换言之，3DGS 在此前所有方法中服务于"训练"（生成训练数据供策略学习），而在这里服务于"重建"（从已采集的真实数据中还原三维结构）。

具体地，框架包含三个核心模块：\textbf{结构感知}从 LiDAR 点云地图 $\mathcal{M}_g$ 中用布料模拟滤波（CSF）分离地面/屋顶，结合 RANSAC 法向量拟合和欧氏聚类识别柱体 $\mathcal{O}_{\mathrm{column}}$ 与墙面 $\mathcal{O}_{\mathrm{wall}}$，为后续路径规划提供结构级先验；\textbf{实例感知扫描规划}根据相机 FOV 和拍摄距离 $D$，对柱体生成螺旋扫描路径（半径 $r_{\mathrm{spiral}}=r+D$，圈数由 $\lceil(h_{\max}-h_{\min})/\mathrm{FOV}_{\mathrm{height}}\rceil$ 决定），对墙面生成基于最小包围矩形的完全覆盖路径规划（CCPP）；\textbf{探索-执行策略}在正式扫描前先对每个结构实例执行探索飞行，通过射线投射体素建模更新局部任务地图，达到探索率阈值 $\tau$ 后再重规划碰撞-无障碍路径。MAV 定位采用 ESKF 融合 LiDAR（Fast-LIO2）和 IMU 测量，运动规划基于 B 样条轨迹优化（含碰撞惩罚 $J_c$、动态平滑性 $J_s$ 和可行性 $J_d$ 三项），控制器通过 PD 增益将期望加速度映射为姿态角和推力指令发送至 PX4 飞控。在 4000 m$^2$ 地下设施中，该框架实现 100\% 扫描路径生成和执行成功率，结构识别 F1 达 0.962，墙面提取覆盖率 85\%，轨迹跟踪 RMSE 在 1.0 m/s 扫描速度下为 0.122 m，3DGS 重建 PSNR 为 30.8 dB。

如果说 Pan 等人将 3DGS 从"策略训练数据生成器"重新定位为"下游重建工具"，那么视觉-环路仿真（2025.02）~\cite{2502_05409} 则进一步将 3DGS 仿真的用途从"训练任何网络"转向"验证已有算法"——不再训练策略网络或位姿估计网络，而是构建一个完整的视觉-环路测试框架，用 3DGS 渲染的照片级图像作为位姿估计算法的输入，在室内安全环境中验证深度单目位姿估计算法在海洋场景中的鲁棒性。这一角色转变的设计动机是现实的：海上实验需要动用研究船，成本高昂且可用窗口有限；而 3DGS 渲染的虚拟海洋环境可以在室内复现多种光照条件下的真实视觉外观，使算法验证不再依赖高代价的实地测试。

具体地，该方法用 GoPro 搭载八旋翼无人机，在切萨皮克湾对美国海军学院研究船 YP689 的飞行甲板及周边结构进行多日、多时段、多船航向的图像采集（0.5 秒间隔定时拍摄，5546$\times$4851 像素），将不同光照条件的图像集分别训练为多个 3DGS 虚拟海洋环境。在仿真框架中，Vicon 动捕系统以 200 Hz 追踪室内无人机的位姿 $(x, R) \in \mathbb{R}^3 \times \mathsf{SO(3)}$，将该位姿映射到 3DGS 虚拟环境中渲染对应视角的照片级图像（640$\times$640 分辨率，75--110 FPS），再由 TNN-MO（Transformer Neural Network Multi-Object）模型从渲染图像实时估计无人机相对于船甲板的 6D 位姿。与 SOUS VIDE 的 FiGS 仿真器直接驱动无人机动力学不同，本框架的渲染图像不用于策略训练，而是作为位姿估计算法的输入——通过将估计位姿与 Vicon 真值对比，在室内即可全面验证飞控软件、硬件和深度位姿估计方案。在两个不同虚拟环境（晴天/阴天）的飞行测试中，TNN-MO 位姿估计的位置 MAE 为 0.089--0.105 m（约占最大作业距离 12 m 的 0.8\%），姿态 MAE 为 1.54--1.78$^\circ$，与真实海洋环境测试结果一致，验证了该框架替代高风险、高成本海上实验的可行性。

\paragraph{多机器人协同重建}

语义引导的多机器人 3DGS 重建（2024.12）~\cite{2412_02249} 沿用 ATLAS Navigator 的"3DGS 在线建图 + 语义引导规划"范式，但将整个系统从单机器人场景扩展到集中式多机器人协同，并在任务生成环节引入实例级语义不确定性约束以缓解局部最优问题。

\textbf{设计动机}。ATLAS Navigator 展示了语言嵌入 3DGS 地图在单机器人导航中的威力，但面临两个扩展性瓶颈：其一，单机器人扫描效率在大规模室内环境中严重不足——一个 187 m$^2$ 的场景需要反复遍历；其二，ATLAS Navigator 基于表面不确定性的重建任务容易使机器人陷入局部最优（高不确定性点的散乱分布导致冗余视点爆炸，机器人在小区域来回移动而无法覆盖全局）。与此同时，已有的多机器人重建方法（如 ACAMS 的 cluster-and-assign、MS3DSR 的贪心聚类分配）在模式与任务分配时需要反复求解 TSP，任务数量增长后规划时间急剧上升。该方法的核心思想是：用开放词汇语义分割将散乱的表面不确定性约束到实例级，使视角采样聚焦于真正不完整的物体而非无差别的高不确定性区域；同时用全局-局部双层模式分配和改进 K-means 任务聚类替代迭代式 TSP，在保证负载均衡的前提下大幅降低规划开销。

\textbf{感知管线}。系统维护三种互补表示：全局 3DGS 地图（基于 Gaussian-SLAM 在线训练，不使用子图策略以加速质量评估）、ASH 并行空间哈希用于场景几何获取和噪声高斯裁剪（距 ASH 表面点超过 0.05 m 的高斯被判定为噪声并移除）、以及 OVIR-3D 开放词汇 3D 实例检索用于在线语义分割。表面不确定性通过损失缓存策略计算——将每帧的投影损失 $L_\text{proj} = L_\text{color} + \lambda_d \cdot L_\text{depth}$（$\lambda_d=0.5$）从图像空间投影到 0.05 m 分辨率的占据哈希体素中，缓存至体素属性中用于后续查询。

\textbf{语义引导的重建任务生成}。与 ATLAS Navigator 在每个高斯上附加 CLIP 语义特征不同，该方法在更高粒度上使用语义信息——通过 OVIR-3D 将多视角 2D 区域提案融合为 3D 实例分割，再结合 ScanNet200 词汇表计算每个实例的物体性分数（objectness score）：对实例 $O_i$ 的 3D 特征 $f_i^{3D}$ 与 200 个类别特征计算相似度 $\Lambda_i$，经缩放因子 $\lambda_e=50$ 放大后 softmax 归一化为分类概率 $C_i = \text{softmax}(\lambda_e \Lambda_i)$。最高概率 $C_i^1$ 即为物体性分数——直觉上，若一个实例被高置信度地识别为某个已知类别（如"椅子"），则它已被充分重建；若置信度很低（可能因 CLIP 标签噪声），则不可靠。系统仅保留 $C_\text{min}=0.2 < C_i^1 < C_\text{max}=0.6$ 的实例作为"重建实例"，聚焦于那些确实存在但尚未完全重建的物体。

对每个重建实例，系统查询其表面点的不确定性 $U_i^\text{rec}$（从全局表面不确定性中提取），以 $N_\text{down}=5$ 倍下采样后提取兴趣点（POI）：贪心地选择不确定性最高的点，要求相邻 POI 间距至少 $d_\text{POI}=1.2$ m，避免冗余采样。信息增益定义为可见实例表面点不确定性的距离衰减加权和：

\begin{equation}
g(v) = \sum_{k=1}^{N_\text{vis}} \sigma_k \, e^{-0.5 \, d_{v,k}}
\end{equation}

其中 $\sigma_k \in U_i^d$ 为可见表面点的不确定性，$d_{v,k}$ 为视点到表面点的距离。对每个 POI，选择信息增益最大的候选视点作为该实例的重建任务。这种实例级约束将散乱的不确定性聚合为以物体为单位的紧凑任务，有效减少了冗余视点生成。

\textbf{层次化多机器人协同规划}。系统将任务分为探索任务（前沿覆盖）和重建任务（语义引导），采用三层层次化规划：

第一层是\textbf{模式分配}。根据全局探索与重建任务的比例确定各模式的机器人数量，再以局部任务统计避免迭代优化：对每个机器人 $R_r$，统计局部范围 $d_\text{local}=6$ m 内的探索任务数 $N_{\text{local},r}^\text{exp}$ 和重建任务数 $N_{\text{local},r}^\text{rec}$，计算差值 $N_{\text{local},r}^\text{ove} = N_{\text{local},r}^\text{exp} - N_{\text{local},r}^\text{rec}$，按差值排序确定优先分配探索模式的机器人——局部探索任务越多的机器人越适合继续探索。与 ACAMS 需要迭代求解 TSP 的模式分配不同，该方法仅需一次排序即可确定所有机器人模式。

第二层是\textbf{改进 K-means 任务聚类与分配}。对每种模式 $M \in \{\text{exp}, \text{rec}\}$ 内的任务，用改进 K-means 聚类到对应模式的机器人，优化目标同时最小化移动代价（任务到聚类中心的欧氏距离 + 机器人到聚类中心的距离）和负载不均衡（任务数偏差 $(N_r^M - \bar{N})^2$ + 内部距离偏差 $|D_r - \bar{D}|$）。与 ACAMS 和 MS3DSR 需要反复计算 TSP 距离不同，该聚类仅基于欧氏距离，消除了 TSP 的迭代计算开销。

第三层是\textbf{单机器人视点路径规划}。对每台机器人的任务集，用 ATSP 求解器求解最短路径，再经 B 样条轨迹优化生成平滑、安全的导航轨迹，以 $t=0.4$ s 间隔沿轨迹采样视点，执行路径长度限制为 $L_\text{exec}=6$ m。

\textbf{实验验证}。在三个虚拟场景（QaLdn 187 m$^2$、Nicut 120 m$^2$、Oyens 80 m$^2$）中使用四台机器人的实验表明：该方法在所有规划方法中取得最高重建质量，且规划时间显著优于 ACAMS 和 MS3DSR（消融实验显示，去除模式分配后探索效率下降导致局部最优，去除负载均衡约束后部分机器人执行路径过长需截断而损失覆盖率）。在真实场景中，两台 Turtlebot（Azure Kinect + Realsense T265）在 12 m $\times$ 8 m $\times$ 4 m 的展厅中约 8 分钟完成探索与重建。

\paragraph{自主驾驶}

GSDrive（2026.04）~\cite{2604_28111} 在 3DGS 仿真 RL 框架的奖励信号环节做出根本性改进——用多模态轨迹探测（trajectory probing）替代稀疏碰撞奖励，将「事后惩罚」转化为「前瞻密集反馈」。与 RAD 等方法仅在碰撞发生时才给出显著信号（但碰撞通常是连续多步错误决策的最终结果，无法将信用合理归因到早期动作）不同，GSDrive 的核心思路是：让策略在每一步决策时，预测 $K$ 条可能的未来轨迹并在 3DGS 环境中前向仿真，取其中最优结果作为当前动作的奖励——策略不再被动等待灾难发生，而是主动预见未来后果并提前修正。

\paragraph{轨迹探测奖励驾驶}

具体地，GSDrive 采用 IL-RL 双阶段训练。IL 阶段用最优传输（OT）引导的 flow matching 训练多模态轨迹预测器：从 3DGS 环境渲染的六相机环绕视图经 ResNet 提取语义特征、压缩为鸟瞰图（BEV）空间特征后，通过交叉注意力融合为查询向量，驱动 flow matching 头预测 $N_{\text{modes}}$ 条轨迹及其模式概率——该预测器本质上是一个「未来奖励探测器」，为 RL 阶段的前瞻性评估做准备。RL 阶段的核心是轨迹探测奖励：策略输出动作后，轨迹预测器生成 $K$ 条候选轨迹，每条在 3DGS 环境中前向 rollout $H$ 步，计算累积物理奖励（包含存活奖励、进度奖励、碰撞惩罚和舒适性惩罚），当前动作获得的探测奖励为 $r_{t}^{\text{probe}} = \max_{i=1\ldots K}[\sum_{h=0}^{H}\gamma^{h}r^{\text{env}}(s_{t+h},a_{t+h})]$——直觉上，策略同时审视 $K$ 条可能的未来，以最乐观的物理评估作为当前决策的奖励信号。为防止 RL 探索偏离 IL 学到的安全先验过远，GSDrive 还引入自适应 KL 散度正则化（EMA 跟踪 KL 偏移量，动态调整惩罚系数），在「探索超越专家演示」与「保持安全归纳偏置」之间取得平衡。在 nuScenes 重建数据集的闭环评估中，GSDrive 的回合奖励达 52.97、碰撞率降至 0.11，超越 RAD（ER 49.24、CR 0.19）等所有仿真 RL 基线。

\paragraph{稀疏奖励闭环驾驶}

与 GSDrive 用多模态轨迹探测将稀疏奖励转化为密集前瞻反馈不同，RAD（2025.02）~\cite{2502_13144} 是首个基于 3DGS 闭环环境的端到端自动驾驶 RL 框架，选择直接面对稀疏奖励挑战，通过解耦动作空间与方向性辅助目标来缓解收敛困难，并将 IL 作为正则化项嵌入 RL 训练循环，使策略在探索安全边界的同时保持人类驾驶行为的基本合理性。

\textbf{动机与问题}。IL 训练的端到端驾驶策略面临两个根本性缺陷：（1）因果混淆——IL 从专家演示中学到的只是观测与动作的相关性而非因果关系，策略可能退化为从历史轨迹外推未来轨迹的捷径学习；（2）开环-闭环鸿沟——IL 以开环方式训练，但实际驾驶是闭环过程，单步轨迹误差随时间累积，将策略推入分布外状态后无法恢复。闭环 RL 训练是自然的解决方案，但需要能与策略交互的驾驶环境。传统仿真器（CARLA 等）基于游戏引擎，传感器模拟不够真实；RAD 的核心思路是：利用 3DGS 从真实驾驶数据重建照片级数字孪生，使策略在大规模试错中学习处理分布外场景，同时通过专门设计的安全奖励和 IL 正则化保证训练有效性与人类对齐。

\textbf{三阶段训练范式}。RAD 采用渐进式训练策略：\textbf{感知预训练}阶段，BEV 编码器将多视角图像特征转换至鸟瞰图空间，地图头（map head）和智能体头（agent head）在真值标签监督下学习矢量化地图要素（车道线、路沿、交通信号等）和动态障碍物运动信息，使实例级 token 隐式编码对应高层语义；\textbf{规划预训练}阶段，冻结感知模块，仅用大规模专家驾驶演示通过双焦点损失（dual focal loss）优化图像编码器和规划头，将动作概率分布初始化为接近人类驾驶的安全先验，避免 RL 训练的冷启动不稳定；\textbf{强化后训练}阶段，RL 与 IL 交替优化策略——RL 通过 PPO 从 3DGS 环境的闭环交互数据中学习安全关键事件的因果响应，IL 作为正则化项约束策略不偏离人类行为过远。训练中仅更新图像编码器和规划头，BEV 编码器、地图头和智能体头始终保持冻结，避免感知与规划目标的梯度冲突。

\textbf{解耦离散动作空间}。RAD 将驾驶动作解耦为横向和纵向两个独立维度，每个维度在 0.5 秒时间窗内定义离散位移：横向 $a^x \in \{d^x_{\min}, \ldots, 0, \ldots, d^x_{\max}\}$（61 个选项，范围 $[-0.75, 0.75]$ m），纵向 $a^y \in \{0, \ldots, d^y_{\max}\}$（61 个选项，范围 $[0, 15]$ m），假设恒线速度和恒角速度。这种解耦将动作空间维度从联合空间的 $61 \times 61 = 3721$ 降为两个独立的 61 分类问题，显著降低探索成本。策略网络输出两个独立的 softmax 分布 $\pi(a^x|s)$ 和 $\pi(a^y|s)$，以及对应的价值函数 $V_x(s)$ 和 $V_y(s)$。在 3DGS 环境中，策略输出动作 $(a^x_t, a^y_t)$ 后，通过运动学自行车模型（kinematic bicycle model）更新自车位姿，3DGS 环境根据新位姿渲染多视角图像作为下一步输入。

\textbf{奖励建模与辅助目标}。奖励函数由四项组成：动态碰撞惩罚 $r_{\text{dc}}$（自车包围盒与标注的动态障碍物包围盒重叠时触发）、静态碰撞惩罚 $r_{\text{sc}}$（与 3DGS 高斯基元的静态障碍物碰撞）、位置偏差惩罚 $r_{\text{pd}}$（偏离专家轨迹超过阈值 $d_{\max}$）、航向偏差惩罚 $r_{\text{hd}}$（航向角偏离专家轨迹匹配航向超过 $\psi_{\max}$）。任一事件触发即终止回合——因为碰撞或严重偏差后 3DGS 环境通常生成噪声传感器数据，不利于后续 RL 训练。奖励按维度分配：横向奖励 $r_t^x = r_t^{\text{sc}} + r_t^{\text{pd}} + r_t^{\text{hd}}$，纵向奖励 $r_t^y = r_t^{\text{dc}}$，对应地使用 GAE 分别计算横向和纵向优势估计 $\hat{A}_t^x$ 和 $\hat{A}_t^y$。为缓解稀疏奖励导致的收敛缓慢，RAD 设计了四项方向性辅助目标：动态碰撞辅助损失根据碰撞位于前方还是后方引导减速/加速；静态碰撞辅助损失根据障碍物在左侧还是右侧引导避让转向；位置偏差和航向偏差辅助损失分别引导横向修正。这些辅助目标将动作概率按方向分解（减速/加速概率、左转/右转概率），根据实时碰撞风险和偏差方向自适应调整，提供比稀疏奖励更密集的梯度信号。

\textbf{3DGS 环境构建}。RAD 从 2000 小时专家驾驶数据中筛选出 4305 个涵盖多样道路类型、交通密度和智能体行为的真实驾驶场景，基于 StreetGaussian 框架对每个场景独立训练 3DGS 模型，并针对闭环训练需求进行优化：路面用网格建模约束高斯几何以确保任意视角的精确路面渲染；天空单独建模防止与前景混淆；前景物体（车辆、行人）在训练中优化其位姿，并以深度和法线一致性作为监督信号增强几何重建精度。每个场景提取固定长度 8 秒的片段，其中 3968 个用于 RL 训练，337 个作为未见过的闭环评估基准。训练时部署 $N$ 个并行 worker，每个 worker 随机采样一个 3DGS 环境执行 rollout，将轨迹数据 $(s_t, a_t, r_{t+1}, s_{t+1}, \ldots)$ 存入共享缓冲区，RL（PPO）与 IL 步骤交替进行（4:1 比例），定期同步策略参数以缓解分布漂移。

\textbf{实验结果}。在 337 个未见 3DGS 环境的闭环评估中，RAD 的碰撞率（CR）比 IL 方法降低 3 倍，在大多数闭环指标上优于 IL 基线。消融实验验证了各设计选择的贡献：纯 IL 偏差低但碰撞率高（CR 0.229），纯 RL 安全性提升但偏离专家行为（ADD 从 0.238 增至 0.345），IL+RL 混合策略取得最佳平衡（CR 0.089, ADD 0.257）。移除动态碰撞奖励项后碰撞率最高（CR 0.238），表明纵向碰撞规避信号对安全性最为关键。一致性分析显示，同一策略在 3DGS 环境和真实环境中的驾驶行为高度一致，验证了 3DGS 闭环评估反映真实世界驾驶性能的有效性。

\paragraph{VLM-规划双系统驾驶}

在 RAD 的基础上，Senna-2（2026.03）~\cite{2603_11219} 沿袭 RAD「3DGS 闭环训练」的仿真范式，但对齐目标从 RAD 的「开环训练 vs 闭环部署」转向「VLM 高层决策 vs E2E 低层规划」的双系统一致性。具体地，Senna-2 采用 VLM（Qwen2.5-VL-3B）生成结构化元动作决策（速度控制：加速/减速/保持/停车；方向控制：直行/左转/右转/变道），经决策适配器（Decision Adapter）融合 VLM 隐状态与可学习类别嵌入为条件特征，通过 AdaLN 注入 DiT 扩散规划器生成轨迹——关键问题在于：若 VLM 输出「减速」而规划器生成的轨迹仍在加速，高层决策对低层规划便失去引导力。

为此，Senna-2 设计三阶段一致性训练范式：第一阶段驾驶预训练分别优化 VLM（QA 监督）和 E2E（扩散去噪损失），再用适配器桥接两者；第二阶段开环对齐用运动学映射函数检测决策-规划不一致样本并选择性修正，双系统一致性 F1 提升 $9.0\%$；第三阶段闭环对齐从 $1{,}300$ 个高风险驾驶片段重建 3DGS 环境，以分层强化学习（HRL）协同优化——底层规划器先在 3DGS 环境中用安全奖励（TTC $< 3\,\text{s}$ 时纵向缩放惩罚）和效率奖励（速度低于导航限速时纵向延伸鼓励）优化轨迹，随后高层 VLM 通过运动学映射函数将优化后轨迹映射回决策类别，若与原始决策不一致则惩罚 VLM 输出分布，确保决策-规划闭环对齐。实验表明，Senna-2 的双系统一致性 F1 较前代 Senna 提升 $19.3\%$，开环规划的最终位移误差（FDE）降低 $5.7\%$，闭环碰撞率（AF-CR）在 3DGS 评估基准上较 RAD 降低 $30.6\%$（从 $0.113$ 降至 $0.077$），验证了「高层决策-低层规划」双系统一致性训练对驾驶安全与效率的协同增益。

\paragraph{视点鲁棒驾驶}

沿袭 3DGS 闭环训练范式，VR-Drive（2025.10）~\cite{2510_23205} 聚焦端到端自动驾驶（E2E-AD）中的相机视点鲁棒性问题——不同车型的相机安装位置与角度各异，导致模型在训练分布外的视点下规划性能严重退化。与 RAD、GSDrive 等 3DGS 主要用于构建仿真环境的思路不同，VR-Drive 将 3DGS 的角色定位为训练时的在线数据增强工具：将前馈 3DGS 场景重建作为辅助任务集成到 E2E-AD 框架中，从稀疏视角以逐像素方式预测高斯基元参数（深度估计推断位置 $\boldsymbol{\mu}$，卷积网络预测缩放 $\mathbf{s}$、旋转 $\mathbf{r}$、不透明度 $\alpha$ 和颜色 $\mathbf{c}$），再通过随机采样的相机外参渲染任意新视点的特征图，无需逐场景优化即可在线合成多样化的训练数据。

为解决合成视点中的特征一致性问题，VR-Drive 引入两个互补模块：\textbf{视点混合记忆库}（viewpoint-mixed memory bank）存储并持续更新来自不同视点的 3D 实例特征，通过交叉注意力实现跨视点的特征交互，使模型学习规范化的 3D 表示；\textbf{视点一致蒸馏}（viewpoint-consistent distillation）以原始视点的高置信度实例特征为教师信号（stop-gradient），通过可学习的 3D 采样点从多相机视图中聚合特征，用 MSE 损失引导新视点特征的语义对齐，防止渲染伪影对下游规划的污染。整个框架以端到端方式联合训练 3D 重建、感知与规划三个分支，循环重建损失（从新视点重建原始视点）确保高斯基元的 3D 一致性。在 nuScenes 和 CARLA 的 5 种未见视点变化（$\pm$俯仰角、$\pm$高度、深度偏移）下，VR-Drive 保持与原始视点相当的规划性能，而 DiffusionDrive 等基线方法的碰撞率在新视点下显著上升。

此外，VR-Drive 发布了首个面向 E2E-AD 视点鲁棒性评估的基准数据集，涵盖多种相机配置下的驾驶场景，为该方向提供了标准化的评测基准（NeurIPS 2025 接收）。

\paragraph{硬件在环驾驶测试}

从策略训练转向硬件测试，VizFlyt（2025.03）~\cite{2503_22876} 将 NavGSim 的"3DGS 场景渲染"思路从策略训练转向硬件在环（Hardware-In-The-Loop, HITL）教学测试——不再用 3DGS 渲染驱动策略网络学习，而是用 3DGS 作为"幻觉传感器"：利用外部动捕系统（Vicon）获取无人机实时位姿，从预训练的 3DGS 场景模型中合成对应视角的照片级 RGB/RGBD 图像流，替代真实机载相机，系统更新率超过 100 Hz。具体地，框架以 ROS2 节点将位姿映射为 3DGS 渲染输出，渲染 PC（RTX 3080）将 RGBD 帧发布至嵌入式板载计算机（Jetson Orin Nano）以模拟真实硬件约束，后者计算控制指令后通过 MAVLink 协议发送至飞行器，形成完整的感知-规划-控制闭环。碰撞检测方面，与 NavGSim 的解析高斯切片方法不同，VizFlyt 将 3DGS 环境导出为点云，构建 3D 占据体素网格，以固定半径球体检测碰撞并自动触发降落。在森林环境的高速避障实验中达到 $5.3\,\text{m/s}$，视觉里程计（DPVO）在幻觉图像上的 ATE 仅为 0.12 m，与真实相机结果一致，验证了 3DGS 幻觉传感器替代真实传感器的有效性。VizFlyt 基于 WPI 多年开设动手航空机器人课程的经验，进一步提出配套的开源开放硬件课程大纲——学生可在安全受控的室内环境中，通过 HITL 框架测试自主感知与控制算法，涵盖视觉里程计、深度估计、位姿估计和路径规划等课程项目，将高风险的户外飞行实验转化为零风险的室内 3DGS 仿真测试（ICRA 2025 接收）。

\paragraph{策略动员}

Mobi-$\pi$（2025.05）~\cite{2505_23692} 将 3DGS 的角色从 RL 训练时的\textit{状态表示}（如 GSRL）拓展到部署时的\textit{位姿搜索引擎}——利用 3DGS 的新视角合成能力，在未见环境中为已训练的操控策略搜索最优机器人基座位姿，从而形式化并求解「策略动员」问题。

\textbf{问题背景与动机}。大多数视觉运动策略仅从有限的机器人基座位姿和相机视点收集数据训练，当部署到移动平台上时，基座位姿的任意变化导致视觉输入和工作空间可达性同时漂出训练分布，策略性能急剧下降。传统解决方案——用大量多视角数据重新训练策略以增强鲁棒性——数据采集成本高昂；或将导航与操控串联、用语言指令或物体位置指定导航目标——但导航目标不考虑策略的视点偏好，到达后策略仍可能处于分布外状态。Mobi-$\pi$ 提出第三条路径：保持操控策略不变，通过优化基座位姿使策略运行在分布内视点上，从而解耦导航与操控策略训练。

\textbf{核心方法}。给定一个预训练操控策略 $\pi$（从固定基座位姿的演示数据训练），以及一个包含目标物体的新场景，Mobi-$\pi$ 的目标是在 $SE(2)$ 空间中搜索最优基座位姿 $p^{*} = \arg\max_{p} J(\pi | (s^{\text{env}}_0, p, q_0))$，使策略的期望回报最大化。方法分为三个步骤：

（1）\textit{场景重建}。机器人进入新场景后，通过驱动环视采集约 1000 张 RGB-D 图像（耗时约 5 分钟），训练一个 3DGS 场景模型。3DGS 的关键优势在于能够从任意候选相机位姿高质量合成新视角——这为后续的位姿评估提供了「预览」能力：无需实际移动机器人到候选位姿，只需在 3DGS 中渲染该位姿对应的相机视角即可评估其适用性。

（2）\textit{混合评分函数}。对每个候选基座位姿 $p$，通过变换 $v = f_{p \to v}(p)$ 获得对应相机位姿，在 3DGS 中渲染图像 $I = \text{GS}(v)$，然后用三重条件评分 $K(p)$ 评估位姿质量。三重条件为：分布内得分 $K_{\text{id}}(p)$——将渲染图像的 DINO 密集描述子与训练演示初始帧的 DINO 特征进行 KNN 匹配，计算平均 L2 距离作为分布内度量，直觉上衡量「从此位姿看到的场景与策略训练时看到的有多相似」；物体可见性 $K_{\text{obj}}(p)$——用视觉语言模型 MiniCPM-v2 检测渲染图像中是否包含目标物体，防止评分函数对「纹理相似但缺少任务物体」的视角给出高分；碰撞检测 $K_{\text{col}}(p)$——查询从深度数据构建的占据图，验证候选位姿是否无碰撞。最终得分 $K(p) = K_{\text{id}}(p)$ 当且仅当 $K_{\text{obj}}(p)=1$ 且 $K_{\text{col}}(p)=1$，否则为 0。

（3）\textit{贝叶斯优化搜索}。由于遮挡、物体可见性变化等因素使评分函数景观非平滑，梯度优化不可靠，Mobi-$\pi$ 采用贝叶斯优化（BO）这一无梯度的采样优化方法：先随机采样 $N^{(\text{BO})}_{\text{init}}$ 个基座位姿并评估得分，再用高斯过程代理模型和 UCB 采集函数迭代提出新候选，经 $N^{(\text{BO})}_{\text{iter}}$ 轮优化后选择最高得分的位姿执行策略。整个过程在 NVIDIA RTX 4090 上约 6 分钟完成。

\textbf{实验验证}。在基于 RoboCasa 的仿真基准（5 个操控任务、10 个未见场景）中，Mobi-$\pi$ 的任务成功率在 5 个任务中的 4 个达到或超过策略在 5 cm 扰动下的表现，显著优于非策略感知基线（LeLaN、VLFM）和策略感知基线（BC w/ Nav，使用 5 倍训练数据仍泛化不良）。在真实杂货店环境的三项任务（拾取薯片、双手抓取纸巾、倒番茄）中，Mobi-$\pi$ 全部实现 100\% 导航成功率，完全无碰撞，且全任务成功率分别达 70\%、100\%、80\%，超越人类直觉基线和端到端基线。此外，该方法支持零样本迁移到未见场景布局、多技能链式执行，以及扩展到 6$\times$10 m 大规模场景。

\paragraph{从单目到碰撞避免}

以上方法（SAFER-Splat、碰撞锥 CBF、SPLANNING、FOCI 等）均将 3DGS 高斯基元的几何信息显式地编码为碰撞约束——无论是欧氏距离、碰撞锥还是概率上界，安全决策都依赖于对高斯基元的直接几何查询。Tordjman-Levavasseur 和 Caron（2025.04）~\cite{2504_06651} 探索了一个截然不同的方向：不查询高斯基元的任何几何属性，而是利用 2DGS 的照片级渲染能力生成训练数据，从单目 RGB 图像中\textbf{隐式学习}碰撞避免策略。与 SAFER-Splat 在已建好的 3DGS 地图上实时计算椭球距离不同，该方法在训练阶段一次性消费 2DGS 渲染的图像来训练策略网络，部署时策略仅依赖单目相机输入，无需加载或查询任何高斯场景模型——这是一个从「在线几何推理」到「离线视觉学习」的根本性范式转移：3DGS/2DGS 的角色从部署时的实时查询对象变为训练时的数据生成器，场景模型在训练完成后即被丢弃。

\textbf{设计动机。}显式碰撞检测方法（如 SAFER-Splat 的 CBF）在安全过滤中表现高效，但它们依赖精确的状态估计——机器人在 3DGS 地图中的位置必须被准确估计，否则安全约束将建立在错误的空间关系上。作者指出，即使环境重建达到超越当前最佳水平一个数量级的精度（厘米级精度、500 Hz 更新率），若状态估计不随之改善，运动性能也不会实质提升——显式方法将高精度要求集中在少数低维状态坐标上，而隐式方法则利用高维隐编码的冗余性分散了对单个坐标精度的依赖。该方法的出发点正是绕过这一瓶颈：与其先建图再在地图上做几何推理，不如让策略直接从视觉观测中学习碰撞避免行为，使安全决策不依赖于状态估计精度。此外，相比于已有隐式方法使用深度图像（深度传感器昂贵且户外受限）或 NeRF（训练缓慢，渲染低于 1 Hz），该方法选择 2D Gaussian Splatting 作为场景表示——2DGS 相比 3DGS 引入法向量一致性和深度失真正则化，能更准确地重建表面几何，同时保持实时渲染速度（$>$100 Hz），直接转化为更快的策略训练。

\textbf{训练管线。}整个方法分为四个阶段，其中视觉编码器与 RL 策略采用\textbf{解耦训练}设计——先独立训练视觉编码器完成 RGB 到深度的映射，再固定编码器训练 RL 策略。这一选择基于两个考虑：便于独立调试视觉模块是否成功收敛（端到端训练中视觉故障与策略故障难以分离），以及适配部署平台 Upkie 轮式双足机器人无专用 GPU 的算力约束——32 维隐编码足以在无加速器的嵌入式平台上实时推理。第一，\textbf{场景采集与 2DGS 重建}：用 GoPro（快门速度 $\geq$ 1/125 s 以消除运动模糊——作者的经验是运动模糊是 3DGS/2DGS 训练失败的首要原因，宽镜头 16 mm）拍摄约 5 分钟的室内视频，通过 FFMPEG 提取关键帧（约 500 张 4K 图像），COLMAP 估计相机位姿后训练 2DGS 模型，并利用已知相机位置进行地理配准以对齐重力方向和尺度。2DGS 训练完成后，提取场景 3D 网格并用碰撞感知凸分解（CoACD）将网格分解为凸子部分以支持仿真中的碰撞检测——纹理较少的表面（如地板）质量较差，需手动替换为平面几何以确保仿真精度。

第二，\textbf{视觉编码器}：在 2DGS 仿真器中采集 65,000 张 128$\times$128 像素的 RGB-深度图像对（数据采集约 20 分钟，深度图由 2DGS 渲染器免费提供），其中 60,000 张用于训练、5,000 张用于测试。训练一个自编码器式 CNN（4 层卷积 + 批归一化）将 RGB 图像映射到深度图的对数空间——选择对数空间是因为近距离深度值对碰撞避免最为关键，对数变换压缩了大深度值的变化范围，使网络更关注近处障碍。编码器输出 32 维隐编码，训练集 MSE 为 0.011、测试集 MSE 为 0.027，权衡了表达力与部署机器人的算力约束。第三，\textbf{RL 碰撞避免策略}：在纯导航环境中用 CrossQ 算法训练 MLP 策略（策略网络 2 层 256 维，评论家网络 2 层 1024 维，均含批归一化）。CrossQ 的选择有明确的动机：由于每步环境交互都需要 2DGS 渲染图像（远比纯物理仿真昂贵），CrossQ 通过在评论家中使用批归一化并移除目标网络，以更低的更新-数据比（UTD）实现高效样本利用，从而减少渲染开销。

观测空间为摇杆输入 + 当前速度 + 视觉隐编码（无历史帧、无循环网络），动作空间为速度修正量。奖励函数设计为：无碰撞时 $R = 1 - \|a\|_1$（鼓励最小修正），碰撞时 $R = -100$（强烈惩罚）。机器人距障碍物 $<2$ cm 即判定碰撞终止；代理在场景中随机初始化（距障碍 $>50$ cm），遵循简化的运动学模型并添加高斯噪声以模拟真实动力学。为增强视觉鲁棒性，训练时施加多种数据增强（对比度、亮度、色调、饱和度及噪声扰动），并用零均值高斯噪声扰动图像俯仰角（约束在 $[-\pi/4, \pi/4]$ 弧度内）以模拟轮式倒立摆运动时的自然倾斜——这些增强使策略在部署时对光照变化和机器人姿态抖动更加鲁棒。训练 500,000 步，在 RTX 4070 Ti SUPER 上耗时约 2.5 小时。

第四，\textbf{底层运动控制}：修正后的速度指令 $(v_x, v_\theta)$ 送入基于线性化倒立摆动力学的模型预测控制器（MPC），以二次规划（QP）形式求解（PROXQP 求解器，支持热启动），100 Hz 运行，偏航速率 $v_\theta$ 通过差分驱动模型直接映射为左右轮速度偏移，前向速度 $v_x$ 积分为参考轨迹后由 MPC 跟踪，输出的车轮速度指令发送至执行器。

\textbf{实验结果。}在 Upkie 轮式双足机器人（搭载 Luxonis OAK-D Lite RGB 相机，640$\times$480 像素裁剪降采样至 128$\times$128）上的实验从三个层面评估策略性能。在仿真纯导航环境中（摇杆全速前进，最长 20 秒），视觉策略平均存活 15 秒（100 次试验），无修正时仅 4 秒；仿真中含底层运动控制时障碍清除率 100\%。在训练数据采集房间的真机实验中，家具位置已改变（墙壁和地板颜色不变），从 5 个不同起始位置各执行 10 次试验（摇杆全速直行，目标为居中放置的沙发），成功率从位置 1 的 90\% 到位置 4 的 40\% 不等，位置 5 成功率为 0\%（策略面对障碍未做显著修正）。在分布外测试中，同楼层的办公走廊环境中策略成功导航窄通道并动态调整轨迹；户外环境中，策略对大型不可移动障碍（墙壁、停放的汽车）表现出谨慎的拒绝前进行为，但对视觉结构复杂的物体（自行车、栅栏门等精细或半透明障碍）未能正确反应。这些结果表明，仅凭单目 RGB 训练的碰撞避免策略在室内结构化环境中具有实用价值，但视觉泛化能力仍是主要瓶颈。

\begin{takeaway}
导航与路径规划展示了 3DGS 在机器人导航中的双重角色：显式方法将高斯基元的几何信息编码为碰撞约束实现安全规划，隐式方法则利用渲染能力生成训练数据，实现从在线几何推理到离线视觉学习的范式转移。
\end{takeaway}

% ======================================================================

\subsection{近期方法脉络：语义导航与规划、主动建图与NBV规划、Object-centric lifelong Gaussian maps等}

这一组近期工作延续 GSFusion、3DGS map、3DGS mapping，主要把改动集中在扩展、应用、任务条件化、增量改进等环节。下面只保留相对既有路线的变化，便于和前文方法直接对照。

\paragraph{LiftNav}
LiftNav（2026.05）~\cite{2605_31376} 延续 GSFusion 在“语义导航与规划”上的基本设定，主要改变是：LiftNav：基于GSFusion的TSDF+GS双地图，结合YOLO检测、TSDF三维提升和B样条轨迹优化。3DGS 的几何-渲染统一表示保持为底层地图。
其设计动机是：提出混合导航框架，在TSDF+GS双表示上进行实时语义提升和轨迹优化，同时改进碰撞惩罚以提高路径平滑性和安全性。
具体而言，方法使用GSFusion双地图：TSDF承担可碰撞几何，GS提供照片级外观；YOLO检测到的2D目标经TSDF提升到3D，规划端使用B样条和hinge-loss碰撞项生成安全轨迹。
适用边界是：依赖检测器和TSDF质量；语义表达不是密集3D语言嵌入；真实机器人泛化仍需更多验证。

\paragraph{Multi-Agent Next-Best-View Optimization for Risk-Averse Planning}
Multi-Agent Next-Best-View Optimization for Risk-Averse Planning（2026.06）~\cite{2606_04158} 延续 3DGS map 在“主动建图与NBV规划”上的基本设定，主要改变是：分布式risk-aware NBV框架，每个机器人维护私有局部3DGS地图，通过C-ADMM协同最大化轨迹掩码区域的信息增益。3DGS 的几何-渲染统一表示保持为底层地图。
其设计动机是：提出通信高效的多智能体NBV优化框架，只交换候选视角、轨迹描述和标量EIG贡献，同时用AVaR建模路径碰撞风险。
具体而言，每个机器人维护局部3DGS地图并估计沿候选轨迹的期望信息增益，团队通过通信图上的Consensus ADMM求解分布式目标，避免共享原始传感器数据。
适用边界是：依赖局部3DGS地图质量和候选视角集合；通信拓扑变化和真实机器人延迟可能影响收敛。

\paragraph{CubifyGS}
CubifyGS（2026.06）~\cite{2606_28720} 可作为“Object-centric lifelong Gaussian maps”中的新锚点候选。相较于 3DGS mapping，主要改变是：将可移动实例建模为可复用 Gaussian assets 并进行检索、刚体变换和剪枝。3DGS 的几何-渲染统一表示保持为底层地图。
其设计动机是：提出 object-level mapping 框架，显式处理物体出现、消失和重排，使 3DGS 地图更适合终身机器人环境。
具体而言，把可移动物体作为 Gaussian assets 存储和复用，检测变化后通过资产检索、刚体变换、显式剪枝和局部补洞来维护地图，而非从头重建。
这一锚点判断来自当前分类中的机制关系：候选锚点：扩展，且相关性 0.91；后续方法可围绕 把 Gaussian map 组织为对象级可维护资产。 继续比较和扩展。
适用边界是：主要适合刚体重排；非刚性变化、复杂接触和错误实例检测可能降低资产复用效果。

\paragraph{ArtiTwinSplat}
ArtiTwinSplat（2026.06）~\cite{2606_24628} 延续 3DGS 在“Interactive digital twins for manipulation”上的基本设定，主要改变是：3DGS 外观重建结合无监督 articulation 发现和运动学估计。3DGS 的几何-渲染统一表示保持为底层地图。
其设计动机是：无需 CAD、仿真资产或人工标注，即可从 RGB-D 视频自动生成具备几何、外观和关节可交互性的物体模型。
具体而言，以 3D Gaussian Splatting 保持几何和光度细节，再通过无监督 articulation discovery 从观察到的运动中恢复部件结构和关节参数，并联合跟踪与优化。
适用边界是：依赖 RGB-D 观测和可见运动，对遮挡严重或未充分运动的部件可能受限；具体速度和失败模式需看全文。

\paragraph{GaussLite}
GaussLite（2026.06）~\cite{2606_30809} 可作为“Task-conditioned Gaussian mapping”中的新锚点候选。相较于 3DGS，主要改变是：用自然语言任务解析驱动 Gaussian 表示密度分配。3DGS 的几何-渲染统一表示保持为底层地图。
其设计动机是：指出均匀分配 Gaussian capacity 会浪费 onboard compute，提出 task-conditioned density allocation。
具体而言，给定 posed RGB-D stream 和任务文本，使用 one-shot LLM parser 提取目标与相关区域，再优先优化任务相关几何和外观，降低无关区域表示密度。
这一锚点判断来自当前分类中的机制关系：候选锚点：任务条件化，且相关性 0.94；后续方法可围绕 用任务语义动态调整 Gaussian density 和在线优化资源。 继续比较和扩展。
适用边界是：依赖任务解析准确性；任务变化频繁或语言歧义可能导致重要区域被低估。

\paragraph{GN0}
GN0（2026.06）~\cite{2606_03682} 延续 3DGS engine 在“导航仿真与评测”上的基本设定，主要改变是：GN0：基于3DGS引擎的大规模导航数据和高保真仿真平台，配合GN-Bench和BAE导航模型。3DGS 的几何-渲染统一表示保持为底层地图。
其设计动机是：提出从3D场景到导航数据、BEV评测和策略学习的统一管线，其中3DGS提供照片级漫游和碰撞感知仿真环境。
具体而言，系统构建GN-Matrix数据集，使用3DGS引擎支持交互式roaming和动态avatar评测，再训练Break and Establish导航基础模型。
适用边界是：3DGS主要作为仿真平台底座；导航策略本身不是3DGS方法创新。

\paragraph{FalconTrack}
FalconTrack（2026.06）~\cite{2606_29783} 延续 Gaussian Splatting simulator 在“Synthetic supervision from Gaussian simulators”上的基本设定，主要改变是：在 Gaussian Splatting simulator 中分离目标高斯并合成随机背景自动生成标签。3DGS 的几何-渲染统一表示保持为底层地图。
其设计动机是：把可编辑 GS 仿真与多头感知、物理感知跟踪结合，实现零样本 sim-to-real 迁移。
具体而言，从短目标视频中隔离 target Gaussians，再与随机背景复合生成带标签的照片级训练数据，训练感知模块并结合物理模型跟踪。
适用边界是：3DGS 是仿真数据生成平台而非核心跟踪算法；目标分离和背景复合的真实性决定迁移效果。

\paragraph{Path Planning in Physically Viable World Models}
Path Planning in Physically Viable World Models（2026.07）~\cite{2607_00673} 延续 3DGS 在“Physically viable Gaussian world models”上的基本设定，主要改变是：在重建 3D Gaussian splat 场景上叠加物理仿真，生成未来地形变化的 what-if 场景。3DGS 的几何-渲染统一表示保持为底层地图。
其设计动机是：让机器人在部署前采集的旧地图上推演未来地形变化，避免使用已失效的静态可行路径。
具体而言，系统以 reconstructed 3D Gaussian splat scenes 为环境表示，结合 physics-based simulation 生成物理修改后的场景，再评估导航路径可行性。
适用边界是：仿真可信度依赖物理参数和环境变化模型；3DGS 几何可碰撞性也会影响规划。

\paragraph{VLK}
VLK（2026.06）~\cite{2606_30645} 延续 3DGS 在“Embodied learning data from Gaussian scenes”上的基本设定，主要改变是：用 3DGS 重建室内环境并渲染 egocentric observations，配合任务语言和运动轨迹。3DGS 的几何-渲染统一表示保持为底层地图。
其设计动机是：填补缺少大规模 egocentric images、language commands 和 robot-compatible trajectories 同步数据的问题。
具体而言，通过 3DGS 重建 metric-scale indoor environments，用 privileged scene information 合成导航和物体交互轨迹，并渲染第一视角观察作为训练监督。
适用边界是：3DGS 是数据生成环境而非机器人算法核心；sim-to-real 差距和场景交互物理真实性仍需验证。

\paragraph{FastBridge}
FastBridge（2026.07）~\cite{2607_01200} 延续 3DGS safety filter 在“Safety filters on Gaussian scene representations”上的基本设定，主要改变是：基于完整四旋翼动力学的高相对阶 collision cone CBF 与 backup CBF。3DGS 的几何-渲染统一表示保持为底层地图。
其设计动机是：缩小 3DGS safety filter 的 model-based realization gap，使快速飞行中的障碍规避更符合真实动力学。
具体而言，基于 analytic collision cone barrier for 3DGS，推导 nonlinear actuator-aware safety filter、高相对阶指数 CBF 和保持 QP 可行性的 backup CBF。
适用边界是：依赖 3DGS 几何和碰撞锥估计准确性；高速动态障碍和感知延迟仍需进一步处理。

\paragraph{PolyMerge}
PolyMerge（2026.06）~\cite{2606_16232} 延续 3DGS 在“可规划高斯地图”上的基本设定，主要改变是：PolyMerge：用凸多面体覆盖过近似3DGS中所有障碍物，并调节多面体数量平衡保守性和计算成本。3DGS 的几何-渲染统一表示保持为底层地图。
其设计动机是：提出从3DGS到控制友好凸几何的转换框架，使原本昂贵的高斯场可用于机载避障和运动规划。
具体而言，方法分析高斯场的占据/障碍区域，以凸多面体联合过近似，并与控制屏障或规划器集成。
适用边界是：过近似会牺牲可通行空间；从照片级高斯到几何障碍的阈值选择影响安全和效率。

\paragraph{SplatlessDF}
SplatlessDF（2026.06）~\cite{2606_13990} 延续 2D Gaussian Splatting 在“可规划高斯地图”上的基本设定，主要改变是：SplatlessDF：用各向异性Gaussian elements在空间域参数化连续Distance Field，并可与2DGS耦合。3DGS 的几何-渲染统一表示保持为底层地图。
其设计动机是：把高斯从图像渲染primitive转为空间距离场元素，支持距离和梯度查询，并可与2DGS表面/渲染形成统一高斯框架。
具体而言，方法直接优化各向异性高斯以恢复可微DF，在导航等下游任务中查询空间距离；同时研究与2DGS耦合以兼顾表面和photometric rendering。
适用边界是：非传统3DGS渲染方法；距离场质量依赖传感器和优化约束。

\paragraph{Geometry-Preserving in 3D Gaussian Splatting for LiDAR-Camera Extrinsic Calibration}
Geometry-Preserving in 3D Gaussian Splatting for LiDAR-Camera Extrinsic Calibration（2026.06）~\cite{2606_20103} 延续 3DGS 在“多传感器标定与地图”上的基本设定，主要改变是：Geometry-Preserving 3DGS calibration：在3DGS优化中加入几何保持约束以服务外参估计。3DGS 的几何-渲染统一表示保持为底层地图。
其设计动机是：指出NVS导向的3DGS会牺牲标定所需几何判别性，提出几何保持机制改善跨模态外参优化。
具体而言，方法在3DGS可微框架中桥接LiDAR和相机，以密集光度监督和几何约束共同优化外参。
适用边界是：依赖场景几何可观测性；弱纹理或动态场景可能影响标定稳定性。

\paragraph{LOGOS}
LOGOS（2026.06）~\cite{2606_21527} 延续 Gaussian splatting 在“可通行性与障碍物感知”上的基本设定，主要改变是：LOGOS：Gaussian Elevation Splatting道路建模和tiny obstacle segmentation。3DGS 的几何-渲染统一表示保持为底层地图。
其设计动机是：提出适合LiDAR扫描的高斯高程表示，用统一框架处理路面起伏和微小障碍区分。
具体而言，方法将LiDAR点投射/聚合为高斯高程primitive，建模道路表面并识别curb、gravel、pothole等微障碍。
适用边界是：不是照片级3DGS渲染方法；侧重LiDAR几何分割。
\section{机械臂操控}
\label{sec:17_2}

机械臂操控是机器人领域最具挑战性的任务之一，要求机器人在非结构化环境中精确地与物体交互。3DGS 在操控中的应用涵盖多个层面：作为场景表示提供几何和语义信息用于抓取规划、作为可渲染环境生成训练数据用于策略学习、作为世界模型预测操作结果用于动作评估。3DGS 在操控中的独特优势在于其\textbf{可编辑性}——可以通过移动、删除或修改特定高斯基元来模拟物体被抓起、放置或推倒后的场景变化，从而为策略学习生成大规模多样化的训练数据，而不需要真实的机器人操作。本节按功能区域推理、抓取与操控策略学习、灵巧操控、可变形物体操控和其他操控方法展开讨论。

\subsection{功能推理与抓取策略}
\label{ss:17_2a}

本节讨论两条技术路线：（1）功能区域推理；（2）抓取与操控策略学习。

\subsubsection{功能区域推理}

本路线介绍基于 3DGS 精细几何表示的功能区域（affordance）推理方法。功能区域（affordance）推理将自然语言指令与 3D 物体的功能区域关联，是机器人实现任务导向操控的关键能力。直观地说，affordance 回答了"物体的哪个部位可以用来做什么"——例如杯子的把手用于抓取、椅子的座面用于坐。3DGS 的连续、稠密表示为精细的功能区域识别提供了理想的几何基础：每个高斯基元既承载了精确的 3D 位置和形状信息，又可以附加语义特征，使得功能区域的预测可以在逐点级别上执行。

\paragraph{大规模功能区域数据集}

\textbf{设计动机与 Baseline 局限}：功能区域推理的现有方法主要依赖稀疏 3D 点云表示——例如 3D AffordanceNet、LASO、PIAD 等基准数据集均以点云为核心模态。然而，点云的离散采样本性导致其对连续曲面和复杂结构的表达能力不足：增大采样密度虽可改善几何保真度，但计算开销显著上升。3DGS 作为高保真连续表示，在保留精细几何和纹理的同时可实现实时渲染，理论上更适合功能区域的精细识别。但 3DGS 在功能区域推理中的应用面临三重障碍：（1）缺乏面向 3DGS 的大规模功能区域标注数据；（2）为点云/图像设计的现有模型无法利用 3DGS 的连续特性；（3）传统的 3DGS 语义嵌入方法（如参数扩展技术）为每个高斯基元静态分配单一语义特征，无法处理同一基元在不同功能上下文中承担不同角色的多属性场景。

\textbf{改动概要}：3DAffordSplat（2025.04）~\cite{2504_11218} 修改了 3DGS 的语义表示与功能推理环节——用动态任务特定描述符替代静态语义嵌入，使每个高斯基元可根据文本查询自适应响应多种功能语义——3DGS 的几何表示和渲染管线保持不变。为此，该方法首先构建首个面向 3DGS 功能区域推理的大规模多模态数据集，包含 23,677 个高斯实例、8,354 个点云实例和 6,631 个手动标注的功能区域标签，覆盖 21 个物体类别和 18 种功能区域类型。每个物体-功能区域组合配备 15 个问答模板，支持语言引导的功能区域推理评估。

\textbf{方法架构}：提出的 AffordSplatNet 以 3DGS 结构特征（位置 $\mu \in \mathbb{R}^3$、缩放 $s \in \mathbb{R}^3$、旋转 $r \in \mathbb{R}^4$）作为输入——功能区域属性主要源自物体的局部结构特征，因此模型仅处理结构特征 $\mathcal{G}_\text{struct} = \{m, s, r\} \in \mathbb{R}^{10}$，外观特征（不透明度和球谐颜色）不参与推理。具体而言，利用 PointNet++ 编码多粒度几何特征，通过 RoBERTa 语言模型提取文本查询的 $\langle\text{Aff}\rangle$ token 嵌入，该 token 同时编码了问题理解和掩码生成的中间语义表示。跨模态融合模块在空间和通道两个层面执行交叉注意力——空间注意力以语言特征为查询、几何特征为键值，通道注意力通过拼接和残差连接融合全局语言上下文与局部几何细节。粒度自适应选择机制通过可学习门控权重 $W_\text{gate}$ 在三个粒度级别上竞争性地分配重要性，最终通过逆距离加权插值上采样至原始高斯密度。解码器基于 Transformer 生成动态卷积核，将多模态特征解码为逐高斯的功能区域掩码——这意味着卷积核不是固定的，而是根据当前文本查询动态生成，从而实现同一高斯基元在不同功能上下文中产生不同的功能区域预测。

\textbf{跨模态预训练}：预训练阶段引入跨模态结构对齐模块，利用结构一致性先验——同一物体类别在不同显式表示中功能区域与整体结构的相对空间关系保持不变——将高斯功能区域和点云功能区域分别编码后投影至共享 $d_\text{consis}$ 维空间，通过多头交叉注意力计算结构亲和矩阵，以 Chamfer 距离加权的余弦损失对齐跨模态的相对结构关系，使模型能够从大规模点云功能区域数据中迁移知识至 3DGS 领域。训练分为预训练（利用大量点云功能区域数据进行跨模态对齐）和微调（使用 3DAffordSplat 数据集的高斯功能区域标注进行监督训练）两个阶段。在已见和未见设置下均展现鲁棒泛化能力。

\paragraph{序列化功能区域推理}

\textbf{改动概要}：SeqAffordSplat（2025.07）~\cite{2507_23772} 修改了 3DGS 功能区域推理的任务定义、推理引擎和训练策略——用 LLM 驱动的自回归序列生成替代独立的单步预测，将``一条指令 $\to$ 一个掩码''映射为``一条指令 $\to$ 有序掩码序列''——3DGS 的几何表示和高斯基元编码方式与锚点保持一致。

\textbf{与锚点 3DAffordSplat 的核心差异}在于任务定义的根本性变化：3DAffordSplat 的每条指令仅映射到单个二值掩码（单步、单物体），而 SeqAffordSplat 引入场景级序列功能区域推理——给定一条简洁的高层指令（如``用微波炉加热碗里的食物''），系统需要同时推理出多步原子动作序列（打开微波炉门 $\to$ 放入碗 $\to$ 关闭门），并在包含多个交互物体和干扰物的复杂场景中逐步定位对应的功能区域掩码。这一设计动机源于现实世界指令的固有简洁性——人类习惯于给出高层复合指令而非逐步原子指令，因此功能区域推理系统必须具备序列分解和跨物体区域转移的能力。为支持这一新任务，SeqAffordSplat 构建了首个场景级序列 3DGS 功能区域基准，包含 1,800+ 复杂场景（通过手动组合 3D-AffordanceNet 中的多物体并施加平移、旋转、缩放变换来模拟真实环境）、14,000+ 功能区域掩码和 8,000+ 序列指令（由 GPT-4o 在场景视觉和文本上下文提示下生成并经人工验证），覆盖 21 个物体类别和 18 种功能区域类型。评估采用三种递增配置：\textit{Single}（单步独立预测）、\textit{Sequential with GT seq}（给定真值序列的掩码定位）、\textit{Sequential}（端到端：从一条指令推理完整动作序列并定位所有掩码），并配套序列化指标 sIoU/sAUC/sSIM/sMAE 以整体评估多步预测质量。

\textbf{SeqSplatNet 架构}采用端到端自回归设计，将任务规划与掩码定位统一于单一生成过程：3DGS 编码器（基于 PointNet++，与锚点一致仅处理几何属性）从高斯基元提取逐点几何特征 $F_\text{geo} \in \mathbb{R}^{N \times d}$；LLM（Qwen-3-0.6B，以 LoRA 微调 q/k/v\_proj 和 lm\_head）作为中央推理引擎处理输入指令，在自回归过程中交替生成语言 token 和特殊 $<$SEG$>$ token——每个 $<$SEG$>$ token 既触发条件可供性解码器生成对应的 3D 掩码，又提供从其隐藏状态导出的动态指令向量 $h_\text{seg} \in \mathbb{R}^d$，该向量通过 LLM 的掩码注意力机制编码了来自原始指令和先前基元的上下文依赖关系。条件可供性解码器以 $h_\text{seg}$ 作为潜在查询，从编码的几何特征中解码出第 $t$ 步的功能区域掩码 $M_t$。整体训练损失联合优化语言建模损失 $\mathcal{L}_\text{lang}$（自回归交叉熵）和逐步分割损失 $\mathcal{L}_\text{mask}$（BCE + Dice），后者在每个 $<$SEG$>$ token 处激活。

\textbf{条件几何重建预训练}解决从零训练 3DGS 编码器的困难：采用双编码器设计，掩码编码器 $\Phi_\text{mask}$ 将真实掩码映射为条件嵌入 $e_\text{mask}$，$e_\text{mask}$ 作为查询对 3DGS 编码器 $\Phi_\text{enc}$ 输出的几何特征 $F_\text{geo}$ 执行注意力聚合，解码器从聚合特征重建掩码。这一任务迫使网络学习从抽象语义概念到对应空间几何的映射——显式地将功能区域概念与几何结构解耦表示——为下游序列任务提供优越的初始化。

\textbf{VFM 语义特征注入}弥补纯几何表示的语义理解不足：利用 3DGS 的高保真渲染能力，从多视角渲染 RGB 图像并通过冻结的 DINOv2（自监督视觉Transformer）/CLIP（对比语言-图像预训练模型）提取 2D 语义特征，再通过逆渲染操作（基于 $\alpha$-混合 权重的加权平均）将 2D 特征提升到每个高斯基元的 3D 语义特征 $f_i^\text{sem}$，在条件可供性解码器的多个尺度上以加法融合方式注入，从粗到细逐级增强分割的语义一致性。消融实验表明，预训练将 sIoU 从 20.3 提升至 24.1（+3.8），DINOv2 特征注入进一步提升至 26.2（+2.1），两者缺一不可。在三种评估配置下：端到端 \textit{Sequential} 设置的 sIoU 达到 26.2，比基于点云的唯一序列基线 SeqAfford（12.1）提升 14.1 个百分点；\textit{Sequential with GT seq} 设置下 sIoU 为 36.0，比最强点云基线 PointRefer（30.3）提升 5.7 个百分点；\textit{Single} 设置下 mIoU 为 37.0，比 3DAffordSplat（30.5）提升 6.5 个百分点。在锚点 3DAffordSplat 数据集的跨域泛化测试中，mIoU 达到 40.2，比原始 3DAffordSplat 基准（30.3）提升 9.9 个百分点，证明架构设计和训练策略在非序列场景下同样有效。

\paragraph{可泛化功能区域学习}

GEAL（2024.12）~\cite{2412_09511} 针对 3D 功能区域学习的泛化性不足和鲁棒性薄弱问题，提出利用 3DGS 建立跨模态一致映射的双分支框架，使 3D 分支继承 2D 基础模型的丰富语义和强泛化能力。该方法还构建了两个噪声鲁棒性基准 PIAD-C 和 LASO-C，涵盖 7 种噪声类型和 5 个严重级别。

\textbf{与锚点 3DAffordSplat 的核心差异}在于跨模态架构设计的哲学不同：3DAffordSplat 直接在 3DGS 结构特征上执行功能区域预测，依赖手工设计的跨模态融合模块；GEAL 则利用 3DGS 的可微渲染能力构建双分支架构——通过 3DGS 将稀疏 3D 点云渲染为真实感 2D 图像（保留语义和深度信息），建立从 3D 到 2D 的桥接，使 3D 功能区域学习能够利用在大规模数据上预训练的 2D 基础模型（冻结的 DINOv2，自监督视觉Transformer）的优越泛化能力。这一设计动机源于一个关键观察：3D 功能区域模型的泛化性能受限于有限的标注数据和仅关注位置/几何编码的骨干网络（如 PointNet++），而 2D 模型受益于海量标注数据和丰富的全局语义捕获。值得注意的是，GEAL 的输入模态是稀疏 3D 点云而非 3DGS 场景——3DGS 在此仅作为 3D-2D 跨模态桥接工具，将点云初始化为高斯基元以实现真实感渲染。

\textbf{3D-2D 映射}是整个框架的基石：将高斯基元的位置 $\mu$ 设为输入点云坐标 $P$，协方差 $\Sigma$ 和不透明度 $\alpha$ 手动调整后固定训练以保留原始几何。利用深度图和预定义颜色图为 $V$ 个相机位姿生成真实感图像 $I \in \mathbb{R}^{V \times 3 \times H \times W}$，同时将功能区域分数 $y \in [0,1]$ 作为灰度颜色赋给高斯（$c = y$），通过 $\alpha$-混合 体渲染公式生成 2D 功能区域掩码 $y_\text{2D}$，建立 3D-2D 的连贯映射。与直接投影不同，3DGS 的 $\alpha$-混合 渲染产生平滑的像素过渡，保留遮挡和深度线索，使渲染图像适合 2D 基础模型进行特征提取。

\textbf{粒度自适应融合模块（GAFM）}通过两个子模块实现多粒度视觉特征与文本指令的对齐：（1）\textbf{柔性粒度特征聚合}——将骨干网络最后 $m$ 层的特征图拼接后，通过带高斯噪声的可学习门控权重 $W = \text{softmax}(f_\text{con} W_g + \sigma \epsilon)$ 自适应融合各层特征，其中 $W_g$ 为可训练权重矩阵，$\sigma$ 控制噪声尺度，$\epsilon \sim \mathcal{N}(0,1)$ 为随机扰动以增强自适应性；（2）\textbf{文本条件视觉对齐}——通过交叉注意力（视觉特征为键值、文本特征为查询）增强文本特征后经 MLP 精炼，再反向查询获得嵌入问题语义的增强视觉特征 $f_\text{enh}$，保持空间结构的同时注入指令相关信息。两个分支的 GAFM 处理方式有所不同：2D 分支（DINOv2 骨干）各层分辨率一致，可直接拼接后计算软权重；3D 分支（PointNet++ 骨干）各尺度空间分辨率和特征维度不同，需先对每尺度独立执行文本条件视觉对齐，再上采样至统一分辨率后执行粒度聚合。此外，2D 分支的输入提示经过视图依赖增强——在每个视角的文本查询后追加``a depth map of a [object] in [view]''以丰富上下文理解，语言模型 RoBERTa 可微调适配。

\textbf{2D-3D 一致性对齐模块（CAM）}实现跨模态知识传播：将 3D 增强特征 $f_\text{enh}^\text{3D}$ 通过两层 Conv1D 降维至 $f_\text{cam}^\text{3D}$，再将每个高斯基元的点特征向量作为其固有属性，利用 3DGS 的 $\alpha$-混合 体渲染公式投影到 2D 像素特征 $f_\text{cam}^{\text{3D} \to \text{2D}}$；2D 增强特征 $f_\text{enh}^\text{2D}$ 通过三层 Conv2D 上采样并降维至 $f_\text{cam}^\text{2D}$，在共享嵌入空间中用 $L_2$ 损失 $L_\text{consis} = \|f_\text{cam}^{\text{3D} \to \text{2D}} - f_\text{cam}^\text{2D}\|_2$ 强制对应空间区域的特征一致，从而将 2D 基础模型的语义知识隐式传播至 3D 分支。

\textbf{解码与训练策略}：共享的 Transformer 解码器以文本特征为查询、增强视觉特征为键值，经过 3 层自注意力+交叉注意力后，用更新后的文本特征作为动态卷积核与视觉特征相乘并经 sigmoid 激活得到功能区域分数预测。损失函数为 BCE + Dice 的组合以同时处理类别不平衡和分割精度。训练采用两阶段策略：第一阶段仅训练 2D 分支；第二阶段冻结 2D 分支（CAM 除外），仅训练 3D 分支，总损失为 $L^\text{3D} = L_\text{BCE}^\text{3D} + L_\text{Dice}^\text{3D} + L_\text{consis}$。推理时仅使用 3D 分支，确保轻量预测。在 PIAD 和 LASO 数据集的已见/未见类别设置下均达到最优，在噪声基准 PIAD-C 和 LASO-C 上全面超越 LASO 基线。消融实验表明：单独使用 2D 或 3D 分支性能相当，引入 CAM 后显著提升，加入 GAFM 的完整模型在已见和未见设置下均取得最佳结果。

\paragraph{语言到灵巧动作}

BLaDA（2026.04）~\cite{2604_08410} 提出从语言到灵巧动作的可解释零样本框架，是首个将语言理解、感知定位和物理执行统一于功能抓取范式的研究。

\textbf{与锚点 3DAffordSplat 的核心差异}在于目标任务的层面不同：3DAffordSplat 聚焦于功能区域的逐点分割——预测每个高斯基元是否属于特定功能区域，输出为逐点的二值掩码；BLaDA 则进一步将功能区域语义桥接到灵巧手的物理执行，建立从``哪个区域可以用来做什么''到``手指如何精确抓握该区域并执行任务''的完整推理链。这一设计动机源于功能抓取与普通抓取的本质区别：灵巧手的功能性抓取不仅需要稳定持握，还要求在正确功能部位执行有目的的交互（如按压喷雾瓶按钮而非仅抓握瓶身），涉及细粒度手-物体接触、位姿级约束和高精度控制。端到端 VLA 方法虽能直接映射语言到动作，但缺乏结构化中间表示导致可解释性弱、对分布外场景泛化能力差；而现有的层次化方法多停留在基础抓取，难以满足功能抓取对语义-位姿耦合和手指级精细控制的要求。

\textbf{知识引导语言解析（KLP）}模块将开放域自然语言指令解析为结构化六元组 $\mathcal{S} = (g^a, g^r, g^t, g^f, \kappa, \tau)$：$g^a$（功能区域可达性，指定物体表面可用区域的空间可达范围）、$g^r$（手指角色分配，明确各手指在接触中的功能逻辑）、$g^t$（抓握姿态类型，直接对应不同粗粒度手型的关节角度值）、$g^f$（交互力度等级，设定底层动力学的交互强度）、$\kappa$（任务意图先验：press/click/open/hold，将动词短语规范化为四种原子任务类型以排除逻辑不一致的抓握组合）和 $\tau$（工具拓扑先验：rod/handle/knob/surface，将物体按几何形状和人类操作习惯分为四类拓扑结构以提供结构化执行约束）。与仅依赖 LLM 通用推理的方法不同，KLP 显式注入领域先验——通过精心设计的三组件 prompt（角色规范定义具身智能执行上下文、结构化注入将抓握分类法和工具拓扑模型以形式化方式提供给 LLM、上下文示例用代表性``指令-推理-元组''范例引导标准化输出），将开放词汇理解与操控领域知识结合。实验表明，KLP 作为即插即用推理模块，在 ChatGPT-4.0 上的平均语言推理准确率从 0.508 提升至 0.753（+21.5\%），其中需要领域专业知识的抓握类型 $g^t$ 准确率从 0.13 跃升至 0.81。

\textbf{三角功能点定位（TriLocation）}模块在 3DGS 场景中定位三个接触关键点：功能指尖接触点 $p_1$（对应食指或拇指的功能性接触）、侧向支撑点 $p_2$（对应小指侧的接触）和腕部支撑点 $p_3$（对应掌根附近的接触），三点构成三角结构直接约束后续手-物体接触位姿。该模块的核心创新包含两个层面。

\textbf{分层语义提取（HSE）}策略解决 GraspSplats 等方法中部件级特征被全局物体语义淹没（语义抑制）和局部裁剪导致全局上下文丢失（语义漂移）的双重问题。具体而言，通过 YOLO 检测候选框经 SAM 生成掩码后，面积比一致性超参数 $\alpha$ 将掩码解耦为物体级集合 $\mathcal{M}_o$（大面积背景或主体区域）和部件级集合 $\mathcal{M}_p$（局部功能组件）。对部件级区域采用上下文感知裁剪：以填充比 $\gamma$ 扩展裁剪边界 $B_p'=[x_1-\gamma w, y_1-\gamma h, x_2+\gamma w, y_2+\gamma h]$ 后经 CLIP 编码，再通过掩码引导投影仅将中心特征映射回对应 SAM 掩码，在保留物体类别上下文的同时提取高纯度部件特征。对不属于任何检测掩码的像素，显式将真值特征向量置零（背景一致性约束），最大化信噪比并防止能量泄漏到背景区域。最终通过 3DGS 体渲染公式将每个高斯基元的低维潜特征渲染为多分支特征图（物体级、部件级、DINOv2），以 $L = L_\text{obj} + L_\text{part} + \lambda L_\text{dino}$ 监督训练。最优超参数配置为 $\alpha=0.1, \gamma=0.4$。

\textbf{三角关键点定位}在构建的多级语义 3DGS 场中执行：首先将 $g^a$ 输入 CLIP 获得查询向量，与每个高斯的语义特征 $f_i$ 计算相似度分布，将相似度高于阈值 $\delta$ 的高斯点在 3D 空间中聚类，取最高置信度簇的质心作为语义锚点 $p_1$。然后根据 $g^r$ 从手三点点模板库中选择功能手模板（含手部三个关键点 $p_{h1}, p_{h2}, p_{h3}$），利用任务-拓扑耦合局部坐标系进行刚体对齐：以 $p_1$ 为原点、由 $\kappa$-$\tau$ 组合确定主任务轴 $\hat{z}$ 和手朝向轴 $\hat{y}$（经 RANSAC 从点云提取结构轴或经 PCA 估计表面法线），通过 $\hat{x} = \hat{y} \times \hat{z}$ 补全右手坐标系，旋转矩阵 $R$ 的列向量对应局部轴 $(\hat{x}, \hat{y}, \hat{z})$ 的世界坐标表示。在此局部坐标系中，$p_3$ 固定沿 $-z$ 方向，$p_2$ 位于 $y$-$z$ 平面内与 $p_3$ 成角 $A_1$，由模板提供的边长 $(L_{12}, L_{13}, L_{23})$ 确定世界坐标 $\mathbf{p}_2 = \mathbf{p}_1 + R[0, -L_{12}\sin A_1, -L_{12}\cos A_1]^T$，$\mathbf{p}_3 = \mathbf{p}_1 + R[0, 0, -L_{13}]^T$。消融实验表明，去除局部坐标系约束后 Press 和 Open 任务的定位成功率均降至 0\%，均值 LSR 从 68.75\% 骤降至 25\%。

\textbf{KGT3D+ 执行模块}从三个关键点构建精确手掌位姿 $(R, T)$，其物理可解释性体现在坐标轴的几何构造：$\vec{z} = (p_1 - p_3)/\|p_1 - p_3\|$ 定义接近方向（从腕部指向功能指尖），$\vec{y} = (p_2 - p_3) \times \vec{z}/\|(p_2 - p_3) \times \vec{z}\|$ 定义手朝向，$\vec{x} = \vec{y} \times \vec{z}$ 补全右手系，$R = [\vec{x}\ \vec{y}\ \vec{z}]$，$T = p_3$。抓握类型 $g^t$ 和力度 $g^f$ 通过预定义功能抓握库（F2F）映射为关节配置 $J$ 和逐指力分布 $F$：$J, F = \text{F2F}(g^t, g^f)$。整个执行流水线避免端到端黑箱映射，执行时序为：先根据 $g^t$ 产生粗粒度手型关节角度实现初始包络，再根据 $g^r$ 收紧非功能手指提升稳定性，最后由 $g^f$ 驱动功能手指执行任务特定操作。在零样本设置下，``持握喷雾瓶''任务成功率 80\%（vs MKA 40\%、Diffusion Policy 50\%），``按压喷雾瓶''任务成功率 30\%（vs MKA 和 DP 均为 10\%），在 66 个真实场景 14 种任务-工具组合上整体 LSR 达 68.75\%（vs MKA 22.5\%）。

\subsubsection{抓取与操控策略学习}

本路线介绍利用 3DGS 渲染能力进行机器人抓取与操控策略学习的方法。3DGS 的照片级真实感渲染能力为机器人操控策略学习提供了关键基础——可以从重建的 3DGS 场景中生成大量多样化的训练数据，大幅降低真实世界示范数据的需求。

\paragraph{零样本抓取}

RobMRAG（2026.03）~\cite{2603_00500} 提出基于 3DGS 增强的多模态检索增强生成框架，构建包含物体接触帧、任务完成帧和位姿参数的多源操控知识库。3D 感知位姿精化模块利用 3DGS 在 3D 空间中对齐参考物体和目标物体的位姿，在 30 类家用物体的测试集上比最佳零样本基线提升 7.76\%。SparseGrasp（2024.12）~\cite{2412_02140} 仅需稀疏视角 RGB 图像即可运行开放词汇机器人抓取系统，利用 DUSt3R 生成稠密点云初始化 3DGS，结合视觉基础模型的语义感知和 PCA 特征压缩实现快速场景更新。

\textbf{与锚点 Splat-MOVER 的核心差异}在于对深度传感器的依赖：Splat-MOVER 假设 RGBD 输入并通过精确的深度引导完成场景重建和语义蒸馏；SparseGrasp 则完全基于纯 RGB 输入运行——利用 DUSt3R 立体匹配基础模型从两张或少量 RGB 图像直接生成稠密点云，无需深度传感器。这样设计的动机来自实际部署中的传感器限制：深度相机在面对透明和金属物体时常常失效，而 SparseGrasp 的纯 RGB 流水线在这些场景下具有天然优势。具体技术差异包括：（1）用 DUSt3R 替代深度传感器提供 3DGS 初始化所需的点云；（2）通过 PCA 压缩 CLIP 高维特征以保持实时语义查询能力；（3）实验验证在透明物体和金属物体组成的场景中成功完成抓取，而依赖深度传感器的方法因无法获取有效点云而完全失败。

\paragraph{可编辑场景操控}

Splat-MOVER（2024.05）~\cite{2405_04378} 修改了 3DGS 的场景表示和渲染管线——用语义--可供性增强高斯替代标准 3D 高斯——基于 tile 的光栅化渲染机制保持不变。系统由三个模块组成：ASK-Splat 负责知识蒸馏，SEE-Splat 负责实时场景编辑，Grasp-Splat 负责可供性对齐的抓取规划。

\textbf{ASK-Splat：语义与可供性知识蒸馏。}
标准 3DGS 仅优化每个高斯的位置 $\boldsymbol{\mu}$、协方差 $\boldsymbol{\Sigma}$、不透明度 $\alpha$ 和球谐参数，缺乏语义理解和抓取知识。ASK-Splat 为此给每个高斯新增两个参数：语义特征 $f_s \in \mathbb{R}^l$ 和可供性分数 $\hat{\alpha} \in [0,1]$。然而，直接将 CLIP 的高维特征（$C=512$ 维 ViT 或 $C=1024$ 维 ResNet）赋给每个高斯会带来巨大的内存和计算开销，破坏 3DGS 的实时渲染速度。因此，ASK-Splat 引入自编码器将 CLIP 特征压缩到低维潜空间，编码器 $g_\phi^{\text{enc}}: \mathbb{R}^{H_f \times W_f \times C} \to \mathbb{R}^{H_f \times W_f \times l}$ 将特征维度从 $C$ 压缩至 $l = 3$，解码器 $g_\psi^{\text{dec}}$ 将其恢复。自编码器通过以下损失函数训练：

\begin{equation}
    \mathcal{L}_g = \lambda_g \frac{1}{|\mathcal{D}|} \sum_{i=1}^{|\mathcal{D}|} \|g_\psi^{\text{dec}}(g_\phi^{\text{enc}}(f_{\text{gt},i})) - f_{\text{gt},i}\|_2^2 + \frac{1}{|\mathcal{D}|} \sum_{i=1}^{|\mathcal{D}|} (1 - \text{sim}(g_\psi^{\text{dec}}(g_\phi^{\text{enc}}(f_{\text{gt},i})), f_{\text{gt},i}))
\end{equation}

即同时最小化重建均方误差和余弦相似度损失，确保压缩后的特征尽可能保留原始语义信息。

在语义特征蒸馏阶段，ASK-Splat 将压缩后的 CLIP 特征蒸馏到每个高斯的语义参数 $f_s$ 中，通过以下损失函数优化：

\begin{equation}
    \mathcal{L}_s = \lambda_s \frac{1}{|\mathcal{D}|} \sum_{i=1}^{|\mathcal{D}|} \|\bar{I}_i^f - g_\phi^{\text{enc}}(f_{\text{gt},i})\|_2^2 + \frac{1}{|\mathcal{D}|} \sum_{i=1}^{|\mathcal{D}|} (1 - \text{sim}(\bar{I}_i^s, g_\phi^{\text{enc}}(f_{\text{gt},i})))
\end{equation}

其中 $\bar{I}_i^s$ 是从高斯 splats 渲染出的 2D 语义特征图，损失函数要求渲染结果与 CLIP 编码器的输出在 MSE 和余弦相似度两个维度上对齐。这意味着每个高斯的语义参数学会在 3D 空间中复现 CLIP 的 2D 语义嵌入。

在可供性蒸馏方面，ASK-Splat 利用大规模人--物交互数据集训练的视觉可供性基础模型 VRB 提取每个像素的抓取可供性分数，再将其蒸馏到每个高斯的参数 $\hat{\alpha}$ 中，优化 MSE 损失 $\mathcal{L}_\alpha = \lambda_\alpha \frac{1}{|\mathcal{D}|} \sum_{i=1}^{|\mathcal{D}|} \|\bar{I}_i^\alpha - I_i^\alpha\|_2^2$，其中 $\bar{I}_i^\alpha$ 和 $I_i^\alpha$ 分别为渲染和真实 2D 可供性图。$\hat{\alpha}$ 通过 sigmoid 激活函数约束在 $[0,1]$ 以提供平滑梯度。对于开放词汇查询，ASK-Splat 计算查询文本的 CLIP 嵌入与每个高斯语义参数之间的余弦相似度，再经成对 softmax 函数转化为概率，阈值为 $0.5$，即可定位场景中与自然语言描述语义相关的物体。实验发现，先训练高斯的非语义参数再训练语义参数（而非同时训练）能获得更好的语义特征局部化效果。

\textbf{SEE-Splat：3D 语义编辑与修复。}
ASK-Splat 赋予了场景语义理解能力，但多阶段操控任务还需要在执行动作后实时更新场景表示。SEE-Splat 正是解决这一问题的模块。它包含三个步骤：首先利用 ASK-Splat 的语义相似度图定位目标物体的高斯集合；然后通过掩码操作提取语义相关的高斯，形成稀疏点云——为了获得足够密集的点云以支持照片级渲染，SEE-Splat 将点云特征从 3D 欧氏空间提升到 7D 特征空间（加入 RGB 颜色和语义分数），使用 KD-tree 查询该 7D 空间中的邻域高斯以增密点云；最后应用变换 $T: \mathcal{G}_s \to \mathcal{G}_s$ 编辑相关高斯的空间和几何属性，实现物体的移动、插入或删除。变换 $T$ 不限于刚体变换，可由启发式规则、大语言模型或人类指令生成。编辑后场景中可能因高斯移除而出现空洞等视觉伪影，SEE-Splat 通过 3D 高斯修复（infilling）在缺失几何区域引入新高斯以恢复照片级渲染质量。这样，SEE-Splat 在虚拟 3D 环境中为操控任务创建了一个可实时更新的"数字孪生"。

\textbf{Grasp-Splat：可供性对齐的抓取规划。}
Grasp-Splat 利用 SEE-Splat 生成的物体稠密点云，通过 GraspNet 模型估计 6 自由度抓取位姿及其质量分数。然而，GraspNet 生成的抓取并不总是理想的——它不区分物体的哪个部位更适合抓取。因此，Grasp-Splat 引入抓取度量 $M: \text{SE}(3) \to \mathbb{R}$，计算每个候选抓取位姿 $X \in \text{SE}(3)$ 处的可供性分数，按 $M(X)$ 对候选抓取重新排序，使抓取位姿优先选择人类更可能抓握的部位，从而提高抓取成功率。

Splat-MOVER 在 Kinova 机器人上完成四个单阶段和四个多阶段操控任务。在三个单阶段任务中，Splat-MOVER 的成功率比 LERF-TOGO 基线提升至少 $2.4$ 倍，比 F3RM* 基线提升 $1.2$--$3.3$ 倍。

上述方法通过语义编辑实现静态场景操控，但多阶段操控任务需要系统能够预测动作的物理后果——推倒物体后场景会如何变化、抓取后物体将出现在哪里。以下方法将 3DGS 提升为动态世界模型，使其能够从当前状态预测未来场景演变。

\paragraph{高斯世界模型操控}

ManiGaussian++（2025.06）~\cite{2506_19842} 将其前代 ManiGaussian（单臂高斯世界模型）扩展至双手操控场景，通过任务导向高斯泼溅区分稳定臂与作用臂，并构建领导者-追随者层次化高斯世界模型编码多体时空动态。

\textbf{与锚点 Splat-MOVER 的核心差异}在于操控范式的根本转变：Splat-MOVER 利用 3DGS 的可编辑性实现静态场景中的语义引导操控——通过编辑高斯基元模拟物体移动，但不预测动作的物理后果；ManiGaussian++ 则将 3DGS 提升为\textbf{动态世界模型}，通过未来场景预测编码多体时空动态，使视觉表示能够捕捉双臂系统中稳定臂和作用臂之间的复杂协作模式。与前代 ManiGaussian 仅编码场景级粗粒度时空动态（单臂设置）不同，ManiGaussian++ 专注于双臂系统的多体交互：领导者变形模型先预测稳定臂运动引起的场景变化，追随者变形模型在此基础上叠加作用臂运动的物理后果，两级级联形成完整的多体未来场景预测。这一设计动机源于双手操控的独特挑战：双臂系统不是两个单臂的简单组合——它需要同时操作和稳定目标物体的特定协作模式，要求视觉表示能够理解两个操控器和目标物体之间的多体交互。

\textbf{具体技术差异}包括：（1）\textbf{任务导向高斯泼溅}：稀疏卷积体素编码器从 RGB-D 观测提取视觉表示后，前馈高斯回归器从视觉表示生成当前高斯辐射场，并在标准高斯参数 $\theta_i^{(t)} = (\mu_i^{(t)}, c_i^{(t)}, r_i^{(t)}, s_i^{(t)}, \sigma_i^{(t)})$ 基础上新增实例级 logits $l_i^{(t)} \in \mathbb{R}^3$，通过 GroundedSAM 从人类指令关键词生成真实实例图作为监督，使每个高斯基元能够区分属于稳定臂、作用臂还是目标物体；（2）\textbf{层次化高斯世界模型}：采用领导者-追随者架构，高斯变形遵循 SE(3) 刚体变换（仅预测位置和旋转变化，保持颜色、缩放、不透明度和实例标签在马尔可夫转移中不变）。

领导者变形模型 $q_{\text{s},\phi}$ 以当前高斯状态 $\theta^{(t)}$、稳定臂动作 $\mathbf{a}_{\text{s}}^{(t)}$ 和视觉表示 $v^{(t)}$ 为输入，预测稳定臂运动引起的位置偏移 $\Delta\boldsymbol{\mu}_{\text{s}}^{(t)}$ 和旋转变化 $\Delta\mathbf{r}_{\text{s}}^{(t)}$；追随者变形模型 $q_{\text{a},\phi}$ 在领导者预测的高斯状态基础上，额外接收作用臂动作 $\mathbf{a}_{\text{a}}^{(t)}$，预测作用臂运动的物理后果 $\Delta\boldsymbol{\mu}_{\text{a}}^{(t)}$ 和 $\Delta\mathbf{r}_{\text{a}}^{(t)}$，两者级联形成 $(\boldsymbol{\mu}^{(t+1)}, \mathbf{r}^{(t+1)}) = (\boldsymbol{\mu}^{(t)} + \Delta\boldsymbol{\mu}_{\text{s}}^{(t)} + \Delta\boldsymbol{\mu}_{\text{a}}^{(t)},\; \mathbf{r}^{(t)} + \Delta\mathbf{r}_{\text{s}}^{(t)} + \Delta\mathbf{r}_{\text{a}}^{(t)})$ 的完整多体未来场景预测；（3）\textbf{多目标联合训练}：当前场景的多视角光度重建损失 $\mathcal{L}_\text{Recon}$ 正则化空间一致性，实例掩码的交叉熵损失 $\mathcal{L}_\text{Task}$ 优化任务导向高斯辐射场，未来场景预测损失 $\mathcal{L}_\text{Pred}$ 将多体时空动态嵌入视觉表示，行为克隆损失 $\mathcal{L}_\text{BC}$ 通过 PerceiverIO 多模态 Transformer 从增强的体素表示和语言指令预测双臂动作（位置、朝向、开合度的离散化候选中分类选择最优动作）。该方法在 RLBench2 的 10 个仿真任务中平均成功率 35.6\%，比前代 ManiGaussian（18.8\%）相对提升 89.4\%，比 PerAct2（15.4\%）绝对提升 20.2 个百分点；在 9 个真实世界双手操控任务中平均成功率约 60\%。

GAF（2025.06）~\cite{2506_14135} 提出高斯动作场（Gaussian Action Field），通过可学习运动属性 $\delta \in \mathbb{R}^3$（位移向量）扩展 3DGS 实现 4D 动态场景和操控动作建模，建立 V-4D-A（视觉-4D-动作）范式。

\textbf{与锚点 ManiGaussian++ 的核心差异}在于 4D 表示的构建方式和动作推理的独立程度：ManiGaussian++ 的未来场景预测由动作条件驱动的变形模型实现——领导者-追随者架构需要显式的机器人动作作为输入来预测高斯状态变化，本质上是``给定动作 $\to$ 预测后果''的前向模拟；GAF 则将运动信息内嵌到每个高斯基元中，仅从视觉输入（稀疏多视角 RGB 图像，无需位姿）直接推断场景的未来演变，再从当前-未来高斯状态的差异中提取初始动作，形成``当前高斯 $\to$ 未来高斯 $\to$ 初始动作''的完整感知-动作管线。这一设计动机源于对现有两类范式的共同局限性的反思：V-A（视觉到动作）范式直接映射像素到动作但缺乏 3D 空间推理，V-3D-A 范式引入 3D 表示但仍以静态快照建模场景——两者均未捕捉场景几何的\textbf{时间演变}，导致静态场景理解与动态动作生成之间存在固有失配。GAF 通过在 3DGS 中显式编码运动属性，将高斯表示从静态 3D 提升到动态 4D，使视觉表示能够同时编码``场景现在是什么样''和``场景未来会变成什么样''，从而为动作推理提供时间感知的几何基础。与 ManiGaussian++ 依赖 RGB-D 输入和已知相机位姿不同，GAF 仅需稀疏 RGB 图像和相机内参即可完成整个 4D 重建，消除了对深度传感器和位姿标定的依赖。

\textbf{具体技术差异}包括：（1）\textbf{运动增强高斯表示}：GAF 将每个高斯基元定义为连续函数 $F_\Theta: \{g(x), t\} \mapsto \{\mu, \delta, f\}$，其中 $\delta \in \mathbb{R}^3$ 是编码时间位移的运动属性，$f = (c, \alpha, r, s)$ 为外观参数（颜色、不透明度、旋转、缩放）。GAF 产生三种相互关联的输出——当前高斯 $\{\mu, f\}$ 表示当前场景，未来高斯 $\{\mu + \delta, f\}$ 表示场景的时间演变，初始动作通过抓取器部分高斯的当前-未来状态的 ICP（迭代最近点）配准估计刚体变换 $T^{t \to t+\Delta t} = \arg\min_T \sum_{k \in \text{gripper}} \|T(\mu_k) - (\mu+\delta)_k\|^2$，该变换矩阵描述末端执行器在时间间隔 $\Delta t$ 内的运动，被插值为离散时间步的变换序列形成初始动作 $a_\text{init}$；（2）\textbf{无位姿前馈重建}：与 ManiGaussian++ 依赖 RGB-D 和精确相机位姿不同，GAF 采用几何无关、无需位姿的方法——ViT 主干从两个无位姿输入视图通过跨视角注意力提取特征，经三个基于 DPT 架构的专用解码器（高斯中心头 $h_\text{center}$ 预测位置、参数头 $h_\text{feature}$ 估计外观参数、运动预测头 $h_\text{motion}$ 预测逐点位移 $\delta$）直接重建运动增强高斯。

重建损失 $\mathcal{L}_\text{GAF} = \mathcal{L}_\text{LPIPS}^t + \mathcal{L}_\text{MSE}^t + \mathcal{L}_\text{LPIPS}^{t+\Delta t} + \mathcal{L}_\text{MSE}^{t+\Delta t}$ 同时约束当前帧和未来帧的渲染质量，实现运动属性的联合优化；（3）\textbf{动作-视觉对齐去噪}：初始动作仅视觉合理但缺乏物理可行性（GAF 模块完全在视觉数据上训练，无真实动作监督），因此采用扩散模型精化。受 R\&D 方法启发，在每一步去噪中将 $a_\text{init}$ 对应的抓取器位姿渲染到当前多视角 RGB 图像上，生成``可操作多视角 RGB 引导'' $R^c = \text{Render}(T_{w2c} \cdot T_g^w \cdot a, K^c)$，将 3D 视觉观测与时间预测动作统一为单一表示，精化损失为 $\mathcal{L}_\text{refine} = L_1(D, D^\text{gt}) + L_1(\epsilon, \epsilon^\text{gt}) + \text{BCE}(g, g^\text{gt})$，其中 $D$ 为抓取器去噪方向、$\epsilon$ 为末端执行器动作噪声、$g$ 为开合二值变量。整个控制闭环（GAF 4D 场景表示 $\to$ 去噪精化 $\to$ 执行 $\to$ 新观测采集）迭代运行直至任务完成，在遮挡和交互不确定性下保持鲁棒。在 RLBench 的 9 个仿真任务中，GAF 比 Diffusion Policy 提升 15.7\% 成功率，重建质量达到 +11.54 dB PSNR、+0.39 SSIM 和 -0.56 LPIPS 的改进，且成功迁移至真实世界环境。

\paragraph{预测感知操控}

Forecast-GS（2026.05）~\cite{2605_11144} 提出预测感知 3DGS 框架，在语言引导的抓取-放置操控中显式预测任务完成状态的 3D 场景。

\textbf{与锚点 Splat-MOVER 的核心差异}在于预测的时间方向：Splat-MOVER 通过 SEE-Splat 在执行动作后实时更新场景表示（事后编辑），而 Forecast-GS 在执行动作前显式预测任务完成后的候选终态（事前预测）。这一设计动机源于抓取-放置任务的核心需求：机器人不仅需要知道``物体当前在哪里''，更需要评估``动作执行后物体会在哪里''，以便在执行前验证候选结果的几何一致性和语义可行性。

\textbf{具体技术差异}包括：（1）\textbf{语义高斯场构建}：从多视角 RGB-D 观测构建 3DGS 表示，将 Detic（一种开放词汇目标检测器）、DINOv2（自监督视觉Transformer）和 CLIP（对比语言-图像预训练模型）的语义特征蒸馏到每个高斯基元 $g_j = (x_j, \Sigma_j, c_j, f_j)$ 的高维语义特征 $f_j$ 中；（2）\textbf{候选终态生成}：给定自然语言指令，先从语义高斯场中利用开放词汇语义查询定位被操控物体和目标区域，然后在可达工作空间内采样平移和绕竖直轴的旋转变换 $T_k$，生成离散的候选终态集合 $\{G_k = T_k(S)\}$；（3）\textbf{候选评估与选择}：每个候选按综合得分 $G^* = \arg\max_{G_k} E(G_k) + \Phi(G_k)$ 排序，其中 $E(G_k)$ 度量语言导出空间约束（接近度、包含、对齐）的违反程度，$\Phi(G_k)$ 度量物理可行性（碰撞、工作空间边界、放置稳定性），自动选择模式在切割器入箱、苹果入碗、海绵入托盘三个任务中分别实现 21/25、23/25 和 16/25 的成功率，超越 ReKep 基线。

\paragraph{工具使用学习}

Tool-as-Interface（2025.04）~\cite{2504_04612} 将人类工具使用知识迁移至机器人，利用 3DGS 新视角合成增强策略对视角变化的鲁棒性，采用分割观测和以工具为中心的任务空间动作实现具身不变视觉运动策略学习。

\textbf{与锚点 Splat-MOVER 的核心差异}在于 3DGS 的应用层面：Splat-MOVER 将 3DGS 用作场景表示和语义编辑界面（在场景级别操作高斯基元）；Tool-as-Interface 则将 3DGS 用作数据增强工具——从双视角人类操控视频通过 MASt3R 重建 3D 场景后，利用 3DGS 的新视角合成能力生成多角度训练图像，使策略对视角变化具备鲁棒性。这一设计动机源于跨具身学习中的视角差异问题：人类演示视频和机器人部署时的相机位姿通常不同，直接用人类视频训练的策略会因视角偏移而失败。

\textbf{具体技术差异}包括：（1）\textbf{具身无关感知对齐}：使用 Grounded-SAM 分别在训练阶段遮蔽人手、在部署阶段遮蔽机器人手臂，确保训练和测试的视觉空间仅保留任务相关信息（工具和物体）；（2）\textbf{以工具为中心的动作表示} $T_\text{tool}^\text{task}$：将动作描述为工具在任务坐标系中的 6D 位姿，独立于人类或机器人形态、相机位姿和基座配置，部署时通过已知变换 $T_\text{eef}^\text{base} = T_\text{task}^\text{base} T_\text{tool}^\text{task} T_\text{eef}^\text{tool}$ 转换为末端执行器指令。该方法比遥操作扩散策略任务成功率提升 71\%，数据收集时间减少 77\%。

\paragraph{主动视觉模仿学习}

ObAct（2025.11）~\cite{2511_18140} 提出观察者-执行者框架。

\textbf{与锚点 Splat-MOVER 的核心差异}在于 3DGS 的角色从离线场景表示转变为在线主动感知工具：Splat-MOVER 先完成完整的 3DGS 场景扫描和训练，再在固定表示上执行操控规划；ObAct 则将 3DGS 构建嵌入操控循环中——观察者臂从三张图像快速构建 3DGS 表示，在该表示上通过虚拟探索（渲染不同候选视角并评估策略置信度）找到最优相机位姿后移动至该位姿，执行者臂利用观察者的最优视角观测执行策略。这一设计动机源于固定相机视角的固有局限：静态相机与任务无关放置，腕装相机视野有限且缺乏全局感知，而两者结合又引入冗余观测。主动视觉策略通过在执行时动态选择最优视角解决这一矛盾。

\textbf{具体技术差异}包括：（1）3DGS 表示用于快速评估候选视角的观测质量（通过新视角合成模拟不同位姿下的图像），而非用于场景编辑或抓取规划；（2）观察者臂不是固定的感知臂，而是可以在观察者和执行者角色之间切换的通用操控器，避免了专用感知臂的硬件浪费和有限可达视角问题；（3）轨迹传递方法通过在 3DGS 表示上训练的策略先验将人类演示轨迹适配到新视角，在无遮挡和有遮挡设置下分别提升 145\% 和 233\%。

\subsection{灵巧操控与特殊操控}
\label{ss:17_2b}

本节讨论三条技术路线：（1）灵巧操控与多指手；（2）可变形物体与特殊操控；（3）其他操控方法。

\subsubsection{灵巧操控与多指手}

本路线介绍基于 3DGS 的灵巧操控与多指手控制方法。灵巧操控需要精确的 3D 理解来处理复杂的任务动态，3DGS 的视觉保真度对训练鲁棒的感知模型至关重要。

\paragraph{灵巧手仿真到真实迁移}

ViserDex（2026.04）~\cite{2604_11138} 修改了 3DGS 的\textbf{渲染管线数据增强环节}——用\textbf{预光栅化增强}（pre-rasterization augmentation）替代传统后处理图像增强——3D 高斯的位置、协方差、不透明度保持不变，仅修改球谐系数以生成视觉多样的训练图像。该方法解决的核心问题是灵巧操控中的仿真到真实迁移：多指手操控物体时，物体在手指间快速运动、频繁被遮挡且伴随运动模糊，导致直接从 RGB 图像估计物体位姿极为困难。传统生成训练数据的策略面临两难——后处理图像增强（如颜色抖动、亮度变化）计算便宜但仅施加全局 2D 变换，无法捕捉 3D 空间中光照和材质的结构化变化；而基于光线追踪的域随机化虽然可以精确模拟光照变化，但计算开销过大，难以在仿真循环中实时运行。

ViserDex 的核心洞察在于：3DGS 的显式高斯属性——特别是球谐（SH）系数——本身就是一个天然的数据增强操作界面。每个 3D 高斯的颜色由球谐系数决定，其中零阶系数 $\text{SH0}$ 编码与视角无关的漫反射基色，高阶系数 $\text{SHN}$ 编码高频镜面效果。论文提出的预光栅化增强直接在高斯属性空间操作，在光栅化之前对 SH 系数施加结构化随机扰动，既保持了物理一致性（相邻高斯协同变化），又避免了全场景光线追踪的计算开销。具体而言，论文设计了四层增强策略，从局部到全局逐层施加：（1）\textbf{随机噪声层}：对每个高斯独立施加随机扰动，模拟高频传感器噪声和像素级伪影；（2）\textbf{空间聚类层}：将高斯按空间位置用 k-means 聚为 $K_s=64$ 个空间簇，对每个簇整体扰动，模拟阴影、磨损等局部一致的外观变化；（3）\textbf{颜色聚类层}：将高斯按 $\text{SH0}$ 基色聚为 $K_c=32$ 个颜色簇，对同一材质区域的漫反射属性施加统一偏移，模拟反照率变化或反射率差异；（4）\textbf{全局偏移层}：将整个高斯场景视为单一簇，施加全局亮度和色温偏移，模拟环境光照变化和相机曝光差异。每层增强以概率 $p_\text{aug}$ 激活，仅扰动比例为 $p_\text{cluster}$ 的簇，扰动幅度在 $[\eta_\text{min}, \eta_\text{max}]$ 范围内均匀采样，多层扰动可叠加产生复合效果。

\textbf{渲染与遮挡合成}。在灵巧操控中，相机固定而物体被手指操控产生运动。为了用静态 3DGS 表示渲染运动中的物体，ViserDex 对相机施加物体变换的逆变换，保持场景静态以满足标准 3DGS 的假设，同时生成反映物体位姿变化的图像。物体渲染后，还面临手指遮挡问题：3DGS 仅表示物体，不包含机器人手。解决方法是利用物理仿真器中的低精度深度光线投射器生成手部深度图 $D_\text{phys}$，与 3DGS 渲染的物体深度图 $D_\text{splat}$ 比较，将手在前方的像素（$D_\text{phys} < D_\text{splat}$）从 RGB 图像中遮蔽，以零额外光线追踪成本恢复物理一致性。

\textbf{策略训练三阶段流水线}。在感知端，利用上述 3DGS 渲染管线生成的照片级真实感增强图像训练 ResNet-34 位姿估计器，回归 9 个关键点的 2.5D 坐标 $(u, v, d)$，再通过刚性 Procrustes 算法恢复 6D 物体位姿。在控制端，采用教师-学生蒸馏架构：教师策略在仿真中以 PPO 算法和完整状态观测（包括物体速度、指尖接触力等特权信息）训练，使用基于连续成功率的性能驱动课程——随智能体表现逐步增加正则化惩罚、动作延迟和时间窗口约束，替代传统广泛域随机化的计算开销。学生策略通过循环信念编码器从噪声本体感知和噪声位姿观测中推断系统隐状态，以复合损失 $L = L_\text{BC} + L_\text{recon}$ 端到端训练，其中行为克隆损失 $L_\text{BC}$ 最小化学生与教师动作的 L2 距离，重建损失 $L_\text{recon}$ 迫使潜在表示能够恢复教师的特权观测。在线 DAgger 算法以初始混合比 $\beta=0.9$ 随机选择教师或学生动作执行，指数衰减 $\beta$ 逐步将控制权移交学生，缓解协变量偏移。在配备腕装 Intel RealSense D435i 相机的 16 自由度 Allegro Hand 上，策略以 30 Hz 推理、300 Hz PD 控制运行，在五种不同几何和物理属性的物体（立方体、球体、药瓶、3D 打印玩具、橡皮鸭）上，在正常和对抗性光照条件下均实现了鲁棒的连续物体重定向。

D-REX（2026.03）~\cite{2603_01151} 提出可微 Real-to-Sim-to-Real 框架，利用高斯泼溅构建视觉-物理双重建：2D Gaussian Splats 提供碰撞几何，3D Gaussian Splats 提供视觉外观，两者共同生成 MJCF 格式的仿真场景。

\textbf{与锚点 ViserDex 的核心差异}在于 3DGS 的用途和仿真到真实的桥接策略：ViserDex 利用 3DGS 的显式高斯属性（球谐系数）作为数据增强界面，在渲染管线内部缩小仿真到真实的视觉域间隙；D-REX 则利用 3DGS 作为物理仿真的几何基础——2D Gaussian Splats 的表面法线信息提供碰撞检测几何，3D Gaussian Splats 的体密度提供视觉外观，两者共同构建可导入 MuJoCo 仿真器的完整物理场景。设计动机是机器人操控中物理参数（特别是物体质量）的不确定性——视觉观测可以恢复几何和外观，但质量参数无法从图像中直接获取，而错误的物理参数会导致仿真策略在真实世界失败。

\textbf{具体技术差异}包括：（1）\textbf{可微物理引擎驱动的参数优化}：通过可微物理引擎（基于 Brax + MuJoCo）在仿真和真实世界执行相同的机器人动作，以轨迹损失 $\mathcal{L}_\text{traj}(m) = \sum_{t=0}^T \|s_t^\text{sim}(m) - s_t^\text{real}\|_2^2$ 优化物体质量参数，无需真实质量标签或接触点标注——直觉上，当仿真中的物体质量与真实物体一致时，相同动作下的轨迹将最吻合；（2）\textbf{力感知混合控制}：在策略学习阶段，将人类手部演示通过 HaMeR + Dex-Retargeting 迁移为机器人轨迹，策略网络同时预测关节位置和与质量相关的抓持力，克服了传统纯位置策略忽略重力补偿导致的抓持不稳定问题。

\paragraph{持续物体跟踪}

POGS（2025.03）~\cite{2503_05189} 将三类特征嵌入紧凑的高斯泼溅表示——基于 Detic 分割的对比学习物体分组特征（用于物体实例分割）、CLIP 语言特征（用于自然语言查询）和 DINOv2 自监督视觉特征（用于跟踪优化）——实现从场景扫描到持续跟踪的完整流水线。

\textbf{与锚点 ViserDex 的核心差异}在于 3DGS 的使用模式和跟踪策略：ViserDex 在仿真环境中构建物体的静态 3DGS 模型，通过逆变换渲染不同位姿下的图像作为数据增强，物体位姿由物理仿真器直接提供；POGS 则在真实世界中从多视角扫描构建场景级 3DGS 表示，并通过在线特征对齐优化持续跟踪物体在人类或机器人操控下的位姿变化——物体位姿不是给定的，而是需要从视觉观测中估计的。

\textbf{具体技术差异}包括：（1）\textbf{三类特征协同}：分组特征通过对比学习将属于同一物体掩码的高斯在嵌入空间中吸引、不同掩码的高斯排斥，再用 HDBSCAN 聚类实现物体实例分割；语言特征通过 CLIP（对比语言-图像预训练模型）文本编码器支持开放词汇查询定位目标物体；DINOv2（自监督视觉Transformer）视觉特征经 PCA（主成分分析）降维至 $d=64$ 维后直接蒸馏到每个高斯的可训练特征向量中，用于跟踪阶段的特征对齐；（2）\textbf{在线跟踪机制}：将每个物体表示为一簇高斯基元，对每一帧新观测，同时渲染模型侧的 DINO 特征图和深度图，与实时提取的真实观测特征图和深度图计算损失，通过梯度下降优化逐物体位姿参数（6D 刚体变换），各物体独立并行优化；（3）\textbf{机器人抓取集成}：语言查询定位目标物体簇后，将该簇的高斯均值作为点云输入 Contact-GraspNet 生成抓取候选，利用 3DGS 的完整 3D 几何信息而非仅从深度图反投影的部分遮挡点云。该系统利用单立体相机即可运行，无需重新扫描或 CAD 模型，实现最多 12 次连续物体重置和 80\% 的抓持中工具扰动恢复。

\paragraph{机器人自建模}

高保真机器人自建模（2025.03）~\cite{2503_05398} 利用 3D 高斯表示机器人的静态形态和纹理，将 3D 高斯通过 K-means 聚类为神经椭球骨骼，每块骨骼由中心点、半径和旋转四元数定义，再通过运动学神经网络将关节角度映射为骨骼级变换。

\textbf{与锚点 ViserDex 的核心差异}在于建模对象和变形机制：ViserDex 建模的是被操控物体的 3DGS 表示，物体变形通过施加逆变换于相机来实现（保持 3DGS 静态）；机器人自建模则建模机器人自身（而非外部物体）的 3DGS 表示，需要处理关节运动引起的非刚性变形。设计动机是机器人需要了解自身形态以进行精确的碰撞检测、运动规划和视觉位姿估计——传统方法依赖人工测量或 CAD 模型，但机器人形态会因磨损、意外损伤而变化。

\textbf{具体技术差异}包括：（1）\textbf{神经椭球骨骼表示}：将 3DGS 的高斯点通过 K-means 聚类为骨骼段，每块骨骼由中心点、半径和旋转四元数参数化，运动学神经网络将关节角度映射为骨骼级 SE(3) 变换；（2）\textbf{混合蒙皮权重}：线性混合蒙皮（LBS）将骨骼变换传递到每个高斯，蒙皮权重由基于马氏距离的原始权重和 MLP 学习的增量权重相加后经 softmax 归一化得到——原始权重提供物理合理的初始化，增量权重补偿关节附近的非刚性变形；（3）\textbf{课程学习训练}：先在静态姿态训练 3DGS 并冻结高斯参数，再以课程学习方式从小范围关节运动逐步扩展到全范围运动（总计 90 万步），配合等距约束（相邻高斯间距不变）、局部刚性约束（局部高斯簇保持刚性）和循环一致性损失确保变形的物理合理性。

\subsubsection{可变形物体与特殊操控}

本路线介绍利用 3DGS 处理可变形物体与特殊操控场景的方法。可变形物体的操控对场景表示提出了更高要求——需要同时建模物体的几何变形和视觉外观变化。

\paragraph{可变形线性物体}

DLO-Splatting（2025.05）~\cite{2505_08644} 修改了 3DGS 的\textbf{场景表示与优化环节}——用\textbf{节点链 + 球形高斯的耦合表示}替代标准 3DGS 的自由高斯集合——3DGS 的 $\alpha$-混合 渲染和可微光栅化管线保持不变。

\textbf{问题背景。}可变形线性物体（Deformable Linear Object, DLO，如绳索、电缆）的操控是机器人 surgery 和工业自动化中的基础任务，但其形状跟踪极具挑战：绳索在密集打结等拓扑复杂场景中频繁自遮挡，仅靠视觉追踪无法区分交叉点的上下层关系。此前基于布料的方法（如 DeformGS、ClothSplatting）在预测步使用图神经网络（GNN）学习动力学，需要大量训练数据；DLO-Splatting 则用纯物理模型替代数据驱动的动力学预测，消除了对训练数据的依赖。

\textbf{预测步：基于位置动力学（PBD）。}绳索被离散化为 $N$ 个节点组成的链 $\mathbf{X}^t \in \mathbb{R}^{N \times 3}$。给定夹持器动作，首先定位被夹持节点 $\mathbf{x}_g^t = \arg\min_i \|\mathbf{x}_i^t - \mathbf{p}_g^t\|$（即距夹持器中心最近的节点），然后通过 Verlet 速度积分预测下一状态：

\begin{equation}
    \mathbf{x}_i^{t+1} = \mathbf{x}_i^t + (\mathbf{x}_i^t - \mathbf{x}_i^{t-1})\,\Delta t + \frac{1}{2}\,\hat{\mathbf{f}}_i^t\,\Delta t^2
\end{equation}

直觉上，这就是``惯性 + 外力''的经典运动预测——当前速度（由前两帧位置差估计）加上外力（重力、法向力、摩擦力）的加速度。每个节点质量为 $m_i = m / N$（总质量均分），外力 $\hat{\mathbf{f}}_i^t = \mathbf{f}_{i,g}^t + \mathbf{f}_{i,N}^t + \mathbf{f}_{i,f}^t$ 包含重力、法向力和库仑摩擦力。随后施加长度约束以保持相邻节点间距恒等于段长 $L$——位置修正量 $\Delta\mathbf{x}_i^{t+1}$ 沿相邻节点连线方向施加，使预测后的绳索既符合物理运动，又保持不可伸缩的绳索长度约束。整个过程无需任何训练数据或学习过程。

\textbf{更新步：3DGS 渲染损失驱动的状态修正。}预测状态 $\mathbf{X}_{\text{PBD}}^{t+1}$ 存在累积误差，需要用视觉观测修正。DLO-Splatting 在绳索中心线上等间距分布 $M$ 个球形 3D 高斯（每段 $d$ 个），协方差矩阵统一设为 $\boldsymbol{\Sigma} = \sigma^2 \mathbf{I}_{3 \times 3}$，其中 $\sigma$ 对应绳索直径——所有高斯大小相同且为球形，因此旋转参数在变形过程中可以忽略，大幅简化了优化。由于 DLO 不呈现非朗伯效应（如金属反射），该方法用 RGB 颜色代替标准 3DGS 的球谐系数，初始颜色由绳索初始分割的平均值确定。

渲染损失定义为多视角观测与渲染图像之间的 $L_2$ 距离：

\begin{equation}
    \mathcal{L}_{\text{obs}} = \|\mathbf{I}_k^t - \hat{\mathbf{I}}_k^t\|_2
\end{equation}

即``渲染出来像不像真实照片''。渲染过程 $\hat{\mathbf{I}}_k^t = h(\mathbf{X}_{\text{PBD}}^{t+1}, \pi_k)$ 将节点链上的高斯投影到第 $k$ 个相机视角，通过 $\alpha$-混合 聚合（与标准 3DGS 完全一致）。关键设计：高斯位置直接由节点链决定（耦合表示），因此优化节点位置就等于优化高斯位置。最终状态更新为：

\begin{equation}
    \mathbf{X}^{t+1} = \mathbf{X}_{\text{GS}}^{t+1} = \mathbf{X}_{\text{PBD}}^{t+1} + \Delta\mathbf{X}^{t+1}, \quad \Delta\mathbf{X}^{t+1} = \arg\min_{\Delta\mathbf{X}}\,\mathcal{L}_{\text{obs}}(\mathbf{X}_{\text{PBD}}^{t+1} + \Delta\mathbf{X})
\end{equation}

直觉上，就是在物理预测的基础上，通过梯度下降微调节点位置，使多视角渲染结果与真实观测一致。这种预测-更新的贝叶斯滤波式框架，使 DLO-Splatting 能够融合多视角信息追踪视觉上无法区分的拓扑结构。

该方法使用三台相机（Intel RealSense D435/D405 + Luxonis Oak-D Pro，30 Hz）在 ABB IRB120 机器人打结场景中验证，比纯视觉方法 TrackDLO 更准确地跟踪绳索尖端抬离平面的运动。该方法使用 gsplat 作为光栅化后端以提升渲染速度和内存效率。

\paragraph{颗粒介质操控}

Gaussian Splatting Visual MPC（2024.10）~\cite{2410_09740} 在高斯泼溅表示上学习视觉动力学模型，通过梯度 MPC 操控颗粒介质（豆类、坚果、大米）。

\textbf{与锚点 DLO-Splatting 的核心差异}在于动力学建模和操控策略：DLO-Splatting 针对可变形线性物体（绳索）采用物理模型（基于位置动力学 PBD）进行预测，并通过 3DGS 渲染损失修正预测状态，是一种物理先验驱动的跟踪方法；Gaussian Splatting Visual MPC 则针对颗粒介质（无数微小粒子的集合）采用数据驱动的图神经网络（GNN）学习动力学——因为颗粒介质的粒子间交互过于复杂，难以用解析物理模型精确描述，但 GNN 的消息传递机制天然适合建模局部交互的传播。

\textbf{具体技术差异}包括：（1）\textbf{初始化策略}：从 RGBD 观测初始化高斯泼溅后，经最远点采样降采样以控制高斯基元数量，移除背景和高透明度的高斯，仅保留表示颗粒材料的部分；（2）\textbf{GNN 动力学模型}：将高斯基元构建为图结构（$L_2$ 距离小于阈值 $\epsilon$ 的节点间连边），通过 SAGE 层编码器提取节点特征，经 $R$ 步迭代消息传递捕获多跳邻居的交互影响，解码器预测每个高斯的平移和旋转位移，以 Chamfer 距离 $L_\text{dyna}$ 训练；（3）\textbf{密度场代价函数}：规划阶段利用高斯泼溅的密度场 $d(x) = \sum_{n_i} \alpha_i \exp(-\frac{1}{2}(x-\mu_i)^T \Sigma_i^{-1}(x-\mu_i))$ 定义占用代价，在基于梯度的 MPC 循环中采样 $K$ 条初始动作序列并迭代优化，实现零样本迁移到新环境。

\paragraph{关节物体操控}

ArtGS（2025.07）~\cite{2507_02600} 将 3DGS 扩展至关节物体的视觉-物理建模，包含三个阶段：先用 3DGS 重建静态关节物体和机器人臂的分别高斯表示；再用微调的 InternVL-2.5 VLM 从正视图推理所有可移动部件的边界框和关节类型（旋转/平移），通过深度图将像素坐标映射到世界坐标并经 PCA 提取关节轴方向；最后在操控过程中利用阻抗控制驱动机器人，以线性混合蒙皮（LBS）将刚性骨骼变换应用到高斯，通过可微渲染优化关节参数和物体状态。

\textbf{与锚点 DLO-Splatting 的核心差异}在于目标物体的类型和变形建模方式：DLO-Splatting 处理可变形线性物体（绳索），其变形通过节点链的位置更新实现，不涉及关节结构；ArtGS 处理关节物体（柜门、抽屉、微波炉），其变形由预定义的关节参数（旋转轴/平移轴 + 原点）驱动，通过线性混合蒙皮（LBS）和可学习蒙皮权重 $W_i \in \mathbb{R}^{K+1}$ 将刚性骨骼变换传递到每个高斯基元。设计动机是关节物体在日常生活和工业环境中的普遍性及其对操控的挑战——关节参数的精确估计直接决定操控轨迹的可行性。

\textbf{具体技术差异}包括：（1）\textbf{VLM 驱动的关节初始化}：利用 3DGS 的新视角合成生成正视图，微调的 InternVL-2.5 VLM 推理所有可移动部件的边界框和关节类型，通过深度图映射和 PCA 提取关节轴方向，为后续优化提供初始参数；（2）\textbf{操控-重建循环优化}：关键区别在于 ArtGS 在操控过程中迭代优化——每次执行操控后用新的 RGB-D 观测通过可微渲染更新关节骨骼参数和物体状态 $\theta_t$，逐步提高建模精度，即使 VLM 初始估计存在显著误差（如关节轴误差超过 $20^\circ$），迭代优化仍能有效修正；（3）\textbf{跨具身重建}：利用修改的 Denavit-Hartenberg（MDH）正运动学分别重建机器人各连杆的高斯表示，支持 Franka、xArm7 等不同机器人臂的高保真数字孪生。

\paragraph{物体扫描与检测}

Omni-Scan（2025.08）~\cite{2508_00354} 利用双手机器人（ABB YuMi）进行全方位物体扫描。

\textbf{与锚点 DLO-Splatting 的核心差异}在于 3DGS 的使用目的和训练范式：DLO-Splatting 在操控过程中在线构建物体的 3DGS 表示用于跟踪，输入是已知相机位姿的多视角观测；Omni-Scan 则解决的是物体扫描本身的问题——如何在夹持器遮挡和运动相机的条件下构建高质量的全方位 3DGS 模型，其核心挑战是反转传统 3DGS 的假设（传统假设静态场景 + 运动相机，而 Omni-Scan 是静态相机 + 运动物体）。设计动机是工业零件检测等应用对全方位高保真 3D 模型的需求。

\textbf{具体技术差异}包括：（1）\textbf{双臂交接扫描策略}：左臂抓取物体并绕腕部旋转 $360^\circ$ 扫描（100 张图像，约 6 分钟），然后通过双臂交接将物体传递给右臂，从新抓取姿态再次扫描以覆盖先前被夹持器遮挡的区域——双臂交接相比单臂重新放置提供更可靠的两次扫描间相对变换；（2）\textbf{3DGS 训练管线修改}：引入物体不透明度损失（$L_1$ 损失约束累积不透明度与物体掩码匹配，消除浮动伪影）和夹持器无关损失（在夹持器遮挡区域将逐像素损失置零，使两个扫描数据集互补训练——一个数据集中被夹持器遮挡的区域可从另一个数据集中学习）；（3）\textbf{模型合并}：两个独立 3DGS 模型通过双臂交接时的手眼变换 $T_{lh}^{-1} T_{rh}$ 对齐（左抓取器到右抓取器的变换），再经标准 ICP 和彩色 ICP 精细化后，以合并的彩色点云初始化训练最终全方位模型，在 12 种工业和家用物体的缺陷检测中平均准确率 83\%。

\paragraph{力感知操控}

DexFruit（2025.08）~\cite{2508_07118} 结合光学触觉感知（DenseTact 传感器）和 3DGS 实现脆弱水果的自主操控与损伤评估。

\textbf{与锚点 DLO-Splatting 的核心差异}在于 3DGS 的双重角色和感知模态：DLO-Splatting 仅利用 3DGS 的多视角渲染损失驱动状态跟踪，感知模态是纯视觉的；DexFruit 则为 3DGS 赋予了两个新角色——（1）作为操控策略的训练环境，通过新视角合成增强策略鲁棒性；（2）作为损伤评估的量化工具，通过比较操控前后的 3DGS 模型变化来量化擦伤。此外，DexFruit 引入了光学触觉感知（DenseTact 传感器）与视觉的融合。

\textbf{具体技术差异}包括：（1）\textbf{自适应视觉-触觉切换}：操控策略采用 Diffusion Policy，以腕部 RGB 相机和双 DenseTact 触觉图像为输入，核心设计是当触觉接触能量超过阈值时关闭视觉输入——这避免模型过度依赖视觉而忽略触觉信号，在脆弱水果操控中尤为重要，因为视觉无法感知抓持力度而触觉可以；（2）\textbf{FruitSplat 损伤评估}：在每个高斯上扩展两个参数 $s_i \in [0,1]$（草莓概率）和 $b_i \in [0,1]$（擦伤概率），通过微调的 YOLOv12 和 SAM2 生成分割掩码监督训练，对操控前后的 3DGS 模型比较擦伤点云比例变化来量化损伤——实现 92\% 的抓取策略成功率和最高 20\% 的视觉擦伤减少。

\subsubsection{其他操控方法}

本路线介绍 3DGS 在机器人操控的若干新兴场景中的应用。除上述按任务类型划分的子方向外，3DGS 还在机器人操控的若干新兴场景中展现出独特价值，涵盖可微渲染、仿真到真实迁移、新视角抓取、语言引导操控、手术机器人等 20 条技术路线。

\paragraph{可微机器人渲染}

可微机器人渲染（2024.10）~\cite{2410_13851}（Dr.\ Robot）将机器人运动学感知变形模型与高斯泼溅集成，使机器人身体的视觉外观直接对控制参数可微。其核心由三个组件构成：以高斯泼溅建模标准姿态外观、隐式线性混合蒙皮（LBS）根据目标姿态重定位 3D 高斯、姿态条件外观变形模型自适应调整球谐系数和协方差矩阵。这种全链路可微表示允许通过梯度反向传播将视觉基础模型的信号直接转化为控制指令，避免无梯度优化的低效性，在从野外视频重建机器人位姿任务上以较大优势超越此前最优方法。

\paragraph{运动-语义-几何统一}

MSGField（2024.10）~\cite{2410_15730} 修改了 3DGS 的\textbf{基元类型与属性表示环节}——用\textbf{2D 高斯}替代标准 3D 高斯，并在每个基元上同时编码语义特征和运动参数，以实现运动-语义-几何（Motion-Semantics-Geometry, MSG）的统一表示——高斯泼溅的可微实时渲染管线保持不变。

\textbf{与锚点可微机器人渲染的核心差异}在于建模对象、基元类型和运动建模方式：可微机器人渲染用 3D 高斯表示机器人自身的静态外观，通过运动学驱动的线性混合蒙皮（LBS）将姿态参数传递到每个高斯——运动来源是已知的机器人关节角度，变形是刚性骨骼变换；MSGField 则用 2D 高斯表示场景中的所有物体（包括柔性、极小物体），通过可学习运动基的加权组合建模任意非刚性运动——运动来源是视觉观测，需要从稀疏视角图像中优化得到。

\textbf{设计动机}：可微机器人渲染解决的是``已知控制参数→渲染视觉外观''的正问题；MSGField 解决的是其逆问题——``已知视觉观测→估计物体运动''——在语言引导操控中，机器人需要同时理解物体的身份（语义）、位置（几何）和运动方式（动力学），才能规划抓取动作。选择 2D 高斯而非 3D 高斯有两个原因：（1）2D 高斯通过深度失真损失约束基元贴合表面，避免 $\alpha$ 混合深度落入基元之间的``悬浮''伪影，提供更精确的几何表示；（2）$\alpha$ 混合的语义特征会失去含义（如混合红色和绿色得到黄色是合理的，但混合``杯子''和``盘子''的语义特征是无意义的），2D 高斯的表面贴合特性使得每个像素仅由少量基元主导，减轻了语义特征的混合退化。

\textbf{具体技术差异}可以按四步读。
第一步是对象级语义场蒸馏：YOLO-World 和 SAM 先给出对象掩码，射线-泼溅交叉再判断每个 2D 高斯是否属于某个对象。
交叉比超过阈值的基元被标记为对象成员，背景伪影再用梯度阈值过滤。
最终，同一对象的所有基元共享一个 CLIP 特征 $F_\text{obj}$，并用一维标签属性索引节省存储。
相对只编码外观颜色的可微机器人渲染，MSGField 给基元额外绑定了语义身份。

第二步是运动基分解。
每个基元从时刻 $0$ 到时刻 $t_i$ 的运动不再独立学习，而是写成 $B$ 个可学习运动基的加权组合：

\[
\Delta_i = \sum_{b=1}^B w^{(b)} \cdot (\delta^{(b)}_{0 \to t_i}, \delta q^{(b)}_{0 \to t_i}).
\]

其中平移和旋转基由所有基元共享，权重 $\{w^{(b)}\}_{b=1}^B$ 由每个基元私有。
这个分解把运动自由度从 $N \times T \times 6$ 压缩为 $B \times T \times 6 + N \times B$，使两个相机视角也能优化非刚性运动。

第三步是 Dice 损失。
稀疏帧采样时，相邻帧视角重叠不足，逐像素 $L_1/L_2$ 监督容易失效。
MSGField 用 $L_\text{Dice} = 1 - \text{Dice}(I_\text{opacity}, I_\text{mask})$ 监督运动掩码，更关注整体区域重叠，而不是每个像素是否完全对齐。

第四步是语言引导操控。
给定语言指令后，系统通过 CLIP 文本-特征匹配定位目标对象的基元簇，再在 NDC（Normalized Device Coordinates）空间做视锥裁剪，选择遮挡较少的视角渲染给机械臂。
在包含柔性物体和极小物体的数据集上，静态环境成功率为 79.2\%，动态环境成功率为 63.3\%，指定物体抓取成功率为 90\%。

\paragraph{仿真到真实操控学习}

与 MSGField 在真实场景中直接编码运动基以表示动态场景不同，RL-GSBridge（2024.09）~\cite{2409_20291} 改变了 3DGS 在仿真管线中的角色——将高斯泼溅从场景表示提升为 RL 仿真的\textbf{可编辑照片级渲染器}，构建 Real2Sim2Real 闭环：先用消费级相机采集多视角图像重建 3DGS 场景模型（Real2Sim），再在物理仿真器中通过动力学驱动的 GS 编辑生成真实感训练图像供 RL 策略学习（Sim），最后零样本迁移至真实机器人（Real）。

\textbf{设计动机。}传统 RL 仿真使用低多边形网格渲染，与真实世界的视觉域间隙（domain gap）大——仿真中训练的策略在真实环境中性能骤降。缩小域间隙的常规手段是域随机化（随机化纹理、光照、相机参数以覆盖真实世界的多样性）或域适应（学习仿真与真实的共享特征空间），但这些方法要么需要大量增强数据，要么依赖大型学习模型，对特定任务效率低下。RL-GSBridge 的核心洞察是：用辐射场重建（3DGS）创建视觉逼真的仿真环境即可从根本上缩小域间隙——仅用消费级手机拍摄 1--2 分钟视频即可构建照片级仿真场景，无需昂贵的 3D 扫描设备或 CAD 建模专长。

\textbf{软绑定约束：突破硬绑定的灵活性瓶颈。}在 Real2Sim 建模阶段，RL-GSBridge 需要在保持几何可编辑性的同时实现高质量渲染。此前 GaMeS 方法将高斯位置硬约束为三角网格顶点的凸组合 $\mu^i = \alpha_1^i \mathbf{v}_1 + \alpha_2^i \mathbf{v}_2 + \alpha_3^i \mathbf{v}_3$（$\sum \alpha_j^i = 1$），确保高斯严格贴合网格面——但这导致两个问题：（1）当网格模型本身不够精确时（如 openMVS 从消费级图像重建的 mesh），高斯被``锁死''在不精确的表面上，无法自由优化以补偿网格误差；（2）对于非结构化物体（如毛绒玩具、柔性包装），硬绑定使高斯分布过于僵硬，无法拟合表面细节。RL-GSBridge 的软绑定在凸组合基础上引入沿面片法向量方向的浮动分量：

\begin{equation}
    \mu^i = \alpha_1^i \mathbf{v}_1 + \alpha_2^i \mathbf{v}_2 + \alpha_3^i \mathbf{v}_3 + \alpha_n^i \mathbf{v}_n, \quad \alpha_n^i \in [-1, 1]
\end{equation}

其中 $\mathbf{v}_n$ 为面片法向量，$\alpha_n^i$ 为可学习参数。直觉上，这允许每个高斯在法向方向上``浮起来或沉下去''——仍然与网格保持关联（通过 $\mathbf{v}_1, \mathbf{v}_2, \mathbf{v}_3$ 的凸组合），但拥有沿法向微调的自由度，使高斯分布更贴合真实表面，渲染质量显著提升（PSNR 提高约 1--2 dB），甚至为非刚性物体的处理提供了可能性。

\textbf{物理动力学驱动的 GS 编辑：仿真-渲染同步。}获得 GS 模型后，RL-GSBridge 通过 RANSAC 平面回归和手动对齐将 GS 模型与仿真器中的几何模型注册到同一坐标系。在 RL 训练的每一步，仿真器提供每个物体的实时位姿变化（旋转四元数 $q_k$ 和齐次变换矩阵 $T_k$），GS 编辑模块据此执行刚性变换更新每个高斯的位置和旋转：

\begin{equation}
    \hat{\mu}_k^i = T_k\,\mu_k^i, \quad \hat{r}_k^i = q_k\,r_k^i
\end{equation}

这一过程确保渲染结果与物理交互过程完全同步——机器人在仿真中抓取物体时，GS 渲染器输出的图像真实地反映了物体的位移和旋转，消除了视觉感知与接触几何之间的不一致。

\textbf{Sim2Real 迁移策略。}RL 策略基于 SAC（Soft Actor-Critic）训练，采用 SACwB（SAC with Baseline）策略——以概率 $\eta$ 执行基线控制器的动作而非策略网络输出，并通过行为克隆损失监督演员网络，使策略在学习过程中保持稳定。在抓取和拾取-放置任务上，仿真中训练的策略直接零样本迁移至真实环境，成功率变化极小（如抓取任务仿真 93\% 对真实 87\%），证明 GS 渲染的真实感足以有效缩小视觉域间隙。

\paragraph{新视角抓取}

与锚点 RL-GSBridge（2024.09）~\cite{2409_20291} 将 3DGS 作为仿真管线的可编辑渲染器不同，基于辐射场的抓取生成（2025.05）~\cite{2505_11467} 改变了 3DGS 在抓取管线中的角色——不修改高斯基元本身（既不变形也不编辑），而是利用其\textbf{新视角合成}能力扩展抓取检测器的输入视野。

\textbf{核心差异}：RL-GSBridge 的 3DGS 场景在仿真中由物理动力学驱动，高斯基元随物体运动实时变换（$\hat{\mu}_k^i = T_k \mu_k^i$），3DGS 本身是被操控的对象；而本方法中 3DGS 场景是静态的桌面重建，高斯基元在优化后保持不变，其价值在于从真实相机无法到达的视角合成高质量的颜色和深度图像，为预训练抓取检测器提供额外的观测输入。

\textbf{设计动机}：眼在手（eye-in-hand）机器人获取多视角图像需要移动机械臂到不同位姿，时间开销随视角数增长且受可达性约束；而辐射场一旦从少量真实视角构建完成，即可零成本地从任意新视角渲染图像——核心假设是这些合成视角能提供真实视角无法覆盖的物体几何信息（如被遮挡的底面或侧面），从而生成额外的力闭合抓取。

\textbf{技术实现}：从 Graspnet-1billion 数据集中每个场景的 256 个视角中选择 $M=3$ 个真实视角（模拟眼在手相机的最小运动场景），利用 SplaTAM（一种基于 3DGS 的 RGB-D SLAM 系统）优化场景表示——对每帧 RGB-D 数据执行高斯致密化（在致密化掩码区域初始化新高斯）和地图更新（最小化所有关键帧的 RGB 和深度误差），由于数据集提供精确相机位姿，无需相机跟踪步骤。优化完成后从 $N=16$ 个接近真实视角的新视角渲染颜色和深度图像——新视角选择靠近真实视角且朝向相似，原因是预训练抓取检测网络在训练分布内的视角上表现更好。通过将优化后的高斯投影至新视角生成点云，送入 Graspnet-1billion 预训练抓取检测网络推理抓取位姿。后处理采用三路并行策略：（1）pose-NMS 合并平移距离 $< 0.03\text{m}$ 且旋转距离 $< 15^\circ$ 的冗余抓取；（2）聚类 + top-grasp 过滤；（3）无后处理保留全部抓取。

最终以 Dex-net 2.0 的力闭合指标（五种摩擦系数 $\mu \in \{0.2, 0.4, 0.6, 0.8, 1.0\}$，任一 $\mu$ 下力闭合即为成功）评估抓取质量。实验验证了核心假设：新视角不仅贡献了额外的力闭合抓取（$G_\text{nvs}$ 中存在 $G_\text{real}$ 中未出现的有效抓取），还增加了场景中可有效计算抓取位姿的物体数量（即抓取覆盖率），改善了整体抓取覆盖范围。重建质量方面，190 个场景的平均 PSNR 为 30.608、MS-SSIM 为 0.984、LPIPS 为 0.053、深度 L1 损失为 0.105，确认了合成视角的几何保真度。

\paragraph{人体到机器人递送}

H2RH-SGS（2025.07）~\cite{2507_08726} 和相关方法（2025.08）~\cite{2508_09855} 利用稀疏视角 3DGS 重建人体到机器人递送场景，生成包含图像-动作对的机器人示范数据。具体而言，H2RH-SGS 使用深度图初始化点云获得度量尺度，通过 GraspNet 生成安全抓取位姿并过滤与人体手部碰撞的候选，在重建场景中球面采样初始位姿并经三阶段轨迹规划（对齐-接近-精化）生成仿真示范，最终以监督学习训练策略直接在真实环境中部署，无需真实机器人训练数据。

\paragraph{手-物交互生成}

FlowHOI（2026.02）~\cite{2602_13444} 提出基于流的语义接地手-物交互生成框架，以自我中心观测、语言指令和 3DGS 场景重建为条件生成时序一致的手-物交互序列。在 GRAB 和 HOT3D 基准上实现最高动作识别精度和 1.7 倍于最强扩散基线的物理仿真成功率，推理速度提升 40 倍。

\paragraph{视觉编码器对齐}

VEGA（2026.05）~\cite{2605_10485} 将 VLA 视觉编码器输出与来自 DINOv2-FiT3D（使用多视角一致 3DGS 监督微调）的空间感知特征对齐，在视觉编码器输出层面进行对齐以避免语言语义纠缠，推理时无额外计算开销。

\paragraph{结构化潜在点}

结构化潜在点学习（2026.05）~\cite{2605_21258} 在点云自编码器的潜在空间插入逐点变分自编码器，联合正则化逐点特征和坐标趋向高斯先验，产生的紧凑潜在表示保留了粗结构倾向，结合轻量 3DGS 渲染管线在 RLBench、ManiSkill2 和真实机器人平台上展现一致增益。

\paragraph{空中操控}

AirVLA（2026.03）~\cite{2603_25038} 研究将操控预训练 VLA（$\pi_0$）迁移至空中拾取-放置任务，利用 3DGS 管线合成导航训练数据，在 460 次真实世界实验中验证合成数据的关键作用——解锁导航任务 100\% 成功率。

\paragraph{操控仿真}

GSWorld（2025.10）~\cite{2510_20813} 提出 GSDF（Gaussian Scene Description File）格式，将高斯-网格表示与机器人 URDF 融合，支持零样本仿真到真实操控策略学习、自动 DAgger 数据收集和可复现策略基准评测。DISCOVERSE（2025.07）~\cite{2507_21981} 提出首个统一、模块化、开源的 3DGS 仿真框架，通过全面 Real2Sim 管线合成超逼真几何和外观。

\paragraph{接触点跟踪}

基于 3DGS 的物体和接触点跟踪（2024.11）~\cite{2411_03555} 从视频示范中提取触摸交互点并跟踪物体运动，为交互式模仿学习提供关于交互位置和方式的详细知识。该方法首先用 3DGS 重建包含目标物体的场景，再通过 RAFT 光流和 SAM2 从示范视频中分割物体，经 GS2Mesh 提取网格后由 FoundationPose 跟踪 6-DoF 物体位姿，最终利用深度图和手部掩码的距离阈值定位接触点，为铰接物体操控提供精确的交互位置信息。

\paragraph{机器人自建模操控}

MATT-GS（2025.03）~\cite{2503_19330} 利用掩码注意力 3DGS 增强工业和智能工厂环境中的机器人感知和物体检测。该方法先用 U2-Net 背景移除将原始图像中的目标物体与杂乱背景分离，确保模型仅处理相关数据；再在 3DGS 框架中集成 Sobel 滤波器注意力机制，增强边缘响应以捕获螺丝、电线和精细纹理等工业场景关键特征。在 L1 损失、SSIM 和 PSNR 指标上，背景移除与注意力机制联合使用均显著优于原始 3DGS 基线。

\paragraph{物理信息世界模型}

PIN-WM（2025.04）~\cite{2504_16693} 利用可微物理仿真从视觉观测端到端识别 3D 刚体动力学系统，通过 3DGS 诱导的观测损失学习，无需状态估计。物理感知数字表亲（Digital Cousins）通过扰动物理和渲染参数生成多样变体，实现鲁棒非预hensile操控的仿真到真实迁移。

\paragraph{反应式移动操控}

ReMoSPLAT（2025.12）~\cite{2512_09656} 基于二次规划的反应式控制器利用高斯泼溅表示进行碰撞避免，协调移动操控器的基座和手臂运动。该方法将机器人几何近似为球体集合，比较两种距离计算策略：球-椭球纯几何方法计算高效但忽略不透明度信息，在低不透明度高斯场景中可能产生虚警；基于多视角深度光栅化的虚拟传感器方法通过六方向虚拟相机聚合不透明度信息，在复杂场景中保持鲁棒。两者均通过速度阻尼器不等式约束和主动避碰代价项集成至 QP 控制器。

\paragraph{语言引导链式思维操控}

ManualVLA（2025.12）~\cite{2512_02013} 在 MoT 架构上统一多模态手册生成和动作执行，通过 3DGS 构建的高保真数字孪生工具包自动生成规划专家训练的手册数据。在乐高组装和物体重排任务上平均成功率比层次化 SOTA 基线高 32\%。

\paragraph{具身视觉表示学习}

几何接地 3D 视觉表示（2026.01）~\cite{2601_22988} 提出统一表示-策略学习框架，利用点云重建和前馈高斯泼溅在多视角监督下学习整体几何表示，通过多步蒸馏将预训练几何理解有效迁移至操控技能。在 12 个 RLBench 任务上平均成功率超越当时最优方法 12.7\%。

\paragraph{从人类示范学习}

WARPED（2026.04）~\cite{2604_10809} 从自我中心人类示范视频合成真实腕部视角观测训练视觉运动策略，利用视觉基础模型初始化交互场景，通过手-物交互管线跟踪和重定向轨迹，最后通过 3DGS 合成照片级真实感腕部视角。该方法在五个桌面操控任务上达到遥操作示范数据策略可比的成功率，数据收集时间减少 5--8 倍。

\paragraph{机器人示教感知}

GraspSplat（2025.10）~\cite{2510_19200} 利用 3DGS 生成手-物交互的高保真新视角，通过光度损失精化灵巧抓取预测。其核心思想是将抓取合成与分析-合成范式的逆过程结合：先用 Vision Transformer 和编码器-解码器从 RGB 图像预测 MANO 手部参数，再将手部 3DGS 模型与物体 3DGS 模型拼接后从随机视角渲染，光度损失（$L_1$ + SSIM）使误差可微反向传播至网络预测，在 YCB 数据集上比现有纯图像方法提升高达 36.9\% 的抓取成功率。

\paragraph{灵巧操控特征蒸馏}

LensDFF（2025.03）~\cite{2503_03890} 通过语言增强特征融合策略高效地将视角一致的 2D 特征蒸馏至 3D 点，实现单视角少样本灵巧抓取泛化。针对 CLIP 视觉特征跨视角不一致的问题，该方法利用语言特征方向更稳定的特性，将每个点的视觉特征投影至对应语言特征方向：$f_i^{\text{aligned}} = \langle f_{\text{vis}}(x_i), f_{\text{lan}} \rangle / \|f_{\text{lan}}\|^2 \cdot f_{\text{lan}}$，无需额外训练即可实现跨视角对齐。结合五种抓取原语（钩握、柱握、捏握、三指握、蚂状握）的本征抓取降维，在真实灵巧手上实现稳定且高度灵巧的抓取。

\paragraph{手术机器人}

arg-VU（2026.03）~\cite{2603_26814} 将物理感知功能区域推理引入手术机器人，利用 3DGS 重建手术场景并转换为时序跟踪表面表示。与传统仅依赖视觉线索的方法不同，该方法在跟踪几何上嵌入基于 XPBD 的变形约束（静水压约束保体积、偏差约束抗剪切），通过约束灵敏度雅可比矩阵构建各向异性刚度度量 $K_i = J_i^\top J_i + \epsilon I$，量化工具运动在组织各方向上的机械兼容性。由此定义物理感知柔性得分（PACS），将末端执行器运动映射至局部组织响应，在可变形手术环境中实现可靠的功能区域推理。

\begin{takeaway}
机械臂操控中 3DGS 发挥三重作用：作为几何表示提供抓取规划约束，作为可渲染环境生成多样化训练数据，作为可编辑场景模拟物体操作后的变化，语言增强特征蒸馏则实现了少样本灵巧抓取泛化。
\end{takeaway}

% ======================================================================

\subsection{近期方法脉络：具身仿真与策略学习、高斯世界模型、操作数字孪生}

这一组近期工作延续 3DGS、3D Gaussian Splatting，主要把改动集中在应用、适配、扩展等环节。下面只保留相对既有路线的变化，便于和前文方法直接对照。

\paragraph{LEGS}
LEGS（2026.05）~\cite{2606_01458} 延续 3DGS 在“具身仿真与策略学习”上的基本设定，主要改变是：LEGS：在3DGS重建背景上叠加mesh机器人/物体，合成无遥操作示范并校准到真实相机。3DGS 的几何-渲染统一表示保持为底层地图。
其设计动机是：提出Embodied Gaussian Splatting混合仿真环境，把真实场景3DGS背景、mesh前景和程序化动作原语结合起来，用于VLA策略微调和sim-to-real。
具体而言，系统从手持采集重建部署环境3DGS，将机器人、物体和道具作为mesh前景合成，并通过两阶段颜色校准对齐真实机器人相机。
适用边界是：物理交互主要由mesh前景承担，3DGS背景不可交互；任务类型和场景范围有限。

\paragraph{Learning Action-Conditional and Object-Centric Gaussian Splatting World Models for Rigid Objects}
Learning Action-Conditional and Object-Centric Gaussian Splatting World Models for Rigid Objects（2026.06）~\cite{2606_01950} 延续 3D Gaussian Splatting 在“高斯世界模型”上的基本设定，主要改变是：MRO-GWM：用物体中心canonical Gaussians表示刚体对象，并用时空Transformer预测未来刚体运动。3DGS 的几何-渲染统一表示保持为底层地图。
其设计动机是：提出面向刚体对象的动作条件Gaussian世界模型，兼顾任意形状对象表示、遮挡下部分观测和未来动作驱动预测。
具体而言，每个对象在规范坐标系中由一组高斯表示，未来动作通过时空Transformer预测对象SE(3)运动，再把高斯变换到未来状态用于渲染和推理。
适用边界是：主要针对刚体对象；复杂非刚体、接触物理和真实机器人长时任务仍需扩展。

\paragraph{ManiSplat}
ManiSplat（2026.06）~\cite{2606_10645} 延续 3DGS 在“操作数字孪生”上的基本设定，主要改变是：ManiSplat：图结构解耦高斯表示，将机器人、物体和背景组织为独立可优化Gaussian subfields。3DGS 的几何-渲染统一表示保持为底层地图。
其设计动机是：提出Graph-Structured Disentangled Representation，把复杂接触交互场景拆成机器人、物体和背景子高斯场，并用任务导向约束提高操作轨迹合成稳定性。
具体而言，系统从单目ego-view视频中恢复场景图和多个Gaussian subfields，各子场可独立优化和控制，便于生成或重放机器人操作轨迹。
适用边界是：单目输入下接触和遮挡仍难；图结构和分割质量影响解耦；复杂可变形物体未充分覆盖。
\section{具身感知与交互}
\label{sec:17_3}

具身感知与交互是连接机器人感知和行动的桥梁。3DGS 在这一领域的应用涵盖遥操作、人机协作、主动感知和多模态记忆构建。本节从遥操作系统、数字孪生与混合现实、多模态感知与记忆、以及其他具身应用四个方面展开。

\subsection{遥操作与数字孪生}
\label{ss:17_3a}

本节讨论两条技术路线：（1）遥操作与远程呈现；（2）数字孪生与混合现实。

\subsubsection{遥操作与远程呈现}

本路线介绍利用 3DGS 渲染能力实现沉浸式遥操作与远程呈现的方法。遥操作要求操作者能够准确感知远程环境并进行精确控制。3DGS 的照片级真实感渲染和实时性能为沉浸式遥操作提供了理想的视觉表示。

\paragraph{沉浸式遥操作}

Reality Fusion（2024.08）~\cite{2408_01225} 修改了 3DGS 的场景表示环节——用「预构建 3DGS 全局场景 $+$ 实时深度点云」的双表示融合替代传统的单一 3D 表示（点云/网格）——3DGS 的泼溅渲染管线保持不变，新增的是将实时点云投影至同一世界坐标系的融合管线。

\textbf{传统做法与局限}：在 Reality Fusion 之前，沉浸式遥操作系统主要依赖两类 3D 表示来呈现远程环境。其一为实时点云，即直接将机载 RGBD 传感器的深度数据投影到 VR 空间——优点是延迟低、反映实时变化，但受限于传感器有限视场（通常水平 85$^\circ$、垂直 54$^\circ$），离散点会产生孔洞和遮挡，操作者难以建立对全局环境的完整心理模型。其二为基于动态 SLAM（如 Voxblox）的在线网格重建——可以逐步构建环境几何，但重建精度对复杂几何（如工业管道、加速器隧道内的密集设备）不够准确，且 SLAM 引入数秒级的重建延迟，不适合需要即时干预的遥操作场景。此外，多相机或 360$^\circ$ 全景相机方案虽能扩大视场，但高带宽视频流引入显著传输延迟，在 VR 中易导致运动眩晕。

\textbf{设计动机}：在受限工业环境（如 CERN 粒子加速器隧道）的机器人巡检场景中，操作者同时需要两类互补信息：全局空间感知（用于任务规划和路径决策）和实时视觉反馈（用于精确控制和即时干预）。上述传统方案均无法同时满足——3DGS 提供照片级真实感的全局环境渲染却无法呈现实时变化，实时深度点云提供即时反馈但缺乏环境全局感知。Reality Fusion 的核心思想是：在半静态工业场景中，环境的全局结构不会频繁变化，因此可以用预构建的 3DGS 模型离线提供高保真全局上下文，同时仅需单个轻量化深度传感器提供实时局部反馈，两者在统一坐标系中融合即可兼顾质量与实时性。

\textbf{技术实现}：Reality Fusion 的融合管线将两种异构 3D 表示统一到同一 VR 空间坐标系中。对于实时感知，系统通过立体相机的视差估计获取深度图 $d$——即利用左右视图的像素对应关系反算每像素的深度值。获取深度后，利用立体投影变换 $P_C = K^{-1} P_d$ 将 2D 像素反投影至 3D 相机坐标，即把图像平面上的每个像素还原为相机坐标系中的 3D 空间点。再通过全局变换
\begin{equation}
    P_W = M_\text{VP} \times T_\text{offset} \times M_\text{camera} \times P_h
\end{equation}

将点云从相机坐标系变换至世界坐标系，其中 $M_\text{camera}$ 描述相机在全局世界空间中的位置和朝向（由内置 IMU 提供 6-DoF 追踪），$M_\text{VP}$ 为 Unity 虚拟相机的视图投影矩阵，$T_\text{offset}$ 通过手动标定补偿机器人追踪中心与实际相机位置之间的平移偏移。直观地说，该变换将机器人前方看到的实时点云精确放置到 VR 世界中的对应位置。对于全局环境，系统预构建的 3DGS 模型通过标准泼溅渲染提供任意视角的高分辨率场景视图。关键在于两种表示共享统一的世界坐标系——3DGS 模型中的机器人位姿通过 ICP 算法与 CAD 模型对齐，确保实时点云与预构建场景无缝衔接。在实际部署中，系统在 RTX 2080 上对包含约 $622{,}335$ 个高斯基元的工业场景达到 40--45 fps（单眼 $900 \times 960$），在 RTX 3090 上对约 $727{,}019$ 个高斯基元的房间级场景达到 30--35 fps（单眼 $1536 \times 1440$），端到端点云流延迟约 $153 \pm 33$ ms。

\textbf{效果}：这种融合方案仅需单个立体相机（ZED Mini）作为实时传感器（满足轻量化平台需求 G1），同时实现低延迟渲染（满足实时控制需求 G3），并支持 6-DoF 视角切换和外视角/内视角机器人控制（满足态势感知需求 G2 和无运动眩晕需求 G4）。24 人用户研究表明，Reality Fusion 在用户性能（任务完成时间显著缩短）、态势感知（SART 得分显著提高）和用户偏好（12/24 偏好外视角融合模式，8/24 偏好内视角融合模式）方面显著优于纯点云基线（C1）。与 Reality Fusion 将预构建 3DGS 场景与实时深度点云在线融合不同，Boehringer 等人的沉浸式 locomanipulation 遥操作框架（2025.04）~\cite{2504_15229} 改变了 3DGS 场景的交互模式——将 3DGS 作为 VR 操控阶段的\emph{全静态}场景副本，操作者不再通过外部相机观察场景，而是直接「置身于」高斯泼溅重建的 3D 环境中，通过游戏式的「抓取-拖拽」交互控制机械臂末端执行器。

\textbf{设计动机}：Reality Fusion 的融合方案需要实时传感器持续提供深度流，适合半静态工业巡检场景；而在移动操控平台（locomanipulation）场景中，机器人需要先导航至目标位置、再执行精细操控——这两个阶段对感知的需求截然不同（导航阶段需广角实时视野，操控阶段需高保真局部几何），因此该框架采用分阶段策略：导航阶段使用传统 2D 视频流（投影至 VR 内固定屏幕以减少运动眩晕），到达目标后切换至操控阶段，离线构建高斯泼溅场景。

\textbf{技术实现}：3DGS 场景构建方面，机械臂移动至预规划位姿序列采集重叠图像，经 SfM 估计相机位姿和三角化 3D 点后训练高斯泼溅模型，在 Unity VFX 图中实现为节点化粒子系统。

操控交互方面，不同于现有系统将控制器锚定在操作者第一人称视角坐标系（限制自由视角移动），该框架将 3DGS 重建的机械臂模型和操控场景直接放置在操作者前方，操作者可在 VR 空间中自由走动观察场景，通过 VR 控制器「抓取」虚拟末端执行器并「拖拽」至目标位姿，对应笛卡尔坐标通过 ROS 反馈回路传输至真实机械臂控制器。为弥补 3DGS 静态重建无法反映抓取后物体状态变化的问题，基座和末端执行器的实时相机画面以画中画形式叠加在 VR 视图中作为参考。

\textbf{具体差异}：相比 Reality Fusion 的双表示在线融合管线（预构建 3DGS $+$ 实时点云，需持续深度流和 ICP 对齐），该框架的 3DGS 重建是完全离线的一次性快照，不依赖实时传感器；交互方式从 Reality Fusion 的外视角/内视角切换（通过 VR 头显自由选择观察位置）升级为直接操控——操作者不是「看着」场景进行控制，而是「站在」场景中拖拽虚拟机器人手臂，真实机械臂的关节角度同步回传至 Unity 镜像虚拟模型运动。

\textbf{效果}：15 人用户研究（遮挡场景中的拾取-放置任务，对比双相机摇杆基线）表明，93\% 参与者整体偏好高斯泼溅界面，66\% 参与者更快完成任务（平均时间减少 43\%）；量化评分中泼溅界面在易用性（5.53 vs.\ 4.33）、视觉反馈清晰度（5.80 vs.\ 4.87）、态势感知（6.13 vs.\ 4.13）、沉浸感（5.87 vs.\ 4.33）和认知负荷（5.60 vs.\ 4.40）上均显著优于基线（7 分制），100\% 参与者推荐未来使用。在真实机器人（Robotnik Summit-XL 移动底盘 $+$ Franka Emika 机械臂）上验证了导航-操控两阶段流程的可行性。

\paragraph{辐射场遥操作}

Wilder-Smith 等人（2024.07）~\cite{2407_20194} 系统性地将在线辐射场（NeRF 和 3DGS）替代传统遥操作管线中的重建-可视化组件。该系统将遥操作流程简化为机器人与传感器、重建方法和可视化系统三个可替换模块，每种模块均可根据部署场景灵活配置。系统在静态机械臂、移动四足机器人和带臂四足机器人三种配置上进行了验证，可视化端同时支持 2D RViz 屏幕窗口和 VR 头显中的 2.5D/3D 视图。相比基于 Voxblox 的传统网格重建基线，辐射场方法提供了更高质量的场景表示和更灵活的视角控制。

\paragraph{实时照片级映射}

基于高斯泼溅 SLAM 的实时照片级映射（2025.09）~\cite{2509_06433} 提出模块化高效 GPU 集成方案，将高斯泼溅 SLAM（基于 CaRtGS）与在线地图遥操作系统结合。系统的核心挑战在于优化器更新高斯参数与遥操作员交互式查看地图需要同时渲染，传统交替渲染方案会导致两者性能互相干扰。该方案利用 CUDA 进程间通信（IPC）实现优化器和交互式查看器对 GPU 上 GS 地图的并行访问，彻底解耦地图生成与交互渲染。交互式查看器基于 Unity 构建，支持 VR 头显多视角同时渲染和游戏手柄双摇杆控制（分别控制机器人目标位置和虚拟相机视角），真实世界四旋翼飞行器实验表明该方法显著改善决策速度和环境交互准确性。

\paragraph{安全监控}

TAIL-Safe（2026.05）~\cite{2605_01195} 利用 3DGS 构建高保真数字孪生系统收集失败数据，训练 Lipschitz 连续 Q 值函数映射状态-动作对至长期安全分数。与 Physically Embodied GS 的物理仿真闭环不同，TAIL-Safe 改变了 3DGS 的应用环节——不用于实时视觉-物理交互，而是作为数据生成引擎从约 100 张 RGB 图像（约 20 分钟采集加 15 分钟重建）快速构建数字孪生，用于在仿真中系统性探索策略失败模式。设计动机在于：学习安全集合的边界本质上需要失败轨迹数据，但在真实机器人上收集此类数据既危险又低效。TAIL-Safe 的核心创新是任务无关的安全启发式函数 $h(s,a)$，融合可见性（目标物体在视场内）、可识别性（观测与专家分布的 Mahalanobis 距离）和可抓取性（末端执行器与抓取候选的对齐度）三个标准，通过可学习权重网络动态融合为标量安全分数，并训练 Reach-Avoid Q 函数传播轨迹中的最低安全分数。当名义策略提出不安全动作（$Q < 0$）时，基于 Nagumo 定理的恢复机制利用 Q 函数的梯度上升（通过谱归一化保证 Lipschitz 连续性）引导策略回到安全集合，在桌面操控任务中实现 100\% 恢复成功率和 99.3\% 状态级准确率。

\subsubsection{数字孪生与混合现实}

本路线介绍基于 3DGS 构建数字孪生与混合现实环境的方法。3DGS 构建的数字孪生为机器人在虚拟环境中进行安全测试和训练提供了高保真平台。与基于网格的传统数字孪生不同，3DGS 数字孪生可以从真实照片直接重建，无需人工建模，且渲染质量接近照片级真实。

\paragraph{物理感知交互}

传统的 3DGS 仅作为静态场景的视觉表示，无法编码碰撞、重力等物理先验，因此既不能预测物体在受力后的未来状态，也无法在动态场景中保持物理一致性。已有的动态 3DGS 方法（如 D3DGS）将物理约束作为辅助损失融入高斯优化，但这一策略存在两个根本局限：一是无法作为前向模型预测``如果施加力会发生什么''，二是需要每帧提供前景掩码，在仅有 3 个相机的实时设定下难以鲁棒运行。

Physically Embodied Gaussian Splatting（2024.06）~\cite{2406_10788} 修改了 3DGS 的场景表示与状态更新机制——用``粒子-高斯''双表示替代纯高斯表示，通过物理引擎驱动的预测-校正闭环替代直接梯度优化——高斯基元的可微泼溅渲染流程保持不变。

该方法的\textbf{核心创新}在于引入\textbf{``高斯-粒子''（Gaussian-Particle）双表示}，用一个刚体变换（bond，即记录高斯基元相对于父粒子的局部坐标偏移与朝向）将每个高斯基元挂载到其父粒子上，从而在物理仿真和可微渲染之间建立桥梁。直观地说，每个高斯基元就像用刚性连杆悬挂在粒子上的灯笼——灯笼负责向观察者展示物体的外观，连杆则保证灯笼始终跟随粒子的物理运动轨迹。

具体而言，每个场景包含两类实体：$N$ 个\textbf{粒子}负责捕获物体的几何与物理状态（位置、速度、质量等），由 Position-Based Dynamics（PBD，一种以位置为直接求解变量的实时物理仿真方法，通过迭代满足约束来模拟刚体碰撞、变形等行为）物理引擎驱动；$M$ 个\textbf{高斯基元}负责从任意视角渲染图像，其位置随父粒子刚性移动。一个典型桌面场景约有 1000 个粒子和 10000 个高斯基元。

系统的运行分为预测和校正两个阶段，形成``物理仿真 $\to$ 渲染比较 $\to$ 视觉力反馈''的闭环。

在\textbf{预测阶段}，PBD 物理引擎以 30\,Hz 对粒子执行前向仿真：首先根据机器人正运动学设定机器人粒子的位置，然后通过 Jacobi 迭代求解器同时处理三类约束——地面约束（防止粒子穿透地面平面）、碰撞约束（防止粒子间相互穿透）和形状匹配约束（保持刚体或可变形体的结构），每步物理仿真约耗时 5\,ms。随后，挂载的高斯基元随父粒子刚性移动至预测位置。

在\textbf{校正阶段}，系统从已知相机视角渲染当前状态，将渲染图像与真实相机观测进行逐像素比较，通过光度损失 $\mathcal{L}_{\text{rgb}} = \sum_{\boldsymbol{u}} |C_{\text{rgb}}(\mathbf{u}) - C_{\text{gt}}(\mathbf{u})|$ 反向传播更新高斯基元的位置。关键的是，高斯基元\emph{不会}被直接移动到优化后的位置——系统记录每个高斯基元优化前后的位置差，通过如下\textbf{视觉力}公式将位移转化为施加在父粒子上的外力。直觉上，这一过程类似于弹簧系统：渲染误差``拉扯''高斯基元偏离预测位置，而视觉力就是弹簧的反作用力——它将拉扯量按高斯不透明度加权后放大，传递给父粒子，使其``感知''到视觉观测的纠偏信号：

\begin{equation}
    \boldsymbol{f}_i = K_p \sum_{j} o_j \left( \boldsymbol{g}_j - \boldsymbol{g}_j^0 \right)
\end{equation}

其中 $\boldsymbol{g}_j$ 和 $\boldsymbol{g}_j^0$ 分别是挂载到粒子 $i$ 上的高斯基元 $j$ 在优化后和优化前的位置，$o_j$ 是该高斯的不透明度（作为置信度权重——越不透明的高斯``说话越算数''），$K_p = 60$ 是比例增益系数。换言之，每个粒子受到的视觉力等于其所有子高斯``建议位移''的不透明度加权平均再乘以增益 $K_p$。校正阶段执行 5 步梯度优化（以满足 30\,Hz 实时约束），随后高斯基元被重置回原始位置，视觉力交由下一帧的 PBD 物理引擎统一求解。

这种设计确保系统始终处于物理可行状态——所有位置更新都由物理引擎协调，视觉观测仅作为外力信号融入仿真循环，不会绕过碰撞、地面等物理约束。整个系统仅需 3 个 RGB 相机即可以 30\,Hz 实时运行，支持刚体推动、物体拾取、多物体碰撞和可变形体（如绳索）操控等多种动态交互场景。

物理仿真闭环为机器人的自主交互提供了基础，但许多场景还需要人类操作者的直接参与——例如在危险环境中远程操控机器人，或在虚拟环境中训练操控策略。以下方法将 3DGS 的实时渲染能力与 VR 技术结合，构建沉浸式人机交互体验。

\paragraph{VR 交互}

GS-Verse（2025.10）~\cite{2510_11878} 直接将物体网格与高斯泼溅表示集成，通过网格-高斯映射 $\Phi$ 建立高斯基元与网格顶点之间的关联，在 VR 环境中实现更精确的表面近似和高度真实的变形交互。18 人用户研究证实该方法在物理感知拉伸操控上统计显著优于 VR-GS 基线，同时在系统可用性（SUS）和认知负荷（NASA-TLX）方面表现优异。VR-Doh（2024.12）~\cite{2412_00814} 定制物质点方法（MPM）模拟手诱导大变形，通过中轴变换（MAT）近似手部几何实现高效碰撞处理，结合解耦外观与物理表示和均匀高斯体积正则化技术增强 3DGS 无缝渲染，支持新手用户直观的接触式和手势式 3D 建模体验。

\paragraph{绿色通信}

GSMR（2025.08）~\cite{2508_08624} 修改了 3DGS 在遥操作中的通信环节——用「3DGS 记忆系统 $+$ 内容切换编码」替代传统的全量图像上传——3DGS 的泼溅渲染管线保持不变，新增的是基于位姿键值的服务端按需渲染和联合优化内容切换与功率分配的跨层框架。

\textbf{传统做法与局限}：在 GSMR 之前，机器人混合现实（RoboMR）系统的通信瓶颈在于：机器人需要将每帧高分辨率图像（例如 $67.2$\,KBytes/帧，即 $537.6$\,Kbits）全部上传至远程仿真服务器，服务器将真实图像与虚拟世界融合后生成混合现实（MR）图像反馈给机器人用于控制。这种全量上传方案的数据速率需求在数百 KBytes/帧以上，对带宽和功率构成巨大压力。已有的资源分配方法（如注水功率分配、最大-最小公平）虽然可以优化通信层面的功率分配，但忽视了两个关键问题：一是未能利用先前 RoboMR 执行中积累的已训练 3DGS 模型所蕴含的场景记忆信息，二是不考虑应用层图像编解码与物理层通信设计之间的跨层耦合关系。

\textbf{设计动机}：GSMR 的核心洞察是：在 RoboMR 场景中，仿真服务器已持有环境的 3DGS 模型，因此服务器可以直接根据机器人位姿渲染出照片级虚拟视角——该位姿仅需 6 个浮点数（$S = 192$\,bits），相比完整图像的 $I = 537.6$\,Kbits 降低了约三个数量级。然而，由于光照变化和环境动态变化，3DGS 渲染的虚拟视角与真实环境之间存在``记忆偏差''（memory bias），不能完全替代真实图像。因此关键问题转化为：如何在每帧自适应地决定``上传真实图像''还是``仅发送位姿让服务器用 3DGS 渲染''，并联合优化上行传输功率分配，在给定功率预算下最小化 MR 图像失真。

\textbf{技术实现}：GSMR 的通信管线在传统 RoboMR 的四阶段循环（感知$\to$通信$\to$融合$\to$控制）中重构了通信和融合两个阶段。在通信阶段，系统将编码器函数设计为内容切换机制：$c_t = E_t(r_t) = x_t r_t + (1 - x_t) s_t$，其中 $x_t \in \{0, 1\}$ 是待优化的切换变量——$x_t = 1$ 时上传真实图像（数据量 $I = 537.6$\,Kbits），$x_t = 0$ 时仅发送位姿键（数据量 $S = 192$\,bits），后者不足前者的 $0.04\%$。在融合阶段，服务器端解码后通过掩码融合生成 MR 图像：$\hat{m}_t = d_t \hat{r}_t + (1 - d_t) v_t$，其中 $d_t$ 是提取真实背景的掩码，$v_t$ 是提取虚拟对象的掩码。当 $x_t = 0$ 时，服务器使用 3DGS 模型从位姿 $s_t$ 渲染虚拟视角 $\Psi(s_t)$ 替代真实图像进行融合。

基于上述管线，GSMR 提出跨层优化框架 GSCLO，将应用层的内容切换决策与物理层的功率分配联合建模为混合整数非线性规划（MINLP）：在平均功率预算 $\frac{1}{T}\sum_{t} p_t \leq P$ 和每帧传输速率约束 $R_t(p_t) \geq C_t(E_t)$ 下，最小化平均 MR 图像失真（由 L1 损失与 DSSIM 加权组合度量）。失真函数的关键结构为 $L_t(1 - x_t)$，即仅当选择位姿传输（$x_t = 0$）时才产生以 GS 渲染误差为上界的失真——这意味着优化器应将有限的功率预算集中分配给 GS 渲染质量差的帧（即``记忆偏差''大的帧），而非均匀上传所有图像。

由于 MINLP 问题可由分支定界法求解但复杂度随帧数 $T$ 指数增长，GSMR 提出加速惩罚优化（APO）算法：首先将二值约束 $x_t \in \{0,1\}$ 松弛为连续约束 $0 \leq x_t \leq 1$，然后通过惩罚项 $\frac{1}{\rho}\sum_t x_t(1 - x_t)$ 促进二元解，最终利用差分凸（DC）规划将非凸惩罚项分解为凸函数之差，通过序列凸优化迭代求解。APO 算法的初始解由排序策略提供：按 GS 失真 $L_t$ 降序排列，优先为失真最大的帧分配 $x_t = 1$（上传图像），在等信道增益条件下该排序解具有最优性。此外，GSMR 还扩展到信道不确定场景（通过 BILS 算法保证中断概率约束）、功率最小化场景（推导出闭合形式的 GS 功率节省因子）和多机器人场景（通过迫零接收消除共信道干扰）。

\textbf{效果}：在轮式机器人（LIMO，288 帧评估）和四足机器人（Unitree Go2，370 帧评估，$1.73$\,MBytes/帧）上的 ROS 实验表明：GSMR 将通信开销降至每帧不足 200\,bits（相比传统 RoboMR 的 $537.6$\,Kbits 降低约三个数量级），在功率预算 10\,mW 下实现功耗降低 75\%（即节省 6\,dB），同时 PSNR 保持 30\,dB 以上（相比无跨层优化的基线提升至少 5\,dB）。APO 算法相比 Mosek 分支定界法将计算时间降低 100 倍以上，同时性能接近最优。在动态场景中（如办公室内播放视频的环境变化），APO 通过检测``记忆偏差''自适应触发图像上传，验证了``数据混合''策略对增强 3DGS 在动态场景中鲁棒性的有效性。

GSRMR（2025.04）~\cite{2504_13697} 修改了 RoboMR 系统的上行通信环节——用「GS 记忆模型 $+$ 按需内容切换」替代「逐帧上传完整图像」——机器人的传感与控制环路保持不变，改变的是与仿真服务器之间的数据传输策略：利用先前执行中采集的历史图像训练 GS 模型，使服务器能根据机器人位姿直接渲染照片级视角，机器人只需上传位姿 key（仅数十字节）而非完整图像（数十千比特），从源头将通信数据量降低数个数量级。由于 GS 渲染与真实环境之间存在光照变化和物体移动导致的偏差（discrepancy），GSRMR 进一步提出 GS 跨层优化（GSCLO）框架，联合优化二值内容切换变量 $x_t \in \{0,1\}$（决定每帧上传图像还是仅上传位姿）和上行传输功率 $p_t$，在满足 MR 图像质量约束下最小化总通信能耗 $\tau \sum_{t=1}^{T} p_t$。这是一个混合整数非线性规划（MINLP）问题，传统分支定界法的复杂度随帧数 $T$ 指数增长。GSRMR 设计了加速惩罚优化（APO）算法，通过惩罚函数 $\varphi(\mathbf{x}) = \frac{1}{\beta}\sum_t x_t(1-x_t)$ 将二值约束松弛为连续约束，再利用差凸（DC）分解将非凸惩罚项线性化，每步迭代求解一个复杂度为 $\mathcal{O}((2T)^{3.5})$ 的凸问题，配合基于 GS 损失排序的贪心初始化加速收敛。在 $4\,\text{m} \times 4\,\text{m}$ 的 RoboMR 平台上（LIMO 机器人 $+$ CARLA 仿真器），GSRMR 将通信能耗降低超过 10 倍（在松驰质量阈值下可达 100 倍），APO 算法的 PSNR 至少比次优方案高 6\,dB，收敛仅需约 10 次迭代。

\subsection{多模态感知}
\label{ss:17_3b}

本节讨论两条技术路线：（1）多模态感知与记忆；（2）其他具身应用。

\subsubsection{多模态感知与记忆}

本路线介绍将 3DGS 与基础模型结合构建多模态空间记忆的方法。将 3DGS 与基础模型结合构建多模态空间记忆，是具身智能的前沿方向。核心挑战在于：基础模型（如 CLIP、LLM）的特征维度通常高达数百甚至数千维，直接存储在每个高斯基元上会导致巨大的内存开销和渲染延迟；而简单的特征压缩（如线性降维到 16--64 维）又会丢失信息，使蒸馏后的特征难以被基础模型解码利用。以下方法从不同角度解决了这一难题。

\paragraph{3D 多模态记忆}

M3（2025.03）~\cite{2503_16413} 修改了 3DGS 的特征存储与渲染环节——用「记忆库 $+$ 注意力检索」替代「逐高斯特征蒸馏」——3DGS 原有的几何、颜色和 $\alpha$-混合渲染管线保持不变。

\textbf{Baseline 做法与局限。}先前的特征泼溅方法（如 Feature 3DGS、Feature Splatting）将基础模型的 2D 特征图通过可微分渲染蒸馏到每个高斯基元上。这一范式面临两个根本问题：其一，基础模型（如 CLIP、DINOv2、LLaMA3）的特征维度 $d$ 通常为 $512 \sim 4096$ 维，若为每个高斯存储完整特征，内存开销呈爆炸式增长；因此这些方法被迫将特征线性压缩至 $16 \sim 64$ 维，造成严重的信息瓶颈——从 1024 维压缩到 16 维仅保留不到 $2\%$ 的信息量。其二，原始 2D 特征图本身未必具有 3D 一致性，在蒸馏过程中强制施加 3D 约束会导致蒸馏特征与原始特征之间的错位（misalignment），使得下游基础模型无法正确解码这些特征。

\textbf{核心思路：分而治之。}M3 的设计动机来源于一个关键观察：不同视角的基础模型特征存在大量冗余——相邻像素的特征高度相似，跨视角的重叠区域也共享相似特征。与其在每个高斯上存储压缩后的低维近似，不如把原始高维特征集中存储在一个紧凑的记忆库中，然后让高斯通过注意力机制按需检索。这就像图书馆（记忆库）与读者（高斯基元）的关系：不需要每人随身携带一本书（逐高斯蒸馏），只需每个人持有索书号（principle query），通过检索即可获取完整内容。

\textbf{主场景成分（Principle Scene Components, PSC）。}PSC 是对场景全部原始特征执行相似性约简后得到的紧凑特征集合，也称为记忆库（memory bank）。具体而言，M3 先用多个基础模型（CLIP、SigLIP、DINOv2、LLaMA3/v、SEEM）提取 $n$ 个视角的 2D 特征图 $\mathbf{E}_i \in \mathbb{R}^{[h \times w, d]}$，将所有视角的特征展平为 $\mathbf{R} \in \mathbb{R}^{[n \times h \times w, d]}$。随后通过贪心式相似性约简算法：逐块计算特征间的余弦相似度，保留与已选成分不重复（相似度低于阈值 $\theta$）的特征，标记其余相似特征为冗余。最终得到 $\mathbf{PSC} \in \mathbb{R}^{[t, d]}$，其中 $t \ll n \times h \times w$。直觉上，这一过程类似主成分分析（PCA）——从海量冗余特征中提取出构成场景的「关键成分」，但 PSC 保留的是原始维度 $d$ 的完整特征而非低维近似。

\textbf{原理查询与渲染。}每个高斯基元新增一个可优化属性——原理查询（principle query）$\mathbf{q} \in \mathbb{R}^{l}$，其中 $l$ 的维度可灵活设置以适配不同基础模型。渲染时，原理查询与颜色、不透明度等属性一样，通过 3DGS 标准的 $\alpha$-混合管线进行光栅化：

\begin{equation}
\mathbf{Q}_p = \sum_{i \in \mathcal{N}} q_i \alpha_i T_i, \quad \text{where}\ T_i = \prod_{j=1}^{i-1}(1 - \alpha_j)
\end{equation}

即每个像素的查询值由沿射线排序的高斯加权混合得到，渲染机制与 3DGS 原始的颜色渲染完全一致。

\textbf{高斯记忆注意力（Gaussian Memory Attention）。}渲染得到的逐像素原理查询 $\mathbf{Q}_p^{\mathbf{V}_*} \in \mathbb{R}^{[H, W, l]}$ 并非直接可用作特征，而是作为「索书号」去检索记忆库。具体通过高斯记忆注意力机制实现：

\begin{equation}
\hat{\mathbf{R}} = \text{Softmax}(\mathbf{Q}_p^{\mathbf{V}_*} \times \mathbf{W}_m \times \mathbf{PSC}^T) \times \mathbf{PSC}
\end{equation}

其中 $\mathbf{W}_m \in \mathbb{R}^{[l, d]}$ 是可学习的记忆投影矩阵，将低维查询投影到与 PSC 相同的特征空间以计算注意力权重；Softmax 产生的权重决定了每个像素应从记忆库中检索哪些关键成分；最终输出 $\hat{\mathbf{R}}$ 是 PSC 的加权组合，维度与原始基础模型特征完全一致（$\mathbb{R}^{[H, W, d]}$）。因为 $\hat{\mathbf{R}}$ 保持原始维度且直接从记忆库中的真实特征组合而来，无需经历降维-升维的信息瓶颈，所以可以被基础模型直接解码：判别式模型（CLIP、SigLIP、SEEM）的特征直接用于检索和定位任务，生成式模型（LLaMA3/v）的特征可直接用于场景描述和视觉问答。

\textbf{效率与实验。}M3 以 35M 参数（其中每个基础模型分配 $l{=}32$ 维的原理查询，共 $32 \times 6{=}192$ 个查询维度）同时支持 6 个基础模型，而 Feature 3DGS 使用 61M 参数仅蒸馏单一低维近似。在特征距离指标上，M3 在 CLIP/SigLIP 等判别式模型上与 F-3DGS 持平或更优，在下游检索任务（I2T@1）上从 F-3DGS 的 55.2\% 大幅提升至接近真实标签水平；在四足机器人平台上，M3 成功实现了基于文本查询的目标定位与抓取演示。训练采用基于点采样的损失函数（每张图像采样 2000 个点），避免了逐像素损失的高内存开销，全 6 个基础模型的端到端训练仅需约 45 分钟。

\paragraph{统一辐射场}

URF-GS（2026.01）~\cite{2601_19216} 提出统一射频-光学辐射场框架。与 M3 在视觉语言模态上构建多模态记忆不同，URF-GS 改变了 3DGS 的属性表示和渲染环节——为每个高斯基元扩展材质属性（散射反照率、金属度、粗糙度）和法向量，使同一组高斯同时支持光学渲染和射频信号传播仿真。设计动机在于：机器人在复杂环境中导航和规划时不仅依赖视觉信息，还需要预测无线信号的空间分布（如 Wi-Fi 接入点部署优化），而传统方法将视觉和无线建模割裂处理，无法共享场景几何信息。URF-GS 通过两阶段管线解决这一问题：第一阶段利用单目深度和法向量先验正则化 3DGS 优化以获取精确几何；第二阶段基于物理的逆向渲染将射频信号传播分解为表面交互（散射+反射，遵循 Cook-Torrance BRDF 模型）和自由空间传输（自由空间路径损耗模型），从多模态观测联合优化材质属性和辐射模式。空间谱准确率提升高达 24.7\%，样本效率提升 10 倍，并展示了在 Wi-Fi AP 部署和机器人路径规划任务中的应用。

\paragraph{语义不确定性}

连续语义泼溅（2024.11）~\cite{2411_02547} 在 3DGS 上实现概率语义更新和光栅化。与 M3 通过注意力机制蒸馏基础模型特征不同，该方法改变了 3DGS 的语义编码和渲染环节——利用 3DGS 的 $\alpha$ 混合权重作为贝叶斯核推理（BKI）的核函数，直接在 3D 椭球体上执行概率语义更新，而非通过额外训练学习特征渲染。设计动机在于：传统语义建图方法（如体素化 BKI）将语义信息离散化至固定分辨率体素，存在精度和内存的矛盾，而 3DGS 的椭球体天然提供连续、自适应分辨率的场景覆盖，可以作为 BKI 核函数的直接载体。具体而言，该方法在每个椭球体上维护 Dirichlet 共轭先验分布，通过 $\alpha_i^n = \alpha_i^{n,0} + \sum_j w(x_j, x_n) y_j^c$ 累积语义观测浓度参数，其中 $w(x_j, x_n)$ 即 3DGS 的 $\alpha$ 混合权重，在满足 BKI 核函数约束（$0 \leq w \leq 1$、$\sum w = 1$）的同时无需额外训练。概率光栅化通过 $E[\pi_i] = \sum_n w(x_i, x_n) E[\pi_n]$ 和方差传播 $V[\pi_i] = \sum_n w(x_i, x_n)^2 V[\pi_n]$ 返回可量化不确定性的逐像素分割预测，将连续语义建图从体素扩展至椭球体。

\paragraph{长期服务机器人}

GS-LTS（2025.03）~\cite{2503_17733} 是首个面向长期服务机器人的 3DGS 系统。与 M3 专注于静态场景的多模态记忆构建不同，GS-LTS 改变了 3DGS 的场景维护环节——将 3DGS 从一次性重建工具转变为可增量更新的长期场景表示，支持动态变化的室内环境。设计动机在于：服务机器人在长期部署中面临物体搬移、新增和移除等场景变化，传统 3DGS 重建无法适应这些动态变化，需要从零重新训练。GS-LTS 提出三阶段管线：首先通过双分支变化检测（像素级绝对差异截断 + EfficientSAM 特征级余弦相似度）判断场景是否发生变化并分类为新增、移除或搬移；然后通过规则策略（沿目标区域切线方向左右移动收集多视角观测）自主收集变化区域的多视角数据；最后通过高斯编辑（新增时实例化新高斯并继承近邻属性、移除时剪枝并修补孔洞、搬移时执行删除-重建策略）高效更新场景表示，仅需少量微调即可恢复重建质量。实验表明 GS-LTS 在重建和导航方面具有优势，场景更新比图像训练基线更快且质量更高。

\paragraph{多智能体 3D 推理}

数据图方法（2024.07）~\cite{2407_10743} 利用图结构允许 LMM 迭代查询大型环境的较小部分。与 M3 将全局特征嵌入每个高斯基元不同，该方法改变了 3DGS 场景与大型多模态模型（LMM）的交互方式——将大型 3DGS 场景分解为节点（局部 3DGS 子场景）和边（空间邻接关系）组成的图结构，使 LMM 能够按需查询相关子场景而非一次性处理整个场景。设计动机在于：LMM 的上下文窗口有限，无法一次性处理包含数百个视角的大规模 3DGS 场景，而机器人搜索救援等任务需要在大型环境中进行 3D 语言推理（如``找到红色背包旁边的伤员''）。该方法通过图遍历算法优先查询最相关位置，结合 LMM 的空间推理能力逐步缩小搜索范围，在搜索救援任务中实现 3D 场景语言任务的规模化处理。

\subsubsection{其他具身应用}

本路线介绍 3DGS 在事件相机、人形机器人、农业与建筑机器人等跨领域具身应用中的方法。3DGS 的应用远不限于导航、操控和遥操作，还延伸至事件相机、人形机器人、农业与建筑机器人、声纳感知、XR 自适应等 14 个跨领域方向。

\paragraph{事件相机}

Event3DGS（2024.06）~\cite{2406_02972} 利用事件相机的卓越时间分辨率，在高速自我运动下重建高保真 3D 结构和外观。重建质量提升 3 dB，同时计算成本降低 95\%。高速动态 3D 成像（Sensor Fusion Splatting）（2025.02）~\cite{2502_04630} 融合 RGB、深度和事件相机三种互补模态，利用可变形 3D 高斯联合优化高斯参数和时间变形场。该方法通过共享场景表示跨模态渲染并最小化多传感器观测差异，在合成和真实高速场景（如蜂鸟振翅、牛顿摆）上均显著优于 Deformable 3DGS 和 4DGS 等基线，即使在小基线和低光照条件下仍保持鲁棒重建。

\paragraph{人形机器人导航}

MIF（2026.05）~\cite{2605_21935} 提出面向人形机器人的多模态交互场，集成置信度感知语义 3DGS、差异触发的空间记忆更新和任务驱动几何重建。在 Unitree-G1 人形机器人上，非静态环境重定位成功率从 12\% 提升至 94\%，语义记忆占用减少 91.4\%。

\paragraph{遥操作语义}

Zero-Splat TeleAssist（2025.12）~\cite{2512_08271} 集成 Monte Carlo dropout CLIPSeg 视觉语言分割、MiDaS 单目深度估计和加权 PCA 位姿提取，将商品级 CCTV 视频流转换为多边遥操作的共享 6-DoF 3DGS 世界模型，无需基准标记或深度传感器。系统以 10 Hz 进行位姿估计、1 Hz 重规划，全部计算可在 Nvidia Jetson Orin 上实时运行，加权方案消融显示移除置信度权重后方向误差从 8.7$^\circ$ 升至 11.5$^\circ$。

\paragraph{计算机视觉数据集生成}

基于 3DGS 的数据集生成管线（2025.12）~\cite{2512_13411} 创建大规模、高度真实、自动标注的数据集。两遍渲染技术结合泼溅真实感和代理网格生成的阴影贴图，像素完美分割掩码自动生成并直接用于 YOLO 等检测模型。混合训练策略（少量真实图像+大量合成数据）产生最佳检测和分割性能。

\paragraph{建筑机器人}

共享自主建筑机器人系统（2025.11）~\cite{2511_09695} 使用高斯泼溅进行在线 3D 重建并引入运动参数建模运动物体。与本章其他方法聚焦于桌面操控或导航场景不同，该方法改变了 3DGS 的建图场景——面向建筑工地等非结构化环境，需要在人机协作中处理动态障碍物（如移动的工人或设备）。设计动机在于：建筑工地的环境持续变化，传统静态地图无法实时反映动态障碍物，而共享自主系统需要同时保证人类安全和机器人任务执行效率。该方法通过运动参数将 3DGS 中的运动物体建模为时变表示，配合神经配置空间障碍方法计算安全规划走廊，在动态建筑环境中实现人机安全协作。

\paragraph{农业机器人}

GreenhouseSplat（2025.10）~\cite{2510_01848} 构建首个温室级辐射场仿真数据集。与本章其他方法聚焦于工业或室内场景不同，GreenhouseSplat 改变了 3DGS 的重建场景——面向农业温室中的高密度植物排列，解决叶片严重遮挡和重复纹理导致的重建困难。设计动机在于：农业机器人在温室中执行采摘、监测等任务需要精确的植物 3D 模型作为导航和操作的参考，但温室环境的高湿度、强光照变化和高密度植物排列对传统 3D 重建方法构成挑战。该数据集包含 82 株黄瓜植物的多行配置，支持基准标记定位等机器人任务，为农业机器人的感知算法研发提供了标准化的评测基准。Botany-Bot（2025.10）~\cite{2510_17783} 利用 3D 分割高斯泼溅模型构建植物活体详细标注数字孪生。与 GreenhouseSplat 提供静态数据集不同，Botany-Bot 改变了 3DGS 与机器人的交互模式——通过旋转转台多视角捕获和 SAM2 掩码预处理解决植物重建中背景干扰和光照不一致的问题，再利用 GARField 多层次 3D 分割将植物分解为叶片、茎干等独立组件。设计动机在于：植物叶片的背面和茎芽等关键区域（病虫害检测的重要部位）在静态 3D 重建中因遮挡而无法观测。Botany-Bot 利用 3DGS 的显式表示特性计算每个叶片的主轴方向和中心位置，自动规划机器人叶片提升/下压轨迹，操纵叶片拍摄被遮挡细节（茎芽、叶片背面）的高分辨率图像，叶片分割准确率 90.8\%。

\paragraph{声纳感知}

SonarSplat（2025.04）~\cite{2504_00159} 提出面向成像声纳的高斯泼溅框架，每个高斯编码声学反射率和饱和度属性，开发高效光栅化方法生成符合声学成像模型的距离/方位图像，PSNR 提升 3.2 dB。

\paragraph{LiDAR-相机可微标定}

在 SonarSplat 扩展高斯属性类型的基础上，鲁棒 LiDAR-相机标定方法（2025.04）~\cite{2504_00525} 利用 2DGS 从相机图像序列重建几何信息解决多传感器标定问题。与 SonarSplat 扩展高斯的属性类型不同，该方法改变了 2DGS 在标定管线中的角色——利用 2DGS 的可微渲染特性将外参标定转化为可微优化问题。设计动机在于：LiDAR 和相机之间的精确外参标定是多传感器融合的基础，但传统方法依赖人工选取对应点或特定标定板，在无纹理场景中容易失败。该方法先用 LiDAR 点云重建无色 2DGS 获取几何骨架，再通过最小化光度损失同时更新高斯颜色和相机外参参数，实现无标定板、端到端的自动标定。

\paragraph{基准标记泼溅嵌入}

基准标记泼溅（2025.08）~\cite{2508_17012} 提出将基准标记（如 AprilTag）高效嵌入 3DGS 场景的混合框架。设计动机在于：3DGS 渲染的高质量图像中缺少可检测的基准标记，导致无法在渲染视角中直接估计相机位姿，而重新在场景中放置物理标记又会破坏重建。该框架在渲染管线中虚拟嵌入标记图案，在密集温室等复杂环境中实现位姿估计精度和效率的双重提升。

\paragraph{XR 自适应}

SAGE（2025.03）~\cite{2503_16747} 通过语义分割驱动的自适应 LOD 减少内存和计算开销。与本章其他方法聚焦于 3DGS 的场景内容不同，SAGE 改变了 3DGS 的优化环节——将语义分割信息融入高斯基元数量的自适应控制过程，按语义类别独立优化高斯基元的数量而非全局统一处理。设计动机在于：标准 3DGS 对所有场景区域使用统一的密度控制策略，导致重要语义区域（如机器人需要交互的物体）和次要区域（如远处墙壁、地面）使用相同的高斯密度，造成内存浪费和渲染开销过大。SAGE 首先通过语义分割为每个 3D 点分配类别标签（多数投票跨视角融合），然后为每个语义类别独立拟合质量-距离参数模型 $\text{SSIM}_{l,i}(d_{\min,l})$，该模型描述在不同相机距离下达到目标 SSIM 所需的最小优化迭代次数，从而预测每个语义类别在何处停止高斯基元的分裂和加密。实验验证了该模型在跨场景设置中的可迁移性——在某个场景中学习到的语义质量参数可以直接应用于新场景的同类语义区域，在保持目标视觉质量的同时显著减少高斯基元数量。

\paragraph{建筑数字孪生}

数字孪生建筑框架（2025.02）~\cite{2502_05769} 连接云端地图平台、多智能体 LLM 数据分析和基于高斯泼溅的网格提取管线。与其他方法聚焦于机器人操作场景不同，该方法改变了 3DGS 的数据来源和重建目标——不依赖自采集的多视角照片，而是利用 Google Map API 的街景卫星图像作为输入，自动重建建筑的 3DGS 模型和语义描述。设计动机在于：城市环境中的建筑机器人或服务机器人需要快速获取目标建筑的 3D 结构信息，但传统方法依赖人工拍摄和多视角重建，无法规模化。该框架通过多智能体 LLM 协调数据检索、视觉分析和模型生成流程，从地址、邮编或坐标检索建筑 3D 模型和视觉描述，为建筑机器人的任务规划提供场景先验。

\paragraph{轻量级视觉直觉}

L-VITeX（2024.10）~\cite{2410_07872} 在资源受限机器人和集群上实现轻量级视觉直觉系统。与本章其他方法在 GPU 服务器上运行 3DGS 不同，L-VITeX 改变了 3DGS 的前端感知环节——在极端资源受限的边缘设备（如 ESP-32 Cam，单核微控制器）上运行轻量级目标检测模型（FOMO，基于截断 MobileNet-V2），作为 3DGS 重建的注意力引导机制。设计动机在于：小型移动机器人（如探索管道的微型机器人）的计算能力极其有限，无法在板上运行完整的目标检测或语义分割模型，但又需要在 3D 建图过程中关注感兴趣区域（如裂缝、异物），而非无差别采集所有视角。L-VITeX 的 FOMO 模型峰值 RAM 占用小于 50 KB、推理时间小于 200 ms，通过 8 位整数量化进一步压缩模型权重，在检测到感兴趣区域时触发``强调功能''（如靠近目标获取更详细观测），从而引导 3DGS 重建聚焦于环境中的关键结构。在 TinyTurtle 微型机器人上验证了集成 FOMO 的 3DGS 重建比无引导基线具有更高的场景覆盖完整性。

\paragraph{4D 数据收集}

ExoGS（2026.01）~\cite{2601_18629} 提出无机器人 4D Real-to-Sim-to-Real 框架。与本章其他 Real-to-Sim 方法依赖真实机器人采集数据不同，ExoGS 改变了示范数据的采集环节——使用自设计的被动外骨骼 AirExo-3（与目标机器人运动学同构、末端定位误差 $<1$ mm）替代真实机器人捕获操作示范，无需机器人硬件即可获得运动学一致的轨迹。设计动机在于：高质量机器人示范数据的采集受限于机器人硬件的可用性和成本，而外骨骼设备成本低廉且可由人类直接佩戴操作。ExoGS 将机器人、物体和环境重建为可编辑 3DGS 资产后，通过四种数据增强策略（相机视角扰动、高斯颜色/光照抖动、背景替换、物体位姿/尺度变化）指数级扩展训练数据多样性，并引入 Mask Adapter 轻量模块（基于 3DGS 的像素级语义掩码引导 ViT 编码器的注意力聚焦于交互相关区域），缩小仿真到真实的视觉域间隙。

\paragraph{开放词汇定位}

WoMAP（2025.06）~\cite{2506_01600} 利用 3DGS Real-to-Sim-to-Real 管线可扩展生成数据。与其他方法将 3DGS 用于策略训练数据的视觉渲染不同，WoMAP 改变了 3DGS 的应用环节——将 3DGS 场景作为世界模型的训练数据生成器，从开放词汇物体检测器（GroundingDINO）蒸馏密集奖励信号至潜在世界模型。设计动机在于：基于世界模型的主动感知方法需要大规模多样化的训练数据来学习环境动态和奖励预测，但传统世界模型依赖图像重建目标函数导致训练不稳定，同时真实世界数据采集难以覆盖足够的环境多样性。WoMAP 的核心创新是奖励蒸馏策略——在 3DGS 渲染的每帧图像上计算检测置信度并按边界框大小缩放作为密集奖励信号，替代传统的图像重建目标函数，实现数据高效的奖励预测器训练。在规划阶段，WoMAP 利用 VLM 的常识推理（链式思维提示）生成高层动作提议，再通过模型预测控制（MPC）在世界模型的潜在空间中优化具体动作序列，成功率比 VLM 和扩散策略基线分别高 9 倍和 2 倍。

\paragraph{主动视觉搜索基准}

E3VS-Bench（2026.04）~\cite{2604_17969} 构建具身 3D 视觉搜索基准。与其他方法将 3DGS 用于机器人策略训练不同，E3VS-Bench 改变了 3DGS 的应用目标——利用 3DGS 构建标准化评估环境，系统性测试具身智能体的主动视觉搜索能力。设计动机在于：现有的 3D 问答基准大多基于预给定视角的被动式问答，无法评估智能体主动探索未知场景、搜索信息并回答问题的能力，而这一能力是搜索救援、目标检测等实际机器人任务的核心需求。该基准包含 99 个高保真 3DGS 场景和 2014 个问题驱动回合，要求智能体在 5-DoF（3D 平移 + 水平/垂直旋转）中控制视角，主动收集视角依赖证据进行问答，评估指标同时考虑回答准确率和视角搜索效率。

\paragraph{自我中心动态场景}

EgoGaussian（2024.06）~\cite{2406_19811} 首次从仅 RGB 自我中心输入同时重建 3D 场景并动态跟踪 3D 物体运动。与本章其他动态 3DGS 方法（如 Physically Embodied GS）依赖外部物理引擎不同，EgoGaussian 改变了 3DGS 的场景分解和优化策略——利用手-物体交互的时序结构将视频分割为静态片段（无物体运动的交互准备/后续阶段）和动态片段（物体被移动的阶段），分别重建背景和动态物体。设计动机在于：自我中心视角下的手-物体交互视频存在大量遮挡和场景不一致性（物体被拿起后原位置被遮挡），直接对全序列进行 3DGS 优化会产生大量浮游伪影。EgoGaussian 利用高斯泼溅的离散特性，首先在静态片段中训练完整场景的 3DGS 并通过相邻动态片段的物体掩码识别动态物体的高斯基元，然后将物体高斯作为整体进行刚体位姿估计（6D 旋转表示 + 平移），在动态片段中逐帧优化物体位姿并同步精化物体形状，最终将背景模型和动态物体模型合成为完整的 4D 重建结果。

\begin{takeaway}
具身感知与交互将 3DGS 从静态重建推向动态理解：自我中心动态场景重建利用手-物体交互时序结构分解静态与动态，3DGS 构建的标准化评估环境则为主动视觉搜索等能力提供了客观评测平台。
\end{takeaway}

% ======================================================================
\section{仿真与真实迁移}
\label{sec:17_4}

仿真到真实迁移（sim-to-real transfer）是机器人学习最核心的挑战之一。策略在仿真中训练后部署到真实世界，需要克服视觉外观差异（渲染质量、纹理、光照）和物理属性差异（动力学、摩擦、质量）。3DGS 的照片级真实感渲染能力使其成为弥合仿真-真实视觉鸿沟的理想工具——与传统仿真器使用低多边形网格渲染不同，3DGS 可以从真实世界照片直接重建场景，生成的合成图像在视觉上与真实照片几乎无法区分，从而大幅缩小视觉域间隙。本节从 Real2Sim2Real 框架、数据增强与合成、以及评估基准三个方面展开。

\subsection{Real2Sim2Real框架}
\label{ss:17_4a}

本节讨论两条技术路线：（1）Real2Sim2Real 框架；（2）数据增强与合成。

\subsubsection{Real2Sim2Real 框架}

本路线介绍基于 3DGS 的 Real2Sim2Real 框架，将真实场景重建为高保真仿真环境以训练和迁移机器人策略。Real2Sim2Real 框架将真实世界场景重建为 3DGS 模型，在仿真中训练策略后迁移回真实世界。3DGS 的关键优势在于重建的场景在视觉上与真实世界几乎无差异，极大降低了视觉域迁移的难度。

\paragraph{高保真数字孪生}

基于 3DGS 的高保真数字孪生（2026.01）~\cite{2601_03200} 修改了 3DGS 的初始化与后处理环节——用 InstantSplat 的预训练几何先验（MASt3R）直接估计替代传统 SfM 逐帧匹配初始化，并在渲染输出之上增加了可见性感知语义融合和三阶段几何过滤后处理管线——3DGS 的核心可微光栅化渲染机制保持不变。系统提出从稀疏 RGB 输入到物理仿真就绪的完整数字孪生生成管线。

\textbf{基线与动机}：传统 Real-to-Sim 方法依赖 NeRF 重建（耗时常达数十分钟）和人工标注分割，难以实现快速闭环机器人操控的数字孪生。设计动机在于：机器人操控需要快速构建环境的数字孪生以支持"扫描-规划-执行"闭环，而 NeRF 的长优化时间和手工标注流程使得单次扫描到规划的时间远超实际任务需求（桌面级任务中物体通常在 5--30 秒内保持静止）。该框架的关键创新在于将 3DGS 的高效重建（分钟级）与自动化语义-几何处理管线结合。

\textbf{管线架构}：系统包含三个核心阶段。

第一阶段是\textbf{高保真场景重建}，其采用 InstantSplat 方法利用预训练几何先验（MASt3R）直接估计初始点云和相机位姿，绕过耗时的 SfM 步骤，通过自监督光度误差联合优化高斯参数与相机位姿。直觉上，传统 3DGS 需要先由 COLMAP 完成特征匹配和稀疏重建才能开始高斯优化，而 InstantSplat 用 MASt3R 的立体对应关系直接"猜"出点云和相机位置，省掉了最耗时的前处理。每个 3D 高斯由位置 $\mu$、协方差矩阵 $\Sigma$、颜色和不透明度定义，其未归一化密度为 $G(\mathbf{x}) = \exp\!\bigl(-\frac{1}{2}(\mathbf{x}-\mu)^\top \Sigma^{-1}(\mathbf{x}-\mu)\bigr)$，通过最小化渲染图像与输入视图之间的光度误差联合优化高斯参数与相机位姿。这一优化绕过了原始 3DGS 的自适应密度控制步骤，使收敛速度大幅提升。

第二阶段是\textbf{可见性感知语义融合}，针对 2D 分割模型（如 SAM）在物体边界处产生的"渗透效应"（bleeding effect，即背景像素被误包含进前景掩码），系统采用基于深度聚类（DBSCAN）的空间隔离策略——对每个视角的 2D 掩码投影至 3D 后按深度值聚类，保留最大聚类作为真实物体几何，其余为伪影。多视角语义融合通过置信度加权投票实现，其中权重 $W_{i,j}$ 由点是否属于主深度聚类决定（属于则 $w_\text{cluster} \approx 1$，否则 $\approx 0$）。直觉上，如果一组高斯在深度方向紧密聚集（间距 $< 2\,$cm），它们大概率属于同一物体表面；而深度方向上远离主簇的离散点则更可能是分割噪声。最终语义标签由超过共识阈值 $\tau_\text{consensus} = N/\alpha$（$\alpha=1.5$）的投票确定，迭代精化（3 轮）逐步收紧语义边界。该共识阈值的设定使投票要求随视角覆盖度增长而提高，同时容忍偶发的分割失败。

第三阶段是\textbf{物理就绪几何重建}，通过三步过滤将噪声视觉表示转化为清洁碰撞几何：固有属性过滤（移除低透明度 $\alpha < 0.1$ 的半透明"薄雾"基元和针状拉伸基元）、语义引导连通性剪枝（保留最大连通组件消除浮游伪影）、Alpha Shapes 网格化（紧贴点分布保留尖锐特征，如盒子棱角和手柄边缘）。生成的碰撞模型通过 ROS2 桥接与 Unity-ROS2-MoveIt 物理引擎无缝集成，实现从感知到规划到验证的完整 Sim2Real 工作流。系统整体平均重建时间约 229 秒（$<$4 分钟），视觉保真度达 37.03\,dB PSNR 和 0.9821 SSIM，相比 NeRF 基线实现约 5 倍加速；消融实验中几何清理管线对人工标注真值的 F1 分数达 0.9989；真实 Franka 机器人上的长程重排任务成功率为 90\%。

\paragraph{物理接地仿真}

GaussTwin（2026.03）~\cite{2603_05108} 将基于位置的动力学（PBD）与离散 Cosserat 杆公式结合实现物理接地仿真。与高保真数字孪生方法的静态重建-网格提取管线不同，GaussTwin 改变了 3DGS 的动态更新环节——通过将高斯锚定至物理基元（球体粒子）并强制一致 SE(3) 更新，实现 25\,Hz 实时的预测-校正闭环。设计动机在于：刚性物体和可变形体（如绳索）的操控需要同时具备物理一致性（碰撞、接触力）和视觉真实性，而传统方法中物理仿真和视觉渲染使用独立表示，无法实现两者的闭环耦合。

\textbf{对象-高斯绑定。}场景初始化阶段，GaussTwin 从多视角 RGB-D 图像中用 SAM2 提取实例掩码，在物体包围盒内均匀填充半径 5\,mm 的球体粒子，通过深度投影剔除掩码外、地面以下或深度不一致的粒子，剩余粒子近似物体几何。由粒子密度积分计算物体的质量、质心、位姿和惯性矩阵。3D 高斯从分割后的物体点云初始化，经 1000 步 Adam 优化后锚定至最近的球体粒子，建立高斯外观与物理基元的固定绑定关系。对于绳索等可变形线性物体（DLO），先训练其 3D 高斯表示，再从俯视渲染图中提取 2D 骨架并沿中心线排序生成球体粒子序列。

\textbf{预测-校正闭环。}在线追踪以 25\,Hz 运行，分为两个阶段。预测阶段：将机器人连杆位姿设为正运动学计算的当前关节配置，PBD 物理引擎（基于 NVIDIA Warp，单步 0.1\,ms）执行前向仿真，驱动球体粒子执行刚体碰撞、剪切-拉伸和弯曲-扭转约束求解，挂载的高斯随父粒子刚性移动。校正阶段：渲染当前高斯的三个视角图像，用 EfficientTAM 提取物体分割掩码，对物体高斯的均值和协方差施加全局变换 $T \in \text{SE}(3)$，通过 6 步 Adam 优化最小化渲染图像与真实图像的均方误差（约 10\,ms）；优化完成后将校正力 $f_i = K_p(p_j T(p_j) - p_j) / N_i$（其中 $K_p$ 为可调增益，$N_i$ 为每个粒子绑定的平均高斯数）反馈至 PBD 仿真，确保所有状态更新均通过物理引擎协调。总延迟约 40\,ms（分割 24\,ms + 位姿优化 10\,ms + 仿真与 I/O 6\,ms）。这一设计的关键差异在于：校正力不直接修改高斯位置（如 PEGS 方法对每个高斯独立优化），而是将视觉误差转化为施加于物理粒子的力和力矩，再由 PBD 求解器统一处理——物理一致性始终由仿真引擎保证，视觉观测仅作为外力源参与校正。

\textbf{实验验证。}在 Franka Research 3 平台上，GaussTwin 在刚体追踪精度上一致优于基于形状匹配的 PEGS 和基于刚体动力学的 RBD 基线，且方差更低。在绳索推动任务中，Cosserat 杆公式通过连续介质力学显式表示弯曲和扭转，比形状匹配方法产生更物理一致的变形和接触反应。系统还支持基于推动的下游规划任务。

\paragraph{视角泛化增强}

GenSplat（2026.03）~\cite{2603_29192} 提出前馈 3DGS 框架，通过新视角渲染增强视觉运动策略的视角泛化能力。与高保真数字孪生方法将 3DGS 用于场景重建不同，GenSplat 改变了 3DGS 的用途——不重建特定场景，而是为策略训练批量生成视角增强数据。设计动机在于：主流 2D 视觉运动策略对训练时未见过的相机视角严重退化，而大规模机器人数据集通常缺乏可靠的相机位姿标注，无法直接使用传统 3DGS 进行新视角合成。

\textbf{置换等变架构。}GenSplat 采用前馈置换等变 Transformer 架构，从稀疏未标定图像直接预测高斯参数和相机位姿，无需逐场景优化。每张图像先经 DINOv2 提取 patch token，通过交替的帧内注意力和全局自注意力层聚合多视角上下文，确保输出与输入视角顺序无关。token 随后分别路由至相机位姿解码器（预测局部到全局的 SE(3) 变换）和高斯参数解码器（预测位置、颜色、不透明度和缩放）。这一架构的关键差异在于：与传统 3DGS 依赖 COLMAP/SfM 逐场景重建不同，GenSplat 在单次前向传播中完成从任意数量未标定视图到完整 3DGS 表示的映射，推理速度比逐场景优化快数个量级。

\textbf{3D 先验蒸馏策略。}纯光度监督在稀疏视图、无位姿条件下容易导致几何坍塌和渲染伪影。GenSplat 通过从预训练几何基础模型 PI3D 提取密集点图监督和仿射不变位姿约束，正则化 3DGS 优化。直觉上，PI3D 为场景提供了一个粗略但合理的几何"骨架"，防止优化过程陷入退化的局部最优（如所有高斯退化为平面或零厚度壳体）。

\textbf{视角增强管线与实验。}对于每条专家示范，GenSplat 重建 3DGS 场景后，在相机绕世界轴旋转 $[-30^\circ, 30^\circ]$ 并沿 $+Z$ 轴平移 $[0.5, 6]$\,cm 的范围内采样新视角渲染合成轨迹，严格保持时间对齐和动作标签不变。在 Franka Research 3 平台上的 6 个操控任务中，GenSplat 增强后的策略在大扰动下将 Diffusion Policy 的成功率从 27.78\% 提升到 43.33\%（相对提升 56.0\%），且在无扰动设定下不降低性能。该策略对策略架构（Diffusion Policy 和 ACT）均有效，实现了架构无关的视角泛化增强。

\paragraph{端到端 Real-to-Sim}

Splatting Physical Scenes（2025.06）~\cite{2506_04120} 提出混合场景表示 SplatMesh——将 3DGS 的照片级真实感渲染与适合物理仿真的显式物体网格融合在单一表示中。与高保真数字孪生方法的静态重建-网格提取管线不同，Splatting Physical Scenes 改变了 3DGS 与网格几何的耦合方式——通过将高斯基元的均值约束在网格表面（加权重心坐标初始化 $\mu = w_1 v_1 + w_2 v_2 + w_3 v_3$），建立高斯外观与底层网格几何的可微映射，使梯度可以从 RGB 损失通过高斯反向传播至网格顶点位置。设计动机在于：3DGS 擅长照片级渲染但缺乏碰撞检测所需的精确几何，显式网格适合物理仿真但渲染质量差，两者需要在一个可微框架中联合优化。端到端优化管线利用可微渲染和可微物理（MuJoCo）联合精化所有场景组件，采用 Euclidean Distance Transform 平滑二值分割掩码解决梯度稀疏问题，并引入 Laplacian 正则化促进网格表面平滑。该系统在 ALOHA 2 双臂操控器上验证，同时实现高保真物体网格重建、照片级新视角和免标注机器人位姿标定。

\paragraph{软体仿真}

SoMA（2026.02）~\cite{2602_02402} 提出面向软体操控的 3DGS 仿真器，在统一潜在神经空间耦合可变形动力学、环境力和机器人关节动作，实现端到端 Real-to-Sim 仿真。与 GaussTwin 通过 PBD 物理引擎驱动高斯基元不同，SoMA 改变了软体动力学建模的方式——采用层次图结构建模机器人-物体-环境交互，将环境效果建模为显式全局力（重力和支持力），机器人交互通过交互图中的隐式力计算，两者直接作用于高斯基元层级。设计动机在于：现有的神经高斯动力学模型主要基于状态回归，将物体运动与接触效果纠缠在一起，在多样化交互模式下难以泛化。SoMA 通过力驱动的动力学解耦（$\dot{v}_i, \dot{\omega}_i = f_\theta(g_i^{t-1}, g_i^{t-2}, f_i)$，其中 $f_i = f_i^\text{env} + f_i^\text{rob}$），使模型学习力到状态的映射而非直接的状态转移，从而在未见过的交互模式中保持泛化能力。多分辨率训练策略沿时间（粗到精细的时间步长）和图像（超分辨率重建几何+原始分辨率训练动力学）两个维度优化长时域稳定性，并结合遮挡感知图像监督（仅在物体可见区域施加光度损失）和动量一致性正则化（约束相邻层次节点间的动量守恒 $\mathcal{L}_\text{mom} = \sum_l \|m_{c_l}\dot{v}_{c_l} - \sum_{i \in C_{l-1}} m_i \dot{v}_i\|_2^2$）处理部分可观测性问题。在真实世界机器人操控上重仿真准确率和泛化能力提升 20\%，支持复杂长时序布料折叠等任务的稳定仿真。软体数字孪生评估框架（2025.11）~\cite{2511_04665} 从真实世界视频构建软体数字孪生，验证仿真的回放与真实执行性能强相关。

\paragraph{自动化资产生成}

Scalable Real2Sim（2025.03）~\cite{2503_00370} 提出全自动 Real2Sim 管线，仅使用机器人关节扭矩传感器和外部相机生成仿真就绪资产。与高保真数字孪生方法依赖外部分割模型和语义处理管线不同，Scalable Real2Sim 改变了资产获取的数据来源和物理参数辨识方式——利用机器人自身执行预设计的激励轨迹（有限傅里叶级数参数化），通过半定规划（SDP）从关节扭矩和运动学状态同时辨识机器人连杆的惯性参数（质量、质心、转动惯量）和关节摩擦系数，然后通过差分法从空载和持物两种状态下的参数差获取物体的完整惯性属性。设计动机在于：传统仿真资产生成需要人工标注几何和物理属性，无法规模化；而机器人自身的关节扭矩传感器可以作为物理参数辨识的``天然标定工具''，无需额外传感器。在几何重建方面，alpha 透明训练方法将背景像素替换为迭代相关的随机颜色，迫使模型在背景区域分配零密度，解决了以物体为中心重建中的密度渗透问题；DINO 特征空间的帧选择策略确保下采样后的训练图像保持最大多样性。最终生成的资产包括纹理网格、凸分解碰撞几何和物理属性，可直接集成至现有拾取-放置系统。

\paragraph{跨域迁移}

Semantic 2D Gaussian Splatting（S2GS）（2025.12）~\cite{2512_04731} 提取以物体为中心的域不变空间特征，通过特征级高斯泼溅将多视角 2D 语义场投射至统一 3D 空间。S2GS 改变了策略输入的语义表示方式——采用层次语义提取机制，通过 SAM 生成部件级分割掩码，利用 CLIP 分别提取全局语义特征（掩码内平均池化）和局部语义特征（裁剪后独立提取），再通过浅层 MLP 解码器将渲染的低维紧凑特征映射回高维语义空间，以余弦相似度损失监督蒸馏过程。设计动机在于：策略学习需要域不变的物体状态表示，但原始 RGB 图像包含大量无关背景信息且随视角变化剧烈，而 3DGS 中以物体为中心的空间特征天然具备视角一致性和背景无关性。S2GS 还提出动态场景更新机制——假设物体刚性运动，通过优化 SE(3) 变换参数最小化渲染图像与观测图像的 L1 损失，结合已知机器人操作作为强运动先验，实现实时场景更新。语义过滤机制移除无关背景内容，确保策略学习的干净一致输入，在 ManiSkill 仿真和真实世界部署中显著提升仿真到真实可迁移性。

\paragraph{视觉-物理双重对齐迁移}

在 S2GS 通过特征蒸馏实现域不变表示的基础上，TwinAligner（2025.12）~\cite{2512_19390} 通过 SDF 重建和可编辑 3DGS 渲染实现像素级视觉对齐。与 S2GS 不同，TwinAligner 改变了 Real-to-Sim 的对齐环节——通过三步闭环对齐（Mesh-GS 数字孪生构建、视觉-动态 Real-to-Sim 对齐、Sim2Real 策略学习）同时缩小视觉外观和物理动力学两个维度的仿真-真实差距。设计动机在于：视觉渲染质量和物理动力学参数（摩擦、质量、质心、机器人控制器参数）的不匹配均会导致策略迁移失败，但现有方法通常只关注其中之一。TwinAligner 通过``Control-Hit-Slide''交互协议同时暴露多种物理属性——机器人控制阶段对齐控制器参数、撞击阶段确保脉冲一致、滑动阶段通过刚体运动方程优化摩擦系数和质心位置——实现物理参数的联合辨识，并通过 3DGS 渲染和梯度无关优化器（粒子群算法）实现像素级视觉对齐，最终实现策略零样本泛化。

\paragraph{MLLM 自动资产生成}

RoboSimGS（2025.10）~\cite{2510_10637} 将多视角真实世界图像转换为可扩展、高保真、物理交互的仿真环境。与 SplatSim 依赖手动分割和 ICP 对齐不同，RoboSimGS 改变了资产生成的自动化程度——采用解耦重建策略（3DGS 用于静态背景的高保真渲染、显式网格用于可交互物体的物理仿真），并通过 MLLM（GPT-4o）分析视觉数据推断物体的运动学结构（关节类型、运动轴、运动极限）和物理属性（密度、杨氏模量、泊松比），自动生成可交互的关节物体 URDF 资产。设计动机在于：传统仿真环境构建需要大量人工建模和参数调优，而 MLLM 的常识推理能力可以从多视角图像推断物体的物理和运动学属性，实现从真实照片到物理仿真环境的全自动转换。系统还通过全场景增强（物体位姿/尺度随机化、相机视角扰动、高斯颜色/光照抖动、轨迹中间点随机化）进一步缩小域间隙，实现零样本仿真到真实迁移。

\paragraph{物理参数在线适应迁移}

Phys2Real（2025.10）~\cite{2510_11689} 结合 VLM 推理的物理参数先验与交互式不确定性感知融合，在平面推动任务中实现 100\% 成功率（底部加权 T 块），超越域随机化基线的 79\%。与 RoboSimGS 通过 MLLM 推断物理属性用于仿真构建不同，Phys2Real 改变了物理参数在策略部署中的使用方式——将快速运动适应（RMA）的两阶段训练范式从隐式潜在空间扩展到显式可解释的物理参数空间，使在线估计的物理参数可以直接与 VLM 的视觉推理结果通过逆方差加权融合。设计动机在于：接触丰富的操控任务中物体物理属性可能在接触前无法观测，而纯交互历史在非预抓取操控中可能提供无信息量的反馈。Phys2Real 首先用 SuGaR 从 3DGS 重建中提取水密网格作为仿真资产，然后在仿真中训练物理参数条件化的 RL 策略，最后训练由 10 个模型组成的集成预测物理参数及其认知不确定性和偶然不确定性。融合机制在交互历史不确定时更多依赖 VLM 先验，在 VLM 置信度低时更多依赖在线适应，实现鲁棒的物理参数感知策略部署。

\paragraph{3DGS 渲染仿真}

SplatSim（2024.09）~\cite{2409_10161} 修改了 3DGS 的场景表示与渲染环节——用按刚体分割的高斯泼溅替代仿真器中的传统网格图元——标准 3DGS 的可微光栅化渲染管线保持不变。传统机器人仿真器（如 PyBullet）使用低多边形网格作为场景表示，渲染出的图像缺乏纹理细节和光照保真度，导致在此类合成图像上训练的策略部署到真实世界时出现显著的视觉差距。SplatSim 的核心洞察是：如果场景中每个刚体（机器人连杆、物体、夹持器）的 3DGS 能够被独立分割出来，并与仿真器坐标系对齐，那么就能利用仿真器提供的正向运动学直接驱动这些高斯子集渲染出任意关节角度下的照片级图像。

具体实现分三个阶段。首先是\textbf{坐标系对齐}：从真实场景的 3DGS $S_\text{real}$ 中手动分割出属于机器人的 3D 高斯，取其均值形成点云，再用迭代最近点（Iterative Closest Point, ICP，一种通过最小化两组点云距离来估计刚体变换的算法）将此点云与仿真器中的真值点云对齐，得到从泼溅坐标系到机器人坐标系的变换矩阵 $T_\text{splat}^\text{robot}$。其次是\textbf{连杆分割与变换}：利用机器人 CAD 模型的轴对齐包围盒将 $S_\text{real}$ 中的 3D 高斯分配到各连杆 $l$，得到子集 $S_\text{real}^l$；仿真器根据关节角 $q_t$ 计算正向运动学变换 $T_\text{fk}^l$，然后通过相似变换将高斯位置和协方差映射到新姿态：

\begin{equation}
    T^l = (T_\text{splat}^\text{robot})^{-1}\, T_\text{fk}^l\, T_\text{splat}^\text{robot}
\end{equation}

直觉上，这个公式先用 $T_\text{splat}^\text{robot}$ 把高斯从泼溅坐标系变换到仿真器的机器人坐标系，再用 $T_\text{fk}^l$ 施加正向运动学驱动，最后逆变换回泼溅坐标系以保持与其他场景元素的一致性。对于\textbf{自由物体}，采用类似的 ICP 对齐获取变换 $T_{\text{obj}}$，然后根据仿真器给出的物体位姿 $(p_t^k, R_t^k)$ 计算变换并渲染。对于\textbf{关节物体}（如平行爪夹持器），由于其连杆不与坐标轴对齐，标准包围盒无法精确分割，因此改用在仿真器标注点云上训练的 $K$ 近邻分类器推断每个高斯所属的连杆类别。

\paragraph{四模块 Real-to-Sim2Real 工具链}

最后是\textbf{策略训练}：给定专家在仿真器中生成的轨迹 $\tau_E = \{(s_1, a_1), \ldots, (s_T, a_T)\}$，SplatSim 对每个状态 $s_t = (q_t, x_t^1, \ldots, x_t^n)$ 渲染照片级图像 $I^\text{sim}$，得到演示数据集 $\tau_G = \{(I_1^\text{sim}, a_1), \ldots, (I_T^\text{sim}, a_T)\}$。策略采用 Diffusion Policy（2023）~\cite{chi2023diffusion} 从合成图像预测动作。尽管高斯泼溅大幅缩小了视觉仿真-真实差距，残余差异（如缺乏阴影、柔性线缆的刚性假设导致渲染失真）仍然存在，因此训练时叠加图像增强（高斯噪声、随机擦除、亮度对比度调整）以提升鲁棒性。操控策略在 SplatSim 中训练后以零样本方式部署至真实世界，6 项操控任务平均成功率 86.25\%，接近真实数据训练策略的 97.5\%。RoboGSim（2024.11）~\cite{2411_11839} 修改了 3DGS 仿真的管线完整度——将 SplatSim 的单阶段渲染扩展为四模块 Real-to-Sim2Real 工具链（高斯重建器、数字孪生构建器、场景合成器、交互引擎），实现仿真数据模型在真实机器人上的零样本性能。

与 SplatSim 专注于单任务场景、依赖手动分割和 ICP 对齐不同，RoboGSim 改变了资产生成和数据合成的自动化程度——数字孪生构建器通过 SAM 自动分割场景高斯中的物体资产，场景合成器支持在新视角、新物体配置和新场景布局下批量生成多样化训练数据，交互引擎则提供在线、可复现且安全的策略评估环境。设计动机在于：SplatSim 虽然通过高斯泼溅缩小了视觉差距，但每次只能对单一场景的单条轨迹渲染，无法规模化生成跨任务的训练数据，而遥操作采集真实示范的成本极高且难以扩展。RoboGSim 通过轨迹重定向将单次人类示范适配至不同物体位姿和相机视角，使策略在新视角和新场景上的表现甚至超过真实数据训练的策略。Real2Sim 和 Sim2Real 双向迁移实验均显示出纹理和物理上的高一致性，验证了该 Real2Sim2Real 闭环的有效性。

\paragraph{城市级导航仿真评测}

Wanderland（2025.11）~\cite{2511_20620} 修改了 3DGS 仿真的场景尺度与几何精度——从桌面操控扩展到城市级室内外开放世界环境，并通过几何接地重建解决大规模场景中渲染质量与导航策略可靠评估的矛盾。与 SplatSim 和 RoboGSim 聚焦于桌面操控不同，Wanderland 改变了 3DGS 仿真的应用范式——不再用于策略训练数据生成，而是构建可复现的闭环评估测试平台。设计动机在于：开放世界具身 AI（如视觉导航）的可复现闭环评估一直是瓶颈——虽然近期的视频-3DGS 方法降低了场景捕获成本，但渲染与几何的 sim-to-real 差距仍然过大，无法作为可信的评测基准。Wanderland 系统性地揭示了三个关键问题：纯图像管线在大规模场景中扩展性差、几何质量直接影响新视角合成质量、两者共同对导航策略学习和评估可靠性产生负面影响。为此，该框架采用多传感器捕获（LiDAR、RGB-D、全景相机）和可靠重建构建精确的 3DGS 几何表示，配合鲁棒视图合成和物理交互引擎，不仅支持导航、搜索和探索等长时域具身 AI 任务的标准化评测，其丰富的原始传感器数据还可作为 3D 重建和新视角合成模型的独立基准。

\paragraph{端到端视觉导航迁移}

GaussFly（2026.04）~\cite{2604_05062} 通过对比表示学习从渲染 RGB 图像提取紧凑、抗噪的潜在特征，实现视觉运动策略的鲁棒零样本迁移。与 SplatSim 直接将 3DGS 渲染图像用于策略训练不同，GaussFly 改变了 3DGS 渲染结果与策略学习之间的衔接方式——在渲染图像和策略网络之间插入对比学习预训练的视觉编码器，通过 InfoNCE 损失最大化同一渲染图像在不同数据增强下的潜在表示一致性，使编码器学习到对视觉扰动（光照变化、纹理噪声、相机偏移）不变的紧凑特征。设计动机在于：直接在 3DGS 渲染图像上训练的策略容易过拟合于渲染特有的视觉特征（如高斯泼溅的模糊边界），在迁移至真实世界时性能下降。GaussFly 采用解耦重建策略（背景和前景资产独立重建后自由组合）和面约束 3DGS 变体（约束高斯基元最小尺度轴对齐表面法向量，通过射线-平面相交计算精确深度），在 Isaac Sim 中以 30\,Hz 实时渲染训练数据，PPO 策略在冻结编码器提取的特征上训练后直接部署至真实无人机。

SplatGym（2024.10）~\cite{2410_19564} 提出开源神经仿真器，从单个视频创建照片级虚拟环境。与 SplatSim 依赖机器人 URDF 模型和 ICP 对齐不同，SplatGym 改变了仿真环境的构建方式——通过 nerfstudio 工具链从单视频自动训练 3DGS 模型并提取点云，再将点云转换为 Octree 碰撞表示（使用 Flexible Collision Library），无需任何机器人 CAD 模型即可实现照片级新视角渲染和碰撞检测。设计动机在于：传统仿真器（如 PyBullet、MuJoCo）依赖人工建模的低多边形网格，渲染质量差且建模成本高，而完整的 3DGS 仿真环境需要同时解决视觉渲染和物理碰撞两个问题。SplatGym 将 3DGS 的渲染能力（比 NeRF 推理快 291 倍、训练快 45\%）与 Octree 碰撞检测（$O(\log n)$ 时间复杂度）结合，支持自我相机视图生成、碰撞检测和虚拟物体修复，使 RL 智能体在照片级环境中学习导航策略。

\subsubsection{数据增强与合成}

本路线介绍利用 3DGS 的可渲染性进行大规模数据增强与合成数据生成的方法。3DGS 的可渲染性使其成为大规模数据增强和合成数据生成的强大工具。关键优势在于：通过简单移动、旋转或删除特定高斯基元即可模拟不同的物体配置和视角，而无需重新拍摄，这使得从少量真实示范指数级扩展训练数据成为可能。

\paragraph{高保真示范生成}

基于 DMP 的示范合成（2026.05）~\cite{2605_01232} 将专家轨迹用动态运动基元（DMP）建模并重定向至新目标、物体配置和视角，产生阶段一致、形状保持的运动。分析式障碍感知 DMP 公式直接在 3DGS 诱导的连续密度场上操作实现碰撞避免。RoboSplat（2025.04）~\cite{2504_13175} 直接操控 3D 高斯生成多样化、视觉真实的示范，覆盖六种泛化类型的五种增强技术，实现单样本设置下 87.8\% 的平均成功率。与 SplatSim 通过仿真器正向运动学驱动高斯渲染不同，RoboSplat 改变了帧对齐和示范生成的精度——利用 3DGS 可微渲染替代传统 ICP 进行精细帧对齐，通过最小化渲染分割掩码与 URDF 掩码的像素差异优化变换矩阵，解决高斯基元尺度属性导致的点云坐标与外观不匹配问题。设计动机在于：传统 ICP 仅基于点坐标对齐，忽略了 3DGS 特有的尺度属性导致的系统性偏移。RoboSplat 覆盖六种泛化维度（物体位姿、物体类型、相机视角、机器人本体、场景外观、光照条件）的系统化增强策略中，物体类型增强通过 GPT-4 生成可抓取物体名称并使用 3D 内容生成模型创建高斯资产，场景外观增强利用 COCO 数据集和重建 3DGS 场景替换背景，光照增强通过随机缩放、偏移和高斯噪声调制每个高斯的漫反射颜色。

\paragraph{大规模数据生成}

Real2Render2Real（R2R2R）（2025.05）~\cite{2505_09601} 仅需智能手机捕获的物体扫描和单个人类示范视频，渲染数千个高视觉保真、机器人无关的示范数据。与 RoboSplat 在 3DGS 空间中直接操控高斯基元不同，R2R2R 改变了数据生成的技术路线——先用 3DGS 重建物体资产再通过 GARField 分割转换为纹理网格，利用 4D 可微部件建模（4D-DPM）从单目人类示范视频中提取部件 6-DoF 运动轨迹，再通过差分逆运动学求解器将物体轨迹转换为机器人关节空间轨迹。设计动机在于：传统数据生成需要机器人遥操作采集示范，成本高昂且难以规模化，而人类通过智能手机完成示范的自然操作几乎零成本。R2R2R 引入轨迹插值策略（空间归一化+球面线性插值 Slerp）将单个示范轨迹适配至任意新的起始-终止位姿配置，并采用抓取姿态采样从 3DGS 均值点构建粗网格进行解析式反向抓取采样，支持双手协调任务。

基于 IsaacLab 的 GPU 并行渲染在单张 RTX 4090 上以 51 个示范/分钟的速率生成数据（对比人类遥操作的 1.7 个/分钟，加速 30 倍），实验表明单个人类示范生成的数据可匹配 150 次人类遥操作示范的效果。RoboPearls（2025.06）~\cite{2506_22756} 基于增量语义蒸馏和 3D 正则化 NNFM 损失支持多种场景操控操作，结合 LLM 自动化仿真制作和 VLM 分析关闭仿真回路。与 R2R2R 将 3DGS 转换为网格渲染不同，RoboPearls 改变了 3DGS 场景编辑的语义粒度——通过增量语义蒸馏（ISD）机制在运行时按需细化分割粒度，当用户请求的物体粒度（如炉灶上的小按钮）超出 SAM 初始掩码的分辨率时，系统先用 G-DINO 验证当前掩码是否匹配目标，不匹配时用边界框作为 SAM 的精细提示生成新掩码，仅微调相关高斯的身份编码而非重新训练整个场景。设计动机在于：SAM 的多层级掩码无法覆盖所有物体粒度，而重新训练整个场景的语义嵌入成本过高。

RoboPearls 还引入 3D 正则化 NNFM 损失限制风格迁移仅作用于目标物体高斯并保持背景不变，多 LLM 智能体协作自动化整个仿真制作流程。

\paragraph{合成数据集}

基于 3DGS 的合成数据生成（2025.06）~\cite{2506_05092} 在 Unreal Engine 中利用照片级高斯泼溅快速生成标注合成数据。与 R2R2R 利用 3DGS 进行轨迹级数据增强不同，该方法改变了 3DGS 在数据生成管线中的角色——将 3DGS 场景作为 Unreal Engine 中的可渲染资产，在游戏引擎内系统性生成大量带自动标注（边界框、分割掩码）的合成训练数据。设计动机在于：机器人感知模型的训练需要大量标注数据，但人工标注成本高昂且耗时，而在纯虚拟环境中生成的合成数据又往往因渲染不够真实而无法有效迁移至真实世界。3DGS 从真实照片重建的场景提供了近乎照片级的渲染质量，同时保留了游戏引擎对相机位姿、物体位置和光照条件的精确控制能力，使生成的合成数据在视觉真实性和标注精度之间取得平衡。验证实验表明，合成图像训练的检测器在机器人足球比赛场景中与人工标注真实图像训练的检测器性能相当。

\paragraph{视角增强}

Splat2Real（2026.03）~\cite{2603_10638} 提出 CN-Coverage 视角选择课程，核心设计动机在于：朴素增加渲染视角数量并不单调提升下游深度估计质量，因为随机采样容易在低价值或外推视角上浪费预算。该方法在每步贪婪选择中联合评估覆盖增益 $|V(T) \setminus U|$（尚未被已有视角覆盖的体素数量）与新颖度惩罚 $\exp(-d(T, T_\text{train})/\tau)$（与训练视角的距离），其中距离 $d$ 综合相机中心平移和偏航角差异，默认权重 $\alpha=0.20$、$\tau=0.35$。此外引入高斯观测层（GOL）安全机制，根据验证集 3DGS 重渲染质量（LPIPS 阈值 0.20–0.80）动态混合 GS 渲染与网格回退，避免低质量场景中盲目使用 GS 观测导致性能回退。在 20 个 TUM RGB-D 序列上的步匹配缩放实验表明，CN-Coverage 在中等至高预算（$N \geq 200$）下提供最稳定的 AbsRel 误差，$N=2000$ 时最低尾部误差为 0.348，显著优于纯覆盖策略的 0.419 和随机采样的 0.462。

\paragraph{物理推理框架}

SMS（2025.05）~\cite{2505_14938} 集成 3DGS 场景重建、视觉基础模型语义分割、VLM 材质属性推理和物理仿真，实现可泛化的物理推理和以物体为中心的规划。在台球操控和四旋翼着陆场景中验证了跨域迁移和真实世界实验的鲁棒性能。

\paragraph{多模态策略学习}

GS-Playground（2026.04）~\cite{2604_25459} 开发高性能并行物理引擎与批量 3DGS 渲染管线集成，在 $640\times480$ 分辨率下实现 $10^4$ FPS 突破性吞吐量。自动化 Real2Sim 工作流重建照片级真实感、物理一致且内存高效的环境。GaussGym（2025.10）~\cite{2510_15352} 将 3DGS 集成至 IsaacGym 等矢量化物理仿真器。与 GS-Playground 的端到端并行管线不同，GaussGym 改变了 3DGS 的数据输入和渲染优化方式——利用 VGGT（视觉接地几何 Transformer）从多种来源（已标定数据集、手机扫描、甚至视频生成模型的原始 RGB 序列）估计相机参数和密集点云，将点云直接初始化高斯基元以加速收敛并提高几何保真度。设计动机在于：大规模视觉强化学习需要同时在数千个不同场景中训练，但传统仿真器要么渲染质量差（低多边形网格），要么渲染速度慢（光线追踪），无法同时满足高保真和高吞吐量。GaussGym 的核心创新是将 3DGS 作为仿真器的``即插即用渲染器''（drop-in renderer），通过多线程 PyTorch 内核实现跨环境的批量泼溅渲染，并将渲染频率与控制频率解耦（控制 50\,Hz、相机 10\,Hz）以节省计算资源，同时引入运动模糊模拟（沿相机速度方向偏移渲染多帧后 $\alpha$-混合）提升视觉真实感。单张 RTX 4090 可同时渲染 4096 个环境、128 个独特场景，在消费级 GPU 上实现超过每秒 10 万步，同时保持高视觉保真度。

\paragraph{VLA 平台}

RealMirror（2025.09）~\cite{2509_14687} 构建面向人形机器人的综合开源 VLA 平台，集成生成模型和 3DGS 重建真实环境和机器人模型，实现零样本仿真到真实迁移——仅在仿真数据上训练的模型无需微调即可在真实机器人上执行任务。GWM（2025.08）~\cite{2508_17600} 提出高斯世界模型用于机器人操控。与其他方法将 3DGS 仅用于渲染训练图像不同，GWM 改变了 3DGS 的应用层级——将 3D 高斯基元作为世界状态的基本表示，通过 3D 高斯变分自编码器（VAE）将可变数量的高斯编码为固定长度的潜在嵌入，再由潜在扩散 Transformer 学习高斯嵌入空间的动态转换 $p(z_{t+1} | z_t, a_t)$。设计动机在于：传统世界模型在 2D 图像空间预测未来状态，难以保持 3D 一致性和精细的场景级细节（如物体位姿变化、局部形变），而以 3D 高斯作为中间表示可以在保持 3D 结构一致性的同时实现高效的动态预测和场景级未来状态重建。GWM 可无缝集成至离线模仿学习（作为特征编码器）和在线模型强化学习（作为动力学模型），支持多种机器人操控任务。

\paragraph{Robo-GS 物理一致表示}

Robo-GS（2024.08）~\cite{2408_14873} 提出混合表示模型集成网格几何、3DGS 和物理属性。与 SplatSim 将高斯按刚体分割并通过 ICP 对齐不同，Robo-GS 改变了 3DGS 与物理仿真的耦合方式——通过高斯-网格-像素绑定技术建立网格顶点和高斯模型间的同构映射 $\Phi$，使梯度可以从真实世界视频流经高斯反向传播至网格（后向优化），也可以从物理仿真结果经网格正向传播至高斯渲染（前向渲染）。设计动机在于：机器人 Sim2Real 需要同时具备三种能力——高保真视觉渲染（3DGS 擅长）、精确碰撞检测（网格擅长）和物理一致的动态仿真（运动学/动力学方程驱动），但这三种表示通常各自独立维护，无法实现端到端的梯度流通。Robo-GS 采用基于位置的动力学和半隐式欧拉积分，将高斯基元的运动分解为运动学变换（Modified Denavit-Hartenberg 参数驱动机器人连杆）和动力学变换（牛顿-欧拉方程驱动物体运动），统一在可微渲染管线中实现全链路优化。

\paragraph{NeRF 与 GS 互转}

NeRF-to-GS 转换方法（2024.05）~\cite{2405_09717} 在有限视角场景中实现 NeRF 和 3DGS 的双向转换。与本章其他方法从零训练 3DGS 不同，该方法改变了 3DGS 的获取方式——利用预训练 NeRF 模型已学习的场景知识直接初始化和优化 3DGS 参数，无需原始训练图像。设计动机在于：NeRF 在有限视角下具有更好的泛化能力（能合理补全未见区域），但推理速度慢；3DGS 渲染速度快但需要密集视角训练且容易在稀疏视角下产生伪影。该方法兼顾两者的优势——先用 NeRF 的体渲染采样点云初始化高斯基元的位置、颜色和不透明度，再通过 NeRF 渲染的大量视角作为伪训练数据优化高斯参数，转换计算成本远低于从头训练。在仿真到真实迁移场景中，该方法允许先用少量视角训练 NeRF 获取鲁棒的场景理解，再转换为 3DGS 用于策略训练的实时渲染。

\paragraph{世界模型}

DreMa（2024.12）~\cite{2412_14957} 将机器人世界模型重新构想为可学习数字孪生。与 GWM 在高斯潜在空间中学习动态不同，DreMa 改变了 3DGS 与物理仿真的集成方式——将高斯泼溅渲染器和物理仿真器直接嵌入世界模型的前向预测管线，使机器人能够``想象''物体的新配置（如不同摆放位置的桌面）并预测机器人动作的未来视觉后果。设计动机在于：传统世界模型在图像空间或抽象潜在空间中预测未来，缺乏对物体几何和物理规律的显式建模，导致长时域预测不准确且无法支持物理一致的轨迹规划。DreMa 集成高斯泼溅和物理仿真器，从单次示例学习策略即可实现真实机器人任务。

\paragraph{大规模视频世界模型}

在 DreMa 的基础上进一步扩展规模，GigaWorld-0（2025.11）~\cite{2511_19861} 集成大规模视频生成和 3DGS 重建、可微物理系统辨识和可执行运动规划，VLA 模型在无真实世界交互训练下实现强真实世界性能。

\paragraph{4DGS 数据引擎世界模型}

EnerVerse（2025.01）~\cite{2501_01895} 采用多视角视频表示结合 4DGS 数据引擎形成自强化数据循环——通过 4DGS 从机器人交互数据中重建动态场景，生成的新视角数据反过来用于训练更强大的视频生成模型，形成数据质量和模型能力的持续提升闭环。

\begin{takeaway}
仿真与真实迁移是 3DGS 在机器人学习中最具变革性的应用：照片级真实感渲染从根本上缩小视觉域间隙，4DGS 数据引擎形成自强化数据循环，大规模视频世界模型使策略在无真实交互训练下实现强真实世界性能。
\end{takeaway}

% ======================================================================

\subsection{近期方法脉络：具身仿真与场景构建、具身仿真与数据增强、机器人仿真与数字孪生}

这一组近期工作延续 3DGS、Gaussian Splatting，主要把改动集中在应用、适配等环节。下面只保留相对既有路线的变化，便于和前文方法直接对照。

\paragraph{GASE}
GASE（2026.06）~\cite{2606_17520} 延续 3DGS 在“具身仿真与场景构建”上的基本设定，主要改变是：GASE：Gaussian Splatting-Based Automated System，从多视角视频重建仿真场景并提取前景对象。3DGS 的几何-渲染统一表示保持为底层地图。
其设计动机是：提出面向机器人学习的自动仿真场景构建系统，改进数据采集效率、视觉质量和前景对象提取。
具体而言，系统从多视角视频重建3DGS环境，分离/处理前景对象并导出可用于 embodied simulation 的场景资产。
适用边界是：场景物理属性和可交互性仍依赖后处理；复杂动态环境可能难自动化。

\paragraph{One Demo is Worth a Thousand Trajectories}
One Demo is Worth a Thousand Trajectories（2026.06）~\cite{2606_19586} 延续 Gaussian Splatting 在“具身仿真与数据增强”上的基本设定，主要改变是：Action-View Augmentation：面向广角鱼眼相机的Gaussian Splatting重建与编辑，加轨迹优化生成新视角动作序列。3DGS 的几何-渲染统一表示保持为底层地图。
其设计动机是：从一个eye-in-hand演示出发，联合生成视觉观测和物理可行轨迹，显著扩展visuomotor policy训练覆盖。
具体而言，方法针对wide-FoV fisheye cameras改造Gaussian Splatting，重建并编辑3D场景中的未知物体，再用轨迹优化产生平滑无碰撞且适合渲染的动作轨迹。
适用边界是：重建质量、鱼眼标定和轨迹优化约束共同影响数据真实性；可变形/接触丰富任务更难。

\paragraph{QuadVerse}
QuadVerse（2026.06）~\cite{2606_07118} 延续 3DGS 在“机器人仿真与数字孪生”上的基本设定，主要改变是：QuadVerse：用几何约束3DGS重建真实场景，支持照片级ego-view渲染和可碰撞语义mesh提取。3DGS 的几何-渲染统一表示保持为底层地图。
其设计动机是：提出同时对齐视觉感知、物理接触和执行器动态的集成框架，其中3DGS提供真实场景渲染，mesh提取提供碰撞和摩擦校准。
具体而言，系统从RGB视频重建几何约束3DGS，再提取语义mesh用于接触校准，通过轨迹后验搜索优化空间变化摩擦先验，并处理执行器动态差异。
适用边界是：mesh提取质量影响物理校准；3DGS本身不直接处理可变形交互；场景采集和校准仍有成本。
\section{强化学习与场景表示}
\label{sec:17_5}

将 3DGS 作为强化学习智能体的环境表示，是 3DGS 在机器人学习中的基础性应用方向。好的场景表示直接影响强化学习的样本效率和策略泛化能力——其核心挑战在于：传统 3DGS 需要逐场景优化，无法满足 RL 中实时编码观察的需求，因此需要设计\textbf{可泛化}（generalizable）的 3DGS 框架，使网络能够从任意新输入直接预测高斯参数。

\subsection{强化学习方法}
\label{ss:17_5a}

本节讨论将 3DGS 作为强化学习状态表示的设计方法。

\subsubsection{可泛化高斯泼溅表示设计}

本路线介绍将 3DGS 改造为可泛化表示以服务于强化学习训练的方法。

\paragraph{可泛化高斯泼溅表示}

GSRL（2024.04）~\cite{2404_07950} 修改了 3DGS 的场景构建环节——用前馈可泛化的图像条件化估计器替代逐场景优化——3D 高斯的渲染管线（$\alpha$-混合光栅化）和属性结构（位置、旋转、缩放、颜色、透明度）保持不变，从而使 3DGS 首次能够作为强化学习的实时状态表示。

\textbf{Baseline 做法与局限}。视觉强化学习中，传统状态表示主要分为三类：（1）图像表示——将多视角图像沿通道拼接后经 ResNet 编码，但 2D 图像无法感知 3D 空间结构，策略难以获得几何感知能力；（2）点云与体素——点云通过深度图反投影获得空间坐标，体素通过三线性插值将点云离散化为 $64^3$ 网格，但受分辨率限制，无法描述复杂的局部几何形状（如物体凹面、细长结构）；（3）NeRF 隐式表示——将场景编码为连续神经场，但通常需要前景掩码区分目标物体，且用单一向量表示整个场景导致无法泛化到新场景。更根本的是，传统 3DGS 虽然兼具显式结构与可微渲染的优势，但其逐场景优化范式需要在每个新场景上迭代数千步调整高斯参数——在 RL 中，智能体每一步都接收新的观察并需要即时编码，逐场景优化的时间开销完全无法满足这一实时性需求。

\textbf{核心洞察}。RL 任务中，同一任务的不同场景和观察在结构上具有高度相似性（例如机械臂操作任务中，桌面、物体、机械臂的几何拓扑始终类似），因此存在从 2D 图像块到 3D 局部几何的共享映射。3D 高斯的属性——位置描述全局坐标、旋转和缩放描述局部形状、颜色和透明度描述外观——天然适合作为这种局部结构的参数化载体。可泛化 3DGS 框架的本质，就是让网络在预训练阶段学习这种从 2D 图像到 3D 局部结构的先验映射，而非在每个新场景上重新优化。

\textbf{两阶段框架}。第一阶段训练图像条件化的 3DGS 估计器网络，包含三个模块。

\textit{深度估计模块}——从立体图像对桥接 2D 到 3D。给定两幅源视图 $I_{s1}, I_{s2}$，首先通过共享参数的特征提取器获得语义特征 $F_{s1}, F_{s2}$，再通过单应变换构建代价体（cost volume）预测视差图。直觉上，代价体衡量了两幅图像在不同视差假设下的匹配程度，类似于双目视觉中的视差感知。预测的归一化对数视差通过以下公式转换为绝对深度：

\begin{equation}
    D_\text{pred} = \exp(D_\text{max} \cdot \sigma(\mathcal{D}(F_{s1}, F_{s2})))
\end{equation}

其中 $\mathcal{D}$ 为视差预测网络，$\sigma$ 为 Sigmoid 激活，$D_\text{max}$ 为数据集最大视差值。指数变换确保深度值为正且覆盖从近到远的完整范围。有了深度后，每个像素通过 $\mathbf{X} = P_t \cdot D_\text{pred} \cdot K_t^{-1} \cdot \mathbf{u}$（$\mathbf{u}=\{u,v,1\}$ 为齐次像素坐标）反投影至 3D 空间，作为高斯的位置参数。

\textit{高斯回归模块}——预测每个高斯的形状与外观。将预提取特征 $F_s$、源图像 $I_s$ 和预测深度 $D_s$ 拼接后送入 UNet 编码器得到融合特征图 $\mathcal{F}_R = E_\phi(D_s \oplus F_s \oplus I_s)$，再通过三个独立的 $1\times1$ 卷积头分别预测旋转 $\mathbf{R}$（归一化激活确保四元数单位长度）、缩放 $\mathbf{S}$（指数激活确保正值）和透明度 $o$（Sigmoid 激活确保 $[0,1]$ 范围），颜色则直接复用源图像像素值 $\mathbf{c} = I_s$ 以简化流程。这意味着每个像素生成一个完整的 3D 高斯基元，其属性全部由网络前馈预测，无需任何迭代优化。

\textit{高斯精化模块}——平滑跨视角不一致的噪声。由于不同视角可能存在颜色偏差等固有不一致，初始预测的高斯属性含有高频噪声。该模块采用 GNN 自编码器，以 $K=16$ KNN 邻域在 3D 空间构建子图，编码器通过 3 层 MLP 提取 128 维节点特征，解码器恢复原始属性。GNN 的局部平滑特性天然充当正则化项，消除邻近视角差异导致的高频噪声。

第一阶段的总训练损失为渲染损失与 GNN 重建损失的加权和：

\begin{equation}
\mathcal{L}_\text{total} = \underbrace{\|R_\theta(I_{s1,s2}|C_t) - I_t\|_2^2}_{\text{渲染损失}\,\mathcal{L}_r} + \lambda \underbrace{\|D_\phi(E_\psi(G)) - G\|_2^2}_{\text{GNN 重建损失}\,\mathcal{L}_\text{recon}}
\end{equation}

其中渲染损失确保预测的高斯云能正确渲染目标视图，GNN 重建损失约束自编码器不丢失有效信息，平衡系数 $\lambda=0.15$。训练时先单独训练深度估计模块至收敛，再冻结深度模块联合训练回归和精化模块。

\textbf{RL 集成}。第二阶段将预训练的 3DGS 估计器完全冻结并集成至 RL 环境的观察管线：每一步，环境输出的多视角图像观察被实时编码为 3D 高斯云（前馈推理，无需优化），RL 策略网络（采用 SetTransformer 处理无序高斯点集）直接在该表示上训练。高斯点数默认为 8192，消融实验表明降至 4096 时性能几乎不受影响，降至 2048 时简单任务不受影响但困难任务分别下降 12\% 和 26\%。

\textbf{优势分析}。3DGS 继承了点云的空间结构特性，同时通过旋转和缩放参数提供更丰富的局部几何描述——点云仅有三维坐标，而每个高斯额外携带 4 维旋转和 3 维缩放，共 14 维属性（含颜色和透明度），表达能力显著更强。由于从多视角图像构建，3DGS 天然具有 3D 一致性和遮挡鲁棒性。在 RoboMimic 的四个任务（Lift、Can、Square、Transport）上，与图像、点云、体素三种基线表示在 BCQ、IQL、IRIS 三种 RL 算法下对比，该方法在大多数组合中取得最高成功率。在最难的 Transport 任务上，IRIS 算法下分别比图像、点云、体素提升 3\%（$36.0 \pm 2.0$ vs $33.0 \pm 3.0$）、11\%（$36.0$ vs $25.0$）和 4.5\%（$36.0$ vs $31.5$）；在 Square 任务上，BCQ 算法下相比图像、点云、体素分别提升 6.5\%、20.5\% 和 6.5\%。此外，渲染质量消融实验表明，3DGS 重建的 PSNR 从 25.9 提升至 33.7 时，Square 任务成功率从 42.0 提升至 48.5，验证了高质量几何表示对 RL 性能的正向贡献。

\paragraph{基于查询的语义高斯场}

Query-based Semantic Gaussian Field（QSGF）（2024.06）~\cite{2406_02370} 修改了 3DGS 的语义 grounding 与渲染驱动环节——在高斯基元上附加层次语义特征，并用基于表面点的查询机制替代逐场景优化来驱动高斯参数预测——标准 3DGS 的 $\alpha$-混合光栅化渲染管线保持不变。与 GSRL 仅利用几何先验（深度反投影）构建可泛化 3DGS 表示不同，QSGF 改变了场景表示的语义丰富度和渲染参数化方式——增加层次语义编码使每个高斯携带部件级语义信息，并采用基于查询的特征泼溅替代直接回归以实现从全局潜在向量到高斯参数的可泛化映射。设计动机在于：NeRF 基的 RL 场景表示（如 SNeRL）沿射线均匀采样，大量采样点落在空白区域导致几何感知能力弱，同时对自由空间和占据空间均匀建模使语义信息冗余且不精确；而 3DGS 的离散点表示天然集中在表面区域，但此前（如 GSRL）仅编码几何属性，缺乏细粒度语义信息——例如抓取锤子时需要区分手柄和锤头，这种部件级语义对操控策略至关重要。

该方法包含两大模块。

\textbf{层次语义编码}（Hierarchical Semantics Encoding, HSE）负责将细粒度语义特征 grounding 至 3D 高斯基元。具体而言，首先利用 SAM 为每幅输入图像生成部件级分割掩码（而非物体级掩码），然后对每个掩码区域分别执行 CLIP 特征提取——通过掩码平均池化（Masked Average Pooling, MAP）将掩码裁剪图像的 CLIP 特征聚合至像素级。由于一个像素可能同时属于多个层次（如一个像素既属于整个橱柜，又属于橱柜的把手），HSE 将该像素所属所有层次的部件特征取平均作为最终语义表示，从而保留层次化的语义信息。然而，CLIP-ViT 的 512 维特征空间对于 RL 任务中少量物体类别（机械臂、水龙头、橱柜等）高度冗余，直接将其作为高斯的附加特征会破坏 3DGS 的实时渲染效率。因此，HSE 训练一个轻量自编码器（MLP 结构）将 512 维 CLIP 特征压缩至 3 维紧凑空间 $l=3$，通过重建损失 $\mathcal{L}_{ae} = \sum_{i \in I}\|\mathcal{D}_\theta(\mathcal{G}_\phi(\mathcal{A}(I))) - \mathcal{A}(I)\|_2^2$ 监督，其中 $\mathcal{A}$ 表示从图像到层次语义特征图的完整映射。3 维语义特征作为高斯基元的附加属性，与颜色一起通过 $\alpha$-混合渲染为语义特征图，从而维持 3DGS 的实时渲染速度。

\textbf{基于查询的可泛化特征泼溅}（Query-based Generalizable Feature Splatting, QGFS）负责从全局潜在向量渲染场景，无需逐场景优化。与 GSRL 的多视角深度估计+直接回归范式不同，QGFS 首先用多视角编码器将 $V$ 个视角的 RGBD 观察编码为单一全局潜在向量 $\mathbf{z} = \Omega(o^{1:V}, K^{1:V})$，然后利用深度观测的几何先验将 RGBD 图像反投影至 3D 空间得到表面点集 $\mathbf{x} = P \cdot K^{-1} \cdot (u, v, \mathcal{D}(u,v))$，以这些表面点坐标作为查询键（query key）从全局潜在向量中检索对应信息——具体地，将点坐标映射为查询向量后通过交叉注意力机制与全局潜在特征交互，再经轻量 MLP 解码器预测每个点的高斯参数（旋转 $\mathbf{R}$、缩放 $\mathbf{S}$、透明度 $o$）和 3 维语义特征 $f_i$。核心区别在于：NeRF 体渲染沿射线均匀采样，大部分采样点位于空白区域且不利用几何先验；QGFS 仅在几何表面上执行查询，避免了对空白区域的冗余计算，因此相比 NeRF 具有更强的几何感知能力和数据效率。渲染时采用标准 3DGS 的 $\alpha$-混合同时输出颜色图和语义特征图，语义特征图的自编码器重建损失提供语义监督。

\textbf{RL 集成}。与 GSRL 类似，QSGF 也采用两阶段范式：第一阶段离线训练 3DGS 编码器（含 HSE 和 QGFS），第二阶段冻结编码器参数，RL 策略直接在全局潜在向量 $\mathbf{z}$ 上训练。不同之处在于，QSGF 的潜在向量同时编码了几何和语义信息——语义特征通过蒸馏从 HSE 注入全局潜在向量，使策略能够获得部件级的场景理解。在 ManiSkill2（6 个任务）和 Robomimic（4 个任务）共 10 个不同 RL 任务上，QSGF 在 8 个任务中取得最佳成功率，其余 2 个任务取得第二。与 NeRF 基线 SNeRL 相比，关键优势来自 HSE 的细粒度语义编码——消融实验表明移除 HSE 后性能显著下降，验证了部件级语义对操控策略的积极作用；同时 QGFS 的表面点查询机制比 NeRF 的射线采样更高效地利用几何先验，直接反映在更少的训练步数和更高的最终性能上。

\paragraph{RL 后训练增强导航}

S2E（2025.07）~\cite{2507_22028} 提出从观看到体验（Seeing-to-Experiencing）的混合学习框架，通过 RL 后训练扩展导航基础模型的交互能力。与 SplatGym 在 3DGS 环境中从零训练 RL 策略不同，S2E 改变了 RL 在导航基础模型中的使用方式——采用两阶段混合范式（离线视频预训练 + 仿真环境 RL 微调），先在大规模真实世界导航视频（35 小时，覆盖多种机器人和环境）上学习可泛化的视觉运动表示，再在仿真环境中通过 RL 注入避障和行人避让等反应行为，而非将 RL 作为唯一的训练手段。设计动机在于：纯离线视频预训练虽然能学习到广泛的视觉泛化能力，但视频数据缺乏物理交互和因果反馈，模型无法理解动作后果（如在人群中穿行、避让移动障碍物），仅增加离线数据量会遭遇严重的边际收益递减（从 100k 扩展到 200k 样本仅提升 3\% 成功率）；而纯 RL 从零训练又受限于采样效率和缺乏归纳先验，难以获得跨场景的泛化导航能力。S2E 通过两阶段学习克服各自的局限。

\textbf{锚点引导分布匹配}（Anchor-Guided Distribution Matching）：针对导航轨迹固有的多模态性（同一观测下多种合理动作并存），S2E 采用锚点引导高斯混合模型（GMM）表征动作分布。具体地，在机器人前方均匀采样 $M$ 个意图点 $\bm{p}_a \in \mathbb{R}^{M \times 2}$ 作为锚点，每个锚点对应一个高斯模态并学习选择概率 $q_m$，动作分布表示为 $\bm{q}(w_t | \bm{o}_{t-k+1:t}) = \sum_{m=1}^{M} q_m \cdot \mathcal{N}_m(w_x - \mu_x^m, \sigma_x^m; w_y - \mu_y^m, \sigma_y^m; \rho^m)$，其中 $\mu, \sigma, \rho$ 分别由第 $m$ 个轨迹头预测。每个模态还预测速度标量 $v \in \mathbb{R}^+$ 以建模绝对运动幅度。训练采用负对数似然（NLL）损失——选择与真值方向最匹配的锚点模态 $h$ 进行优化：$\mathcal{L}_\text{nll} = -\log \mathcal{N}_h(\hat{w}_x - \mu_x^h, \sigma_x^h; \hat{w}_y - \mu_y^h, \sigma_y^h; \rho^h) - \log(q_h)$——加上速度回归损失 $\mathcal{L}_\text{reg} = \|\hat{v} - v\|_2^2$。与离散动作、单模态高斯（如 ViNT）和扩散策略相比，锚点引导 GMM 兼具结构化多模态表达力和训练稳定性：不同锚点专门化不同高层意图（直行、转向、让行），模态内保留适度的方差以表达不确定性。相比 ViNT 的单模态匹配，S2E-BC 在障碍物+行人场景中成功率提升 33\%、碰撞次数减少 2.11 次。

\textbf{残差注意力模块}（Residual Attention Module, RAM）：在仿真中用 PPO（256 个并行环境，每个含随机分布障碍物的 $30 \times 30$ 区域）微调时，RAM 解决了直接全参数微调导致预训练知识被仿真数据覆盖（分布漂移）的问题。核心设计是：冻结视觉编码器和自注意力层（对域偏移敏感），仅适配交叉注意力层（显式建模智能体-环境交互，跨域一致性好）。具体地，冻结原始交叉注意力参数 $\Theta$，创建可训练副本 $\Theta_l$，用两个零初始化线性层 $\mathcal{Z}$ 包裹，输出为 $\bm{y}' = \bm{y} + \mathcal{Z}(\psi_D(\mathcal{Z}(\bm{y}); \bm{K}, \bm{V}; \Theta_l))$——零初始化确保训练开始时适配分支对原始输出无影响，之后逐渐学习调制交互相关特征。PPO 的奖励函数为 $R = R_\text{vel} + R_\text{dis} + R_\text{arrive} + R_\text{collision}$，分别鼓励朝目标运动、惩罚每步位移、到达奖励和碰撞惩罚。消融实验表明 RAM 在 NavBench-GS-Obstacle 上成功率 0.76、碰撞次数 0.56，优于全参数微调（0.74/0.83）和纯 SFT（0.71/0.77）。

\textbf{NavBench-GS 基准}：基于照片级 3DGS 重建（源自 Vid2Sim）构建的端到端评估基准，包含 26 个场景 $\times$ 4 种任务配置（空场景、静态障碍物、移动行人、障碍物+行人），以成功率、路径完成率和碰撞次数评估导航基础模型的泛化性和安全性。与 Wanderland 用于开放世界具身 AI 评测不同，NavBench-GS 专注于城市导航场景中策略的避障和交互能力评估。S2E-Full 在 NavBench-GS 上相比纯预训练模型（S2E-BC）在障碍物场景成功率提升 21\%，在真实世界轮式机器人和四足机器人上均实现零样本迁移。此外，RL 后训练相比 SFT 后训练在增加训练成本时保持鲁棒性能，而 SFT 出现严重过拟合。

\paragraph{动作条件世界模型}

ChronoDreamer（2025.12）~\cite{2512_18619} 修改了 3DGS 的渲染输入与语义通道——不再渲染场景外观，而是将 3D 接触力渲染为与相机对齐的图像——标准高斯泼溅的光栅化管线和各向同性核结构保持不变，但语义通道从 RGB 颜色重新定义为力方向和力大小的物理量编码。与 DreMa 将 3DGS 嵌入世界模型用于渲染未来视觉外观不同，ChronoDreamer 改变了高斯泼溅在世界模型中的角色——从渲染"场景看起来怎样"变为渲染"接触力在哪里、有多大、朝哪个方向"，使接触信息以图像原生格式直接输入视觉骨干网络。

\textbf{深度加权高斯泼溅接触编码}的设计动机在于：接触丰富的操控任务中，3D 接触力本质上是稀疏的（仅发生在物体表面微小区域），但现有世界模型无法有效处理这种稀疏 3D 物理量——要么将其简化为标量信号丢失空间信息，要么需要在 3D 空间中维护昂贵的显式接触场。ChronoDreamer 借用 3DGS 的泼溅思想，将每个接触点视为一个"接触高斯基元"：通过针孔投影将 3D 接触位置 $\mathbf{p}_i$ 映射至图像平面，力方向（沿力向量投影后的 2D 位移归一化并映射至 $[-1,1]$）编码到绿蓝通道、力大小（截断归一化）编码到红色通道和各向同性高斯核半径 $r_i = r_{\min} + (r_{\max} - r_{\min})(m_i/m_{\max})^\gamma$。关键差异在于引入深度衰减权重 $w_i^{\text{depth}} = \exp(-X_i / \tau_{\text{depth}})$，使距离相机更近的接触点通过类软 z-buffer 的归一化加权混合 $\mathbf{I}(u,v) = \tilde{\mathbf{I}}(u,v) / W(u,v)$ 获得更高影响力——这与标准 3DGS 的深度无关 $\alpha$-混合排序不同，是专为物理接触场景设计的深度优先策略，确保机器人末端执行器附近的接触信号不被远处环境接触淹没。

\textbf{空间-时间 Transformer}采用 24 层因子化时空注意力块，通过 Cosmos DI8 编码器将 RGB 和接触泼溅图像经有限标量量化（FSQ，词表大小 $V \approx 65{,}536$）为 $32 \times 32$ 离散 token 网格，以 MaskGIT 迭代去掩码训练目标联合预测未来视频帧 token、接触分布图 token 和关节角度。动作指令与关节角观测先分别线性投影后加性融合为每帧的控制 token，再与视觉 token 拼接为增强序列输入 Transformer。推理阶段，预测的 rollout 序列通过视觉-语言模型（VLM）进行碰撞推理，实现不安全动作的拒绝采样。在 DreamerBench 基准（基于 Project Chrono 物理引擎生成的同步 RGB、接触泼溅、本体感知和物理标注数据集，覆盖刚体和可变形物体场景）上验证了接触感知世界模型在操控规划中的有效性。

\paragraph{RoboTidy 基准}

RoboTidy（2025.11）~\cite{2511_14161} 提供统一的家庭整理语言引导基准，包含 500 个照片级 3DGS 家庭场景、6.4k 操控示范轨迹和 1.5k 导航轨迹，支持 VLA 和 VLN 的训练和评测。与 ChronoDreamer 作为动作条件世界模型用于操控规划不同，RoboTidy 改变了 3DGS 的应用角色——从规划辅助工具扩展为完整的语言引导整理评测平台，首次在 3DGS 场景中联合评估高层语义推理（物体-容器匹配）和低层反馈驱动控制（操控执行与跨房间导航）。设计动机在于：现有家庭整理基准（如 TidyBot）要么依赖用户显式指定每件物品的目标容器导致可扩展性不足，要么通过跨用户平均学习非个性化规则忽略个体差异，且普遍缺乏物理仿真和照片级渲染能力，难以全面评估语言到动作的端到端能力。

\textbf{方法框架}：基于 NVIDIA Isaac Sim 构建，系统采用四模块化架构。（1）\textbf{物体分类模块}：利用 Qwen2.5-VL 提取物体类别和属性语义，通过 LLM 将观察到的物体-容器对应关系归纳为个性化整理规则，并设计四种最小化操控动作原语——拾取放置（Pick and Place）、拾取抛投（Pick and Toss）、开容器和关容器。分类规则覆盖属性（材质/尺寸）、功能（用途场景）、安全性（危险/易碎）和卫生（污染风险）四个维度，当多个准则冲突时按预设优先级选择主容器，其余作为等效备选。（2）\textbf{3DGS 场景构建}：基于 InteriorGS 构建 500 个覆盖客厅、卧室和厨房的家庭场景。关键设计是 3DGS-网格混合表示——3DGS（通过 3DGUT 导出为 USDZ）负责照片级渲染，艺术家创建的三角网格经 CoACD 凸分解生成逐物体碰撞体作为不可见刚体，使场景同时具备高保真视觉和精确物理接触。（3）\textbf{操控模块}：IK 求解器与运动规划器在关节限制、速度/加速度约束和碰撞约束下生成无碰撞轨迹，夹持器采用闭环宽度阈值判据确定抓取和释放时机，同步记录多视角 RGB 图像、关节配置和夹爪状态。（4）\textbf{导航模块}：A* 路径规划器在 InteriorGS 的 2D 语义地图上规划参考路径，PID 控制器追踪轨迹并记录完整导航数据。评估指标包括物体放置准确率（OPA）和有效整理成功率（VSSR），后者要求容器预测正确且操控动作全部完成。

\paragraph{最终米数导航}

POINav（2026.05）~\cite{2605_28237} 修改了 3DGS 在导航基准中的角色定位——从 RoboTidy 以 3DGS 作为家庭场景的静态评测环境，转向以 3DGS 商业街区重建作为兴趣点（Point of Interest, POI）目标导航的闭环评测与训练平台，并针对最终米数到达问题提出 Brain-Action 双系统框架。与 RoboTidy 在家庭整理场景中评估物体放置准确率不同，POINav 关注的是真实世界导航的最后一公里挑战：当机器人已到达 POI 附近后，如何精确识别并物理抵达目标商铺的入口。设计动机在于：现有 POI 目标导航基准要么依赖合成轨迹数据存在仿真-真实差距（如 BridgeNav 的 AI 生成视频），要么以粗粒度街道段为单位评测（如 CitySeeker），无法评估四足机器人在高密度商业区域中以米级精度到达特定商铺入口的能力。

\textbf{POINav-Bench} 是首个面向真实世界 POI 目标导航的闭环评测基准，包含 11 个商业区域的高保真 3DGS 重建（采用高精度 LiDAR 和摄影测量采集，覆盖 126,398\,$m^2$ 和 163 个独立 POI），集成至 NVIDIA Isaac Sim 平台进行闭环仿真评测。场景标注包括可通行区域标注（基于表面法向量分割可通行地面并投影为 2D 占据栅格）、水平包围盒入口定义（两个齐平点之间的水平边界框覆盖入口全宽）、A* 最短可行路径真值轨迹和仅提供 POI 名称（而非固定指令模板）的灵活评测协议。所有起始点经过人工验证，确保从四足机器人视角至少可见部分招牌视觉线索。

\textbf{Brain-Action 框架}将导航解耦为语义推理与几何执行两个互补模块。

\textbf{语义大脑模块}将 POI 接地建模为层次依赖推理——将目标分解为 POI 招牌 $b_{id}$（语义锚）和物理入口 $b_{geo}$（几何定位）两级视觉参考，通过条件概率链 $P(\mathcal{G}_t|V_t,L) \propto P(b_{id}|V_t,L) \cdot P(b_{geo}|b_{id},V_t,L)$ 强制入口定位以招牌为条件，避免将入口误定位至无关商铺。为支持该推理，作者构建 POINav-Dataset（70K 真实世界招牌-入口配对），通过多阶段自动化流水线（极端视角过滤、几何约束验证、物理可通行性检查）从街景序列中提取高质量训练样本。

\textbf{几何动作模块}通过弹性时序采样保持全局运动记忆——对完整观测历史 $\mathcal{H}_{raw}=\{V_0,\ldots,V_t\}$ 按时间膨胀索引函数均匀采样 $K$ 帧，构建稀疏但全局一致的上下文集 $\mathcal{H}_{ctx}$，在保持常量计算复杂度的同时覆盖整个回合的时空信息。可学习潜在查询 $Q_{nav} \in \mathbb{R}^{N \times d}$ 经交叉注意力从视觉语言骨干提取空间特征，轻量 MLP 解码为连续路径点序列 $\mathcal{W}_{future}=\{\Delta p_1,\ldots,\Delta p_H\}$。两模块通过可见性条件化的动态视觉线索耦合——当入口 $b_{geo}$ 在视野内可解析时裁剪入口区域，否则回退至招牌区域 $b_{id}$ 提供方向引导（$C_t = \mathbb{I}[b_{geo} \in V_t] \cdot \Phi(V_t,b_{geo}) + (1-\mathbb{I}[b_{geo} \in V_t]) \cdot \Phi(V_t,b_{id})$），实现从远距离语义锚定到近距离几何对齐的确定性视觉交接。在 POINav-Bench 上，POINav 以 SR(2m)\,42.33\%、SPL\,40.29\% 显著超越 OmniNav 的 34.36\%/31.52\% 和 ViNT 的 19.01\%/18.24\%（均使用 Qwen2.5-VL-3B 骨干），验证了双系统框架在最终米数到达任务上的有效性。

\begin{takeaway}
强化学习与场景表示将 3DGS 作为 RL 智能体的环境表示，核心挑战在于设计可泛化的 3DGS 框架以支持实时编码，3DGS 商业街区重建则为最终米数导航等精细任务提供了闭环评测与训练平台。
\end{takeaway}

% ======================================================================

\section{近期研究（2026 年 8--11 月）}
\paragraph{以可刷新的语义高斯地图连接移动操作的各阶段}
2026.08~\cite{2608_10756}。该系统将可更新的语义高斯作为主动观测、语言目标定位、障碍几何、可达站位与视觉语言动作模型的共享接口。3DGS 在此承担可查询地图，定位结果还需经几何与可达性处理后才能用于操作。

真实机器人实验考察长流程任务、高度变化、易混淆干扰和杂乱环境。重建、语义定位与低层动作各有误差来源，结果对应论文验证的装置与任务，不能等同于对开放环境的通用操作能力。

\paragraph{SpotlessGS：把车载或机载局部照明与表面反射分开}
2026.08~\cite{2608_14713}。在暗环境中移动自带光源会改变同一表面的外观，普通高斯难以分辨光照变化与场景属性。SpotlessGS 在 DarkGS 基础上给高斯加入法向和可学习反射属性，用轻量 BRDF 网络建模非朗伯行为，同时估计球谐环境光以解释照明。

论文用合成和自采有参考的图像测试重照明，在自采数据上还比较原图与重照明图的 COLMAP 重建。下游结果针对该采集条件，不能等同于已验证长期自主机器人跟踪。

\paragraph{GaussMemory：以任务查询读取并更新持久高斯记忆}
2026.08~\cite{2608_14986}。GaussMemory 把多视角 RGB 经前馈高斯编码转换为观测标记，并在统一注意力中让任务查询读取持久场景记忆。另一输出分支学习把新观测匹配、插入或裁剪记忆原语，使被操作物体的状态能随时间更新。

作者在 LIBERO 的空间、对象、目标和长程任务以及 VLABench 比较成功率和进展。能够处理这类任务依赖多视角感知与训练设置，表中的部分基线还使用额外腕部相机和本体状态，比较时应保留输入差异。

\paragraph{OccamView：在相同采集帧预算下优先补足物体附近几何}
2026.08~\cite{2608_16499}。OccamView 从在线 RGB-D 测量和目标检测维护物体记忆，将附近未观测体素作为物体遮挡区的代理。候选视角按考虑遮挡的代理覆盖率重新排序，但只有原始几何收益满足阈值的视角才允许重排，从而保留原有探索目标。

在 Replica 与 MP3D，作者按相同观测帧数而非规划轮数比较重建，报告完整度收益。遮挡代理不是真实未观测表面，检测遗漏和深度支持不足会影响它的质量。

\paragraph{GS-VLA：在冻结机器人策略前校正相机视点}
2026.08~\cite{2608_19066}。GS-VLA 接在冻结视觉语言动作策略前，先利用标定深度把可见内容几何映射到规范视点，只用局部高斯模块补足新露出的区域。策略权重保持不动，修正发生在图像观测空间。

作者在 LIBERO 四个套件对比多种策略与相机扰动幅度。方法需要单视图的标定深度，在源视图完全看不到的内容不能由重投影可靠复原，因此不能把有界扰动的结果推成大幅换视角鲁棒性。

\paragraph{ChainSplat：以开链动力学连接可变形线体的高斯外观}
2026.08~\cite{2608_28570}。ChainSplat 将线状物的运动表示为基于旋量理论的低维开链，同时让每段链节携带高斯外观。从多视图 RGB 先联合恢复外观、几何与关节状态序列，再识别动力学参数；这个物理模型可用于后续状态估计和模型驱动动作优化。

论文在桌面接触和空中悬挂两类机器人操作中评估，分别对应摩擦和重力主导的条件。开链和多视角观测是方法假设，不能直接推广到拓扑变化或任意软体。

\paragraph{RIDE：用高斯重定位的米制锚点校准单目视频深度}
2026.09~\cite{2609_11079}。RIDE 从既有米制高斯场执行重定位，保留通过 PnP-RANSAC 的匹配，用这些稀疏米制深度约束冻结的视频深度模型。在线分开估计全局尺度、有限空间修正和时序记忆，使少量锚点不必单独决定全部像素深度。

论文按场景划分公共 RGB-D 数据，并测试五个场景中的 27 条真实机器人路线。部署不需额外深度传感器来提供每帧锚点，但前提是已有米制高斯地图且能找到有效重定位匹配。

\paragraph{Neverwhere：从真实场景重建构建视觉跑酷基准}
2026.09~\cite{2609_16443}。Neverwhere 从多视图影像计算相机位姿和多视立体网格，再采样网格初始化高斯。网格承载仿真场景几何，高斯提供视觉观察。构建场景时仍需手工对齐坐标并标注跑酷路线点。

基准把跨障碍、越间隙、斜坡和楼梯拆成逐类任务，检查视觉运动策略成功率和路线进展。这里重建与真实环境之间的差异仍需通过实机试验衡量，基准上的泛化不能直接推出部署表现。

\paragraph{SemSafe-3DGS：用语义和地图不确定性调整导航安全约束}
2026.09~\cite{2609_19330}。SemSafe-3DGS 给地图原语增加语义风险权重，把碰撞后果和几何不确定性一起用于构造控制屏障。主动观测根据当前轨迹调整感知方向，以改善真正影响控制的局部地图。

实验分别考察组合屏障效率、轨迹相关主动感知，以及 Ackermann 实机部署。安全约束依赖地图、语义和所用车辆动力学，观察到的导航收益不等同于任意场景的无碰撞保证。

\paragraph{AirSplan：在高斯场中评估飞行路径的碰撞风险}
2026.09~\cite{2609_21226}。AirSplan 用一组随轨迹变化的体积包络近似飞行器占据，在归一化高斯地图上估计这些包络的碰撞风险，并把风险作为多候选路径优化的约束。这里高斯地图是风险查询的输入。

论文分别在仿真中比较碰撞查询和端到端规划，作者明确指出实机验证仍缺，当前还假定地图准确且能追踪计划轨迹，并需要外部计算。结果不能视作实飞安全保证。

\paragraph{Splat-CBF：在同一控制问题中处理避障与主动观测}
2026.09~\cite{2609_23100}。Splat-CBF 把风险感知的安全屏障与沿轨迹的信息感知屏障放进同一个控制二次规划。发生冲突时，用自适应松弛降低感知要求，并保留避障约束。

四场景的受控模拟中，两种配置观测到的安全率均为 99.5\%，与 SAFER-Splat 相同。论文还区分室内导航、Isaac Sim 机械臂与实物 Ackermann 车辆。所用动力学、地图与软硬约束有别，不能直接交换三组结果。

\paragraph{PlenoCI：用全光场导数区分几何和外观变化}
2026.09~\cite{2609_28930}。两次建图的颜色差可能是物体真的变了，也可能是观察方向改变。PlenoCI 将高斯场用作全光场，用可解析的空间和方向导数刻画原语在不同射线下的响应。实例高斯提供对象组织，导数特征再用于发现变化和区分几何、外观两类原因。

论文将 CL-Splats 和 PASLCD 的二类变化检测与 PASLCD 的多类评价分开，并用未变化独立采集考察误报。对照属于所测静态与半静态场景，不能从几何与外观的区分推断任意动态光照或运动都能正确分类。

\paragraph{Robot-GST：将重建外观和交互几何联合用于机器人行为预演}
2026.09~\cite{2609_33872}。Robot-GST 用短 RGB-D 记录和 SAM3D 对齐可渲染高斯及对象级几何代理。高斯用于提供照片级外观，显式代理承担碰撞、接触等查询。语言任务先变成以三维关键点和关系表示的多段行为，再在重建环境里比较候选结果。实机执行使用估计的目标状态规划动作。

作者用刚性、软质和可形变物体任务检查仿真筛选与实机表现。遮挡、场景变化及接触模型近似会影响预演可靠性，照片级渲染也不等于已经验证相应物理参数。

\paragraph{CollisionSplatting：用高斯空间不确定性评估避碰代价}
2026.09~\cite{2609_35619}。CollisionSplatting 把场景高斯的空间不确定性纳入距离和避碰代价，利用可调保守系数扩大对潜在障碍的避让。GPU 并行碰撞检查作为模块接入采样规划和模型预测控制，视觉渲染特征可同时用来定义目标。

论文给出仿真中的吞吐量、显存和路径对比，也在真实机械臂和移动机器人上展示了受视觉目标约束的任务。保守度提高通常会带来更多绕行，图中以高斯中心的最近距离做代理指标，实验结果不能当作对新场景无碰撞的保证。

\paragraph{从带状态的模拟教师蒸馏仅用 RGB 的导航安全滤波}
2026.09~\cite{2609_36520}。该工作用高斯渲染生成的交互仿真为教师提供动态障碍状态。训练时，控制屏障函数教师修改给定导航指令，学生从相机图像序列与控制输入学习类似校正。部署时不用查询三维高斯地图，地图在这里主要服务训练和模拟。

作者在模拟和实机动态避障试验中统计成功率、无碰撞率与安全距离。学生的表现由这些试验衡量，不能把教师使用的屏障函数形式自动解释成实机学生的形式安全证明。

\paragraph{在返航时间约束下主动采集三维高斯重建视角}
2026.10~\cite{2610_04013}。系统用 RGB-D 更新在线高斯地图供渲染重建，并另保留 TSDF 与占据数据作为碰撞查询和三维 Voronoi 路网的基础。下一视角选择同时考虑信息量、行程和剩余飞行时间，并保留返航所需的时间预算。

作者比较室内照片级仿真中的覆盖、重建误差和返航结果，Gibson 与 HM3D 用场景平均，ReplicaCAD 用物体实例平均。周期性重建路网和遇障时的反应式调整仍有时间开销，返航表现只在所测试的仿真预算与航行模型下得到评价。

\paragraph{FOCUS：从既有高斯地图分辨物体保留、替换和移除}
2026.10~\cite{2610_05639}。同一位置换成外形相似的新物体时，只比对占据体积容易漏检。FOCUS 用带位姿的 RGB-D 序列建立参考图，在无变化的首次采集回放中估计各物体的正常观测残差。后续访问将几何一致性与高斯渲染出的视角相关外观作为两类似然，连续更新保留、替换和移除三种状态。

作者在 Isaac Sim 仓库中报告物体替换 F1 从概率基线的 0.20 升至 0.84，并以相同参数评测真实 TorWIC。消融显示，几何通道能识别移除但难以识别原地替换，单用外观又更易误报。结果依赖参考建图、分割和重访观测，不足以保证换季或长期光照变化下仍能保持原有校准。

\paragraph{GS-Pool：在两次高斯建图中汇总逐对象的变化证据}
2026.10~\cite{2610_06688}。GS-Pool 需要对同一空间两次访问分别训练高斯场，并使相机坐标一致。方法将 SAM2 的二维对象掩码提升到有渲染贡献的高斯，在对象内汇总几何、颜色、语义与光度变化线索。光度线索来自把另一次访问的照片作为对照，将图像误差经渲染器回传到高斯。共享未变对象提供每个场景自己的参照。

在 PASLCD 的共同坐标设置下，采用 FastPGSR 的一组结果为 mIoU 0.7509、F1 0.8461，作者报告 mIoU 相对既有最佳结果提高约 17\%。换用原始 3DGS 训练器时，论文报告相对提升约 12\%，两者是不同配置下的 mIoU 对照。分割未产生单独掩码的变化无法在后续找回，单独光照变化也可能造成误报。

\paragraph{SURGE：融合视觉和声呐定位后进行目标级高斯重建}
2026.10~\cite{2610_07472}。水下单目视觉能提供目标成像方向，但缺少可靠米制尺度。SURGE 利用图像门控从多帧二维成像声呐中获取距离线索，在因子图里联合估计拖缆机器人的轨迹和目标中心，并让声学残差容许主要回波来自近侧表面。随后用米制位姿训练目标范围内的声呐高斯。

实验使用淡水中环绕单个目标的记录，按离线流程处理，报告相对光学基线更集中的目标表示。声呐俯仰歧义、视角覆盖有限和缺少独立测量真值限制了结论，结果并未验证精确表面网格或全场景重建。

\paragraph{2DGS-Planner：直接在可渲染地图中查询路径净空}
2026.10~\cite{2610_11752}。2DGS-Planner 不把整张高斯地图先转换成密集占据数组，而对候选路线调用特定视点的光栅化。多视图结构证据分配路线图节点，候选边附近的几何查询形成可缓存的局部净空数据，再用于不同机器人半径下的在线路径搜索。

原文对路线图构建、净空保真和完整路径规划分开评价。定量环境包括拥有真实参考网格的合成 Stonehenge 与只提供配准代理网格的 Tokyo，两处几何参照可信度不同。系统的净空判断仍取决于重建覆盖与地面估计，实验主要测已有静态地图。

\paragraph{NovaSplat：从手机视频建立无 CAD 对象位姿模型}
2026.10~\cite{2610_13309}。NovaSplat 用目标物体短视频进行一次建模，经 SfM、多视立体和高斯训练获得几何模板和可微渲染表示。查询阶段先做候选定位与外观、几何匹配，再固定高斯参数，通过光度和几何损失联合优化物体相对相机的六自由度位姿。

HOPEv2 和 IndustryShapes 分别按 BOP 的检测和定位协议度量，不能横向比较两者的 AP 与召回。作者报告在 HOPEv2 上与使用 CAD 模型的方法比较有优势，在 IndustryShapes 上仍有差距，尤其是低纹理和强反光物体会影响重建与匹配。

\paragraph{YUBI-Real2Sim：检验高斯外观缺陷会不会改变策略排名}
2026.10~\cite{2610_17521}。YUBI-Real2Sim 用实机排名作为参照，在物理、物体、初值和随机种子保持一致时改动高斯背景，分开比较空洞、漂浮伪影和减少采样的设置。这样观察到的排名变化与外观条件有较清楚的因果联系。

清洁重建的仿真在九个场景任务组合中有八个恢复实机排名，而缺陷条件下更少。排名即使不变，也可能失去统计区分。依据来自所测场景和任务，不能单凭高图像分数保证任意策略排序。

\paragraph{I2O2I-Nav：将室内外连续导航拆成阶段并单独评价交接}
2026.10~\cite{2610_18307}。I2O2I-Nav 把出楼、室外移动、入楼和开放词汇目标搜索串成连续任务。场景用高斯提供视觉，用配准几何判断碰撞。其模块化智能体协调原有子策略，分别处理完成确认、阶段交接和有限次数恢复。

实验分别评估单阶段、一次交接及完整任务，在二十个场景的 1374 个回合用受碰撞约束的运动学智能体运行。人工布景和重建伪影会影响基准覆盖与感知，结果尚不包含关节动力学条件的实机验证。

\paragraph{GS-PRM：用高斯渲染构建和查询视觉路网}
2026.10~\cite{2610_19541}。GS-PRM 以高斯渲染的深度、累计不透明度及覆盖度判断候选位置和路径是否得到地图支持，在通过检查的区域建立可复用路网。查询时用目标图像在已有地图中定位，将起终点接入路网，并为局部策略渲染途中目标。

作者在仿真里用真实碰撞几何检查渲染判据，并在真实差速机器人上评价视觉导航。论文给出的完备性和物理无碰撞结论以覆盖、局部可连接性及保守表面渲染假设为前提，低置信地图区域不可据此视作已验证安全。

\paragraph{借助可微高斯场联合设计编码孔径和深度感知}
2026.10~\cite{2610_28004}。该工作把高斯模型作为可微拍摄环境，在候选相机参数下渲染针孔 RGB 与深度，再用随深度变化的点扩散函数模拟编码孔径。恢复全焦图和估计深度的任务损失可反传给相机和孔径参数，使传感器设计与感知网络一起优化。

论文在单目和双目散焦深度任务比较，并补充原型实拍。结果针对所用光学近似、虚拟场景覆盖与评价任务，不能将图像保真度直接解释为所有下游任务的传感器收益。

\paragraph{Real2Sem2Real：从一个视频获得可迁移的物体级动作数据}
2026.10~\cite{2610_28880}。Real2Sem2Real 用带语义特征的高斯表示跟踪示范视频中的物体状态，再将抓取和放置等操作对齐到同类但外形不同的物体，生成可用于仿真的动作数据。流程还需要示范物体扫描和同类物体参考，单个视频提供的是动作示范。

作者比较四类任务和实机回合，也通过相同视频与同规模样本比较是否引入同类物体变体。流程要求示范物体具有可靠米制信息，这项条件限制了从纯单目视频直接执行完整流程。

\paragraph{AOV-GS：早期用语义新颖度选视，后期回到几何细化}
2026.10~\cite{2610_29374}。AOV-GS 在探索初期用冻结 SAM 与 CLIP 提取候选掩码的语义新颖度，让机器人较早观察更多对象；随后回到几何优先的选视策略。探索结束后再用采集到的嵌入训练场景专属的 LangSplatV2 语言场，因此探索与最终文本查询分成两个阶段。

论文在 Replica 和 Matterport3D 的 Habitat 仿真比较。方法使用相同的真值相机位姿，实验隔离了选视策略的影响，没有同时验证带 SLAM 位姿误差时的效果。平均语义收益也不表示每个场景都改善。

\paragraph{AirPlayground：用重建外观和几何统一比较飞行策略}
2026.10~\cite{2610_29699}。AirPlayground 将每个环境的三维高斯用于相机外观，将同坐标系的米制点云转为用于采样与碰撞的几何。通过统一接口和适配器，不同导航策略可使用相同航迹训练，并在相同闭环条件下执行。

论文把 94 个环境划为 70 个训练、8 个验证及 16 个未见测试场景，并分析离线预测与闭环执行的差异。场景的飞行空间仍由点云、占据和人工标注的范围决定，重建仿真得分不能代替真实飞行测量。

\paragraph{SceneAgent：将拍摄记录转换为可交互仿真场景}
2026.10~\cite{2610_34745}。SceneAgent 调度不同的重建与物理建模工具，将物体、背景、关节和接触等表示组装在仿真器里，再用渲染检查识别并修复部分不一致。对于动态视频，还将模拟运动拟合到观测。高斯等视觉重建给出外观线索，但交互所需的几何与物理属性要另外建模。

论文在静态场景和 4DCodeBench 比较重建，并给出初步实机迁移。视觉检查可以发现部分外观或几何问题，却无法确认未观测结构与物理参数。反复调用模型和外部计算也限制复杂场景的转换速度。

\paragraph{PrecisionReach-GS：在机器人坐标下生成腕部观察与接近动作}
2026.10~\cite{2610_35626}。PrecisionReach-GS 把真实工作区的高斯模型对齐到机器人基座坐标，用可行末端位姿渲染腕部相机观测，并和对应的平滑接近轨迹配成视觉动作样本。部署时由腕部相机驱动微调后的视觉语言动作策略，并用指令相关的验证器判断何时停止。

实验在 Franka FR3 的主板场景验证三个部件的接触前接近。中等程度的位置变化与完整布局变化是不同测试，后者可能要重新重建和生成数据；是否能到达目标也不直接证明接触后拆解或抓取的成功。


\section*{本章小结}
\addcontentsline{toc}{section}{本章小结}
\label{sec:17_summary}

\begin{futurebox}{本章总结与展望}{ch17_summary}
本章系统梳理了 3DGS 在机器人导航与具身智能领域的 170 篇论文，覆盖五个核心方向。

导航与路径规划（56 篇）是论文最多的方向，3DGS 已成为导航地图表示的重要选择。从安全规划角度看，SAFER-Splat、SPLANNING 和 FOCI 分别从 CBF、风险上界和高斯重叠积分三个角度将高斯基元转化为可计算的碰撞约束。从视觉导航角度看，IGL-Nav、SplatSearch 和 3DGSNav 展示了 3DGS 可渲染性在图像目标导航和语义搜索中的独特价值。无人机导航是最活跃的子方向，SOUS VIDE、SINGER、GRaD-Nav 等系统展示了从 3DGS 仿真到真实硬件零样本迁移的完整路径。自动驾驶方向，GSDrive、RAD 和 Senna-2 利用 3DGS 环境进行闭环 RL 训练，显著降低碰撞率。

机械臂操控（42 篇）涵盖了从抓取规划到灵巧操控的完整技术栈。功能区域推理方面，3DAffordSplat 构建了首个大规模 3DGS 功能区域数据集。策略学习方面，ManiGaussian++ 和 GAF 分别从双手操控和 4D 动态场景角度扩展了高斯世界模型。灵巧操控方面，ViserDex 和 D-REX 展示了 3DGS 在灵巧手仿真到真实迁移中的关键作用。可变形物体操控方面，DLO-Splatting 和颗粒介质操控展示了 3DGS 在非刚性物体操控中的潜力。

具身感知与交互（34 篇）展示了 3DGS 在遥操作、数字孪生和多模态记忆中的广泛应用。Reality Fusion 和 VR 遥操作框架（2025.04）~\cite{2504_15229} 证明了 3DGS 增强遥操作的用户体验优势。M3 和 URF-GS 展示了将 3DGS 与基础模型结合构建多模态空间记忆的前景。Event3DGS 展示了事件相机与 3DGS 结合在高速机器人自我运动中的价值。

仿真与真实迁移（32 篇）是 3DGS 在机器人领域最具变革性的应用方向。SplatSim 以 86.25\% 的零样本成功率展示了 3DGS 渲染替代传统网格渲染的可行性。R2R2R 证明了单个人类示范加 3DGS 重建可匹配 150 次遥操作示范的数据效果。GS-Playground 在 $640\times480$ 分辨率下实现 $10^4$ FPS 的突破性吞吐量，为大规模视觉 RL 铺平了道路。SoMA 在软体操控仿真中展示了 20\% 的重仿真准确率提升。

强化学习与场景表示（6 篇）探索了将 3DGS 作为 RL 环境表示的基础性方向。GSRL 和 Query-based Semantic Gaussian Field 首次验证了可泛化 3DGS 在 RL 中的优越性，S2E 展示了 RL 后训练对导航基础模型的增强效果。

展望未来，3DGS 在机器人导航与具身智能领域仍面临多个开放问题：（1）\textbf{闭环具身智能}——当前大多数系统将 3DGS 用于开环训练数据生成或静态地图构建，如何在闭环操控中实时更新 3DGS 表示并与策略学习耦合仍是核心挑战；（2）\textbf{多模态统一表示}——将视觉、触觉、力觉、语言等多模态信息统一嵌入 3DGS 表示，实现真正的多模态具身智能；（3）\textbf{动态场景长期维护}——在场景持续变化的长期部署中维护高质量的 3DGS 表示；（4）\textbf{实时物理仿真}——将 3DGS 的照片级渲染与高效物理仿真深度集成，支持闭环仿真到真实的无缝迁移。随着这些问题的逐步解决，3DGS 有望成为具身智能的核心基础设施，支撑从家庭服务机器人到工业自动化、从搜索救援到太空探索的广泛应用。
\end{futurebox}
